\documentclass[11pt]{article}
\usepackage[utf8]{inputenc}
\usepackage[T1]{fontenc}
\usepackage[a4paper,margin=1in]{geometry}
\usepackage{lmodern}
\usepackage{microtype}
\usepackage{amsmath,amssymb,mathtools}
\usepackage{mathrsfs}
\usepackage{tikz-cd}
\usepackage{enumitem}
\usepackage{longtable,booktabs,array}
\usepackage[hidelinks]{hyperref}
\setlength{\parindent}{0pt}
\setlength{\parskip}{0.55em}
\emergencystretch=3em
\hbadness=10000
\newcommand{\K}{K}
\newcommand{\Gr}{\operatorname{Gr}}
\newcommand{\Fil}{\operatorname{Fil}}
\newcommand{\ch}{\operatorname{ch}}
\newcommand{\Todd}{\operatorname{Todd}}
\newcommand{\cl}{\operatorname{cl}}
\newcommand{\Parf}{\operatorname{Parf}}
\newcommand{\Tor}{\operatorname{Tor}}
\newcommand{\Spec}{\operatorname{Spec}}
\newcommand{\RRR}{\mathrm{RRR}}
\providecommand{\codim}{\operatorname{codim}}
\providecommand{\Inv}{\operatorname{Inv}}
\providecommand{\Pic}{\operatorname{Pic}}
\providecommand{\Z}{\mathbb Z}
\providecommand{\Q}{\mathbb Q}
\providecommand{\R}{\mathbb R}
\providecommand{\Ahat}{\widehat A}
\providecommand{\Ahp}{\widehat A^{+}}
\providecommand{\That}{\widehat T}
\providecommand{\rad}{\operatorname{rad}}
\providecommand{\Gl}{\operatorname{Gl}}
\providecommand{\GL}{\operatorname{GL}}
\providecommand{\ch}{\operatorname{ch}}
\providecommand{\Todd}{\operatorname{Todd}}
\providecommand{\Gr}{\operatorname{Gr}}

% Core macro additions
\providecommand{\Atil}{\widetilde A}
\providecommand{\Ctil}{\widetilde C}
\providecommand{\K}{\mathrm K}
\providecommand{\cC}{\mathcal C}
\providecommand{\cF}{\mathscr F}
\providecommand{\eps}{\varepsilon}
\providecommand{\ellc}{\ell}
\providecommand{\Tor}{\operatorname{Tor}}
\providecommand{\Coker}{\operatorname{Coker}}
\providecommand{\Ker}{\operatorname{Ker}}
\providecommand{\Supp}{\operatorname{Supp}}
\providecommand{\codim}{\operatorname{codim}}
\providecommand{\Gr}{\operatorname{G}}
\providecommand{\rank}{\operatorname{rank}}
\providecommand{\rang}{\operatorname{rang}}
\providecommand{\cO}{\mathcal O}

\hypersetup{hypertexnames=false}
\providecommand{\N}{\mathbb N}
\begin{document}
\begin{center}
{\Large Seminar on Algebraic Geometry of Bois Marie}\par
{\large 1966--67}\par
\vspace{1.5em}
{\Large Intersection Theory and the Riemann--Roch Theorem}\par
{\large (SGA 6)}\par
\vspace{1.5em}
a seminar directed by\par
P. Berthelot, A. Grothendieck, L. Illusie\par
with the collaboration of\par
D. Ferrand, J.-P. Jouanolou, O. Jussila, S. Kleiman, M. Raynaud, J.-P. Serre
\end{center}

\section*{Preface}
The present edition of \emph{Séminaire de Géométrie Algébrique 6} reproduces, up to minor changes, the first edition, distributed by the I.H.E.S. in the form of mimeographed exposés. The corrections made concern essentially printing errors, except in Exposés I and III, where slight modifications were made by L. Illusie to a few statements which were incorrect in the primitive version. Footnotes have also been added, especially in Exposé XIV by A. Grothendieck, in order to point out recent progress on questions that were still open at the time of the seminar.

Finally, note that Exposé XI, by D. Ferrand, on determinant theory for perfect complexes, was not written up and therefore will not be published; to avoid errors in references, we have left unchanged the numbering of the following exposés.

\begin{flushright}June 1971\end{flushright}

\section*{Introduction}
The purpose of the present seminar is to develop a global theory of intersections, and a Riemann--Roch formula, for arbitrary schemes. The reader will find a description of the programme of the seminar in Exposé 0. The possibility in principle of proving a Riemann--Roch formula for general regular schemes, along the lines of the well-known Borel--Serre report of 1958, was already clear at that time, at least for a projective morphism. The extension to the general case is less evident; the programme in which it belongs, and which is partly carried out in this seminar, goes back to 1960. As was also the case for the duality theory of coherent sheaves, an essential tool for formulating a satisfactory theory is Verdier's theory of derived categories; familiarity with it is indispensable for understanding the seminar.

Apart from this theory, the study of the seminar requires little more than a general knowledge of the foundations of algebraic geometry as set out in EGA, chapters I, II, III. We shall only occasionally refer to EGA IV, for certain facts concerning flat or smooth morphisms, which for the most part also appear in the first exposés of SGA 1. Finally, to develop some results with the full generality desirable for applications, we sometimes use the notion of a ringed site and of a ringed topos, for which we refer to SGA 4, Exposés I to IV. A reader unfamiliar with the language of sites and topoi may everywhere replace these objects by ordinary topological spaces, the objects of the topos then being replaced by open subsets of these spaces; nevertheless, we recommend that the reader assimilate the language of topoi, which provides an extremely convenient unifying principle.

The fundamental notion for the theory presented here is that of a perfect complex of modules, together with its various variants, developed in Exposés I, II, III. It seems clear that these notions, also important in other contexts in algebraic geometry, notably for formulas of Lefschetz--Verdier type in cohomology with discrete coefficients (SGA 5) or coherent coefficients, will also have a role to play outside algebraic geometry, especially in the formulation of a complex analytic variant of the Riemann--Roch theorem presented here, or of suitable variants of the Atiyah--Singer index theorem. Some indications in this direction will be given in Exposé II. Unfortunately, at present there is still no statement, even heuristic, which would encompass these two theorems, the first of which remains conjectural for the moment. A new idea seems to be missing, just as one is also missing in algebraic geometry for proving Riemann--Roch beyond projective hypotheses (cf. Exposé XIV, no. 2).

We shall make no use in this seminar of the local theory of intersections, set out in J.-P. Serre, \emph{Algèbre locale. Multiplicités} (Lecture Notes in Mathematics no. 11, Springer), and merely point out here that Serre's course had an evident influence on the genesis of the theory in 1957. Nor shall we use the theory of the Chow ring developed in the Chevalley Seminar of 1958, Exposés 2 and 3 (École Normale Supérieure). That theory is essentially tied to nonsingularity hypotheses, whereas the aim of our seminar is, on the contrary, to develop a theory of intersections on general schemes, and even on general locally ringed topoi. On this matter see the comments in Exposé XIV, nos. 4 and 8, which give the relations between our theory and that of the Chow ring and raise the question of a generalization of the latter. Thus one may say that the seminar presents a coherent and self-contained theory of intersections, but one which is by no means exhaustive of the different points of view useful, indeed indispensable, in intersection theory; consequently it should be complemented by the cited exposés of Serre and Chevalley.

We have appended to the seminar, as an appendix to Exposé 0, Grothendieck's 1957 report, ``Classes of sheaves and the Riemann--Roch theorem'', which served as the basis for the Borel--Serre seminar at Princeton in the same year and for their already cited article. That report sketches two proofs of the Riemann--Roch theorem for nonsingular quasi-projective varieties. The first, valid for the time being only in characteristic zero but giving, in return, a slightly more precise result in the case of an immersion, has not yet appeared in the literature, apart from the work of Borel--Serre and the present seminar. The report presupposes, of course, none of the other exposés of the seminar, and can even serve, like Exposé 0, as an introduction to the study of the latter, especially for a reader not yet familiar with Borel--Serre. The proof just alluded to applies essentially unchanged to the case of a projective morphism of complex analytic spaces, and will doubtless also apply in arbitrary characteristic once the problem of the operations ``exterior powers'' in the derived category is solved (cf. Exposé XIV, nos. 1, 2, 3). The absence of a study of this question is no doubt the most serious gap in this seminar, from the point of view we have adopted.

One will notice the absence, in the table of contents of the present seminar, of two of the participants mentioned on the cover page. One of them, J.-P. Jouanolou, took an active part in the technical elaboration of the first part of the seminar, but was prevented from taking part in the oral exposés; in particular, he is responsible for the linear-independence assertion contained in the important statement VI 1.1, giving the structure of the ring $K^\circ$ of classes of locally free modules on a projective bundle, which had previously been proved only under the assumption that the base scheme has an ample invertible module. The other, J.-P. Serre, gave two oral exposés which were logically independent of the rest of the seminar, and therefore preferred to publish them in the form of a separate article (cf. J.-P. Serre, ``Grothendieck groups of split reductive group schemes'', to appear in \emph{Publications Mathématiques}, no. 34).

As in the other parts of the \emph{Séminaire de Géométrie Algébrique du Bois Marie}, the sigla SGA 1 to SGA 6 refer to the different parts of that seminar; the Roman numeral following the siglum indicates the number of the exposé, and the Arabic numerals following it correspond to the internal numbering of the exposé in question. For internal references within the present seminar the siglum SGA 6 is omitted, and references of the form I 4.9 are used, meaning: Exposé I, statement 4.9. The siglum EGA refers to the \emph{Éléments de Géométrie Algébrique} of J. Dieudonné and A. Grothendieck.

\section*{Table of contents}
\begin{longtable}{@{}p{0.16\linewidth}p{0.64\linewidth}r@{}}
\toprule
Exposé & Title and author & Page \\
\midrule
0 & Outline of a programme for a theory of intersections on general schemes, by A. Grothendieck & 1 \\
Appendix & Classes of sheaves and the Riemann--Roch theorem, by A. Grothendieck & 20 \\
I & Generalities on finiteness conditions in derived categories, by L. Illusie & 56 \\
II & Existence of global resolutions, by L. Illusie & 160 \\
III & Relative finiteness conditions, by L. Illusie & 190 \\
IV & Grothendieck groups of ringed topoi, by L. Illusie & 222 \\
V & Generalities on $\lambda$-rings, by P. Berthelot & 297 \\
VI & The $K^\circ$ of a projective bundle: computations and consequences, by P. Berthelot & 365 \\
VII & Regular immersions and the computation of $K^\circ$ of a blow-up, by P. Berthelot & 416 \\
VIII & The Riemann--Roch theorem, by P. Berthelot & 466 \\
IX & Some computations of $K_\bullet$-groups, by P. Berthelot & 498 \\
X & Formalism of intersections on proper algebraic schemes, by O. Jussila, with an appendix by A. Grothendieck, ``Specialization in intersection theory'' & 519 \\
XI & Not written up & -- \\
XII & A relative representability theorem for the Picard functor, by M. Raynaud, written by S. Kleiman & 595 \\
XIII & Finiteness theorems for the Picard functor, by S. Kleiman & 616 \\
XIV & Open problems in intersection theory, by A. Grothendieck & 667 \\
\bottomrule
\end{longtable}

\section*{Exposé 0}
\begin{center}
{\large Outline of a programme for a theory of intersections on general schemes}\par
by A. Grothendieck
\end{center}

The present exposé is introductory in nature, and its reading is not logically indispensable for the study of the seminar. It is addressed more particularly to readers familiar with the provisional version of the Riemann--Roch theorem, as it is presented in the Borel--Serre report or in Grothendieck's report already mentioned in the introduction, cited below as [RRR], and reproduced as an appendix at the end of the present exposé.

\subsection*{1. Reminder on the Riemann--Roch formula}
Recall the Riemann--Roch formula for a proper morphism
\[
f:X\longrightarrow Y
\]
of quasi-projective smooth schemes over a field $k$, and a coherent sheaf $F$ on $X$, whose class in the group $K(X)$ of classes of coherent sheaves on $X$ is denoted by $\cl(F)$:
\begin{equation}\tag{1.1}
\Todd(T_Y)\,\ch_Y\bigl(f_!(\cl(F))\bigr)=f_*\bigl(\Todd(T_X)\,\ch_X(F)\bigr).
\end{equation}
This formula is valid in $A(Y)\otimes\mathbb Q$, where $A(Y)$ is the Chow ring of $Y$. The $f_*$ on the right hand side is obtained by tensoring with $\mathbb Q$ the homomorphism ``direct image of cycles''
\[
f_*:A(X)\longrightarrow A(Y),
\]
while the term on the left is the Euler--Poincaré characteristic of $F$ relative to $f$:
\[
f_!(\cl(F))=\sum_i(-1)^i\cl\bigl(R^if_*(F)\bigr).
\]
The symbols $\ch_X$, $\ch_Y$ denote the Chern characters on $X$ and on $Y$, and $T_X$, $T_Y$ the tangent bundles of $X$ and $Y$. As is known, $\Todd(-)$ and $\ch(-)$ are universal polynomials with coefficients in $\mathbb Q$ in the Chern classes of the argument. Since the constant term of $\Todd(-)$ is $1$, it is an invertible element for every value of the argument. Thus, after multiplying by $\Todd(T_Y)^{-1}$, formula (1.1) takes the equivalent form, more convenient for our purposes,
\begin{equation}\tag{1.2}
\ch_Y\bigl(f_!(\cl(F))\bigr)=f_*\bigl(\Todd(T_f)\,\ch_X(F)\bigr),
\end{equation}
where
\begin{equation}\tag{1.3}
T_f=T_X-f^*T_Y\in K(X).
\end{equation}
Thus $T_f$ plays the role of a virtual relative tangent bundle of $X$ over $Y$. When $f$ is smooth, one simply has $T_f=T_{X/Y}$, whereas when $f:X\to Y$ is an immersion, one obtains
\[
T_f=-\check N_{X/Y},
\]
where $N_{X/Y}$ denotes the normal sheaf of $X$ in $Y$.

One of the main aims of the seminar is to generalize (1.2) simultaneously in two directions:
\begin{enumerate}[label=\alph*)]
\item to get rid of the hypothesis that a base field $k$ exists;
\item to replace the regularity hypotheses on $Y,X$ by a hypothesis of ``local regularity'' on $f$.
\end{enumerate}
Finally, along the way, we shall also deal with the following problem:
\begin{enumerate}[label=\alph*),start=3]
\item to eliminate the quasi-projectivity hypotheses which, in the absence of a base field, are expressed by the existence of ample invertible modules on $X$ and $Y$.
\end{enumerate}

\subsection*{2. Removing the base field}
Let us first examine generalization a), while retaining the hypotheses of regularity and of the existence of ample invertible modules on $X$ and on $Y$.

The definitions of $K(X)$, $K(Y)$ and of the homomorphism $f_!:K(X)\to K(Y)$ then present no new problem, thanks to the fact that $X$ and $Y$ are regular. The most natural way to give a meaning to (1.2) would seem to consist in defining Chow rings $A(X),A(Y)$ and a group homomorphism
\[
f_*:A(X)\longrightarrow A(Y),
\]
in establishing a theory of Chern classes, giving maps
\[
c_i:K(X)\longrightarrow A(X),
\]
and similarly for $Y$, and finally in describing an element, the virtual relative tangent bundle,
\[
T_f\in K(X).
\]

\subsubsection*{2.1}
For the latter, a definition presents itself rather evidently. One uses the fact that the morphism $f$, thanks to the hypothesis that ample invertible modules exist on $X$ and on $Y$, may be factored as
\[
X\xrightarrow{i}X'\xrightarrow{f'}Y,
\]
where $i$ is a closed immersion and $f'$ is smooth; for example, one may take $X'=\mathbb P^r_Y$ for suitable $r$. One then sets
\begin{equation}\tag{2.1}
T_f=i^*T_{X'/Y}-\check N_{X/X'}\in K(X),
\end{equation}
where $\Omega^1_{X'/Y}$ is the locally free module of relative $1$-differentials of $X'$ over $Y$, or the relative cotangent module of $X'$ over $Y$, and $N_{X/X'}$ is the conormal module of $X$ in $X'$, by definition equal to $J/J^2$, where $J$ is the ideal on $X'$ defining $X$. The regularity hypotheses on $X$ and $X'$ also imply that $N_{X/X'}$ is locally free. Finally, $\check N$ denotes the dual of $N$. One verifies without difficulty, in Exposé VIII, that the element $T_f$ defined by (2.1) is independent of the chosen factorization of $f$.

\subsubsection*{2.2}
As for a definition, under the conditions considered, of a Chow ring and a corresponding theory of Chern classes, although it is plausible that such a theory should be developable on the model of smooth varieties over a field, it seems that this work has not been done at the present time. We shall therefore adopt another viewpoint, by directly defining a ring which replaces the Chow ring, and whose tensor product with $\mathbb Q$ is indeed isomorphic to $A(X)\otimes\mathbb Q$ in the classical case.

A first natural filtration of $K(X)$, the one considered in [RRR], is the topological filtration by the codimension of the support of coherent sheaves: $\Fil^iK(X)$ is generated by the $\cl(F)$ with $F$ coherent and $\operatorname{codim}(\operatorname{Supp}F,X)\ge i$. A first question is the compatibility of this filtration with the multiplicative structure of $K(X)$. In [RRR], this compatibility is obtained, for $X$ smooth over a field, by means of Chow's moving lemma. Assuming this assertion, one obtains a graded ring $\Gr_{\mathrm{top}}(X)$ and a canonical surjective homomorphism
\[
\varphi:A(X)\longrightarrow \Gr_{\mathrm{top}}(X),
\]
sending the class of a prime cycle $Z$ to $\cl(\mathcal O_Z)$; it then turns out, Exposé XIV 4.2, that the kernel is a torsion subgroup. Moreover, [RRR] shows how the images in $\Gr_{\mathrm{top}}(X)$ of the Chern classes $c_i(F)$ can be computed directly in terms of $\cl(F)\in K(X)$ and the exterior-power operations $\lambda^i$ in $K(X)$.

This shows that, for a noetherian regular scheme $X$ endowed with an ample invertible module, the ring $\Gr_{\mathrm{top}}(X)$ is a convenient substitute for the hypothetical Chow ring $A(X)$. It also has the advantage over the latter that certain computations of [RRR], leading to important formulas in this ring, are performed directly in it--formulas which, even for $X$ projective and smooth over a field, have not at present been established in $A(X)$ itself (Exp. XIV, 4.3 and 6.4). To justify this viewpoint completely, however, one would have to solve the problem mentioned above of the compatibility of the topological filtration of $K(X)$ with its ring structure, a problem which seems essentially connected with Chow's moving lemma, itself the essential technical ingredient of the theory of the Chow ring which we have claimed to be able to avoid.

\subsubsection*{2.3}
To avoid this problem, one endows $K(X)$ with a second filtration, finer than the topological filtration, and describable in purely algebraic terms from the structure of $K(X)$ as an augmented $\lambda$-ring. The definition of this filtration is suggested by the expression given in [RRR] for the Chern classes of an element $x\in K(X)$ in $\Gr_{\mathrm{top}}(X)$:
\begin{equation}\tag{2.3}
c_i(x)=\gamma^i(x-\epsilon(x))\mod \Fil^{i+1}K(X),
\end{equation}
where $\gamma^i$ denotes a certain linear combination of the $\lambda^j$ ($j\le i$), made explicit in loc. cit., and where $\epsilon$ is the augmentation. One shows that $\gamma^i(x-\epsilon(x))$ is indeed of topological filtration at least $i$, which gives meaning to formula (2.3). One then defines the $\lambda$-filtration of $K(X)$ by the ideals $\Fil^i_\lambda K(X)$, where $\Fil^i_\lambda K(X)$ consists of finite sums of expressions of the form
\[
\gamma^{i_1}(x_1-\epsilon(x_1))\cdots \gamma^{i_r}(x_r-\epsilon(x_r)),\qquad i_1+\cdots+i_r\ge i.
\]
The associated graded ring for this filtration will be denoted $\Gr^\bullet(X)$; it is this ring which will replace the Chow ring of $X$ in the seminar. In the case under consideration here, one can moreover show that the canonical homomorphism
\begin{equation}\tag{2.4}
\Gr^\bullet(X)\longrightarrow \Gr_{\mathrm{top}}(X)
\end{equation}
is an isomorphism modulo torsion, justifying the viewpoint that, after tensoring with $\mathbb Q$, this ring can play exactly the same role as the Chow ring tensored with $\mathbb Q$. On the other hand, exactly what was needed has been done so that formula (2.3), with $\Fil_{\mathrm{top}}$ replaced by $\Fil_\lambda$, defines a theory of Chern classes with values in $\Gr^\bullet(X)$; the Chern classes in $\Gr_{\mathrm{top}}(X)$ are then deduced by means of the homomorphism (2.4). It is shown in Exposé V, using the fact proved in Exposé VI that $K(X)$ is a special $\lambda$-ring in the sense of [RRR], that is, a $\lambda$-ring in the sense of Exposé V, that this notion of Chern class satisfies all the usual formal properties reviewed in [RRR].

\subsubsection*{2.4}
One should note a very appreciable advantage of the definition of $\Gr^\bullet(X)$ over that of $A(X)$ and $\Gr_{\mathrm{top}}(X)$: it remains meaningful for an absolutely arbitrary scheme $X$, and more generally for an arbitrary locally ringed topos $X$. One then takes for $K(X)$ the group of classes constructed from locally free modules on $X$, which is again an augmented $\lambda$-ring, hence equipped with a corresponding filtration and an associated graded $\Gr^\bullet(X)$. The price of this generality, however, is that when one is not willing to neglect torsion phenomena, i.e. to tensor with $\mathbb Q$, the properties of $\Gr^\bullet(X)$ are in several respects less satisfactory than those of its two competitors (Exp. XIV no. 4). Thus, assuming again that $f:X\to Y$ is a proper morphism of noetherian regular schemes endowed with ample invertible modules, it is not true in general that the homomorphism
\[
f_*:K(X)\longrightarrow K(Y)
\]
maps $\Fil^i_\lambda K(X)$ into $\Fil^{i-d}_\lambda K(Y)$, where $d$ is the relative dimension of $X$ over $Y$, and hence defines, by passage to associated gradeds, a degree $-d$ homomorphism from $\Gr^\bullet(X)$ to $\Gr^\bullet(Y)$. However, Exposé VII proves that this is true after tensoring with $\mathbb Q$, so that one obtains a degree $-d$ homomorphism
\begin{equation}\tag{2.5}
f_*:\Gr^\bullet(X)\otimes\mathbb Q\longrightarrow \Gr^\bullet(Y)\otimes\mathbb Q,
\end{equation}
which plays the role of the homomorphism ``direct image of cycles'' $A(X)\to A(Y)$.

Having said this, we see that we possess all the structural elements needed to give meaning to formula (1.2) when the existence of a base field is not postulated, with $f:X\to Y$ a proper morphism of noetherian regular schemes admitting ample invertible modules.

\subsection*{3. The case of general schemes}
We now show how one can still give meaning to formula (1.2) when the regularity hypothesis on $X,Y$ is abandoned and replaced by a local hypothesis on the morphism $f$, to be specified below.

\subsubsection*{3.1}
A first difficulty comes from the fact that the observation which served as the starting point of the theory [RRR], namely that the group $K(X)$ of classes of sheaves on $X$ is the same whether it is defined in terms of coherent sheaves or in terms of locally free sheaves, is essentially tied to the regularity hypothesis on $X$. For a general scheme $X$, say noetherian for simplicity, one must distinguish between these two groups; we shall denote them respectively by $K_\bullet(X)$ and $K^\bullet(X)$. The first is covariant in $X$ for proper morphisms, the second contravariant in $X$ for arbitrary morphisms. Notice that $K^\bullet(X)$ is naturally endowed with a ring structure, indeed with an augmented $\lambda$-ring structure, whereas $K_\bullet(X)$ no longer has, in general, a ring structure, but only a module structure over $K^\bullet(X)$. The homomorphism $x\mapsto x\cdot\cl(\mathcal O_X)$ from $K^\bullet(X)$ to $K_\bullet(X)$ is no longer an isomorphism in general. For our formulation of the generalized Riemann--Roch theorem, we shall work exclusively with $K^\bullet(X)$ and $K^\bullet(Y)$, to the exclusion of $K_\bullet(X)$ and $K_\bullet(Y)$. As noted in 2.3, these augmented $\lambda$-rings, equipped with their natural $\lambda$-filtration, give rise to graded rings $\Gr^\bullet(X)$, $\Gr^\bullet(Y)$, and to Chern-class homomorphisms
\[
c_i:K^\bullet(X)\longrightarrow \Gr^\bullet(X),
\]
and similarly for $Y$. The difficulty, however, lies in defining a direct-image homomorphism
\begin{equation}\tag{3.1}
f_*:K^\bullet(X)\longrightarrow K^\bullet(Y),
\end{equation}
which can no longer be defined by the naive formula
\[
f_*(\cl(F))=\sum_i(-1)^i\cl\bigl(R^if_*(F)\bigr),
\]
because even if $F$ is locally free, the modules $R^if_*(F)$ are in general no longer locally free, nor even of finite Tor-dimension, and the preceding formula loses its meaning. A solution to this difficulty is furnished by the viewpoint of derived categories. Indeed, whereas the $R^if_*(F)$ are in general ``bad'' sheaves, say of infinite Tor-dimension, the complex $Rf_*(F)$ on $Y$, the direct image of $F$ in the sense of derived categories, tends to be excellent. More precisely, when $F$ has finite Tor-dimension over $Y$, for instance when $f$ has finite Tor-dimension and $F$ is locally free on $X$, then $Rf_*(F)$ is perfect on $Y$, that is, has coherent cohomology and globally finite Tor-dimension (Exp. III). It follows easily, Exp. IV, that its class in the Grothendieck group of perfect complexes has a meaning in $K^\bullet(Y)$.

\subsubsection*{3.2}
To give a meaning to (1.2), it remains to impose on $f$ a hypothesis making it possible to give a meaning to the homomorphism (2.5), i.e. ensuring that the homomorphism (3.1) just defined is, up to a shift and modulo torsion, compatible with the filtrations of these rings; and finally one must again define an element, the ``virtual relative tangent bundle'',
\[
T_f\in K^\bullet(X).
\]
For the latter, returning to the definition proposed in 2.1, one sees that a hypothesis is needed ensuring that, with the notation of loc. cit., the conormal module $N_{X/X'}$ is locally free. The most natural hypothesis to this end is that $i$ be a ``regular immersion'', i.e. identify $X$ with a subscheme of $X'$ defined by an ideal which locally can be defined by an $\mathcal O_{X'}$-regular sequence of sections of $\mathcal O_{X'}$. This hypothesis, purely local in nature, in fact does not depend on the choice of the factorization of $f$ into an immersion followed by a smooth morphism (Exp. VIII), and, as above, the same is true for the class $T_f$ defined by formula (2.1). When the hypothesis just considered is satisfied, the morphism $f$ is called a ``locally complete intersection'' morphism. Such a morphism is also locally of finite Tor-dimension.

We shall also allow ourselves to omit the word ``locally'' in this terminology, which cannot cause any confusion.

If $f$ is a locally complete intersection morphism, one introduces a natural notion of ``virtual relative dimension'' $d(f)$ by the formula
\[
d(f)(x)=\dim_x\bigl(f^{-1}(f(x))\bigr)-\codim_x(X,X').
\]
It is an integer-valued function, locally constant at the point $x\in X$, hence constant if $X$ is connected; moreover it does not depend on the choice of the factorization (2.1). When this virtual relative dimension $d$ is constant, when $f$ is proper, and when $X$ and $Y$ admit ample invertible modules, one can prove (Expos\'e VIII), although this is far from a trivial result, that
\begin{equation}\tag{3.4}
f_*\bigl(\Fil^i K^\bullet(X)_{\mathbb Q}\bigr)\subset \Fil^{i-d}K^\bullet(Y)_{\mathbb Q},
\end{equation}
whence a homomorphism (3.3) of degree $-d$.

\subsubsection*{3.3}
We have therefore assembled all the structural elements needed to give a meaning to formula (1.2) in the case of a proper locally complete intersection morphism of noetherian schemes $X,Y$, both admitting ample invertible modules. With a little additional care, this formulation in fact retains a meaning independently of any noetherian hypothesis.

In this general form, this Riemann--Roch formula will be proved in Expos\'e VIII. The reader will observe that the proof given in this seminar combines elements of the two proofs which appear in [RRR] and which were discussed in the Introduction. The need to resort to the methods of the first of these proofs, involving more $K$-formalism than the second, followed in Borel--Serre, is due to the need for a proof of the inclusion (3.4), a question which did not arise in the more restricted framework of [RRR] and Borel--Serre [2], where one could work with the Chow ring or the ring $\Gr^\bullet_{\mathrm{top}}$. From the second proof, we retain the introduction of the scheme obtained by blowing up $Y$ along $X$ in the case where $f$ is an immersion. This technical artifice will no doubt disappear, and the proof in the case of an immersion will extend to the case where one no longer postulates the existence of an ample invertible module on $Y$, once the question raised in Expos\'e XIV, no. 1, concerning the introduction of exterior-power operations in the derived category $D^-(Y)$ has been resolved.

\subsection*{4. Removing the ample invertible module hypothesis}
We must finally give a few indications of how one can still give a meaning to (1.2) in the absence of a hypothesis asserting the existence of ample invertible modules on $X$ and $Y$.

\subsubsection*{4.1}
First note that for an arbitrary scheme $X$, or more generally an arbitrary locally ringed topos, it seems beyond doubt that the ``right'' invariant which should replace the ring $K(X)$ of [RRR] is not the ring $K^\bullet(X)$ considered above, defined in terms of locally free modules on $X$, but rather the group $K(\Parf(X))$ considered in 3.1, defined in terms of perfect complexes on $X$. We shall denote this group here by $K^\bullet(X)$, and distinguish it from the ``naive'' group of 3.1, denoted $K^\bullet_{\mathrm{naive}}(X)$. This point of view is forced by the need to define a direct-image homomorphism (3.1) for a proper morphism
\[
f:X\longrightarrow Y
\]
of finite Tor-dimension, say between noetherian schemes; the non-noetherian case requires somewhat stronger local finiteness properties on $f$, cf. Expos\'e III. With the definition of $K^\bullet(X)$ now adopted, this homomorphism $f_*$ is indeed defined in a trivial way by formula (3.2), where now $F$ is an arbitrary perfect complex on $X$, which implies that $Rf_*(F)$ is also perfect.

\subsubsection*{4.2}
This new definition of $K^\bullet(X)$ raises, however, a new problem in homological algebra: that of defining a $\lambda$-structure on $K^\bullet(X)$ compatible with its natural ring structure coming from tensor product in the derived category. If $L$ and $L'$ are two perfect complexes, then so is their total tensor product $L\otimes^{\mathbf L}L'$. More precisely, one would need to define ``exterior-power'' operations in the derived category $D^-(X)$, transforming perfect complexes into perfect complexes. This question has not yet been clarified at present, and, as was said in the Introduction, it constitutes the most serious gap in the seminar. Once this question is solved, and proceeding as in 2.3, one obtains a natural filtration on $K^\bullet(X)$, allowing one to form the associated graded ring $\Gr^\bullet(X)$, which will play the role of the Chow ring, and to define a theory of Chern classes
\[
c_i:K^\bullet(X)\longrightarrow \Gr^i(X).
\]

\subsubsection*{4.3}
To give a meaning to formula (1.2) for a proper locally complete intersection morphism
\[
f:X\longrightarrow Y
\]
of noetherian schemes, one must still include in the statement of this formula the nontrivial inclusions (3.4), where $d$ is the virtual relative dimension, in order to be able to define (3.3); and finally one must define the virtual relative tangent class
\[
T_f\in K^\bullet(X).
\]
When $f$ admits a factorization (2.1) through an immersion into a relative smooth scheme, the definition of $T_f$ presents no new difficulty. It can moreover be reformulated in that case, in a way more consonant with the spirit of derived categories, by considering the well-known canonical homomorphism
\begin{equation}\tag{4.1}
N_{X/X'}\longrightarrow \Omega^1_{X'/Y}\otimes_{\mathcal O_{X'}}\mathcal O_X
\end{equation}
induced by the exterior differential, and the chain complex $L_{X/Y}$ defined from (4.1) by ``extending by zeros'', with the complex equal in degree $1$, respectively $0$, to the source, respectively the target, of the homomorphism (4.1). Formula (2.2) then takes the more satisfactory form
\begin{equation}\tag{4.2}
T_f=\cl\bigl(L_{X/Y}^{\vee}\bigr),
\end{equation}
a formula which makes sense in $K^\bullet(X)$ when $f$ is a locally complete intersection morphism, since in that case the complex $L_{X/Y}$ is manifestly perfect. The assertion that $T_f$ is invariant relative to the chosen factorization of $f$ can then be made precise as an invariance assertion for the complex $L_{X/Y}$: up to canonical isomorphism in the derived category $D(X)$, this complex is likewise independent of the chosen factorization (Expos\'e VIII).

\subsubsection*{4.4}
When $f:X\to Y$ is an arbitrary morphism of schemes, one can still construct a canonical object $L_{X/Y}$ of the derived category $D(X)$, defined up to unique isomorphism, which on each affine open of $X$ lying above an affine open of $Y$ is isomorphic to the complex constructed according to the principle explained in the preceding paragraph. In the general case, however, ``smooth'' should be replaced by ``formally smooth'', so that every algebra $B$ over a ring $A$ occurs as a quotient of a formally smooth algebra, for example a polynomial algebra. For such a definition it is not possible to proceed by simple gluing from the affine constructions, since the objects of the derived category $D(X)$ are, essentially, not glueable.

Here is a general construction of $L_{X/Y}$, valid more generally for every morphism
\[
f:X\longrightarrow Y
\]
of commutatively ringed topoi. Let $B=\mathcal O_X$ and $A=f^{-1}\mathcal O_Y$, so that $A$ is a ring on the space $X$ and $B$ is an $A$-algebra. Let $T\to B$ be a homomorphism of sheaves of sets which generates $B$ as an $A$-algebra, i.e. such that the corresponding homomorphism of $A$-algebras
\[
C=A[T]\longrightarrow B
\]
is an epimorphism of sheaves of sets; here $A[T]$ denotes the commutative $A$-algebra freely generated by $T$, that is, the sheaf associated with the presheaf
\[
U\longmapsto A(U)[T(U)].
\]
Let $J$ be the kernel of $C\to B$, and consider, as in 4.3, the chain complex of length $1$ defined by the natural homomorphism
\[
J/J^2\longrightarrow \Omega^1_{C/A}\otimes_C B.
\]
It is not difficult to see that, up to canonical isomorphism in the category $D(X)$, the complex so constructed is independent of the choice of $T$ and of $T\to B$, so that it deserves the designation of the relative cotangent complex $L_{X/Y}$ of $X$ over $Y$. One also verifies that, on every affine open of $X$ over an affine open of $Y$, it is given, up to canonical isomorphism, by the procedure of 4.3.

In the case where $f$ is locally complete intersection, one sees moreover that the result is a perfect complex, which makes it possible again to define an element $T_f$ of $K^\bullet(X)$ by formula (4.2). Here again, in general, it would not have been possible to define $T_f$ as an element of $K^\bullet_{\mathrm{naive}}(X)$, and nothing permits one to suppose that $T_f$ is globally isomorphic in $D(X)$ to a bounded complex with locally free components of finite type. This is therefore an additional justification for the point of view advanced in 4.1 in favour of the ``sophisticated'' $K^\bullet(X)$, in preference to $K^\bullet_{\mathrm{naive}}(X)$.

\subsubsection*{4.5}
Subject to the question raised in 4.2, we have therefore again assembled all the structural elements needed to give a meaning to (1.2) for a proper locally complete intersection morphism of arbitrary schemes; the noetherian hypothesis implicitly made in the preceding sections can in fact be eliminated without any essential difficulty. Unfortunately, this does not mean that, even subject to this reservation, the formula in question is established, even in the case where $Y$ is the spectrum of an algebraically closed field of characteristic $p>0$ and $X$ is, say, a proper smooth algebraic scheme of dimension $3$. See the comments in Expos\'e XIV, no. 2.

\subsection*{5. Analytic variant}
In the form presented in no. 4, the Riemann--Roch formula (1.2) also has a meaning for a proper locally complete intersection morphism of complex analytic spaces, or of rigid analytic spaces in the sense of Tate. But of course the proof given in [RRR] applies to this case only when one assumes in addition that the morphism $f$ is projective; and, as in the algebraic case, it seems that an essentially new idea is needed to treat the general case. The formula we have in view here is a formula with values in an ``analytic Chow ring'' $\Gr^\bullet(X)$, or rather $\Gr^\bullet(X)_{\mathbb Q}$, constructed according to the principle explained in 4.2, the difficulty encountered there disappearing, moreover, because one is in characteristic zero, as is noted in Expos\'e XIV, no. 1.

When $X$ and $Y$ are nonsingular, it is possible to give another version of the Riemann--Roch formula by proceeding as follows. First let $X$ be an arbitrary complex analytic space, and let $X^{\mathrm{top}}$ be the same space equipped with the sheaf of rings of continuous complex-valued functions. We therefore have a canonical morphism of ringed spaces
\[
X^{\mathrm{top}}\longrightarrow X,
\]
which, by the evident contravariant nature of the functor $K^\bullet(X)$ in the variable ringed space $X$, defines a homomorphism of rings
\begin{equation}\tag{5.1}
K^\bullet(X)\longrightarrow K^\bullet(X^{\mathrm{top}}).
\end{equation}
For simplicity suppose that $X^{\mathrm{top}}$ is compact; then, as in the case of a scheme admitting an ample invertible module, one proves (Expos\'e II) that every perfect complex on $X^{\mathrm{top}}$ is globally isomorphic to a bounded complex, locally free in each degree, i.e. to a complex of complex vector bundles on $X^{\mathrm{top}}$, and hence that $K^\bullet(X^{\mathrm{top}})$ is canonically isomorphic to the well-known group $K(X^{\mathrm{top}})$ of topologists [1], defined in terms of complex vector bundles on the compact space $X^{\mathrm{top}}$. Thus one may interpret (5.1) as a homomorphism from the $K^\bullet(X)$ defined in terms of the analytic structure to the well-known ring $K(X^{\mathrm{top}})$ of topologists. If now
\[
f:X\longrightarrow Y
\]
is a morphism of compact nonsingular analytic spaces, one knows how to associate to it, thanks to Atiyah--Hirzebruch [1], a group homomorphism
\begin{equation}\tag{5.2}
f^{\mathrm{top}}_*:K^\bullet(X^{\mathrm{top}})\longrightarrow K^\bullet(Y^{\mathrm{top}}).
\end{equation}
Using likewise the homomorphism $f_*$ of (3.1), defined by formula (3.2), which here implicitly uses Grauert's finiteness theorem [4] for a proper morphism of analytic spaces, one obtains a square of natural homomorphisms involving only groups of type $K^\bullet$, to the exclusion of rings of Chow-ring or cohomology type:
\begin{equation}\tag{5.3}
\begin{array}{ccc}
K^\bullet(X) & \longrightarrow & K^\bullet(X^{\mathrm{top}}) \\
\downarrow f_* && \downarrow f^{\mathrm{top}}_* \\
K^\bullet(Y) & \longrightarrow & K^\bullet(Y^{\mathrm{top}}).
\end{array}
\end{equation}
The Riemann--Roch type formula to which we are alluding consists in asserting the commutativity of this square. Note the difference in nature between this assertion, which does not neglect torsion phenomena and is situated in a group $K^\bullet(Y^{\mathrm{top}})$ of an essentially discrete nature, and the Riemann--Roch formula of the type considered in no. 4, which is situated in a group of ``Chow'' type no longer of a discrete nature, but which, on the other hand, must be tensored with $\mathbb Q$. Let us note that formula (5.3) is proved at any rate when $X$ and $Y$ are nonsingular projective varieties [7], as is the formula of the type considered in no. 4, proved in this case in [RRR]. What is involved here is a plausible variant of the Atiyah--Singer index theorem, and of its generalizations due to Shih [7] and Atiyah, which obviously calls for a common generalization. Once the commutativity of (5.3) is proved, one would immediately deduce from it, using the ``differentiable Riemann--Roch theorem'' of Atiyah--Hirzebruch [1], a cohomological variant of the complex analytic Riemann--Roch formula, namely the commutativity of the diagram
\begin{equation}\tag{5.4}
\begin{array}{ccc}
K^\bullet(X) & \xrightarrow{\;x\mapsto \ch(x)\Todd(T_X)\;} & H^{2\bullet}(X,\mathbb Q) \\
\downarrow f_* && \downarrow f_* \\
K^\bullet(Y) & \xrightarrow{\;y\mapsto \ch(y)\Todd(T_Y)\;} & H^{2\bullet}(Y,\mathbb Q),
\end{array}
\end{equation}
where the horizontal arrows are deduced from cohomological Chern classes, themselves deduced from the homomorphism (5.1) and from the ordinary Chern-class homomorphism
\[
K^\bullet(X^{\mathrm{top}})\longrightarrow H^{2\bullet}(X,\mathbb Z),
\]
the first vertical arrow is the same as in (5.3), and the second is the Gysin homomorphism in rational cohomology.

\subsection*{6. The covariant invariant $K_\bullet(X)$}
In all that precedes, scarcely anything has been said except about the contravariant invariant $K^\bullet(X)$, to the exclusion of the covariant invariant $K_\bullet(X)$ whose existence was indicated in 3.1. If one regards $K^\bullet(X)$ as giving an analogue of the integral cohomology ring of a topological space, one should regard $K_\bullet(X)$ as giving an analogue of integral homology. The module structure of $K_\bullet(X)$ over $K^\bullet(X)$ then corresponds to the cap-product operation. The fact that, when $X$ is regular, the natural homomorphism $K^\bullet(X)\to K_\bullet(X)$ is an isomorphism should be regarded as the analogue of the Poincar\'e duality theorem for oriented topological varieties, asserting that cap product with the fundamental homology class defines an isomorphism from cohomology to homology. These analogies moreover suggest the usefulness of also defining, for a finite polyhedron $X$, say, a covariant invariant $K_\bullet(X)$ whose relations to integral homology -- spectral sequence, Chern homomorphism, and so on -- would be analogous to those existing between $K^\bullet(X)$ and integral cohomology. Nor does it seem excluded that there might exist, in terms of these groups $K_\bullet(X)$, variants of the differentiable Riemann--Roch theorem valid for spaces more general than compact differentiable manifolds.

Returning to the case where $X$ is an arbitrary noetherian scheme, one should endow $K_\bullet(X)$ with an increasing filtration, $\Fil_iK_\bullet(X)$ being generated by classes $\cl_\bullet(F)$, with $F$ a coherent sheaf on $X$ such that $\dim \operatorname{Supp}F\le i$. One will prove (Expos\'e X) the compatibility of this filtration with the module structure and with the filtration of $K^\bullet(X)$, i.e. the formula
\begin{equation}\tag{6.1}
\Fil^iK^\bullet(X)\cdot \Fil_jK_\bullet(X)\subset \Fil_{j-i}K_\bullet(X),
\end{equation}
which allows the graded group associated with the filtered module $K_\bullet(X)$, denoted $\Gr_\bullet(X)$, to be endowed with a module structure over $\Gr^\bullet(X)$. Its elements may in fact be interpreted as classes of algebraic cycles on $X$, for an equivalence relation less fine than rational equivalence when $X$ is of finite type over a base field $k$. If
\[
f:X\longrightarrow Y
\]
is a proper morphism of noetherian schemes, then
\begin{equation}\tag{6.2}
f_*:K_\bullet(X)\longrightarrow K_\bullet(Y)
\end{equation}
is compatible with the filtrations and defines a homomorphism of graded groups
\begin{equation}\tag{6.3}
f_*:\Gr_\bullet(X)\longrightarrow \Gr_\bullet(Y),
\end{equation}
giving rise to a ``projection formula'', i.e. a formula which says that it is $\Gr^\bullet(Y)$-linear; this formula is obtained by passage to associated gradeds from the analogous projection formula for the homomorphism (6.2), which is $K^\bullet(Y)$-linear.

From the point of view proposed in the seminar, a manageable theory of intersections on noetherian schemes $X$, not necessarily regular, should consist in the combined study of the contravariant invariants $K^\bullet(X),\Gr^\bullet(X)$ and the covariant invariants $K_\bullet(X),\Gr_\bullet(X)$. These play roles of essentially different nature, and their properties complement each other. Thus the covariant groups $K_\bullet(X),\Gr_\bullet(X)$ are easier to compute for certain nonproper fibrations, thanks to the ``homotopy exact sequence'' which one can establish for them ([RRR] and Expos\'e IX). For example, one proves (Expos\'e IX) canonical isomorphisms such as
\[
K_\bullet(X[t])\simeq K_\bullet(X),\qquad \Gr_\bullet(X[t])\simeq \Gr_\bullet(X),
\]
which fail in general for $K^\bullet$ when $X$ is a nonregular scheme. On the other hand, $K^\bullet(X)$ and $\Gr^\bullet(X)$ lend themselves well to computations, thanks to their ring structures.

When one considers schemes $X$ proper over a field, the projection
\[
\pi_X:X\longrightarrow (\operatorname{point})
\]
defines a homomorphism, the ``Euler--Poincar\'e characteristic'',
\begin{equation}\tag{6.4}
\chi_X=\pi_{X*}:K_\bullet(X)\longrightarrow K_\bullet(\operatorname{point})=\mathbb Z,
\end{equation}
which induces on the subgroup $\Gr_0(X)=\Fil_0K_\bullet(X)$ the homomorphism ``degree of zero-cycles''
\begin{equation}\tag{6.5}
\deg_X:\Gr_0(X)\longrightarrow \mathbb Z.
\end{equation}
The multiplication homomorphism
\[
\Gr^i(X)\times \Gr_i(X)\longrightarrow \Gr_0(X),
\]
composed with the degree homomorphism (6.5), then furnishes a pairing
\begin{equation}\tag{6.6}
\Gr^i(X)\times \Gr_i(X)\longrightarrow \mathbb Z,
\end{equation}
which explains that, in certain respects, $\Gr^\bullet(X)$ and $\Gr_\bullet(X)$ play dual roles, just as cohomology and homology do. Thus, if $f:X\to Y$ is a morphism of proper algebraic schemes, the two homomorphisms
\[
f^*:\Gr^\bullet(Y)\longrightarrow \Gr^\bullet(X),\qquad f_*:\Gr_\bullet(X)\longrightarrow \Gr_\bullet(Y)
\]
are formally transposes of one another, in the sense of the pairings (6.6).

The homomorphisms (6.4) and (6.5), and the pairing (6.6) deduced from them, may be regarded as the source of numerical relations in intersection theory on proper algebraic schemes over a given field $k$. This point of view will be developed in some detail in Expos\'e X, which is essentially independent of Expos\'es VI, VII, VIII, centred on the proof of the Riemann--Roch theorem, and of Expos\'e IX. Some notions concerning the behaviour of $K^\bullet(X)$ and $\Gr^\bullet(X)$ under specialization will also be given there.

\subsection*{7. Picard groups and the determinant functor}
As in every intersection theory, it is important to specify the place occupied in the theory by the Picard group, or group of classes of divisors in the classical terminology. Subject to a slight restriction, automatically satisfied when $X$ is for instance quasi-compact, one establishes in Expos\'e X a canonical isomorphism
\begin{equation}\tag{7.1}
c_1:\Pic(X)\xrightarrow{\sim}\Gr^1(X),
\end{equation}
which shows the exact role of the Picard group in the theory we are proposing. This role justifies a more detailed study of the Picard group in Expos\'es XII and XIII, which are logically independent of the text of the seminar. Kleiman, in Expos\'e XIII, studies numerical equivalence and certain finiteness theorems for the Picard group, announced without proof in [6], using purely numerical techniques, not relying on the theory of the Picard functor and its representability, due to himself and Mumford, and markedly simpler than Grothendieck's initial proofs. The only result refractory to Kleiman's methods, and which he is obliged to use, is [6, C-07 (i)]; its proof and certain refinements, including certain representability theorems for the Picard functor, occupy Expos\'e XII. The latter is moreover of a spirit and technique entirely foreign to the rest of the seminar, and is connected with [5]; this expos\'e is logically independent of all the others.

Finally, in Expos\'e XI, of a nature appreciably more general than XII and XIII and closer to the general spirit of the seminar, one studies, on an arbitrary locally ringed topos, a remarkable functor called the ``determinant functor'':
\begin{equation}\tag{7.2}
\det^\bullet:\Parf_{\mathrm{is}}(X)\longrightarrow \Inv(X)
\end{equation}
from the category of perfect complexes on $X$, whose morphisms are only the isomorphisms in the sense of the derived category $D(X)$, to the category of invertible modules on $X$. This functor is described, up to very little, by the requirement that it be of ``local nature'' and that it be given, for a bounded complex $L$ with locally free components, by the formula
\begin{equation}\tag{7.3}
\det^\bullet(L)=\bigotimes_i \det(L_i)^{(-1)^i},
\end{equation}
where $\det(L_i)$ denotes the exterior power of maximal degree of the locally free module $L_i$. Although this expos\'e is also logically independent of the rest of the seminar, except for Expos\'e I, it nevertheless constitutes an important element of the general ``yoga'' proposed in the seminar and developed essentially in Expos\'es I to XI.

\section*{Bibliography}
\begin{enumerate}[label={[\arabic*]},leftmargin=2.5em]
\item M. Atiyah and F. Hirzebruch, \emph{Vector bundles and homogeneous spaces}, Symposia in Pure Mathematics, vol. 3, Differential Geometry, pp. 7--38, 1961.
\item A. Borel and J.-P. Serre, \emph{Le th\'eor\`eme de Riemann--Roch}, Bull. Soc. math. France, t. 86, pp. 97--136 (1958).
\item H. Grauert, \emph{Ein Theorem der analytischen Garbentheorie}, Pub. Math. no. 5, pp. 5--64 (1960).
\item A. Grothendieck, \emph{Classes de faisceaux et th\'eor\`eme de Riemann--Roch}, mimeographed notes, Princeton 1957, cited as [RRR] and reproduced as an appendix to Expos\'e 0 of the seminar.
\item A. Grothendieck, \emph{Technique de descente et th\'eor\`emes d'existence en G\'eom\'etrie Alg\'ebrique}, in \emph{Fondements de la G\'eom\'etrie Alg\'ebrique}, extracts from the Bourbaki Seminar 1957--1962, Secr\'etariat Math\'ematique, Paris.
\item W. Shih, S\'eminaire de Topologie \`a l'I.H.E.S., 1966/67.
\item M. Atiyah and F. Hirzebruch, \emph{The Riemann--Roch theorem for analytic embeddings}, Topology, vol. 1, pp. 151--166 (1962).
\end{enumerate}



\clearpage



\section*{Appendix to Exposé 0: RRR}
\begin{center}
{\Large Classes of Sheaves and the Riemann--Roch Theorem}\par
\vspace{0.5em}
by A. Grothendieck\footnote{This is the literal reproduction of the report cited in the Introduction to the present seminar. The footnotes indicated by the signs (*), (**) were added in October 1967.}
\end{center}

\begin{center}
Demonstration of the Riemann--Roch theorem (1 November 1957).
\end{center}

\subsection*{Chapter I. $\lambda$-rings (formal preliminaries)}
\begin{longtable}{@{}p{0.74\linewidth}r@{}}
1. Definitions & 1\\
2. Examples & 4\\
3. The $\lambda$-ring defined by a graded ring & 8\\
4. The operations $\lambda^p(N,x)$ & 13\\
\end{longtable}

\subsection*{Chapter II. Classes of coherent algebraic sheaves and Chern classes}
\begin{longtable}{@{}p{0.74\linewidth}r@{}}
1. The theory of Chow & 18\\
2. Definition of the Chern classes of coherent algebraic sheaves & 22\\
3. Functorial generalities on $K(X)$ & 27\\
4. Some technical results & 34\\
5. Sheaf-theoretic definition of Chern classes. Application to the study of immersion morphisms & 43\\
6. The Riemann--Roch theorem & 49
\end{longtable}

\section*{Chapter I. $\lambda$-rings (formal preliminaries)}

\subsection*{1. Definitions}

A \emph{$\lambda$-ring}\footnote{In this seminar we rather say \emph{pre-$\lambda$-ring} (V 2.1), reserving the name $\lambda$-ring for those satisfying the supplementary condition of page 3 (cf. V 2.4).} means a commutative ring $K$ with unit, provided with a family of maps
\begin{equation}\tag{1.1}
\lambda^i:K\longrightarrow K\qquad (i\text{ an integer }\geq 0)
\end{equation}
satisfying the following conditions:
\begin{equation}\tag{1.2}
\lambda_x^0=1,\qquad \lambda_x^1=x,\qquad
\lambda^n(x+y)=\sum_{i=0}^n \lambda^i_x\lambda^{n-i}_y .
\end{equation}
For every $x\in K$ put
\begin{equation}\tag{1.3}
\lambda_t(x)=\sum_{n\geq 0}\lambda^n(x)t^n\in K[[t]].
\end{equation}
The preceding relations then say that $x\mapsto \lambda_t(x)$ is a homomorphism from the additive group $K$ to the multiplicative group $1+K[[t]]^+$ of formal series over $K$ with augmentation $1$, lifting the canonical homomorphism
\[
1+\sum_{i>0}x_it^i\longmapsto x_1.
\]

Thus on the ring $\Z$ there exists a unique $\lambda$-structure compatible with its ring structure and such that
\begin{equation}\tag{1.4}
\lambda_t(1)=1+t .
\end{equation}
One then has
\begin{equation}\tag{1.5}
\lambda_t(n)=(1+t)^n,
\end{equation}
whence
\[
\lambda^i(n)=\binom{n}{i}.
\]

The notion of a $\lambda$-homomorphism of $\lambda$-rings is clear; an \emph{augmented $\lambda$-ring} is a $\lambda$-ring provided with a $\lambda$-homomorphism to $\Z$.

To define the notion of a \emph{special} $\lambda$-ring, the following considerations are useful. Let $k$ be a commutative ring with unit and let
\[
K=1+k[[t]]^+
\]
be the group of formal series with coefficients in $k$ and augmentation $1$. We shall introduce on it a $\lambda$-ring law whose additive structure is the multiplication of formal series. The multiplication of this ring will be denoted $f\circ g$. It is entirely determined by the condition that the coefficients $c_n=c_n(f\circ g)$ of $f\circ g$ are given by universal polynomials with integral coefficients in the coefficients $a_i=c_i(f)$ and $b_i=c_i(g)$ of $f$ and $g$, respectively:
\begin{equation}\tag{1.6}
c_n=P_n(a_1,\ldots,a_n;b_1,\ldots,b_n),
\end{equation}
that it be distributive with respect to the ``addition'' $fg$, and be given for linear factors by
\begin{equation}\tag{1.7}
(1+at)\circ(1+bt)=1+abt.
\end{equation}
The $P_n$ are computed by calculations with symmetric functions; $P_n$ is isobaric of weight $n$ in the $a_i$ and isobaric of weight $n$ in the $b_i$ (where $a_i$ and $b_i$ have weight $i$), and moreover is symmetric in $a=(a_i)$ and $b=(b_i)$.

Similarly, the $\lambda^i(f)$ are determined without ambiguity by the condition that the coefficients
\[
c_n\bigl(\lambda^i(f)\bigr)=c_n\bigl(\lambda^i(f)\bigr)
\]
of $\lambda^i(f)$ be calculated by universal polynomials with integral coefficients in the coefficients $c_n(f)=a_n$ of $f$:
\begin{equation}\tag{1.8}
c_n=P_{i,n}(a_1,\ldots,a_{in}),
\end{equation}
that the relations (1.2) be satisfied, and that, for $i>1$, one have
\begin{equation}\tag{1.9}
\lambda^i(1+at)=1
\end{equation}
($1$ is the zero element of the ring $K=1+k[[t]]^+$). The $P_{i,n}$ are again computed by calculations with symmetric functions, and one observes that $P_{i,n}$ is an isobaric polynomial of weight $in$ in the $a_i$. Provided with the preceding operations, $K$ becomes a $\lambda$-ring; its zero element is $1$ and its unit element is $1+t$. Also note the relation
\begin{equation}\tag{1.7 bis}
(1+at)\circ f=f(at).
\end{equation}

Let now $K$ be an arbitrary $\lambda$-ring. We say that $K$ is \emph{special} if the additive homomorphism $\lambda_t$ from $K$ to $1+K[[t]]^+$ is a $\lambda$-homomorphism. Thus this means that one has the formulas
\begin{equation}\tag{1.10}
\lambda_t(1)=1+t,\quad
\lambda_t(xy)=\lambda_t(x)\circ\lambda_t(y),\quad
\lambda_t(\lambda^i(x))=\lambda^i(\lambda_t(x)),
\end{equation}
or, more explicitly,
\begin{equation}\tag{1.11}
\begin{cases}
\lambda^i(1)=0 & \text{if }i>1,\\
\lambda^n(xy)=P_n(\lambda^1x,\ldots,\lambda^n x;\lambda^1y,\ldots,\lambda^n y),\\
\lambda^n(\lambda^i(x))=P_{i,n}(\lambda^1x,\ldots,\lambda^{in}x),
\end{cases}
\end{equation}
where the $P_n$ and $P_{i,n}$ are the universal polynomials considered above. In particular the formulas (1.5) follow. One may show that the $\lambda$-ring $1+k[[t]]^+$ constructed from a commutative ring $k$ is special, but we shall not need this fact.

\subsection*{2. Examples}

1) Let $G$ be a group and $k$ a commutative ring. Consider the group $K(G)$, the quotient of the free group generated by the classes of representations of $G$ in finite-type $k$-modules by the subgroup generated by elements of the form
\[
(\rho\oplus\rho')-(\rho)-(\rho')
\]
(where $\oplus$ denotes the direct sum). Tensor product and exterior powers define on $K(G)$ a $\lambda$-ring structure. One may also pass to the quotient by elements of the form
\[
(\rho)-(\rho')-(\rho''),
\]
where $\rho$ is a representation which is an extension of $\rho'$ by $\rho''$, trivial as an extension of $k$-modules. One obtains the group $K_r(G)$ of reduced classes of representations of $G$, and one checks that the $\lambda$-ring law passes to the quotient.

When $k$ is an algebraically closed field of characteristic $0$,\footnote{It is unnecessary for $k$ to be algebraically closed; cf. VI 3.3 for a still more general result, encompassing all those that will be used in the present report.} one can show that $K(G)$, and a fortiori $K_r(G)$, is a special $\lambda$-ring. This is no longer true if the characteristic is not $0$; however, if $k$ is algebraically closed, $K_r(G)$ is still a special $\lambda$-ring. To prove this one is reduced to the case where $G$ is a product $\Gl(m,k)\times\Gl(n,k)$ and where only rational representations of $G$ are considered; our assertion follows from the explicit determination of $K_r(G)$ to be made below (Example 3). Finally, when $k$ has characteristic $0$, complete reducibility of rational representations of $G$ shows that $K(G)=K_r(G)$.

Variants of the preceding example: $G$ is an algebraic group defined over $k$ and one considers only rational representations of $G$ defined over $k$; one may consider the case of a topological group, or of a Lie group, and continuous representations, etc.\footnote{One could build other examples of $\lambda$-rings; for example, the set of equivalence classes of quadratic forms, the set of representation classes of a Lie algebra, etc. We promise that we shall not need all these examples in order to treat the Riemann--Roch theorem.}

2) Let $X$ be an algebraic variety, irreducible or not, defined over the field $k$. We may consider the quotient group of the free group generated by the classes of vector bundles on $X$ (defined over $k$) by the subgroup generated by the elements of the form
\[
(E\oplus E')-(E)-(E'),
\]
on which the $\lambda$-ring operations are defined by tensor product and exterior powers. If $X$ is complete, this is a free group generated by the indecomposable vector bundles on $X$ (thanks to Krull--Remak--Schmidt); knowledge of this ring and of its basis is equivalent to the complete classification of vector bundles on $X$ and to knowledge of the elementary tensor operations on these bundles. If $k$ is algebraically closed of characteristic $0$, all other tensor operations on classes of vector bundles are then known by means of what follows in 3). One may make a more drastic passage to the quotient by the group generated by the elements
\[
(E)-(E')-(E''),
\]
where $E$ is a vector bundle extension of $E'$ by $E''$. One then obtains a $\lambda$-ring, denoted $K(\xi(X))$, where $\xi(X)$ denotes the additive category of vector bundles on $X$ defined over $k$. Using the results of Example 1), one shows that, if $k$ is algebraically closed,\footnote{Unnecessary condition; cf. the preceding footnote.} the $\lambda$-ring $K(\xi(X))$ is special, and that the same is true of the first $\lambda$-ring defined here if the characteristic of $k$ is $0$, but not in the general case.

3) Determination of $K(G)$ and of $K_r(G)$ in some important cases. Suppose the field $k$ is algebraically closed and that the group $G$ is affine algebraic, connected, and ``globally reductive'' (the radical of $G$ is isomorphic to $(k^*)^s$), and consider only rational linear representations. In characteristic $0$, complete reducibility of rational linear representations of $G$ shows that $K(G)=K_r(G)$ is the free group generated by the classes of irreducible rational representations. If $T$ is a maximal torus of $G$, then
\[
K(T)=K_r(T)=\Z(\That),
\]
the algebra of the group $\That$ of rational characters of $T$; one has $\That\simeq \Z^r$ if $r$ is the rank of $G$ (this is independent of the characteristic). The Weyl group $W$ operates on $T$, hence on $K(T)$, in a way which depends only on the operations of $W$ on the free group $\That$. There is an evident homomorphism of $\lambda$-rings, deduced from the injection of $T$ into $G$:
\begin{equation}\tag{1.13}
K_r(G)\longrightarrow K_r(T)^W=K(T)^W,
\end{equation}
where $K(T)^W$ denotes the set of fixed points of $W$ in $K(T)$. The theory of linear representations of $G$\footnote{Cf. Chevalley Seminar 1956/58, \emph{Groupes de Lie algébriques} (École Normale Supérieure).} then proves the following theorem.

\textbf{Theorem 1.1.} The homomorphism (1.13) is an isomorphism.

Thus one sees that $K_r(G)$, once the ``type'' of $G$ is known (characterized by $W$ operating on a group $\Z^r$), does not depend on the characteristic, and that it is a special $\lambda$-ring. There is a canonical basis of $K_r(G)$ corresponding to the orbits of $W$ in $\That$; another basis is obtained as follows. Let $(\rho_i)_{1\leq i\leq r'}$ be the classes of representations of $G$ corresponding to the fundamental representations of $G/\rad G$, and let $(\sigma_j)_{1\leq j\leq r''}$ be a basis of the group of rational homomorphisms from $G$ to $k^*$. Thus $r'+r''=r$ (the rank of $G$).

\textbf{Corollary.} The ring $K_r(G)$ is identified with the subring
\[
\Z[\rho_1,\ldots,\rho_{r'};\sigma_1,\ldots,\sigma_{r''},\sigma_1^{-1},\ldots,\sigma_{r''}^{-1}]
\]
of the field of rational functions in the $\rho_i$ and the $\sigma_j$ with coefficients in $\Q$.\footnote{The derived group $G'$ of $G$ must be supposed simply connected for this statement to be correct.}

In particular, $K_r(G)$ has no zero divisors. If $G=\prod_i\Gl(n_i,k)$, the ring $K_r(G)$ is made explicit as follows. Let $\rho_i$ be the representation of $G$ deduced from the identical representation of $\Gl(n_i,k)$. If one then considers the field of rational functions over $\Q$ in the indeterminates $\lambda^j(\rho_i)$, $1\leq j\leq n_i$, then $K_r(G)$ is identified with the subring generated by these indeterminates and the inverses of the $\lambda^{n_i}(\rho_i)$. This justifies our assertion that, at least in characteristic $0$, the operations $\lambda^i$ make it possible to reconstruct all tensor operations on vector bundles; this remains true over an algebraically closed field of arbitrary characteristic, provided one takes reduced classes.

\textbf{Remarks.} a) In Example 1), the ring $K(G)$ is augmented, the augmentation being defined by the degree of representations; by passage to the quotient, this augmentation defines an augmentation on $K_r(G)$. Likewise, the $\lambda$-rings considered in Example 2) are augmented when $X$ is $k$-irreducible; in $K(\xi(X))$, the augmentation is defined by the rank of a vector bundle, that is, the dimension of its fibres.

b) When the field $k$ is not algebraically closed, it is plausible that $K(G)$ is a special $\lambda$-ring in characteristic $0$, and likewise for $K_r(G)$ in all characteristics.

c) In non-zero characteristic, Theorem 1.1 was graciously supplied by Chevalley.

\subsection*{3. The $\lambda$-ring defined by a graded ring}

Let $A$ be a graded commutative ring with unit. To simplify, assume that the degrees of $A$ are positive and that $A^0=\Z$. Denote by $\Ahat$ the product of the $A^i$ (this is a commutative ring with unit), and by $\Ahp$ the kernel of the evident augmentation $\Ahat\to A^0=\Z$. Then $1+\Ahp$ is a subgroup of the multiplicative group of invertible elements of $\Ahat$. Consider the abelian group
\begin{equation}\tag{1.14}
\widetilde A=\Z\times (1+\Ahp),
\end{equation}
whose elements will be denoted $[n,x]$, with $n\in\Z$ and
\[
x=1+\sum_{i\geq 1}x^i\in 1+\Ahp\qquad (x^i\in A^i).
\]
We write this group additively, thus:
\begin{equation}\tag{1.15}
[n,x]+[n',x']=[n+n',xx'].
\end{equation}
The zero element of this group is $[0,1]$. We shall introduce on it an augmented $\lambda$-ring structure, compatible with its additive structure, the augmentation being $\varepsilon:[n,x]\mapsto n$. The product $[m,x][n,y]$ is completely characterized by the fact that it is written
\[
[mn,x^n y^m(x*y)],
\]
where $x*y$ is an element of $1+\Ahp$ whose components are expressed as universal polynomials with integral coefficients in the $x^i$ and the $y^i$:
\begin{equation}\tag{1.16}
(x*y)^i=Q_i(x^1,\ldots,x^i;y^1,\ldots,y^i)\qquad (i\geq 1),
\end{equation}
the polynomial $Q_i$ being isobaric of weight $i$. Moreover, $x*y$ must be distributive with respect to the ``addition'' $xy$ and, for ``linear'' factors, must reduce to
\begin{equation}\tag{1.17}
[1,1+x^1][1,1+y^1]=[1,1+x^1+y^1].
\end{equation}
The polynomials $Q_i$ are evaluated by calculations with elementary symmetric functions. This is the calculation of the Chern classes of a tensor product of vector bundles by means of the Chern classes $x^i$ and $y^i$ of the factors and of their respective dimensions $m$ and $n$. From (1.17) one deduces
\begin{equation}\tag{1.17 bis}
(1+x^1)*(1+y^1)=\frac{1+x^1+y^1}{(1+x^1)(1+y^1)},
\end{equation}
and an arbitrary product $x*y$ is computed term by term by decomposing formally
\[
x=\prod_{i=1}^n(1+\alpha_i),\qquad y=\prod_{i=1}^n(1+\beta_i),
\]
where the $\alpha_i$ and $\beta_i$ have degree $1$.

In this way $\widetilde A$ becomes a commutative augmented ring with unit $[1,1]$; it may be interpreted as the ring obtained by adjoining a unit to the ring $1+\Ahp$. Moreover, on $\widetilde A$ one defines a filtration by putting $\widetilde A^0=\widetilde A$ and denoting by $\widetilde A^n$ the set of the $[0,x]$ with $x^i=0$ for $1\leq i<n$. This filtration is compatible with the product, for
\begin{equation}\tag{1.18}
\left(1+\sum_{i\geq m}x^i\right)*\left(1+\sum_{j\geq n}y^j\right)
=1-\frac{(m+n-1)!}{(m-1)!(n-1)!}x^m y^n+\cdots,
\end{equation}
the omitted terms having degree $>m+n$. This formula shows that the graded ring $G(\widetilde A)$ associated with $\widetilde A$ is identified with $A$ from the additive point of view, but not from the multiplicative point of view. One can only say that the map $\eta'$ from $A$ to $G(\widetilde A)$ defined by
\begin{equation}\tag{1.19}
\eta'(x^i)=(-1)^{i-1}(i-1)!\,x^i\qquad (x^i\in A^i)
\end{equation}
is a ring homomorphism.

On $\widetilde A$ introduce operations $\lambda^i$ satisfying the axioms of $\lambda$-rings (Section 1, formulas (1.2)), well determined by the condition that the components of $\lambda^i[0,x]$ be deduced from the components $x^i$ of $x$ by universal polynomials with integral coefficients:
\begin{equation}\tag{1.20}
\lambda^i[0,x]=[0,\lambda^i x],\qquad (\lambda^i x)^{(n)}=Q_{i,n}(x^1,\ldots,x^n),
\end{equation}
where $Q_{i,n}$ is isobaric of weight $n$, and such that
\begin{equation}\tag{1.21}
\lambda^i[1,1+x^1]=0\qquad\text{for }i>1.
\end{equation}
The polynomials $Q_{i,n}$ are evaluated by a calculation with elementary symmetric functions: it is the calculation of the Chern classes of the exterior powers of a vector bundle by means of the Chern classes and the rank of that bundle. One then observes that $\widetilde A$ is a special augmented $\lambda$-ring and that the $\widetilde A^i$ are stable under the operations $\lambda^j$.

The augmented $\lambda$-ring $\widetilde A$ may be considered as a functor in $A$. Consequently, the principal automorphism
\[
x=\sum_{i\geq 0}x^i\longmapsto \check x=\sum_{i\geq 0}(-1)^i x^i
\]
of $A$ defines an automorphism of $\widetilde A$, denoted by the same symbol $\check{\ }$.

\textbf{Proposition 1.2.} Let $x\in\widetilde A$ be of the form
\[
x=[n,1+x^1+\cdots+x^n]
\]
(so $x^i=0$ for $i>n$). Then $\lambda^i(x)=0$ for $i>n$, one has $\lambda^n(x)=[1,1+x^1]$, and finally
\begin{equation}\tag{1.22}
\sum_{i=0}^{n}(-1)^i\lambda^i(x)=[0,1-(n-1)!x^n+\cdots].
\end{equation}
From this proposition, and from the fact that the $\widetilde A^i$ are stable under the operations $\lambda^j$, it follows that formula (1.22) is valid whenever $\varepsilon(x)=n$, whether or not $x^i=0$ for $i>n$.

Formula (1.22) admits the following generalizations. Let $K$ be any $\lambda$-ring, and let $n$ be arbitrary. For $x\in K$, put
\begin{equation}\tag{1.23}
\gamma^n(x)=(-1)^n\sum_{i=0}^{n}(-1)^i\lambda^i(x+n)=\lambda^n(x+n-1)
\end{equation}
(the identity of the last two members follows from an immediate induction). In particular $\gamma^0(x)=1$ and $\gamma^1(x)=x$. Next put
\begin{equation}\tag{1.24}
\gamma_t(x)=\sum_{n\geq 0}\gamma^n(x)t^n\in K[[t]].
\end{equation}
One deduces the formula
\begin{equation}\tag{1.25}
\gamma_t(x)=\lambda_{t/(1-t)}(x),
\end{equation}
which is equivalent to
\begin{equation}\tag{1.25 bis}
\lambda_s(x)=\gamma_{s/(1+s)}(x).
\end{equation}
This proves that the $\lambda^i$ are expressed in terms of the $\gamma^i$, and that
\[
\gamma_t(x+y)=\gamma_t(x)\gamma_t(y).
\]
Consequently the maps $\gamma^i$ define a new $\lambda$-ring structure on $K$. Formula (1.22) then shows that, for $K=\widetilde A$ and every $x\in\widetilde A$, one has
\begin{equation}\tag{1.22 bis}
\gamma^n(x-\varepsilon(x))=[0,1+(-1)^{n-1}(n-1)!x^n+\cdots],
\end{equation}
and hence the component $x^n$ is obtained, up to the factor $(-1)^{n-1}(n-1)!$, by reducing modulo $\widetilde A^{n+1}$ the element $\gamma^n(x-\varepsilon(x))$ of $\widetilde A^n$.

\paragraph{The Chern homomorphism.}
For every graded ring $A$ satisfying the preceding conditions, we shall define a homomorphism of augmented rings
\begin{equation}\tag{1.26}
\ch:\widetilde A\longrightarrow \Ahat\otimes\Q.
\end{equation}
This homomorphism is defined without ambiguity by the condition that it be an additive homomorphism, send $1$ to $1$, and be such that the components of $\ch(x)$ in degree $>0$ are given by universal isobaric polynomials with rational coefficients in the components of $x$, and that
\begin{equation}\tag{1.27}
\ch([1,1+x^1])=\exp x^1=\sum_{n\geq 0}(x^1)^n/n!.
\end{equation}
One immediately obtains the explicit expression of $\ch$. Let $\eta$ be the additive endomorphism of $\Ahat\otimes\Q$ defined, for $x^i$ of degree $i$, by
\begin{equation}\tag{1.28}
\eta(x^i)=\frac{1}{(-1)^{i-1}(i-1)!}x^i .
\end{equation}
Then
\begin{equation}\tag{1.29}
\ch\left([n,1+\sum_{i\geq 1}x^i]\right)
=n+\eta\left(\log\left(1+\sum_{i\geq 1}x^i\right)\right).
\end{equation}
Since this formula can be inverted, at least if $A$ has only finitely many homogeneous components, one sees that in this case one has an isomorphism $\ch$ from the ring $\widetilde A\otimes\Q$ onto the ring $\Ahat\otimes\Q$, which completely elucidates the structure of $\widetilde A\otimes\Q$.

Recall that, according to Hirzebruch [1], a formal series $f\in\Q[[t]]$ defines by universal polynomial formulas an additive homomorphism
\[
1+\Ahp\longrightarrow 1+(\Ahat\otimes\Q)^+,
\]
denoted $\mathfrak C_f$. When $f(t)=t/(1-\exp(-t))$, we write simply $\mathfrak C_f=\mathfrak C$ (the initial of Todd). We extend the definition of $\mathfrak C(x)$ to $x\in\widetilde A$. This said, we have the following result.

\textbf{Proposition 1.3.} Let $N\in\widetilde A$ be of the form
\[
N=[q,1+N^1+\cdots+N^q]
\]
(so $N^i=0$ for $i>q$). If one puts
\[
\lambda_{-1}(N)=\sum_{i\geq 0}(-1)^i\lambda^i(N)=\sum_{i=0}^{n}(-1)^i\lambda^i(N),
\]
then
\begin{equation}\tag{1.30}
\ch(\lambda_{-1}(N))=(-1)^qN^q\mathfrak C(\check N)^{-1},
\end{equation}
which is equivalent to
\begin{equation}\tag{1.30 bis}
\ch(\lambda_{-1}(\check N))=N^q\mathfrak C(N)^{-1}.
\end{equation}

\subsection*{4. The operations $\lambda^p(N,x)$}

Let $A$ be the ring of formal series, with coefficients in $\Z$, in two infinite sequences of indeterminates denoted $\lambda^iN$ and $\lambda^i x$ ($i\geq 1$). Put $\lambda^1N=N$ and $\lambda^1x=x$, and introduce on $A$ the unique special $\lambda$-ring structure for which
\[
\lambda^i(N)=\lambda^iN,
\qquad
\lambda^i(x)=\lambda^i x,
\]
and for which the operations $\lambda^i$ are continuous. If one restricted oneself to polynomials in the indeterminates $\lambda^iN$ and $\lambda^i x$, one would have the $\lambda$-ring freely generated by $N$ and $x$; our ring $A$ is its completion. One may then form
\begin{equation}\tag{1.31}
\lambda_{-1}(N)=\sum_{i\geq 0}(-1)^i\lambda^iN,
\end{equation}
which is a formal series with constant term $1$, hence invertible. Thus an element $\lambda^p(N,x)$ of $A$ can be defined without ambiguity by putting
\begin{equation}\tag{1.32}
\lambda^p(N,x)\lambda_{-1}(N)=\lambda^p(x\lambda_{-1}(N)).
\end{equation}

\textbf{Theorem 1.4.} Let $J_q$ be the closed ideal of $A$ generated by the $\lambda^iN$ with $i>q$. Identify $A/J_q$ with the ring of formal series in the $\lambda^i x$ and the $\lambda^jN$ ($1\leq j\leq q$). Then $\lambda^p(N,x)$ reduced modulo $J_q$ is a polynomial, isobaric of weight $p$ with respect to the variables $\lambda^i x$.

Let $N$ and $F$ be two vector spaces of finite dimensions $q,q'$ over an algebraically closed field $k$ of characteristic $0$. Let $N^{(p)}$ be the kernel of the map
\[
(a_1,\ldots,a_p)\longmapsto a_1+\cdots+a_p
\]
from $N^p$ to $N$. We make the symmetric group $\mathfrak S_p$ operate on $N^p$, hence on $N^{(p)}$, hence on $\Lambda(N^{(p)})$, and then on
\[
\Lambda(N^{(p)})\otimes\bigotimes^p F.
\]
Consider this vector space as a graded module over $\mathfrak S_p\times\Gl(N)\times\Gl(F)$, and take its ``alternating part'', which is a graded module over $\Gl(N)\times\Gl(F)=G$. Then
\begin{equation}\tag{1.33}
\lambda^p(N,F)=E_G\left(\left(\Lambda(N^{(p)})\otimes\bigotimes^p F\right)^{\mathrm{alt}}\right).
\end{equation}

N.B. If $M$ is a complex of $G$-modules, $E_G(M)$ denotes the alternating sum in $K(G)$ of the classes of the components of $M$, that is, the Euler--Poincaré characteristic of $M$. In formula (1.33), the first member $\lambda^p(N,F)$ denotes the element of the ring $K(G)$ obtained by substituting, for the variables $\lambda^i_N$ and $\lambda^i_F$, the classes in $K(G)$ of the $G$-modules $\lambda^i(N)$ and $\lambda^i(F)$, respectively.

\emph{Proof.} Expanding the second member of (1.32) by the second formula (1.11), one sees that the first assertion is proved if one proves the analogous assertion obtained by setting $x=1$, i.e. if one proves that in the $\lambda$-ring $\Z[\lambda^1N,\ldots,\lambda^qN]$, where $\lambda^iN=0$ for $i>q$, the element $\lambda^p(\lambda_{-1}(N))$ is divisible by $\lambda_{-1}(N)$. A priori, the quotient
\[
\lambda^p(N,1)=\lambda^p(\lambda_{-1}(N))/\lambda_{-1}(N)
\]
is a formal series contained in the quotient field $K$ of the ring of formal series in the $\lambda^iN$ for $1\leq i\leq q$. On the other hand, also denote by $N$ a vector space of dimension $q$ over an algebraically closed field $k$ of characteristic $0$, and consider the element
\[
\mathscr L^p(N)=E_G((\Lambda(N^{(p)}))^{\mathrm{alt}})\in K(G),
\]
where now $G=\Gl(N)$. We know (Section 2, Theorem 1.3) that $K(G)$ is identified with the subring of $K$ generated by the $\lambda^i_N$ for $1\leq i\leq q$ and by the inverse of $\lambda_N^q$. Moreover, in this ring one verifies easily that
\[
\mathscr L^p(N)\lambda_{-1}(N)=\lambda^p(\lambda_{-1}(N)).
\]
Since we are in a large field $K$ and $\lambda_{-1}(N)\neq 0$, it follows that $\mathscr L^p(N)=\lambda^p(N,1)$. Both members lie in the intersection of $K(G)$ with the ring of formal series in the $\lambda_N^i$, hence are polynomials in the $\lambda_N^i$.

To prove (1.33) in the general case, one verifies that the second member satisfies in $K(G)$ the formula analogous to (1.32), that is, that its product by $\lambda_{-1}(N)$ is $\lambda^p(F\lambda_{-1}(N))$. Likewise, the product of $\lambda^p(N,F)\in K(G)$ by $\lambda_{-1}(N)$ is $\lambda^p(F\lambda_{-1}(N))$, as follows from the fact that $K(G)$ is special. Since $K(G)$ is integral (cf. Section 2), and $\lambda_{-1}(N)\neq 0$ (same reference), formula (1.33) is proved. QED.

\textbf{Remarks.} a) When the dimension $q'$ of $F$ is at least $p$, formula (1.33) characterizes the polynomial $\lambda^p(N,x)$ in the $\lambda_N^i$ ($1\leq i\leq q$) and the $\lambda_x^j$ ($1\leq j\leq p$).

b) It is plausible that (1.33) remains valid if the field $k$ is not algebraically closed.\footnote{This is in fact proved by the same method.} One would have to resolve this question if one wanted to prove the Riemann--Roch theorem over a non-algebraically-closed field of characteristic $0$ by the method to be explained. The passage from an $\mathfrak S_p$-module to its alternating part has reasonable meaning only when the base field has characteristic $0$.

Let now $K$ be any special $\lambda$-ring and let $N$ be an element of $K$ such that $\lambda^iN=0$ for all sufficiently large $i$. Then for every $x\in K$ and every integer $p\geq 1$ one can consider the elements $\lambda^p(N,x)$, which are expressed by universal polynomials with integral coefficients in the $\lambda_N^i$ and the $\lambda_x^j$. One has $\lambda^1(N,x)=x$ and formula (1.32) in $K$. To express the properties of the operations $x\mapsto\lambda^p(N,x)$, it is convenient to introduce a new $\lambda$-ring, denoted $K_N$. As a commutative ring, $K_N$ is obtained by adjoining a unit to the ring whose underlying additive group is $K$ and whose multiplication is given by
\begin{equation}\tag{1.34}
x_Ny=xy\mu\qquad(\mu=\lambda_{-1}(N)).
\end{equation}
The map $\lambda_t$ from $K_N$ to $K_N[[t]]$ is defined by
\begin{equation}\tag{1.35}
\lambda_t(n,x)=(1+t)^n\sum_{p\geq 0}\lambda^p(N,x)t^p,
\end{equation}
where $\lambda^0(N,x)=(1,0)$ (the unit element of $K_N$). I say that, with these operations, $K_N$ is a special $\lambda$-ring. Thus one has to verify the formulas
\[
\lambda_t(X+Y)=\lambda_t(X)\lambda_t(Y),\qquad
\lambda_t(XY)=\lambda_t(X)\lambda_t(Y),\qquad
\lambda_t(\lambda^i(X))=\lambda^i(\lambda_t(X))
\]
for $X,Y\in K_N$. One may restrict oneself to verifying these formulas when $X$ and $Y$ lie in $K$, and one checks that it suffices to verify them when $K$ is the free special $\lambda$-ring generated by $X,Y$ and $N$ subject to the relations $\lambda^iN=0$ for $i>q$. But the group homomorphism
\begin{equation}\tag{1.36}
K_N\longrightarrow K,
\end{equation}
which sends unit to unit and on $K$ reduces to the homomorphism $x\mapsto x\mu$, is a homomorphism compatible with the ring structures and the maps $\lambda^i$, as one verifies immediately. In the present case, its restriction to $K$ is injective. Since $K$ is a special $\lambda$-ring, it follows easily that the above relations are indeed verified for $X,Y\in K$, completing the proof of our assertion.

According to formula (1.23), it is natural, for $x\in K$ and $n\geq 1$, to put
\begin{equation}\tag{1.37}
\gamma^n(N,x)=\sum_{i=0}^{n-1}\binom{n-1}{i}\lambda^i(N,x),
\end{equation}
in other words, $\gamma^n(N,x)=\gamma^n(x_N)$, where $x_N$ denotes $x$ considered as an element of the $\lambda$-ring $K_N$.

\textbf{Proposition 1.5.} Let $A$ be a graded commutative ring with unit, let $N$ be an element of $\widetilde A$ of the form
\[
N=[q,1+N^1+\cdots+N^q],
\]
and let $x$ be an arbitrary element of $\widetilde A$:
\[
x=[\nu,1+\sum_{i\geq 1}x^i].
\]
Then, for every integer $j\geq 0$, one has $\gamma^{q+j}(N,x)\in\widetilde A^j$, and the component of degree $j$ of $\gamma^{q+j}(N,x)$ is of the form
\begin{equation}\tag{1.38}
\bigl(\gamma^{q+j}(N,x)\bigr)^{(j)}=(-1)^{j-1}(j-1)!\,G_{q,j}(\nu,x^1,\ldots,x^j;N^1,\ldots,N^j),
\end{equation}
where $G_{q,j}$ is a universal polynomial with integral coefficients characterized by the condition
\begin{equation}\tag{1.39}
G_{q,j}(\nu,x^1,\ldots,N^j)\,(N^q)=(-1)^q(x*\lambda_{-1}(N))^{(q+j)}.
\end{equation}

\emph{Proof.} We may suppose that $A$ is the ring of polynomials with integral coefficients in indeterminates $x^i$ and $N^j$ ($1\leq j\leq q$). Since, combining (1.32) and (1.37), one has
\begin{equation}\tag{1.40}
\gamma^n(N,x)\lambda_{-1}(N)=\gamma^n(x\lambda_{-1}(N)),
\end{equation}
and, for $n=q+j$, the second member has filtration $\geq q+j$ (use formula (1.22 bis), taking account of the fact that $x\lambda_{-1}(N)$ has zero augmentation), while $\lambda_{-1}(N)$ has filtration $q$ (formula (1.22)), it follows that $\gamma^n(N,x)$ has filtration $j$, since by (1.18) the associated graded of $A$ is integral here. Using the preceding formulas, one immediately obtains the indicated form for $(\gamma^{q+j}(N,x))^{(j)}$. QED.

\section*{Chapter II. Classes of coherent algebraic sheaves and Chern classes}

\subsection*{1. The theory of Chow}

Throughout this chapter, a base field $k$ is fixed; for simplicity it will be assumed algebraically closed. A large part of what follows, perhaps all of it, is nevertheless valid without this restriction. The algebraic spaces and the regular maps considered will be defined over $k$.

An algebraic space is called \emph{quasi-projective} if it is isomorphic to a locally closed subset of a projective space.

Let $X$ be a nonsingular connected quasi-projective variety of dimension $n$. Denote by $A(X)$, and call the \emph{Chow ring} of $X$, the graded ring of classes of cycles on $X$ modulo rational equivalence. Thus $A^i(X)$ is formed from the classes of cycles of dimension $n-i$. The fact that $A(X)$ is indeed a graded ring follows from the following theorem of Chow.

\textbf{Theorem 2.1.}\footnote{As a result of a misunderstanding on the author's part, the stated theorem is more than is presently known. In fact, one uses only the ``moving lemma'' of Chow, saying that two cycles can be moved in their classes so that their intersection is not excessive; cf. Chevalley Seminar 1958, \emph{Anneau de Chow et applications}.} Every cycle on $X$ is rationally equivalent to a nonsingular cycle. Given elements of $A^i(X)$ and $A^j(X)$, one can find nonsingular representatives of these elements, say $Y^{n-i}$ and $Z^{n-j}$, which meet everywhere transversely.

A cycle is called nonsingular if its components are nonsingular subvarieties. Two cycles are said to meet everywhere transversely if every component of one cuts every component of the other transversely, i.e. if the two secant varieties are nonsingular at their intersection points and their tangent varieties there are in general position.

Let $f$ be a morphism, i.e. a regular map, from $X$ to $Y$ (where $X$ and $Y$ are quasi-projective and nonsingular). Then $f$ defines, using the preceding theorem, a homomorphism of graded rings
\begin{equation}\tag{2.1}
f^*:A(Y)\longrightarrow A(X),
\end{equation}
namely inverse image of cycle classes. Thus $A(X)$ becomes a contravariant functor in $X$.

Let $f$ be a proper morphism from $X^n$ to $Y^m$ (i.e. in Weil's terminology, the graph of $f$ is complete over $Y$, cf. [3]). Then $f$ also defines an additive homomorphism
\begin{equation}\tag{2.2}
f_*:A(X^n)\longrightarrow A(Y^m),
\end{equation}
preserving the dimension of cycles, hence increasing the degree by $m-n$.




% --- Source pages 47--75 ---
Thus, relative to proper morphisms, and ignoring multiplicative structures and gradings, $A(X)$ is a covariant functor in $X$. Let us record the classical formula
\begin{equation}\tag{2.3}
 f_*(x\cdot f^*(y))=f_*(x)\cdot y .
\end{equation}

We shall also use Chow's theory of Chern classes. It consists in attaching to every vector bundle $E$ on $X$ (still assumed quasi-projective and nonsingular) Chern classes
\[
 c^i(E)\in A^i(X)\qquad (i\geq 1),
\]
in such a way as to satisfy the following conditions. Put $c^0(E)=1$ and
\begin{equation}\tag{2.4}
 \Ctil(E)=\left[\rank E,\sum_{i\geq 0}c^i(E)\right]\in \Atil(X)
\end{equation}
(see Chapter I, \S 3 for the definition of $\Atil(X)$). The conditions to be verified by a theory of Chern classes are then expressed as follows:
\begin{equation}\tag{2.5}
 \Ctil(E)=\Ctil(E')+\Ctil(E'')
\end{equation}
if the vector bundle $E$ is an extension of $E'$ by $E''$. If one puts
\[
 C(E)=\sum_i c^i(E),
\]
this condition is expressed by $C(E)=C(E')C(E'')$. Formula (2.5) permits one to extend $\Ctil$ to an additive homomorphism from $K(\xi(X))$ (cf. Chapter I, \S 2, Example 2) to $\Atil(X)$:
\begin{equation}\tag{2.6}
 \Ctil:K(\xi(X))\longrightarrow \Atil(X),
\end{equation}
and one wants $\Ctil$ to be a $\lambda$-homomorphism of $\lambda$-rings. This last condition may also be written
\begin{equation}\tag{2.7}
 \Ctil(E\otimes F)=\Ctil(E)\Ctil(F),\qquad
 \Ctil(\lambda^iE)=\lambda^i\Ctil(E),
\end{equation}
for two vector bundles $E$ and $F$; of course $\Ctil$ is automatically compatible with augmentations. Let $D$ be a divisor on $X$ and let $L(D)$ be the vector bundle it defines. One wants
\begin{equation}\tag{2.8}
 c^1(L(D))=\ellc(D),\qquad c^i(L(D))=0\quad (i>1),
\end{equation}
where, for any cycle $Z$, $\ellc(Z)$ denotes its class in $A(X)$. Finally, one imposes that the $c^i$ be functorial, i.e. that if $f$ is a morphism from $X$ to $Y$ and $E$ is a vector bundle on $Y$, then
\begin{equation}\tag{2.9}
 c^i(f^!(E))=f^*(c^i(E)),
\end{equation}
where $f^!(E)$ is the inverse image of $E$ on $X$. One may also write (2.9) in the following form: the notion of inverse image of vector bundles defines a $\lambda$-homomorphism of augmented $\lambda$-rings
\begin{equation}\tag{2.10}
 f^!:K(Y)\longrightarrow K(X),\footnote{The notation $f^!$ used here is in conflict with that used in duality theory (cf., for example, R. Hartshorne, \emph{Residues and Duality}, Lecture Notes in Mathematics, no. 20, Springer, 1966). This is why in this seminar we replace it by the notation $f^*$.}
\end{equation}
where here one writes $K(X)$ instead of $K(\xi(X))$, and (2.9) means that there is commutativity in the diagram
\[
\begin{tikzcd}[column sep=large,row sep=large]
K(Y) \arrow[r,"f^!"] \arrow[d,"\Ctil_Y"'] & K(X) \arrow[d,"\Ctil_X"]\\
A(Y) \arrow[r,"f^*"'] & A(X).
\end{tikzcd}
\]
Here the space on which $\Ctil$ is considered has been put as a subscript.

\textbf{Theorem 2.2.} There exists a theory of Chern classes satisfying conditions (2.5), (2.7), (2.8), and (2.9).

We shall also admit the following proposition.

\textbf{Proposition 2.3.} Let $Y$ be a closed subset of $X$ and let $U$ be its complement. If $x$ is an element of $A(X)$ inducing $0$ on $U$, then there exists a representative of $x$ which is a cycle carried by $Y$.

\textbf{Corollary.} Let $p=\dim X-\dim Y$. Then, if $i<p$, the homomorphism
\[
 A^i(X)\longrightarrow A^i(U)
\]
is bijective. Every element of the kernel of $A^p(X)\to A^p(U)$ is a linear combination of components of dimension $n-p$ of $Y$; hence it is proportional to $\ellc(Y)$ if $Y$ is irreducible.

\subsection*{2. Definition of the Chern classes of coherent algebraic sheaves}

Let $\cC$ be an abelian category, and let $\cC_0$ be a subclass of $\cC$. Denote by $K(\cC_0)$ the quotient group of the free group generated by the isomorphism classes of elements of $\cC_0$ by the subgroup generated by the elements of the form
\[
(E)-(E')-(E''),
\]
where $E$ is an extension of $E'$ by $E''$ and $E,E',E''\in\cC_0$. We must suppose that the classes of elements of $\cC_0$, for the isomorphism relation, form a set. Denote by $\gamma_{\cC_0}(E)$ the class of an object $E\in\cC_0$ in $K(\cC_0)$, and omit the subscript $\cC_0$ when there is no danger of confusion. The group $K(\cC_0)$ is the solution of a universal problem relative to functions
\[
E\longmapsto f(E)
\]
defined on $\cC_0$, with values in an arbitrary abelian group $A$, which are additive in the sense that $f(E)=f(E')+f(E'')$ whenever $E$ is an extension of $E'$ by $E''$.

Let $F$ be an additive functor from $\cC$ to $\cC'$ and let $\cC'_0$ be a subclass of $\cC'$ satisfying the same condition as $\cC_0$, such that for every $A\in\cC_0$ the $R^qF(A)$ are in $\cC'_0$ and are zero for $q$ sufficiently large. We suppose that injective resolutions exist in $\cC$, so that the derived functors $R^qF$ of $F$ exist. Then $F$ defines a homomorphism from $K(\cC_0)$ to $K(\cC'_0)$, still denoted $F$, by the formula
\begin{equation}\tag{2.11}
 F(\gamma_{\cC_0}(E))=\sum_{i\geq 0}(-1)^i\gamma_{\cC'_0}(R^iF(E)).
\end{equation}
More generally, if one has a cohomological functor $(F^i)$ from $\cC$ to $\cC'$ such that for every $E\in\cC_0$ the $F^i(E)$ are in $\cC'_0$ and are zero except for a finite number of $i$, one obtains a homomorphism from $K(\cC_0)$ to $K(\cC'_0)$. These remarks extend to multifunctors; moreover, the classical spectral sequences give transitivity properties which we shall not make explicit.

In particular, if $F$ is an exact functor, formula (2.11) simplifies to
\begin{equation}\tag{2.11 bis}
 F(\gamma_{\cC_0}(E))=\gamma_{\cC'_0}(F(E)).
\end{equation}
This applies in particular if $\cC=\cC'$, if $F$ is the identity functor, and if $\cC_0\subset \cC'_0$; one has a canonical homomorphism from $K(\cC_0)$ to $K(\cC'_0)$.

\textbf{Theorem 2.3.} Let $\cC$ be an abelian category and let $\cC_0\subset \cC_1$ be two subclasses such that the isomorphism classes of objects of $\cC_1$ form a set. Suppose that the following conditions are satisfied:

\begin{enumerate}[label=(\roman*)]
\item For every exact sequence $0\to A'\to A\to A''\to 0$ in $\cC$ such that $A$ and $A''$ are in $\cC_0$ (respectively $\cC_1$), one has $A'\in\cC_0$ (respectively $A'\in\cC_1$).
\item Let $u:A\to B$ and $v:C\to B$, with $A,B,C\in\cC_1$ and $u$ surjective. Then there exist $L\in\cC_0$ and homomorphisms $u':L\to C$ and $v':L\to A$ such that $u'$ is surjective and $vu'=uv'$.
\item For every $A\in\cC_1$ there exists an integer $d=d(A)$ such that for every left resolution $L$ of $A$ by objects of $\cC_0$ one has $Z_d(L)\in\cC_0$, where $Z_d$ denotes the cycles in dimension $d$.
\end{enumerate}
Under these conditions, the canonical homomorphism from $K(\cC_0)$ to $K(\cC_1)$ is an isomorphism.

\emph{Principle of the proof.} Let $A\in\cC_1$. By (ii) and (iii) there exists a finite resolution $L$ of $A$ by objects of $\cC_0$; put
\[
 \ell(A)=\sum_{i\geq 0}(-1)^i\gamma_{\cC_0}(L_i).
\]
One shows that the element of $K(\cC_0)$ thus constructed is independent of the resolution $L$ chosen. If one has a second resolution $L'$, one ``caps'' $L$ and $L'$ by a third resolution $L''$, with surjective homomorphisms from $L''$ to $L$ and $L'$; this is possible by (i) and (ii). One is then reduced to proving that $E_{\cC_0}(L)=E_{\cC_0}(L'')$. But the difference of the two is $E_{\cC_0}(N)$, where $N$ is the kernel of the surjective homomorphism from $L''$ to $L$; this kernel is in $\cC_0$ by (i). Since $N$ is an acyclic complex in all dimensions, $E_{\cC_0}(N)=0$, whence the conclusion. Here $E_{\cC_0}(L)$, for a finite complex with values in $\cC_0$, denotes the Euler--Poincar\'e characteristic
\[
 \sum_i (-1)^i\gamma_{\cC_0}(L_i).
\]
One proves analogously that the function $\ell:\cC_1\to K(\cC_0)$ thus constructed is additive; hence it defines a homomorphism from $K(\cC_1)$ to $K(\cC_0)$, and it is immediate that this homomorphism is inverse to the canonical homomorphism from $K(\cC_0)$ to $K(\cC_1)$. QED.

\textbf{Remark.} To verify (ii), it is enough to verify that $\cC_0$ is contained in a subclass $\cC_2$ of $\cC$ satisfying the two conditions
\begin{enumerate}[label=(ii \alph*)]
\item every $A\in\cC_1$ is isomorphic to a quotient of an $L\in\cC_2$;
\item if $u:A\to B$, with $A,B\in\cC_2$, then $\Ker u\in\cC_1$.
\end{enumerate}

\textbf{Corollary.} Let $X$ be a quasi-projective algebraic space. Let $\cF(X)$ be the category of coherent algebraic sheaves on $X$, let $\xi(X)$ be the category of coherent locally free algebraic sheaves, equivalently vector bundles on $X$, and let $\cF_0(X)$ be the category of coherent algebraic sheaves $F$ such that, for every $x\in X$, the stalk $F_x$ is an $\cO_x$-module of finite cohomological dimension. Thus $\xi(X)\subset\cF_0(X)\subset\cF(X)$. Then the canonical homomorphism
\begin{equation}\tag{2.12}
 K(\xi(X))\longrightarrow K(\cF_0(X))
\end{equation}
is an isomorphism.

We verify the conditions of Theorem 2.3. Conditions (i) and (iii) follow from the definition of cohomological dimension and from the fact that a finite-type projective module over a noetherian local ring $\cO_x$ is free. To verify (ii), we shall prove that $\cC=\cF(X)$ satisfies conditions (ii a) and (ii b) of the remark. The second is trivial; therefore it remains to prove the following result.

\textbf{Proposition 2.4 (Serre).} Let $X$ be a quasi-projective algebraic space. Then every coherent algebraic sheaf on $X$ is isomorphic to the quotient of a locally free algebraic sheaf.

This proposition is true if $X$ is a closed subset of projective space by [2], since, for $n$ large enough, $F(n)$ is generated by its sections, which means that $F$ is a quotient of $\cO(-n)^k$. In the general case, $X$ is an open subset of a closed subset of a projective space, and it suffices to use the following result, due to Cartier and Serre in particular cases.

\textbf{Proposition 2.5.} Let $F$ be a coherent algebraic sheaf on an open subset $U$ of an algebraic space $X$. Then $F$ is isomorphic to the restriction of a coherent algebraic sheaf on $X$.

By applying Zorn's lemma to the opens containing $U$ on which one can extend $F$, one is reduced to the case where $X$ is affine. Let $\overline F$ be the algebraic sheaf on $X$ whose sections on an open $V$ are the sections of $F$ on $U\cap V$. One checks easily that $\overline F$ is quasi-coherent, i.e. defined on every affine open $V$---hence here on all of $X$---by a module, not necessarily of finite type, over the coordinate ring of that open. This is immediate if $U$ itself is affine, since $\overline F$ is then defined on the open $V$ by $\Gamma(U\cap V,F)$ considered as a module over $\Gamma(V,\cO_X)$. In the general case, $U$ is a finite union of affine opens $U_i$, hence $\overline F$ is the kernel of the evident homomorphism of quasi-coherent sheaves
\[
 \prod_i \overline F\big|_{U_i}\longrightarrow \prod_{i,j}\overline F\big|_{U_i\cap U_j},
\]
and is therefore itself quasi-coherent. Since $X$ is affine, $\overline F$ is generated by its sections. Since $\overline F|_U=F$ is coherent, there exists a finite number of sections of $\overline F$ which generate $F$ on $U$. These sections generate a coherent subsheaf of $\overline F$ whose restriction to $U$ is isomorphic to $F$. QED.

Now suppose that $X$ is nonsingular and quasi-projective. By the theorem on syzygies, one has, with the notation of the corollary to Theorem 2.3,
\[
 \cF_0(X)=\cF(X).
\]
In this case, therefore,
\begin{equation}\tag{2.13}
 K(\xi(X))=K(\cF(X)).
\end{equation}
One may express this isomorphism as follows.

\textbf{Corollary 2 to Theorem 2.3.} Let $X$ be a nonsingular quasi-projective variety. For every abelian group $A$, there is a one-to-one correspondence between additive functions $f$ of vector bundles on $X$ with values in $A$, and additive functions $g$ of coherent algebraic sheaves on $X$ with values in $A$. To every $f$ there corresponds the unique $g$ such that
\begin{equation}\tag{2.14}
 f(E)=g(\cO_X(E))
\end{equation}
for every vector bundle $E$ on $X$, where $\cO_X(E)$ denotes the sheaf of germs of regular sections of $E$ on $X$.

To compute $g(F)$ for an arbitrary coherent algebraic sheaf, one considers a finite resolution of $F$ by sheaves of the form $\cO_X(E_i)$, and one then has
\begin{equation}\tag{2.14 bis}
 g(F)=\sum_i (-1)^i f(E_i).
\end{equation}

In particular, if $X$ is connected, in \S 1 we considered an additive function $\Ctil(E)$ of the vector bundle $E$, with values in the group $A=\Atil(X)$. We deduce from it an additive function of coherent sheaves, say $\Ctil(F)$, and, by passing to components, functions
\[
 F\longmapsto c^i(F)\in A^i(X).
\]
The $c^i(F)$ are called the \emph{Chern classes of the coherent sheaf} $F$. Theoretically they are computed from the Chern classes of vector bundles by (2.14 bis), with the understanding that in the second member one must pass to multiplicative notation. As for the component of $\Ctil(F)$ in $\Z$, it is the rank of the sheaf $F$, corresponding to the rank of a module over an integral domain: if $X$ is affine, to which one can always reduce, $F$ is defined by a finite-type module over the coordinate ring of $X$, and one takes the rank of this module.

\subsection*{3. Functorial generalities on $K(X)$}

If $X$ is an algebraic space, put, for brevity,
\[
 K(X)=K(\cF(X)).
\]
When $X$ is nonsingular and quasi-projective, identify this group with $K(\xi(X))$ by the corollary 2 to Theorem 2.3. Since $K(\xi(X))$ is not only an abelian group but a $\lambda$-ring, it follows by transport of structure that $K(X)$ is also a $\lambda$-ring, and is therefore provided with a multiplicative structure and operations $\lambda^i$. It is important to define these structures directly, without passing through vector bundles; this will be done below.

If $E$ is a vector bundle on $X$, denote by $\gamma(E)$ its class in $K(X)$. If $F$ is a coherent algebraic sheaf on $X$, $\gamma(F)$ denotes in the same way its class in $K(X)$. If $Y$ is an irreducible subvariety of $X$, put $\gamma(Y)=\gamma(\cO_Y)$, where $\cO_Y$ denotes the sheaf of local rings of $Y$, considered as a coherent algebraic sheaf on $X$. By linearity, one defines the symbol $\gamma(Y)$ when $Y$ is any cycle on $X$. When there is risk of confusion, one writes $\gamma_X$ for $\gamma$.

Let $X$ be an arbitrary algebraic space. We introduce on $K(X)$ a filtration, denoting by $K(X)^i$ the subgroup of $K(X)$ generated by the $\gamma(F)$ with
\[
 \dim \Supp F\leq n-i,
\]
where $n$ is the dimension of $X$. From Serre's d\'evissage lemma, in the form given in [3], follows the following result.

\textbf{Proposition 2.6.} The group $K(X)^i$ is generated by the $\gamma(Y)$, where $Y$ runs through the subvarieties of dimension $n-i$ of $X$. If $X$ is irreducible, then $\Gr^0(K(X))=\Z$.

The latter isomorphism is given by the homomorphism from $K(X)$ to $\Z$ defined with the aid of the notion of rank of a coherent algebraic sheaf.

Now suppose that $X$ is nonsingular. Following the developments at the beginning of \S 2, define a multiplication in $K(X)$ by putting
\begin{equation}\tag{2.15}
 \gamma(F)\gamma(G)=\sum_{i\geq 0}(-1)^i\gamma(\Tor_i(F,G)),
\end{equation}
the $\Tor_i$ being taken, of course, with respect to the local rings of $X$. The exact sequence of the $\Tor_i$ shows that formula (2.15) indeed defines a biadditive map from $K(X)\times K(X)$ to $K(X)$; associativity follows from the spectral sequence of multiple Tors. One also verifies in this way that one can compute the product of several factors by the formula
\begin{equation}\tag{2.16}
 \gamma(F_1)\cdots\gamma(F_n)=\sum_{i\geq 0}(-1)^i\gamma(\Tor_i(F_1,\ldots,F_n)),
\end{equation}
where, in the second member, the $\Tor_i$ are simultaneous Tors, computed by simultaneous projective resolution of the arguments $F_i$.

If in (2.15) one supposes that $G$ is locally free, one finds
\begin{equation}\tag{2.15 bis}
 \gamma(F)\gamma(G)=\gamma(F\otimes G),
\end{equation}
which shows in particular that $\gamma(\cO_X)$ is the unit element of $K(X)$, and that the canonical homomorphism from $K(\xi(X))$ to $K(X)$ is a homomorphism of rings.

Formula (2.15) is evidently suggested by Serre's ``Tor formula'' in intersection theory. This formula proves the second part of the following proposition, the first part following immediately from (2.15) and from a local Tor calculation.

\textbf{Proposition 2.7.} Let $X$ be a nonsingular variety, and let $Y$ and $Z$ be two cycles in $X$. If $Y$ and $Z$ meet everywhere transversely, then
\begin{equation}\tag{2.17}
 \gamma(Y)\gamma(Z)=\gamma(Y\cdot Z).
\end{equation}
If $Y$ and $Z$ are homogeneous and of respective dimensions $n-i$ and $n-j$, and if they meet without an excessive component, then
\begin{equation}\tag{2.18}
 \gamma_X(Y)\gamma_X(Z)=\gamma_X(Y\cdot Z)\mod K(X)^{i+j+1}.
\end{equation}
One must be careful: equality does not in general hold in (2.18). We shall see in \S 4 that, if $X$ is nonsingular and quasi-projective, the ring structure of $K(X)$ is compatible with its filtration, and that the associated graded ring $\Gr(K(X))$ is identified with a quotient of $A(X)$ by a subgroup of torsion, perhaps always zero.\footnote{No: for a counterexample see this seminar, XIV 4.7.}

I do not know how to define directly the $\lambda^i(F)$ in terms of sheaves except when the field $k$ has characteristic $0$; this is the reason that has prevented us from proving the Riemann--Roch theorem when $k$ is not of characteristic $0$. Let $F$ be a coherent algebraic sheaf on $X$ on which the group $\mathfrak S_p$ of permutations of $p$ letters operates. Denote by $F^{\mathrm{alt}}$ the set of $x\in F$ such that
\[
 s\cdot x=\varepsilon(s)x
\]
for every $s\in\mathfrak S_p$, where $\varepsilon(s)$ is the signature of $s$. It is evidently a coherent algebraic subsheaf of $F$. This said, one has:

\textbf{Theorem 2.8.} Let $X$ be a nonsingular quasi-projective variety over a field $k$ of characteristic $0$. Let $F$ be a coherent algebraic sheaf on $X$. Then
\begin{equation}\tag{2.19}
 \lambda^p(\gamma(F))=\sum_{i\geq 0}(-1)^i\gamma\left((\Tor_i(F,F,\ldots,F))^{\mathrm{alt}}\right).
\end{equation}
Here it is clear that the group $\mathfrak S_p$ operates on the sheaf $\Tor_i(F,\ldots,F)$.

\emph{Principle of the proof.} One considers a finite resolution $L$ of $F$ by locally free sheaves $L_i$, and uses the following general lemma.

\textbf{Lemma 2.9.} Let $L=(L_i)$ be a graded locally free coherent algebraic sheaf with only finitely many nonzero components. In $K(X)$ one has the formula
\begin{equation}\tag{2.20}
 E\left((\otimes^p L)^{\mathrm{alt}}\right)=\lambda^p(E(L)).
\end{equation}
In this formula, $\otimes^pL$ is considered as a sheaf on which the group $\mathfrak S_p$ acts, taking the gradings into account; the gradings therefore introduce signs in the usual way, cf. Cartan--Eilenberg. Moreover, if $M$ is a graded coherent algebraic sheaf on $X$, put
\[
 E(M)=\sum_i(-1)^i\gamma(M_i).
\]

To prove formula (2.20), one is reduced to proving the analogous formula in $K(G)$, where $G$ is an algebraic group, for instance a product of groups $\operatorname{Gl}(n,k)$, and where $L$ is a finite-dimensional $G$-module over $k$ (cf. Chapter I, \S 2, Example 2, for the definition of $K(G)$). In that case, this formula has already been used in Theorem 1.4 to show that the second member of (1.33) satisfies the same functional equation as the first member; the verification is elementary. QED.

Let $f$ be a morphism from $X$ to $Y$, where $Y$ is an algebraic space. The notion of inverse image of vector bundle defines a homomorphism
\[
 f^!:K(\xi(Y))\longrightarrow K(\xi(X))
\]
of $\lambda$-rings. Identifying vector bundles with locally free coherent algebraic sheaves, the notion of inverse image of vector bundle corresponds to the notion of inverse image of algebraic sheaf, already used in [3, no. 8]. When $Y$ is nonsingular, one defines more generally a homomorphism of groups
\begin{equation}\tag{2.21}
 f^!:K(Y)\longrightarrow K(X)
\end{equation}
by a formula involving an alternating sum of Tors, inspired by the developments at the beginning of \S 2, and which the reader will spell out. One then has commutativity in the following diagram:
\[
\begin{tikzcd}[column sep=large,row sep=large]
K(\xi(Y)) \arrow[r,"f^!"] \arrow[d] & K(\xi(X)) \arrow[d]\\
K(Y) \arrow[r,"f^!"'] & K(X),
\end{tikzcd}
\]
which shows in particular, when $X$ and $Y$ are both nonsingular and quasi-projective, that the two definitions of $f^!$ coincide under the identifications of $K(X)$ with $K(\xi(X))$ and of $K(Y)$ with $K(\xi(Y))$.

Another case where one can define a homomorphism (2.21), without $Y$ necessarily nonsingular, is the case where the functor $F\mapsto f^!(F)$ from $\cF(Y)$ to $\cF(X)$ is exact, i.e. where the local rings of $X$ are flat modules over the local rings of $Y$. It is enough then to put
\[
 f^!(\gamma_Y(F))=\gamma_X(f^!(F)).
\]
This case occurs in particular when $X$ is a locally trivial fibre space over $Y$; to see this, one is reduced to the case of a product $X=Y\times T$. Note that one also verifies that in this case one has $f^!(\cO_Z)=\cO_{f^{-1}(Z)}$ if $Z$ is an irreducible subvariety of $Y$, whence
\begin{equation}\tag{2.22}
 f^!(\gamma_Y(Z))=\gamma_X(f^{-1}(Z))
\end{equation}
if $X$ is flat over $Y$. Without this latter condition, formula (2.22) becomes false in general, but one has the following analogue of Proposition 2.7.

\textbf{Proposition 2.8.} Let $f:X\to Y$ be a morphism of nonsingular varieties, and let $Z$ be a cycle on $Y$. If $Z$ is everywhere transverse to $f$, formula (2.22) is exact. If $Z$ is homogeneous of codimension $i$ and if $f^{-1}(Z)$ has no excessive component, then
\begin{equation}\tag{2.23}
 f^!(\gamma_Y(Z))=\gamma_X(f^{-1}(Z))\mod K(X)^{i+1}.
\end{equation}
The proof is analogous to that of Proposition 2.7.

Finally, let $X$ and $Y$ again be arbitrary algebraic spaces, and let $f$ be a proper morphism from $X$ to $Y$. For every coherent algebraic sheaf $F$ on $X$, the direct image sheaf $f_*(F)$ is coherent algebraic, and the same holds more generally for the sheaves $R^if_*(F)$ (cf. [3]). Recall that, by definition, $R^if_*(F)$ is the sheaf on $Y$ associated with the presheaf which sends an open $U$ to the group $H^i(f^{-1}(U),F)$; in particular
\begin{equation}\tag{2.24}
 R^0f_*(F)=f_*(F),\qquad R^qf_*(F)=0\quad(q>\dim X),
\end{equation}
and moreover
\begin{equation}\tag{2.25}
 \Gamma(U,R^qf_*(F))=H^q(f^{-1}(U),F)
\end{equation}
if the open $U$ is affine. According to \S 2, one defines a homomorphism from $K(X)$ to $K(Y)$, still denoted $f_*$, by the formula
\begin{equation}\tag{2.26}
 f_*(\gamma(F))=\sum_{q=0}^{\infty}(-1)^q\gamma(R^qf_*(F)).
\end{equation}
Thus $K(X)$ becomes a covariant functor in $X$ relative to proper morphisms. The relation $(fg)_*=f_*g_*$, just like the relation $(fg)^!=g^!f^!$ which we forgot to mention in its place, is a particular case of the transitivity relations alluded to in \S 2 and which follow from the spectral sequences of composed functors.

Let us note the following formula, valid for $Y$ nonsingular and particularly easy to verify when $Y$ is also quasi-projective by taking for $y$ the class of a vector bundle on $Y$:
\begin{equation}\tag{2.27}
 f_*(x\cdot f^!(y))=f_*(x)\cdot y\qquad (x\in K(X),\ y\in K(Y)).
\end{equation}
In the particular case where the inverse images of the points of $Y$ by $f$ are finite and contained in an affine open,\footnote{This latter condition is evidently a consequence of the first, something the author still did not know in 1957!} one checks that the functor $F\mapsto f_*(F)$ is exact and that $R^qf_*(F)=0$ for $q>0$. In this case, therefore, $f_*(\gamma_X(F))=\gamma_Y(f_*(F))$. In particular, if $X$ is an algebraic subspace of $Y$ and if $i$ is the canonical injection of $X$ into $Y$, then
\begin{equation}\tag{2.28}
 i_*(\gamma_X(F))=\gamma_Y(F),
\end{equation}
where, in the second member, $F$ is considered as a sheaf on $Y$ which is zero outside $X$.

\subsection*{4. Some technical results}

Let $X$ be an algebraic space. We have seen in \S 3 that the projection $p$ from $X\times k$ to $X$ defines a homomorphism
\begin{equation}\tag{2.29}
 p^!:K(X)\longrightarrow K(X\times k).
\end{equation}

\textbf{Proposition 2.9.} The homomorphism (2.29) is surjective.

We argue by induction on $n=\dim X$. The assertion is trivial if $n<0$; suppose $n\geq 0$. By Proposition 2.6, it is enough to verify that for every irreducible closed subset $Y$ of $X\times k$, $\gamma(Y)$ is in the image of $p^!$. It is enough to continue the reasoning by supposing $X$ irreducible. If $p(Y)$ is not dense in $X$, our assertion follows from the induction hypothesis. Otherwise, either $Y=X\times k$ and the assertion is trivial, or $Y$ is a hypersurface in $X\times k$. Let $U$ be an affine open of $X$ formed of nonsingular points. If $D$ is the divisor on $U\times k$ induced by $Y$, one knows by commutative algebra that it is linearly equivalent to a divisor of the form $D'\times k$, where $D'$ is a divisor on $U$. Then
\[
 \gamma_{U\times k}(D)-\gamma_{U\times k}(D'\times k)
\]
belongs to $K(U\times k)^2$ by the corollary to Proposition 2.10, and by what has already been proved this element is of the form $p^!(x')$ with $x'\in K(U)$. Thus
\[
 \gamma_{U\times k}(D)=p^!(x'+\gamma_U(D')),
\]
but $x'$ is the restriction of an element $y$ of $K(X)$ by Proposition 2.10, and consequently the restriction of $\gamma_{X\times k}(D)-p^!(y)$ to $U\times k$ is zero. By Proposition 2.10, $\gamma_{X\times k}(D)-p^!(y)$ comes from an element of
\[
 K(X\times k-U\times k)=K(X'\times k),\qquad X'=X-U.
\]
The proof is finished by applying the induction hypothesis again. QED.

\textbf{Corollary.} Let $X$ be an algebraic space and let $p$ be the projection from $X\times k^n$ to $X$. Then the homomorphism
\[
 p^!:K(X)\longrightarrow K(X\times k^n)
\]
is surjective; it is even bijective if $X$ is nonsingular.

The first assertion is by induction on $n$. For the second, put $f(x)=(x,0)$, whence $pf=1$ and $f^!p^!=1$, which proves that $p^!$ is injective.

\textbf{Proposition 2.10.} Let $X$ be an algebraic space, let $U$ be an open subset of $X$, and let $Y$ be its complement in $X$. Then the sequence of natural homomorphisms
\begin{equation}\tag{2.30}
 K(Y)\longrightarrow K(X)\longrightarrow K(U)\longrightarrow 0
\end{equation}
is exact.

N.B. Compare with Proposition 2.3.

The fact that $K(X)\to K(U)$ is surjective follows at once from Proposition 2.5, or also from Proposition 2.6. It remains to prove that if $x\in K(X)$ has zero restriction in $K(U)$, then it is in the image of $K(Y)$. For this, return to the explicit definition of $K(X)$ and $K(U)$. There are two coherent algebraic sheaves $F$ and $G$ on $X$ with
\[
 x=\gamma_X(F)-\gamma_X(G)
\]
and
\[
 \gamma_U(F)-\gamma_U(G)=0.
\]
This latter relation means that one can find coherent sheaves $P'$ and $Q'$ on $U$, provided with filtrations whose factors are pairwise isomorphic on $U$, and such that one has an isomorphism on $U$
\[
 P=F+Q'\longrightarrow Q=G+P'.
\]
By Proposition 2.5, the sheaves $P'$ and $Q'$ are restrictions to $U$ of coherent sheaves on $X$, denoted by the same letters; in the same way, their filtrations come from filtrations on the corresponding sheaves on $X$, by Lemma 2.11 below. One then has
\[
 \gamma_X(F)-\gamma_X(G)=(\gamma_X(P)-\gamma_X(Q))+\bigl(\gamma_X(P')-\gamma_X(Q')\bigr).
\]
Introducing the factors $P_i$ and $Q_i$ of the filtrations of $P'$ and $Q'$ on $X$, the last term is the sum of terms $\gamma_X(P_i)-\gamma_X(Q_i)$, where $P_i|_U\simeq Q_i|_U$. We are thus reduced to proving the following particular case: if $R$ and $S$ are coherent sheaves on $X$ such that $R|_U\simeq S|_U$, then $\gamma_X(R)-\gamma_X(S)$ is in the image of $K(Y)$. Let $T$ be the graph of the given isomorphism from $R|_U$ to $S|_U$; it is a coherent subsheaf of $(R\times S)|_U$, hence by Lemma 2.11 it is the restriction to $U$ of a coherent subsheaf, again denoted $T$, of $R\times S$. Then
\[
 \gamma_X(R)-\gamma_X(S)=(\gamma_X(R)-\gamma_X(T))-(\gamma_X(S)-\gamma_X(T)),
\]
and it is enough to prove, for instance, that the first term of the second member is in the image of $K(Y)$. But the projection $u:T\to R$ is an isomorphism over $U$, hence
\[
 \gamma_X(R)-\gamma_X(T)=\gamma_X(\Coker u)-\gamma_X(\Ker u),
\]
and $\Ker u$ and $\Coker u$ have support in $Y$.

It remains, then, to prove the following lemma.

\textbf{Lemma 2.11.} Let $X$ be an algebraic space, let $F$ be a coherent algebraic sheaf on $X$, let $U$ be an open subset of $X$, and let $G$ be a coherent subsheaf of $F|_U$. Then $G$ is the restriction to $U$ of a coherent subsheaf of $F$.

Let $\overline G$ be the subsheaf of $F$ whose sections on an open $V$ of $X$ are the sections of $F$ on $V$ whose restriction to $U\cap V$ is a section of $G$. Everything amounts to proving that this sheaf $\overline G$ is coherent. For this one may suppose $X$ and $U$ affine, because $U$ is a union of affine opens $U_i$ and $\overline G$ is then the intersection of the analogous sheaves $\overline{G|U_i}$. Then the coordinate ring of $U$ is a ring of fractions $A_S$ of the coordinate ring $A$ of $X$, and one immediately verifies that if $F$ is defined by the $A$-module $M$, then $F|_U$ is defined by the $A_S$-module
\[
 M_S=A_S\otimes_A M.
\]
The subsheaf $G$ of $F|_U$ is then defined by a submodule $N'$ of $M_S$, and one sees that $\overline G$ is precisely defined by the submodule $N$ of $M$, inverse image of $N'$ under the canonical homomorphism from $M$ to $M_S$. QED.

\textbf{Corollary 1 to Proposition 2.10.} Let $x\in K(X)$. Then $x\in K(X)^i$ if and only if there exists a closed subset $Y$ of $X$, of dimension $\leq n-i$, such that the restriction of $x$ to $X-Y$ is zero.

\textbf{Corollary 2.} Suppose $X$ is connected of dimension $n$, with no singularities in dimension $n-1$. If $D$ and $D'$ are two linearly equivalent divisors on $X$, then
\[
 \gamma(D)=\gamma(D')\mod K(X)^2.
\]
By Corollary 1, one may suppose $X$ nonsingular. The corollary then follows from the general formula
\begin{equation}\tag{2.31}
 \gamma(D)=1-\gamma(L(-D))\mod K(X)^2,
\end{equation}
where $L(-D)$ denotes the rank-one vector bundle defined by the divisor $-D$. To verify (2.31), one writes it in the form
\begin{equation}\tag{2.31 bis}
 \gamma(L(-D))=1-\gamma(D)\mod K(X)^2,
\end{equation}
and notes that the two members are homomorphisms from the group of divisors to the multiplicative group of invertible elements of the ring $K(X)/K(X)^2$. Here $K(X)^1/K(X)^2$ has square zero, which is very easily verified on the explicit formula (2.15), using Corollary 1. Thus one is reduced to the case where $D$ is an irreducible hypersurface. More generally, if $D$ is a hypersurface whose components are disjoint, one has the equality
\begin{equation}\tag{2.32}
 \gamma(D)=1-\gamma(L(-D)),
\end{equation}
as follows from the exact sequence of sheaves
\[
0\longrightarrow \cO_X(L(-D))\longrightarrow \cO_X\longrightarrow \cO_D\longrightarrow 0.
\]

\textbf{Theorem 2.12.} Let $X$ be a nonsingular quasi-projective variety, and let $Z$ and $Z'$ be two homogeneous cycles of codimension $p$ on $X$ which are linearly equivalent. Then
\begin{equation}\tag{2.33}
 \gamma_X(Z)=\gamma_X(Z')\mod K(X)^{p+1}.
\end{equation}
By hypothesis, one can find a cycle $Y$ of dimension $n-p+1$ on $X\times k$ ($n$ being the dimension of $X$), and two points $a,a'\in k$, such that
\[
 Z=Y\cdot (X\times a),\qquad Z'=Y\cdot (X\times a').
\]
By formula (2.23), one then has
\[
 \gamma_X(Z)=i_a^!(\gamma_{X\times k}(Y))\mod K(X)^{p+1},
\]
and an analogous formula for $\gamma_X(Z')$. It is therefore enough to prove that $i_a^!=i_{a'}^!$, which follows at once from Proposition 2.9. QED.

\textbf{Corollary 1.} The filtration of $K(X)$ is compatible with its ring structure, i.e. one has
\[
 K(X)^iK(X)^j\subset K(X)^{i+j}
\]
for all integers $i,j\geq 0$. Moreover there is a natural surjective homomorphism of graded rings
\begin{equation}\tag{2.34}
 \varphi:A(X)\longrightarrow \Gr(K(X)),
\end{equation}
where the second member denotes the graded ring associated with $K(X)$, provided with the filtration $(K(X)^i)_{i\geq 0}$, defined by the condition
\[
 \varphi(\ellc(Z^{n-p}))=\gamma(Z^{n-p})\mod K(X)^{p+1}.
\]
This corollary follows immediately from Theorem 2.12, Propositions 2.6 and 2.7, and Chow's theorem 2.1, by induction on the integer $i+j$.

\textbf{Corollary 2.} Let $X$ and $Y$ be two nonsingular quasi-projective varieties and let $f$ be a morphism from $X$ to $Y$. Then the homomorphism $f^!$ from $K(Y)$ to $K(X)$ is compatible with the filtrations, and one has a commutative diagram
\[
\begin{tikzcd}[column sep=large,row sep=large]
A(Y) \arrow[r,"\varphi_Y"] \arrow[d,"f^*"'] & \Gr(K(Y)) \arrow[d,"\Gr(f^!)"]\\
A(X) \arrow[r,"\varphi_X"'] & \Gr(K(X)).
\end{tikzcd}
\]
Using the graph of $f$, decompose $f$ as the product of an immersion of $X$ into $X\times Y$ and of a projection of $X\times Y$ onto $Y$. For the latter the corollary is trivial, and for an immersion one proceeds as for Corollary 1, using Proposition 2.8 in place of Proposition 2.7.

\textbf{Proposition 2.13 (``cellular decomposition'').} Let $P$ be an algebraic space, irreducible or not, provided with an increasing family
\[
P^0\subset P^1\subset\cdots\subset P^n=P
\]
of closed subsets, such that for every $i$ with $0\leq i\leq n$, the connected components of $P^i-P^{i-1}$ are algebraic spaces $V_{i,\alpha}$ isomorphic to $k^i$; put $P^{-1}=\varnothing$. Then:
\begin{enumerate}[label=\alph*)]
\item For every algebraic space $X$, the natural map from $K(X)\otimes K(P)$ to $K(X\times P)$, deduced from the tensor product of sheaves on $X$ and $P$, is surjective.
\item The $\gamma_P(V_{i,\alpha})$ form a system of generators of the group $K(P)$.
\end{enumerate}
The proof is by induction on $n$, using the corollary to Proposition 2.9 and Proposition 2.10. N.B. Unfortunately, I do not know whether the $\gamma_P(V_{i,\alpha})$ form a basis of $K(P)$ or whether the homomorphism considered in a) is bijective. We shall not need to know this below.

\textbf{Theorem 2.14.} Let $x_i$ ($1\leq i\leq q$) and $E_j$ ($1\leq j\leq r$) be letters, and let $n_j$ ($1\leq j\leq r$) be positive integers. Consider the symbols $\lambda^k(x_i)$ with $k\geq 1$ and $\lambda^k(E_j)$ with $1\leq k\leq n_j$ as indeterminates, and let $R$ be a polynomial with integral coefficients in these indeterminates. Suppose that, for every graded ring $A$, when one substitutes in $R$ elements $\xi_i$ for the $x_i$ and elements
\[
 \eta_j=[n_j,1+\eta_j^{(1)}+\cdots+\eta_j^{(n_j)}]
\]
of $\Atil$ for the $E_j$, subject to the restrictions $\varepsilon(\xi_i)=0$ and $\eta_j^{(k)}=0$ for $k>n_j$, one has
\[
 R\bigl((\lambda^k(\xi_i)),(\lambda^k(\eta_j))\bigr)\in \Atil^p.
\]
Then, if $X$ is an irreducible nonsingular quasi-projective variety, for arbitrary elements $x'_i$ of $K(X)^1$ and vector bundles $E'_j$ of dimension $n_j$ on $X$ ($1\leq i\leq q$, $1\leq j\leq r$), one has
\begin{equation}\tag{*}
 R\bigl((\lambda^k(x'_i)),(\lambda^k(\gamma(E'_j)))\bigr)\in K(X)^p.
\end{equation}

\emph{Proof.} By \S 2, the $x'_i$ are themselves differences of elements of the form $\gamma(E)$, with $E$ a vector bundle. On the other hand, assuming $X$ locally closed in a projective space $P$, it follows from \S 2 and Proposition 2.5 that, for every $F\in\cF(X)$, the sheaf $F(n)$ is generated by its sections for $n$ sufficiently large. If $F$ is defined by a vector bundle $E$, this proves that $E(n)$ is the inverse image of the canonical vector bundle on a Grassmannian variety $G$ by a morphism from $X$ to $G$. Considering the corresponding morphism from $X$ to $P\times G$, one sees that $F$ is the inverse image of a vector bundle on $P\times G$. Applying this to all vector bundles which occur, and using Corollary 2 to Theorem 2.12, we are reduced to the case where $X$ is a product of Grassmannians. In this case the conclusion will follow from the following lemma.

\textbf{Lemma 2.15.} Suppose that $X$ is a product of Grassmannians. Then one can find a $\lambda$-homomorphism from $K(X)$ to a torsion-free $\lambda$-ring of the form $\Atil$, where $A$ is a graded ring, compatible with filtrations, and such that the corresponding homomorphism $\Gr(K(X))\to\Gr(A)$ is injective.

To prove the lemma, we shall use Chow's theory of Chern classes. If one could prove the lemma, or Theorem 2.14 directly, one would deduce a purely ``sheaf-theoretic'' theory of Chern classes; cf. \S 5.

We shall consider the homomorphism $\Ctil$ from $K(X)$ to $\Atil(X)$ and put $A=A(X)$. It remains to prove that the homomorphism $\psi=\Gr(\Ctil)$ from $\Gr(K(X))$ to $\Gr(A(X))\simeq A(X)$ is injective. The fact that $\Ctil$ is compatible with filtrations is true for every nonsingular quasi-projective variety and follows easily from the corollary to Proposition 2.3. On the other hand, since $X$ is a product of Grassmannians $G_j$, there are vector bundles $E_j$ on $X$, inverse images of the canonical bundles on the factors. We shall use the following known facts:
\begin{enumerate}[label=\alph*)]
\item On $X$, linear equivalence of cycles is equivalent to numerical equivalence, and $A(X)$ is a group of finite rank.
\item The $c^i(E_j)$ generate the ring $A(X)$.
\end{enumerate}
From a) it follows that
\[
 \varphi:A(X)\longrightarrow \Gr(K(X))
\]
is injective, and hence bijective, by the following lemma.

\textbf{Lemma 2.15.} Let $X^n$ be a nonsingular projective variety and let $Z^{n-p}$ be a cycle on $X$ such that $\varphi(Z)\in\Gr^p(K(X))$ is zero. Then $Z^{n-p}$ is numerically equivalent to zero.

To prove this lemma, it suffices to use the additive function
\[
 \chi(F)=\sum_{i\geq 0}(-1)^i\dim H^i(X,F)
\]
and the fact that this function is equal to $1$ for $F=\cO_{(x)}$ for every $x\in X$.

Thus, in the case under consideration, $A(X)$ and $\Gr(K(X))$ have the same finite rank in every degree, and $\Gr(K(X))$ is torsion-free, since this is true of $A(X)$. From assertion b) above it follows, in a formal way, that the map $\Ctil$ from $K(X)$ to $\Atil(X)$ is surjective modulo finite groups. Since $\Gr(K(X))$ and $A(X)$ have the same finite rank in every degree, it then follows easily that the homomorphism $\Gr(\Ctil)$ from $\Gr(K(X))$ to $\Gr(\Atil(X))$ is bijective modulo finite groups in every degree; one uses induction on the degree. Since $\Gr(K(X))$ is torsion-free, this homomorphism $\Gr(\Ctil)$ is therefore injective. QED.

N.B. The homomorphism $\Ctil$ is not surjective in general, even if $X$ is a projective space of dimension $2$.

\subsection*{5. Sheaf-theoretic definition of Chern classes. Application to the study of injection morphisms}

\textbf{Theorem 2.16.} Let $X$ be a nonsingular quasi-projective variety. Then, for every $x\in K(X)$ and every integer $i\geq 0$, the element $\gamma^i(x-\varepsilon(x))$ is of filtration $\geq i$. Let $c^i(x)\in\Gr^i(K(X))$ be the class of this element modulo $K(X)^{i+1}$. If one puts
\[
 \widetilde c(x)=\left[\varepsilon(x),1+\sum_{i\geq 0}c^i(x)\right]\in \widetilde{\Gr(K(X))},
\]
then $\widetilde c$ satisfies conditions analogous to those of Theorem 2.2.

The first assertion follows from Theorem 2.14, as does the fact that $x\mapsto\widetilde c(x)$ is a homomorphism of $\lambda$-rings; formula (1.25) implies $\gamma_t(x+y)=\gamma_t(x)\gamma_t(y)$, which already shows that $\widetilde c$ is an additive homomorphism. The functoriality of $\widetilde c$ is verified immediately, and (2.8) follows from (2.31). QED.

\textbf{Corollary 1.} One has
\[
 c^i(x)=\varphi(c^i(x)).
\]
This follows from a uniqueness theorem for Chern classes, proved as in Hirzebruch by passage to flag varieties. One only has to show that, if $Y$ is a flag variety associated with a vector bundle on $X$, then $\Gr(K(X))\to\Gr(K(Y))$ is injective; this is verified by the explicit determination of $K(Y)$, with all its structures, from $K(X)$. I shall not give the details here.

\textbf{Corollary 2.} Suppose that for every element $Z\in A(X)$ there exists a morphism $f$ from $X$ to a product $X'$ of Grassmannians such that $Z$ is in the image of $f^*$. Then the canonical homomorphisms
\[
 \varphi:A(X)\longrightarrow \Gr(K(X)),\qquad
 \psi:\Gr(K(X))\longrightarrow A(X)
\]
satisfy, in dimension $i$, the relations
\begin{equation}\tag{2.35}
 \psi\varphi=(-1)^{i-1}(i-1)!\,1,
 \qquad
 \varphi\psi=(-1)^{i-1}(i-1)!\,1.
\end{equation}
Since $\varphi$ is surjective, it is clearly enough to prove the first formula, and for this one may suppose that $X$ is a product of Grassmannians. Since then $\varphi$ is injective, it is enough to prove the second formula (2.35), which by Corollary 1 is written
\[
 \gamma^i(x-\varepsilon(x))=(-1)^{i-1}(i-1)!\,x\mod K(X)^{i+1}
\]
for every $x\in K(X)^i$. Since $\psi$ is injective, it is enough to prove the analogous relation for $\Ctil(x)\in A(X)^i$, and this is then a consequence of (1.22 bis).

\textbf{Remark.} Corollary 2 makes plausible the validity of (2.35) in all cases; there should exist a direct proof of it. In every case one proves easily the second formula (2.35) modulo the torsion subgroup of $\Gr(K(X))$, and in characteristic $0$ this same formula is a consequence of the Riemann--Roch theorem. The interest of (2.35), when true, lies in the fact that it shows that $\varphi$ is injective modulo torsion groups, perhaps always zero. Thus one does not lose much by defining Chern classes as elements of $\Gr(K(X))$. If one tensors with $\Q$, as will be done later for the Riemann--Roch theorem, the two definitions of Chern classes would then be identical.

Let $X$ still be nonsingular and quasi-projective. Let in addition $Y$ be a nonsingular subvariety of $X$, and let $i$ be the injection of $Y$ into $X$. For $i_*$ and $i^!$ one then has the following formulas:
\begin{enumerate}[label=(\roman*)]
\item $i^!$ is a homomorphism of $\lambda$-rings.
\item $i_*(x\cdot i^!(y))=i_*(x)\cdot y$, whence in particular
\[
 i_*i^!(y)=\xi y,
\qquad \text{with }\xi=i_*i^!(1)=\gamma_X(N).
\]
\item $i^!i_*(x)=x\cdot\lambda_{-1}(N)$, whence in particular
\[
 i^!(\xi)=i^!i_*(1)=\lambda_{-1}(N).
\]
\item $i_*(x)i_*(x')=i_*(xx'\lambda_{-1}(N))$.
\end{enumerate}
Here, by definition, $N$ is the class in $K(Y)$ of the conormal bundle of $Y$ in $X$. Formulas (i) and (ii) are true for arbitrary morphisms, not necessarily injections, and (iii) and (iv) are verified from the definition (2.15) and the analogous formula for inverse images, taking $x$ and $x'$ to be classes of vector bundles and using the fact that
\begin{equation}\tag{2.36}
 \Tor(\cO_Y,\cO_Y)=\Lambda(N).
\end{equation}
Formula (iv) expresses that $i_*$ is a homomorphism of rings if one introduces on $K(Y)$ a new multiplicative structure, by the definition
\[
 x\cdot_N x'=xx'\lambda_{-1}(N)
\]
(cf. Chapter I, \S 4). It is completed by the formula
\begin{enumerate}[label=(\roman*)]
\setcounter{enumi}{4}
\item $i_*(\lambda^p(N,x))=\lambda^p(i_*(x))$ for $p\geq 1$,
\end{enumerate}
which means that $i_*$ defines a $\lambda$-homomorphism from $K(Y)_N$ into $K(X)$. This formula (v), certainly true in all cases, is for the moment proved only in characteristic $0$, because it follows from the conjunction of Theorems 1.4 and 2.8.

There are analogous formulas for $i_*$ and $i^*$:
\begin{enumerate}[label=(\roman* bis)]
\item $i^*$ is a homomorphism of graded rings.
\item $i_*(x\cdot i^*(y))=i_*(x)y$, whence in particular
\[
 i_*i^*(y)=\xi' y,
\qquad \text{with }\xi'=i_*(1).
\]
\item $i^*i_*(x)=x\cdot (-1)^qN^q$, whence in particular
\[
 i^*(\xi')=(-1)^qN^q.
\]
\item $i_*(x)i_*(x')=i_*(xx'(-1)^qN^q)$.
\item $i_*$ increases degrees by $q$ units.
\end{enumerate}
In these formulas, $q$ is the codimension of $Y$ in $X$, and one puts $N^i=c^i(N)$. Only formulas (iii bis) and (iv bis) require proof. I do not know whether they are well known, or even whether they are true, and I confine myself to noting that they follow by reduction from formulas (ii) and (iv), provided one works in $\Gr(K(Y))$ and $\Gr(K(X))$ and not in $A(Y)$ and $A(X)$. One then uses the fact that $i_*$ and $i^*$ come by passage to the quotient from $i_*$ and $i^!$, noting that $i_*$ shifts the filtration by $q$ units.

All universal formulas relating $i_*$ and $i^!$, respectively $i_*$ and $i^*$, are formal consequences of the preceding formulas, as one can convince oneself without great difficulty.

Formula (v) gives
\[
 i_*(\gamma^p(N,x))=\gamma^p(i_*(x)).
\]
The conjunction of Theorem 2.14 and Proposition 1.5 shows that $\gamma^p(N,x)$, for $p=q+j$, is of filtration $\geq j$. If one denotes by $G_{q,j}(N,x)$ its class in $\Gr^j(K(Y))$, Theorem 2.16 shows that
\[
 i_*(G_{q,j}(N,x))=c^p(i_*(x)).
\]
Finally, Proposition 1.5 joined with Theorem 2.14 permits one to express $G_{q,j}(N,x)$ as a polynomial with integral coefficients in the $c^i(N)$, the $c^i(x)$, and $\varepsilon(x)$. One deduces the following theorem.

\textbf{Theorem 2.17.} Suppose the field $k$ has characteristic $0$, and let $i:Y\to X$ be an injection morphism of nonsingular quasi-projective varieties. Then one has formulas (i) to (v), and moreover, for every integer $j\geq 0$, the formula
\begin{equation}\tag{2.37}
 c^{q+j}(i_*(x))=
 i_*\left(\left(\frac{\widetilde c(x)*\lambda_{-1}(\widetilde c(N))}{(-1)^qN^q}\right)^{(j)}\right).
\end{equation}
The second member of formula (2.37) is an abbreviated way of denoting the polynomial in the components of $\widetilde c(x)$ and $\widetilde c(N)$ described in Proposition 1.5.

\textbf{Corollary.} Under the preceding conditions one has
\begin{equation}\tag{2.38}
 \ch(\widetilde c(i_*(x)))=i_*\left(\ch(\widetilde c(x))\,\mathfrak C(\check N)^{-1}\right),
\end{equation}
where $\ch$ is the Chern homomorphism defined by formula (1.26).

It is enough to use (2.37) written in the form
\begin{equation}\tag{2.37 bis}
 \widetilde c(i_*(x))=1+i_*\left(\frac{\widetilde c(x)*\lambda_{-1}(\widetilde c(N))-1}{(-1)^qN^q}\right),
\end{equation}
the explicit formula (1.29), formula (iv bis), and finally (1.30).

One may write (2.38) in the form
\begin{equation}\tag{2.38 bis}
 \ch_X(F)=i_*\left(\ch_Y(F)\,\mathfrak C(T_{X/Y})^{-1}\right),
\end{equation}
where $F$ is a coherent algebraic sheaf on $Y$ (considered in the first member as a sheaf on $X$ which is zero outside $Y$), where $\ch_X$ and $\ch_Y$ are the Chern homomorphisms taken respectively on $X$ and $Y$, and where $T_{X/Y}$ is the normal bundle of $Y$ in $X$.

\textbf{Remarks.} Contrary to (2.38), formula (2.37) is a ``formula without denominators''. It has been proved, even in characteristic $0$, only by placing oneself in $\Gr(K(X))$ and not in $A(X)$, which may not be the same thing. It is certainly true in every characteristic and in $A(X)$. Let us note that one verifies at once that the images of the two members by $i^*$ are the same; hence, applying $i_*$, the product of their difference by $N^q$ is zero. This suggests a completely different method of proof, by trying to obtain $Y$ as a transverse inverse image of a subvariety $Y'$ in a variety $X'$ such that, in dimension $\leq n=\dim X$, the class of $Y'$ in $A(X')$ is not a divisor of $0$. Unfortunately one cannot hope for $X'$ to be a Grassmannian or something close to it.

Note that, if torsion elements are neglected, formula (2.38) is equivalent to (2.37).




% Additional macro definitions
\providecommand{\cS}{\mathcal S}
\providecommand{\cD}{\mathcal D}
\providecommand{\cE}{\mathscr E}
\providecommand{\cG}{\mathscr G}
\providecommand{\cL}{\mathscr L}
\providecommand{\cM}{\mathscr M}
\providecommand{\cP}{\mathscr P}
\providecommand{\cK}{\mathcal K}
\providecommand{\Db}{\mathrm D}
\providecommand{\Kb}{\mathrm K}
\providecommand{\Trc}{\operatorname{Tr}}
\providecommand{\ob}{\operatorname{ob}}
\providecommand{\Coh}{\operatorname{Coh}}
\providecommand{\Cone}{\operatorname{Cone}}
\providecommand{\im}{\operatorname{Im}}
\providecommand{\Id}{\operatorname{Id}}
\providecommand{\RR}{\mathrm{R.R.}}
\providecommand{\RRR}{\mathrm{R.R.R.}}
\providecommand{\GrR}{\mathrm{G}}
\providecommand{\pr}{\operatorname{pr}}


\section*{6. The Riemann--Roch theorem}

Its statement is as follows: let \(f\) be a proper morphism from a non-singular quasi-projective variety \(Y\) into a variety \(X\) of the same kind. Then one has
\begin{equation}\tag{2.39}
 f_*\bigl(\ch_Y(x)\,\mathfrak C(Y)\bigr)
 =
 \ch_X\bigl(f_!(x)\bigr)\,\mathfrak C(X).
\end{equation}
Both sides are in \(A(X)\otimes \mathbb Q\), and \(\mathfrak C(X)=\mathfrak C(T_X)\).

This formula makes it possible to compute, up to torsion, the Chern classes of \(f_!(x)\) from the Chern classes of \(x\), those of \(T_X\) and \(T_Y\), and the homomorphism \(f_*\). Or, equivalently, it makes it possible to determine the homomorphism
\[
  f_!:K(Y)\otimes \mathbb Q \longrightarrow K(X)\otimes \mathbb Q
\]
when these groups are identified, by means of the homomorphisms \(\ch_Y\) and \(\ch_X\), respectively with \(A(Y)\otimes \mathbb Q\) and \(A(X)\otimes \mathbb Q\) (provided that formula (2.35) is true), from \(f_*\) and the Chern classes of \(T_X\) and \(T_Y\).

I can prove formula (2.39) only in characteristic \(0\), and as an equality in \(\GrR(K(X))\otimes\mathbb Q\). For this, since \(Y\) is locally closed in a projective space \(P\), \(f\) is decomposed as the product of an injection of \(Y\) into \(P\times X\) and a projection morphism from \(P\times X\) to \(X\), and (2.39) is proved for each of these two morphisms.

For an injection morphism, one sees immediately that formulas (2.38) and (2.39) are equivalent. The characteristic \(0\) hypothesis enters only through Theorem 2.8, which permits one to express the classes \(\lambda^i(\gamma(F))\), for a coherent algebraic sheaf \(F\), by local cohomological invariants of \(F\).

For a projection morphism \(f:P\times X\to X\), Proposition 2.13 permits us to suppose that \(x\) is of the form \(y\otimes z\), with \(y\in K(P)\) and \(z\in K(X)\). Then the Künneth formula gives
\[
  f_!(x)=\chi(y)\,z,
\]
where, for a coherent algebraic sheaf \(F\) on \(P\), one puts
\[
  \chi(F)=\sum_{i\geq 0}(-1)^i\dim H^i(P,F).
\]
One is therefore at once reduced to proving the formula
\begin{equation}\tag{2.40}
  \chi(F)=\pi_r\bigl(\ch_p(F)\,\mathfrak C(P)\bigr),
\end{equation}
where \(\pi_r\) denotes the component of degree \(r=\dim P\). By Proposition 2.13, one may suppose that \(F=\mathcal O_Q\), where \(Q\) is a linear subvariety of \(P\). Then \(\chi(F)=1\), since this is the arithmetic genus of \(Q\); it remains only to see that the second member of (2.40) is \(1\).

Indeed, if \(V\) is a hyperplane of \(P\) and if one puts
\[
  \xi=\gamma_P(\mathcal O_V),\qquad L=\ell(V),
\]
then
\[
  1-\xi=\gamma_P(E(-L)),
\]
where \(E(-L)\) is the vector bundle defined by \(-L\). Hence
\[
  C_p(1-\xi)=[1,1-L],
\]
and consequently
\[
  \ch_p(1-\xi)=e^{-L},\qquad \ch_p(\xi)=1-e^{-L},
\]
and therefore
\[
  \ch_p(\xi^p)=(1-e^{-L})^p.
\]
If \(p\) is the dimension of \(Q\) in \(P\), one has \(\xi^p=\gamma_P(Q)\); hence
\[
  \ch_p(F)\,\mathfrak C(P)
  =
  \bigl(L/(1-e^{-L})\bigr)^{r+1}(1-e^{-L})^p.
\]
It remains to verify that the coefficient of \(L^r\) in this formal series in \(L\) is \(1\). This follows from the characteristic property of the formal series \(L/(1-e^{-L})\), namely that the coefficient of \(L^q\) in \(\bigl(L/(1-e^{-L})\bigr)^{q+1}\) is \(1\) (put \(q=r+1-p\)).

\subsection*{Bibliographical references}
\begin{enumerate}[label={[\arabic*]}]
\item Hirzebruch: \emph{Neue topologische Methoden in der algebraischen Geometrie}.
\item Serre: Faisceaux algébriques cohérents, \emph{Ann. of Math.}
\item Grothendieck: Sur les faisceaux algébriques cohérents, in Séminaire Cartan 1956/57.
\end{enumerate}

\section*{Proof of the Riemann--Roch theorem}
\begin{center}
(Grothendieck, 1 November 1957)
\end{center}

\subsection*{Reminder of notation}
These are the notations of the ``Report on Riemann--Roch,'' cited as \(\RRR\) below.

All varieties considered are non-singular and quasi-projective, that is, isomorphic to a locally closed subvariety of a projective space. The base field is algebraically closed, of arbitrary characteristic.

If \(X\) is such a variety, \(K(X)\) denotes the abelian group generated by the classes of coherent sheaves on \(X\), modulo the identification coming from an extension into a sum; the alternating sum of Tor makes \(K(X)\) a ring. If \(F\) is a sheaf, respectively \(E\) a vector bundle, respectively \(Y\) a subvariety of \(X\), one denotes by \(\gamma_X(F)\), respectively \(\gamma_X(E)\), respectively \(\gamma_X(Y)\), the element of \(K(X)\) corresponding to \(F\), respectively to \(E\), respectively to \(\mathcal O_Y\). In fact, one often writes simply \(F,E\) instead of \(\gamma_X(F)\).

Let \(A(X)\) be the graded ring of classes of cycles on \(X\) for rational equivalence (cf. Chow); \(A^i(X)\) is formed by the cycles of codimension \(i\). If \(E\) is a vector bundle, \(C(E)=\sum C^i(E)\), with \(C^i(E)\in A^i(X)\), is its Chern class (defined as inverse image of Schubert cycles). By linearity one deduces the classes \(C^i(y)\), for all \(y\in K(X)\). If
\[
  C(y)=\prod (1+a_i),
\]
one sets, following Hirzebruch,
\[
  \ch(y)=\sum e^{a_i},\qquad
  T(y)=\prod\frac{a_i}{1-e^{-a_i}}.
\]
These are elements of \(A(X)\otimes\mathbb Q\). When \(y\) is the tangent bundle of \(X\), one writes \(T(X)\) instead of \(T(y)\).

If \(G\) is a vector bundle, \(\lambda^pG\) denotes its \(p\)-th exterior power and
\[
  \lambda_{-1}(G)=\sum(-1)^p\lambda^p(G);
\]
this is an element of \(K(X)\).

If \(f\) is a proper regular map from a variety \(Y\) into a variety \(X\), one defines homomorphisms
\[
  f^*:A(X)\longrightarrow A(Y),\qquad
  f_*:A(Y)\longrightarrow A(X).
\]
The first is homogeneous of degree \(0\) and multiplicative. The second is homogeneous of degree \(\dim X-\dim Y\). One also defines homomorphisms
\[
  f^*:K(X)\longrightarrow K(Y),\qquad
  f_!:K(Y)\longrightarrow K(X).
\]
If \(F\) is a coherent sheaf on \(X\), \(f^*(F)\) is, by definition, the alternating sum of the \(\Tor_p^{\mathcal O_X}(\mathcal O_Y,F)\). For vector bundles, this is simply the inverse-image operation.

If \(G\) is a coherent sheaf on \(Y\), \(f_*(G)\) is coherent on \(X\). If \(R^pf_*\) denotes the \(p\)-th derived functor of the functor \(f_*\) so defined, one has, by definition,
\[
  f_!(G)=\sum(-1)^p R^pf_*(G).
\]

\subsection*{Riemann--Roch theorem}
Let \(f:Y\to X\) be a proper morphism (regular map) from a non-singular quasi-projective variety \(Y\) into another such variety \(X\), and let \(y\in K(Y)\). Then
\begin{equation}\tag{1}
  f_*\bigl(\ch(y)T(Y)\bigr)=\ch(f_!(y))T(X)
  \quad\text{in } A(X)\otimes\mathbb Q.
\end{equation}

\emph{Proof} (abridged).

\textbf{Lemma 1.} Consider proper morphisms of varieties
\[
  Z \xrightarrow{g} Y \xrightarrow{f} X.
\]
If the theorem \(\RR\) is true for \(f\) and for \(g\), it is true for \(fg\). If it is true for \(fg\) and for \(g\), it is true for \(f\). If \(f_*\) is injective, and if the theorem \(\RR\) is true for \(fg\) and for \(f\), it is true for \(g\).

This is trivial. In fact, when one supposes \(\RR\) true for \(fg\) and \(g\), one must still suppose that \(g_!\) is surjective in order to be able to assert that \(\RR\) is true for \(f\).

Since a proper morphism is composed of an injection and a projection \(X\times P\to X\), where \(P\) is a projective space, one is reduced to treating these two cases separately. The second, as was seen above, is reduced to the standard calculation, by means of Proposition 2.13 (``cellular decomposition'') of \(\RRR\). Thus suppose \(Y\subset X\), with injection morphism \(i\). Formula (1) becomes
\begin{equation}\tag{2}
  \ch(i_!(y))=i_*\bigl(\ch(y)T(E)^{-1}\bigr),
\end{equation}
where \(E=T_{X/Y}\) is the normal bundle of \(Y\) in \(X\).

\textbf{Lemma 2.} Suppose that there exists on \(X\) a vector bundle \(G\) of rank \(p=\codim_XY\), such that
\begin{equation}\tag{3}
  \gamma_X(Y)=i_*(1)=\lambda_{-1}(G),\qquad
  i^*(G)=E,\qquad
  (-1)^pC^p(G)=Y.
\end{equation}
If then
\begin{equation}\tag{4}
  y=i^*(x),\qquad x\in K(X),
\end{equation}
formula (2) is valid.

The proof is immediate from the fact that \(\ch\) is a homomorphism of rings, and that
\begin{equation}\tag{5}
  \ch(\lambda_{-1}(G))=(-1)^pC^p(G)\,T(G)^{-1}.
\end{equation}

\textbf{Corollary.} \(\RR\) is true if \(Y\) is a hypersurface and if \(y\) is of the form \(i^*(x)\), with \(x\in K(X)\).

\textbf{Lemma 3.} Suppose that \(X=Y\times P\), where \(P\) is a projective space, and that \(i\) is the injection \(u\mapsto (u,v_0)\), \(v_0\in P\). Then \(i_*\) is injective and \(\RR\) is true.

This is pure standard calculation.

Suppose again that \(X,Y,y\) are arbitrary, always with \(Y\subset X\). Let \(X'\) be the variety obtained by blowing up \(X\) along \(Y\), and let \(Y'\) be the hypersurface inverse image of \(Y\). One has a diagram of maps
\[
\begin{tikzcd}
Y' \arrow[r,"j"] \arrow[d,"g"'] & X' \arrow[d,"f"]\\
Y \arrow[r,"i"'] & X .
\end{tikzcd}
\]
The inverse image of \(E\) in \(Y'\) admits a subbundle of rank \(1\), whose dual is denoted \(L\). Moreover \(L^{-1}\) is also the restriction to \(Y'\) of the line bundle \(L(Y')\) defined on \(X'\) by the divisor \(Y'\). Put \(F=E/L^{-1}\). Elementary computations give the following lemma.

\textbf{Lemma 4.} One has the formulas, where \(p=\codim_XY\):
\begin{align}
  f^*i_!(y)&=j_*\bigl(\lambda_{-1}(\check F)y\bigr), \tag{6}\\
  \lambda_{-1}(\check F)&\equiv \gamma^{p-1}(E)\lambda^p(\check E)
    \pmod{1-L}, \tag{7}\\
  f_*(1)&=1,\quad\text{hence }f_*f^*=1, \tag{8}\\
  f^*i_*(\alpha)&\equiv j_*\bigl(\alpha C^{p-1}(F)\bigr)
    \pmod{\Ker f_*}\quad\text{for all }\alpha\in A(Y). \tag{9}
\end{align}
To simplify notation, \(K(Y)\) has been made to operate on \(K(Y')\), and \(A(Y)\) on \(A(Y')\), by the homomorphisms \(g^*\). In formula (7), it seems that \(\gamma^{p-1}(E)\) denotes \(\lambda^{p-1}(E-2)\), that is,
\[
  \sum_{i+j=p-1}(\,\cdots\,)\lambda^i(E);
\]
the congruence takes place in \(K(Y')\).

Formula (6) follows from the equality
\begin{equation}\tag{6 bis}
  \Tor_i^{\mathcal O_X}\bigl(\mathcal O_Y(v),\mathcal O_Y(v')\bigr)
  =
  \lambda^i\check F(v')
  \quad (v'\in Y',\ v=g(v')).
\end{equation}
Formula (7) follows from the formulas
\begin{equation}\tag{10}
  g_*(L^{-i})=0\quad\text{for }0<i\leq p-1,\qquad g_*(1)=1,
\end{equation}
themselves consequences of Künneth and of the cohomology of projective spaces.

Formula (8) is trivial. Finally, (9) follows from \(g_*(C^{p-1}(F))=1\), which is deduced, for example, from the immediate formulas
\begin{equation}\tag{10 bis}
  g_*(\xi^i)=0\quad\text{if }0\leq i<p-1,\qquad
  g_*(\xi^{p-1})=1,
\end{equation}
with \(\xi=C^1(L)\in A^1(Y')\). In fact, (9) is very likely an identity, and not merely a congruence. The unnecessarily complicated formula (7) may be replaced, for example, by VII 4.1.

From formula (7) it follows at once that \(\lambda_{-1}(\check F)\in j^*(K(X'))\), provided that
\begin{equation}\tag{11}
  y\,\gamma^{p-1}(E)\lambda^p(\check E)\in i^*(K(X)).
\end{equation}
Suppose this condition (11) satisfied. By the corollary to Lemma 2, \(\RR\) is true for \((Y',X',y\lambda_{-1}(\check F))\). On the other hand, to prove (2), formula (8) reduces the matter to proving that
\begin{equation}\tag{12}
  f^*\ch(i_!(y))
  \equiv
  f^*\bigl(i_*(\ch(y)T(E)^{-1})\bigr)
  \pmod{\Ker f_*}.
\end{equation}
This last formula is verified immediately, using (6), (9), \(\RR\) for \((Y',X',y\lambda_{-1}(\check F))\), and the formula
\begin{equation}\tag{13}
  \ch(\lambda_{-1}(\check F))=C^{p-1}(F)T(F)^{-1}
\end{equation}
(cf. formula (5)).

Thus \(\RR\) is proved every time \(y\,\gamma^{p-1}(E)=0\), in particular every time \(p-1>\dim Y\), since \(\gamma^{p-1}(E)\) is of filtration \(\geq p-1\) in \(K(Y)\), filtered by the codimension of supports; that is, every time
\[
  \dim Y < \frac{\dim X-1}{2}.
\]
One is reduced to this case by combining Lemma 1 and Lemma 3, taking for \(P\) a projective space of sufficiently large dimension, for example \(r>\dim X+1\). This proves the assertion.

\textbf{Remark.} We have used the fact that \(\gamma^{p-1}(E)\) was of filtration \(\geq p-1\). One can prove this fact directly, without passing through Grassmannians as in \(\RRR\), 2.16, by studying the projective bundle associated to \(E\). This method more generally gives Theorem 2.14 of \(\RRR\).

\subsection*{Extract from a letter of Grothendieck}
``Morally,'' the new method rests on the determination of \(K(X)\) and \(A(X)\) when \(X\) is either the projective bundle associated to a vector bundle, or a variety obtained by blowing up a non-singular subvariety. In this second case I have not made the complete determination; a small equality is still missing. With the notation of the short note, one must prove formula (9), but as an equality:
\begin{equation}\tag{9}
  f^*(i_*(\alpha))=j_*(\alpha C^{p-1}(F)),\qquad \alpha\in A(Y).
\end{equation}
One observes that the two members have the same image by \(f_*\), and the same image by \(j^*\), so that the question is equivalent to
\begin{equation}\tag{9 bis}
  \Ker f_*\cap\Ker j^*=0.
\end{equation}
In any case, (9) is true after reducing to \(G(K(X'))\) (which destroys a little torsion, at most), for it then follows from (6); there is therefore hardly any doubt that it is true.

I also point out that equality (9) follows very easily from (6), without passing through the troublesome (6 bis).

\begin{flushleft}
P.C.C.J-P.S.
\end{flushleft}

\newpage
\section*{Exposé I. Generalities on finiteness conditions in derived categories}
\begin{center}
by L. Illusie
\end{center}

\subsection*{0. Introduction}
The object of Exposés I to IV is to develop, with the degree of generality appropriate in this Seminar, the formalism of finiteness conditions in derived categories of ringed topoi.

As was said in Exposé 0, the need to define, on arbitrary schemes, ``Grothendieck groups'' possessing good variance properties forced one to generalize and to soften the finiteness notions used up to now.

A classical notion such as that of a coherent sheaf on a ringed space \((X,\mathcal O_X)\) makes sense only when \(\mathcal O_X\) is coherent. One might think of replacing it by the notion of finite presentation. But this has the inconvenience that the kernel of an epimorphism of modules of finite presentation is no longer, in general, of finite presentation. A satisfactory notion is reached by observing that, when \(\mathcal O_X\) is coherent, a coherent sheaf \(F\) is not only of finite presentation but of \(n\)-finite presentation for every \(n\in\mathbb N\), where \(n\)-finite presentation means that locally there exists an exact sequence
\[
  L_n\longrightarrow L_{n-1}\longrightarrow \cdots \longrightarrow L_0\longrightarrow F\longrightarrow 0,
\]
with the \(L_i\) free of finite type. When \(\mathcal O_X\) is no longer assumed coherent, say that an \(\mathcal O_X\)-module \(F\) is \emph{pseudo-coherent} if it satisfies the preceding condition, that is, if it is of \(n\)-finite presentation for every \(n\). Then the notion of pseudo-coherence has, with respect to short exact sequences, the same stability property as coherence: if two terms of a short exact sequence of \(\mathcal O_X\)-modules are pseudo-coherent, then so is the third.

The preceding notion is easily generalized to complexes of sheaves. A complex \(F\) of \(\mathcal O_X\)-modules is said to be \emph{pseudo-coherent} if it is \(n\)-pseudo-coherent for every \(n\in\mathbb Z\), where \(n\)-pseudo-coherent means that there locally exists a quasi-isomorphism \(L\to F\), with \(L\) a complex bounded above and with components free of finite type in degrees \(>n\). For a complex \(F\) reduced to degree \(0\), this is the same as saying that \(F\) is pseudo-coherent as a complex or as a module.

The notion of pseudo-coherent complex has excellent stability properties, described in Exposé I, §§1 and 2. First of all it is stable under isomorphism in the derived category \(\Db(X)\) of the category of \(\mathcal O_X\)-modules; better still, if two objects of a distinguished triangle of \(\Db(X)\) are pseudo-coherent, then so is the third. In other words, the full subcategory \(\Db(X)_{\mathrm{coh}}\) of \(\Db(X)\) formed by pseudo-coherent complexes is a triangulated subcategory. Moreover, pseudo-coherence is preserved under inverse image and derived tensor product. In the case where \(\mathcal O_X\) is coherent, to say that a complex \(F\) is pseudo-coherent simply means that \(F\) has locally bounded-above cohomology and that all the \(H^i(F)\) are coherent.

If \(\mathcal O_X\) is a sheaf of regular local rings, every coherent \(\mathcal O_X\)-module \(F\) admits locally a finite resolution to the left by finite-type free modules, that is, there exists locally an exact sequence
\[
  0\longrightarrow L_n\longrightarrow \cdots \longrightarrow L_0\longrightarrow F\longrightarrow 0,
\]
where the \(L_i\) are free of finite type. It is therefore natural, when \(\mathcal O_X\) is an arbitrary sheaf of local rings, to study the modules which have the preceding property. Such modules are called \emph{perfect}. More generally, a complex \(F\) of \(\mathcal O_X\)-modules is said to be \emph{perfect} if there exists locally a quasi-isomorphism \(L\to F\), where \(L\) is a bounded complex with components free of finite type. It is shown in Exposé I, §3, that the full subcategory \(\Db(X)_{\mathrm{parf}}\) of \(\Db(X)\) formed by perfect complexes is, like \(\Db(X)_{\mathrm{coh}}\), a triangulated subcategory, stable under inverse image and derived tensor product. One obviously has
\[
  \Db(X)_{\mathrm{parf}}\subset \Db(X)_{\mathrm{coh}},
\]
but the inclusion is strict in general. Pseudo-coherence can be redefined ``by passage to the limit'' from perfection: a complex \(F\) is pseudo-coherent if and only if, for every \(n\in\mathbb Z\), \(F\) can be locally approximated to order \(n\) by a perfect complex, meaning that there exists locally a distinguished triangle
\[
\begin{tikzcd}
L \arrow[rr] && F \arrow[dl]\\
& R \arrow[ul] &
\end{tikzcd}
\]
with \(L\) perfect and \(R\) acyclic in degrees \(>n\). Conversely, perfection is recovered from pseudo-coherence by means of an additional regularity condition: a complex is perfect if and only if it is pseudo-coherent and locally of finite Tor-dimension (Exposé I, §5). Thus one recovers the fact that, if the local rings are regular, every coherent sheaf is perfect, and more generally every pseudo-coherent complex with locally bounded cohomology is perfect.

Pseudo-coherence and perfection are the two fundamental finiteness notions with which we shall work in this Seminar. Sections 1 to 5 of Exposé I are devoted to their definition and to establishing their elementary stability properties. For this purpose, only two or three local properties of the category of finite-type free modules, relative to the category of all modules, are used. It is therefore convenient and useful, notably in Exposé II, to axiomatize the situation by introducing a notion of pseudo-coherence (respectively of perfection) in a category fibered over a site, relative to a fibered subcategory satisfying suitable conditions.

Sections 6 to 8 of Exposé I generalize to perfect complexes a certain number of notions well known for finite-type locally free sheaves: rank, duality, and the trace of an endomorphism.

In Exposé II, one examines the problem of the existence of ``global resolutions'': under what conditions, on a ringed topos \((X,\mathcal O_X)\), is a pseudo-coherent (respectively perfect) complex globally isomorphic in \(\Db(X)\) to a complex of locally free modules of finite type? The importance of this question lies essentially in the fact that at present one does not know how to generalize to perfect complexes certain usual constructions on locally free modules, such as exterior and symmetric powers or Chern classes. That is why it is useful to have workable sufficient criteria for the existence of global resolutions, which reduce certain questions about pseudo-coherent (respectively perfect) complexes to analogous questions about locally free sheaves of finite type. Such criteria are given in Exposé II.

In particular, one shows that global resolutions exist in the following cases: \(X\) is the Zariski topos of an affine scheme, or more generally of a quasi-compact scheme possessing an ample invertible sheaf, or even an ample family of invertible sheaves in the sense of Kleiman, for example a regular scheme; or again \(X\) is the topos of sheaves of sets on a compact topological space and \(\mathcal O_X\) is the sheaf of continuous functions with complex values. By way of illustration, and to show the flexibility of the notions introduced, the appendix gives a ``purely sheaf-theoretic'' definition of the index of a family of elliptic operators.

Exposé III studies the stability of finiteness conditions under derived direct image. To obtain reasonable statements, one must relativize the notions of pseudo-coherence and perfection. We place ourselves here in the framework of ordinary schemes, which is sufficient for the Seminar, but there is no doubt that one will sooner or later need to develop an analogous theory for relative analytic schemes or spaces. Let \(p:X\to S\) be an \(S\)-scheme locally of finite type, and let \(F\) be a complex of \(\mathcal O_X\)-modules. One says that \(F\) is pseudo-coherent (respectively perfect) relative to \(p\) if one can locally extend \(X\), by a closed immersion, into a smooth \(S\)-scheme \(X'\), so that \(F\), extended by zero to \(X'\), is pseudo-coherent (respectively perfect). The notion of pseudo-coherence relative to \(p\) is especially interesting when \(S\) is not locally noetherian; in the opposite case it coincides with the ordinary notion of pseudo-coherence. Moreover, as expected, to say that \(F\) is perfect relative to \(p\) is equivalent to saying that \(F\) is pseudo-coherent relative to \(p\) and locally of finite Tor-dimension relative to \(p\) (i.e. relative to the sheaf of rings \(p^{-1}(\mathcal O_S)\)).

The central theorem of Exposé III is the following finiteness theorem, stated below in a slightly more precise form: if \(f:X\to Y\) is a proper morphism of \(S\)-schemes locally of finite type, the functor \(Rf_*\) transforms complexes pseudo-coherent relative to \(S\) into complexes pseudo-coherent relative to \(S\). This theorem is in fact a conjecture, but it is nevertheless proved in the two particular cases \(S\) locally noetherian and \(f\) projective. It seems likely, unfortunately, that the extension to the general case is of the same order of difficulty as the analogous theorem in analytic geometry, Grauert's theorem. Combining the finiteness theorem with an essentially trivial formula called the projection formula, one obtains workable criteria for the stability of relative perfection under direct image. As a corollary one recovers Grauert's continuity and semicontinuity theorems (EGA III 7.6).

Exposé IV translates the results of Exposés I to III into the language of ``Grothendieck groups.'' On a ringed topos \((X,\mathcal O_X)\), write \(K^*(X)\) (respectively \(K_\bullet(X)\)) for the Grothendieck group of the category of perfect complexes and complexes of finite Tor-dimension (respectively pseudo-coherent complexes with bounded homology). The group \(K^*(X)\) is a ring, and \(K_\bullet(X)\) is a module over \(K^*(X)\). The functor \(K^*\) is contravariant, as suggested above, while \(K_\bullet\) is covariant for proper pseudo-coherent morphisms of schemes, or of analytic spaces, subject to the availability of the finiteness theorem. There are also unusual variants: covariance of \(K^*\), contravariance of \(K_\bullet\), for morphisms satisfying suitable regularity hypotheses. Finally, the globalization criteria of Exposé II make it possible to link these groups with the ``Grothendieck groups'' defined naïvely from coherent sheaves or finite-type locally free sheaves.

\section*{1. Preliminary definitions}

\subsection*{1.1. Fibered categories with additive fibers (respectively abelian, respectively triangulated fibers)}

\textbf{1.1.1.} Let \(S\) be a category, and let \(\mathcal C\) be an \(S\)-category fibered in categories with additive (respectively triangulated) fibers (SGA 1 VI 6.1). We shall say that \(\mathcal C\) is an \emph{additive} (respectively \emph{triangulated}) \(S\)-category, or that \(\mathcal C\) is additive (respectively triangulated) over \(S\), if, for every arrow \(f:X\to Y\) of \(S\), the inverse-image functor
\[
  f^*:\mathcal C_Y\longrightarrow \mathcal C_X
\]
(determined up to unique isomorphism) is additive (respectively exact). We shall say that \(\mathcal C\) is an \emph{abelian} \(S\)-category if \(\mathcal C\) is an additive \(S\)-category with abelian fibers, and that \(\mathcal C\) is a \emph{flat abelian} \(S\)-category if, moreover, the inverse-image functors are exact.

\textbf{1.1.2.} Let \(\mathcal C\) be a triangulated \(S\)-category. One defines, for example with the aid of a cleavage of \(\mathcal C\) (SGA 1 VI 7.1), a fibered \(S\)-category \(\Trc(\mathcal C)\) whose fiber over each \(X\in \ob(S)\) is the category \(\Trc(\mathcal C_X)\) of distinguished triangles of \(\mathcal C_X\).

\textbf{1.1.3.} Let \(\mathcal C\) and \(\mathcal C'\) be additive (respectively triangulated) \(S\)-categories, and let \(F:\mathcal C\to\mathcal C'\) be a functor. We say that \(F\) is an additive (respectively exact) \(S\)-functor if \(F\) is a Cartesian \(S\)-functor that is additive (respectively exact) on each fiber.



% Additional macro definitions
\providecommand{\cC}{\mathcal C}
\providecommand{\cD}{\mathcal D}
\providecommand{\cK}{\mathcal K}
\providecommand{\cO}{\mathcal O}
\providecommand{\cS}{\mathcal S}
\providecommand{\Cpx}{\mathrm C}
\providecommand{\Kpx}{\mathrm K}
\providecommand{\Dpx}{\mathrm D}
\providecommand{\D}{\mathrm D}
\providecommand{\Ho}{\mathrm H}
\providecommand{\Ob}{\operatorname{ob}}
\providecommand{\Coker}{\operatorname{Coker}}
\providecommand{\Ker}{\operatorname{Ker}}
\providecommand{\Im}{\operatorname{Im}}
\providecommand{\Hom}{\operatorname{Hom}}
\providecommand{\Cone}{\operatorname{Cone}}
\providecommand{\id}{\operatorname{Id}}
\providecommand{\qis}{\xrightarrow{\sim}}


\textbf{1.1.3 (continued).} If \(\mathcal C\) and \(\mathcal C'\) are triangulated \(S\)-categories and if \(F\) is an exact \(S\)-functor, then \(F\) defines a Cartesian \(S\)-functor
\[
  \Trc(\mathcal C)\longrightarrow \Trc(\mathcal C').
\]

\textbf{1.1.4.} Let \(\mathcal C\) be a triangulated \(S\)-category, and let \(\mathcal C'\) be a full fibered sub-\(S\)-category of \(\mathcal C\). We say that \(\mathcal C'\) is a \emph{triangulated sub-\(S\)-category} of \(\mathcal C\) if, for every \(X\in\Ob S\), \(\mathcal C'_X\) is a triangulated subcategory of \(\mathcal C_X\) \([V]\), Chap. I, §1, 2.3.

\textbf{1.1.5.} Let \(\mathcal C\) be an additive \(S\)-category. One defines (SGA 4 XVII 2.2, 2.4) an additive (respectively triangulated) \(S\)-category
\[
  \Cpx(\mathcal C)\qquad\text{respectively}\qquad \Kpx(\mathcal C),
\]
whose fiber over \(X\in\Ob S\) is the category \(\Cpx(\mathcal C_X)\) of complexes with differential of degree \(+1\) (respectively the category \(\Kpx(\mathcal C_X)\) of complexes of \(\mathcal C_X\) ``up to homotopy''). If \(\mathcal C\) is a flat abelian \(S\)-category, one defines, loc. cit., a triangulated \(S\)-category \(\Dpx(\mathcal C)\), whose fiber over \(X\in\Ob S\) is the derived category \(\Dpx(\mathcal C_X)\); if \(f:X\to Y\) is an arrow of \(S\), we shall write
\[
  f^*:\Dpx(\mathcal C_Y)\longrightarrow \Dpx(\mathcal C_X)
\]
for the inverse-image functor, obtained from
\[
  f^*:\Kpx(\mathcal C_Y)\longrightarrow \Kpx(\mathcal C_X)
\]
by passage to the quotient. One defines full triangulated sub-\(S\)-categories
\[
  \Kpx^*(\mathcal C)\qquad\text{respectively}\qquad \Dpx^*(\mathcal C)
  \quad (*=+,-,b)
\]
by the usual degree conditions.

\subsection*{1.2. Terminology}

Let \(S\) be a site, let \(\mathcal C\) be a flat abelian \(S\)-category (1.1), and let \(\mathcal C_0\) be a strictly full fibered sub-\(S\)-category.

Let \(F\in\Ob \mathcal C_X\), for \(X\in\Ob S\). The set of arrows \(f:U\to X\) in \(S/X\) such that there exists an epimorphism
\[
  E\longrightarrow f^*F,\qquad E\in\Ob \mathcal C_{0,U},
\]
is a sieve on \(X\). If this sieve is covering, we shall say that \(F\) is of \(\mathcal C_0\)-\emph{finite type} (or simply of finite type when no confusion is to be feared).

We shall say that \(\mathcal C_0\) is \emph{locally stable under kernels of epimorphisms} if the following condition is satisfied: for every \(X\in\Ob S\) and every epimorphism
\[
  u:E\longrightarrow F
\]
of \(\mathcal C_X\), with \(E,F\in\Ob\mathcal C_{0,X}\), the kernel of \(u\) is locally in \(\mathcal C_0\); in other words the set of \(f:U\to X\) in \(S/X\) such that
\[
  f^*\Ker(u)\in\Ob\mathcal C_{0,U}
\]
is a refinement of \(X\).

By the standard process of decomposing a long exact sequence into short exact sequences, one sees that the preceding condition is equivalent to the following apparently stronger one: for every \(X\in\Ob S\) and every exact sequence in \(\mathcal C_X\)
\begin{equation}\tag{*}
  0\longrightarrow E^0\longrightarrow E^1\longrightarrow\cdots\longrightarrow E^n\longrightarrow 0
\end{equation}
with \(E^i\in\Ob\mathcal C_{0,X}\) for \(1\leq i\leq n\), the object \(E^0\) is locally in \(\mathcal C_0\).

We shall say that \(\mathcal C_0\) is \emph{locally quasi-stable under kernels of epimorphisms} if, for every exact sequence such as \((*)\), \(E^0\) is of \(\mathcal C_0\)-finite type.

On the other hand, we shall say that \(\mathcal C_0\) is \emph{locally liftable in} \(\mathcal C\) if, for every \(X\in\Ob S\), every epimorphism
\[
  u:E\longrightarrow F
\]
of \(\mathcal C_X\) with \(F\in\Ob\mathcal C_{0,X}\) locally admits a section. Equivalently, for \(\mathcal C_0\) to be locally liftable in \(\mathcal C\), it is necessary and sufficient that, for every diagram of \(\mathcal C_X\), \(X\in\Ob S\),
\[
\begin{tikzcd}
& H \arrow[d]\\
E \arrow[r,"u"'] & F
\end{tikzcd}
\]
where \(u\) is an epimorphism and \(H\in\Ob\mathcal C_0\), the set of \(U\) over \(X\) for which there exists a commutative diagram
\[
\begin{tikzcd}
& H_U \arrow[d]\\
E_U \arrow[r,"u_U"'] \arrow[ur] & F_U
\end{tikzcd}
\]
is a refinement of \(X\).

We shall say that \(\mathcal C_0\) is \emph{locally quasi-liftable in} \(\mathcal C\) if the following weaker condition is satisfied: for every \(X\in\Ob S\), and every epimorphism
\[
  u:E\longrightarrow F
\]
of \(\mathcal C_X\), with \(F\in\Ob\mathcal C_{0,X}\), the set of arrows \(f:U\to X\) of \(S/X\) for which there exists
\[
  v:F'\longrightarrow f^*E,\qquad F'\in\Ob\mathcal C_{0,U},
\]
such that \(f^*(u)v\) is an epimorphism, is a refinement of \(X\). This condition is equivalent to the following: for every diagram in \(\mathcal C_X\)
\[
\begin{tikzcd}
& H \arrow[d]\\
E \arrow[r,"u"'] & F ,
\end{tikzcd}
\]
where \(u\) is an epimorphism and \(H\in\Ob\mathcal C_0\), the set of \(U\) over \(X\) such that there exists a commutative diagram
\[
\begin{tikzcd}
G \arrow[r,"v"] \arrow[d] & H_U \arrow[d]\\
E_U \arrow[r,"u_U"'] & F_U
\end{tikzcd}
\]
with \(G\in\Ob\mathcal C_0\) and \(v\) an epimorphism, is a refinement of \(X\).

If \(S\) is a chaotic site, for example the punctual site (one object, one morphism), we shall omit the adverb ``locally'' in the preceding definitions.

\subsection*{1.3. Examples}

\textbf{a)} Let \(S\) be the punctual site, let \(A\) be a ring, and let \(\mathcal C\) be the category of left \(A\)-modules. The full subcategory \(\mathcal C_0\) of \(\mathcal C\) formed by the projective \(A\)-modules of finite type is stable under kernels of epimorphisms and is liftable in \(\mathcal C\); an \(A\)-module is of \(\mathcal C_0\)-finite type if and only if it is of finite type in the ordinary sense. The full subcategory \(\mathcal C_1\) of \(\mathcal C\) formed by the \(A\)-modules of finite type is quasi-liftable in \(\mathcal C\), but is not liftable in general; if \(A\) is noetherian, \(\mathcal C_1\) is stable under kernels of epimorphisms. The full subcategory \(\mathcal C'_0\) of \(\mathcal C\) formed by the free \(A\)-modules of finite type is liftable in \(\mathcal C\) and quasi-stable, though not stable in general, under kernels of epimorphisms.

\textbf{b)} Let \((S,\mathcal O_S)\) be a space ringed by local rings, let \(S\) be the site of open subsets of \(S\), and let \(\mathcal C\) be the fibered \(S\)-category, even split (SGA 1 VI 9), defined by the functor
\[
  U\longmapsto \{\text{category of left }\mathcal O_U\text{-modules}\}.
\]
Then \(\mathcal C\) is a flat abelian \(S\)-category (1.1). Let \(\mathcal C_0\) be the strictly full additive sub-\(S\)-category of \(\mathcal C\) defined by
\[
  U\longmapsto \{\text{category of finite-type free }\mathcal O_U\text{-modules}\}.
\]
An \(\mathcal O_S\)-module \(M\) is locally of \(\mathcal C_0\)-finite type if and only if \(M\) is of finite type in the ordinary sense (EGA \(0_{\mathrm I}\), 5.2). From EGA \(0_{\mathrm I}\), 5.2.7 and 5.4.9, it follows that \(\mathcal C_0\) is locally stable under kernels of epimorphisms and locally liftable in \(\mathcal C\).

\textbf{c)} Let \(G\) be a group, let \(S\) be the topos of \(G\)-sets (SGA 4 V 3.11), and let \(k\) be a field, regarded as the constant ring on \(S\). Let \(\mathcal C\) be the split \(S\)-category defined by
\[
  U\longmapsto \{\text{category of left }k_U\text{-modules}\};
\]
then \(\mathcal C\) is a flat abelian \(S\)-category. Let \(\mathcal C_0\) be the strictly full additive sub-\(S\)-category of \(\mathcal C\) defined by
\[
  U\longmapsto \{\text{category of finite-type free }k_U\text{-modules}\}.
\]
Since localization in \(S\) is equivalent to forgetting the operations of \(G\), it is clear that a \(G\)-\(k\)-module \(M\) is locally of \(\mathcal C_0\)-finite type if and only if \(M\) is finite-dimensional as a vector space over \(k\), and that \(\mathcal C_0\) is locally stable under kernels of epimorphisms and locally liftable in \(\mathcal C\).

The reader is left to look at the ``Galois'' variant of c), where \(G\) is a profinite group and \(S\) is the topos of sets on which \(G\) operates continuously.

In fact, examples b) and c) are particular cases of the following example.

\textbf{d)} Let \((S,\mathcal O_S)\) be a ringed topos (cf. SGA 4 IV 2.10), and let \(\mathcal C\) be the split \(S\)-category defined by
\[
  U\longmapsto \{\text{category of left }\mathcal O_U\text{-modules}\};
\]
then \(\mathcal C\) is a flat abelian \(S\)-category. Let \(\mathcal C_0\) be the strictly full additive sub-\(S\)-category of \(\mathcal C\) defined by
\[
  U\longmapsto \{\text{category of finite-type free }\mathcal O_U\text{-modules}\}.
\]
An \(\mathcal O_S\)-module locally of \(\mathcal C_0\)-finite type will simply be called \emph{of finite type}. The category \(\mathcal C_0\) is locally quasi-stable under kernels of epimorphisms and locally liftable in \(\mathcal C\), as follows from:

\textbf{Lemma 1.3.1.} Let \(P\) be an \(\mathcal O_S\)-module. The following conditions are equivalent:
\begin{enumerate}[label=(\roman*)]
\item \(P\) is of finite type and the functor \(\Hom(P,\cdot)\) is exact;
\item \(P\) is of finite type and every epimorphism \(M\to P\) locally admits a section;
\item \(P\) is locally a direct factor of a free module of finite type.
\end{enumerate}
The proof is immediate and is left as an exercise to the reader.

\subsection*{1.4. An ``extension'' lemma (cf. EGA \(0_{\mathrm{III}}\), 11.9.1)}

\textbf{Proposition 1.4.1.} Let \(S\) be a site, let \(\mathcal C\) be a flat abelian \(S\)-category (1.1), and let \(\mathcal C_0\) be a strictly full additive sub-\(S\)-category of \(\mathcal C\). Assume that \(\mathcal C_0\) is locally quasi-liftable in \(\mathcal C\) (1.2). Let \(S\in\Ob S\), let
\[
  u:E\longrightarrow F
\]
be a morphism, of degree \(0\), of complexes of \(\mathcal C_S\) (differentials of degree \(+1\)), let \(G\) be the cone of \(u\) (cf. [V], Chap. I, §1, 2.1), let \(n\in\mathbb Z\), and assume that \(\Ho^n(G)\) is of \(\mathcal C_0\)-finite type (1.2). Then the set of \(X\in\Ob(S/S)\) satisfying property \((P)\) below is a refinement of \(S\).

\((P)\): there exist an object \(M\in\mathcal C_{0,X}\) and arrows of \(\mathcal C_X\)
\[
  r:M\longrightarrow Z^{n+1}(E_X)=
  \Ker\!\left(E_X^{n+1}\xrightarrow{d^{n+1}}E_X^{n+2}\right),
  \qquad
  s:M\longrightarrow F_X^n
\]
such that the following diagram of \(\mathcal C_X\)
\[
\begin{tikzcd}[column sep=large]
\cdots \arrow[r] & E_X^{n-1}
 \arrow[r,"{\binom{d_E^{n-1}}{0}}"]
 \arrow[d,"u_X^{n-1}"'] &
E_X^n\oplus M
 \arrow[r,"{(d_E^n,r)}"]
 \arrow[d,"{(u_X^n,s)}"'] &
E_X^{n+1} \arrow[r] \arrow[d,"u_X^{n+1}"] &
\cdots\\
\cdots \arrow[r] & F_X^{n-1} \arrow[r,"d_F^{n-1}"'] &
F_X^n \arrow[r,"d_F^n"'] &
F_X^{n+1} \arrow[r] &
\cdots
\end{tikzcd}
\]
is a morphism of complexes \(u':E'\to F_X\), that is,
\[
  d_F^n s=u_X^{n+1}r,
\]
and such that
\begin{enumerate}[label=(\roman*)]
\item \(\Ho^{n+1}(u')\) is a monomorphism, with
\[
  \Im\Ho^{n+1}(u')=\Im\Ho^{n+1}(u_X);
\]
\item \(\Ho^n(u')\) is an epimorphism.
\end{enumerate}

In pictorial terms, one may say that, after modifying \(E\) and \(u\) locally in degree \(n\) only by adding, in degree \(n\), an object of \(\mathcal C_0\), one can locally make \(\Ho^{n+1}(u)\) injective, without changing its image, and \(\Ho^n(u)\) surjective.

\emph{Proof.} The procedure is as follows: first one shows that \(r\) and \(s\) may be chosen locally so as to make \(\Ho^{n+1}(u)\) injective; one verifies that the finiteness hypothesis on \(\Ho^n(G)\) is preserved; a new local correction in degree \(n\) then makes it possible to obtain \(\Ho^n(u)\) surjective.

The exact cohomology sequence
\[
  \cdots\longrightarrow \Ho^n(G)\longrightarrow
  \Ho^{n+1}(E)\longrightarrow
  \Ho^{n+1}(F)\longrightarrow \cdots
\]
shows that the kernel of \(\Ho^{n+1}(u)\) is of \(\mathcal C_0\)-finite type. In other words, the set of \(X\) over \(S\) for which there is an exact sequence
\[
  M\longrightarrow \Ho^{n+1}(E_X)
  \xrightarrow{\Ho^{n+1}(u)}
  \Ho^{n+1}(F_X),
  \qquad M\in\Ob\mathcal C_{0,X},
\]
is a refinement of \(S\). Since \(\mathcal C_0\) is locally quasi-liftable in \(\mathcal C\), one may moreover impose that the arrow
\[
  M\longrightarrow \Ho^{n+1}(E_X)
\]
factor through \(Z^{n+1}(E_X)\):
\[
  M\xrightarrow{r} Z^{n+1}(E_X)\longrightarrow \Ho^{n+1}(E_X).
\]
The composite
\[
  M\xrightarrow{u_X^{n+1}r} Z^{n+1}(F_X)\longrightarrow \Ho^{n+1}(F_X)
\]
is zero by construction, so \(u_X^{n+1}r\) factors through
\[
  B^{n+1}(F_X)=\Im\!\left(F_X^n\xrightarrow{d_F^n}F_X^{n+1}\right).
\]
Applying quasi-liftability once more, one may impose that \(u_X^{n+1}r\) factor through \(F_X^n\):
\[
  M\xrightarrow{s}F_X^n\longrightarrow B^{n+1}(F_X)\longrightarrow Z^{n+1}(F_X).
\]
Then \(u'\), defined as in (1.4.1), is plainly a morphism of complexes satisfying (i).

Let \(G'\) be the cone of \(u':E'\to F_X\). We show that \(\Ho^n(G')\) is of \(\mathcal C_0\)-finite type. By construction there is an exact sequence
\[
  0\longrightarrow E_X\xrightarrow{v}E'\longrightarrow M[-n]\longrightarrow 0,
\]
with \(u'v=u_X\). To compare \(G_X\) and \(G'\), one may, for instance, form the octahedron in \(\Kpx(\mathcal C_X)\):
\[
\begin{tikzcd}[column sep=small,row sep=small]
& G_X \arrow[rr] \arrow[dr] && M[-n+1] \arrow[dl]\\
E_X \arrow[ur,"+1"] \arrow[rr,"v"'] && E' \arrow[ur] \arrow[rr,"u'"] && F_X .
\end{tikzcd}
\]
The exact cohomology sequence of the distinguished triangle marked by the octahedron shows that
\[
  \Ho^n(G_X)\longrightarrow \Ho^n(G')
\]
is an epimorphism; hence \(\Ho^n(G')\) is indeed of \(\mathcal C_0\)-finite type.

The preceding correction therefore reduces us to the case where \(\Ho^{n+1}(u)\) is a monomorphism. We suppose this from now on. The exact cohomology sequence of \(E\to F\to G\) then gives an exact sequence
\[
  \Ho^n(E)\xrightarrow{\Ho^n(u)}
  \Ho^n(F)\longrightarrow \Ho^n(G)\longrightarrow 0.
\]
Since \(\Ho^n(G)\) is of \(\mathcal C_0\)-finite type, the set of \(X\) over \(S\) for which there exists an epimorphism
\[
  M\longrightarrow \Ho^n(G_X),\qquad M\in\Ob\mathcal C_{0,X},
\]
is, by definition, a refinement of \(S\). Using the quasi-liftability of \(\mathcal C_0\) in \(\mathcal C\), one may impose that this epimorphism factor through \(Z^n(F_X)\):
\[
  M\xrightarrow{s}Z^n(F_X)\longrightarrow \Ho^n(F_X)\longrightarrow \Ho^n(G_X).
\]
It is clear that \(u'\), defined as in (1.4.1) with \(r=0\), is a morphism of complexes and that
\[
  \Ho^n(u'):\Ho^n(E_X)\oplus M\longrightarrow \Ho^n(F_X)
\]
is an epimorphism. This completes the proof of (1.4.1).

\textbf{Lemma 1.4.2.} Let \(\mathcal C\) be an abelian category, let \(u:E\to F\) be a morphism of complexes of \(\mathcal C\) (respectively an arrow of \(\Dpx(\mathcal C)\)), and let \(G\) be the cone of \(u\). For a given \(n\in\mathbb Z\), the following conditions are equivalent:
\begin{enumerate}[label=(\roman*)]
\item \(\Ho^i(G)=0\) for \(i\geq n\);
\item \(\Ho^i(u)\) is an isomorphism for \(i>n\) and an epimorphism for \(i=n\).
\end{enumerate}
The proof is immediate from the exact cohomology sequence.

\textbf{Definition 1.4.3.} We shall say that \(u\) is an \(n\)-\emph{quasi-isomorphism} (respectively an \(n\)-\emph{isomorphism}) if \(u\) satisfies the equivalent conditions of (1.4.2).

That being so, one deduces at once from (1.4.1) the following:

\textbf{Corollary 1.4.4.} Let \(S,S,\mathcal C,\mathcal C_0\) be as in (1.4.1). Let \(u:E\to F\) be an \((n+1)\)-quasi-isomorphism of complexes of \(\mathcal C_S\), with \(n\in\mathbb Z\) fixed, such that \(\Ho^n(G)\) is of \(\mathcal C_0\)-finite type. Then the set of \(X\) over \(S\) for which there exist \((M,r,s)\) as in (1.4.1) such that \(u'\) is an \(n\)-quasi-isomorphism is a refinement of \(S\).

In pictorial terms, one may, using the finiteness condition on \(\Ho^n(G)\), locally ``gain one notch,'' that is, locally modify \(u\) in degree \(n\) by adding to \(E^n\) an object of \(\mathcal C_0\) in such a way as to make \(u\) an \(n\)-quasi-isomorphism.

\textbf{Corollary 1.4.5.} Let \(\mathcal C\) be an abelian category, let \(n\in\mathbb Z\), and let \(u:E\to F\) be an \(n\)-quasi-isomorphism of complexes of \(\mathcal C\). There exists a commutative diagram in \(\Cpx(\mathcal C)\)
\[
\begin{tikzcd}
E \arrow[r,"v"] \arrow[d,"u"'] & E' \arrow[dl,"u'"]\\
F &
\end{tikzcd}
\]
where \(u'\) is a quasi-isomorphism and \(v\) is a monomorphism inducing the identity in degrees \(\geq n\), that is,
\[
  E'^i=E^i,\qquad v^i=\id(E^i)\qquad\text{for }i\geq n.
\]

\emph{Proof.} Applying (1.4.4) in the case where \(S\) is the punctual site and \(\mathcal C_0=\mathcal C\), one obtains a commutative triangle in \(\Cpx(\mathcal C)\)
\[
\begin{tikzcd}
E \arrow[r,"v_{n-1}"] \arrow[d,"u"'] & E'_{\,n-1} \arrow[dl,"u'_{n-1}"]\\
F &
\end{tikzcd}
\]
where \(u'_{n-1}\) is an \((n-1)\)-quasi-isomorphism and \(v_{n-1}\) is a monomorphism inducing the identity in degrees \(\geq n\). Iterating the process, one finds, for \(i\leq n-1\), a sequence of commutative triangles in \(\Cpx(\mathcal C)\)
\[
\begin{tikzcd}
E'_i \arrow[r,"v_{i-1}"] \arrow[d,"u'_i"'] & E'_{\,i-1} \arrow[dl,"u'_{i-1}"]\\
F &
\end{tikzcd}
\]
where \(u'_i\) is an \(i\)-quasi-isomorphism and \(v_{i-1}\) is a monomorphism inducing the identity in degrees \(\geq i\). The corollary follows by passing to the limit. More precisely, writing \(d'_i\) for the differential of \(E'_i\), define \(E'\) by
\[
  E'^i=E^i,\quad d'^i=d^i\quad (i\geq n),
\]
and
\[
  E'^i=(E'_i)^i,\quad d'^i=(d'_i)^i\quad (i<n),
\]
and then define \(v\) and \(u'\) by
\[
  v^j=(v_i)^j,\qquad u'^j=(u'_i)^j\qquad (i\leq j).
\]
One immediately verifies that \(E'\), \(v\), and \(u'\), so defined, answer the question.

In abbreviated form: one can extend an \(n\)-quasi-isomorphism into a quasi-isomorphism ``without changing anything given in degree \(>n\).''

\section*{2. Pseudo-coherent complexes}

\textbf{2.0.} In this section, \(S\) denotes a site, \(S\) an object of \(S\), \(\mathcal C\) a flat abelian \(S\)-category (1.1), and \(\mathcal C_0\) a strictly full additive sub-\(S\)-category of \(\mathcal C\). We assume that \(\mathcal C_0\) is locally quasi-stable under kernels of epimorphisms (1.2) and locally quasi-liftable in \(\mathcal C\) (1.2).

\textbf{Definition 2.1.} Let \(X\in\Ob S\), let \(E\) be a complex of \(\mathcal C_X\), and let \(n\in\mathbb Z\). We shall say that \(E\) is \emph{strictly \(\mathcal C_0\)-\(n\)-pseudo-coherent} (respectively \(\mathcal C_0\)-pseudo-coherent, respectively \(\mathcal C_0\)-perfect) if \(E^i=0\) for \(i\) sufficiently large and \(E^i\in\Ob\mathcal C_{0,S}\) for \(i\geq n\) (respectively for every \(i\), respectively for every \(i\) and \(E^i=0\) for almost all \(i\)).

When no confusion is to be feared, we shall sometimes say ``strictly \(n\)-pseudo-coherent'' (respectively, and so on) instead of ``strictly \(\mathcal C_0\)-\(n\)-pseudo-coherent''.

\textbf{Lemma 2.1.1.} Let \(E\in\Ob\Cpx(\mathcal C_S)\) and \(n\in\mathbb Z\). If \(E\) is strictly \(n\)-pseudo-coherent and if \(\Ho^i(E)=0\) for \(i>n\), then \(\Ho^n(E)\) is of finite type.

\emph{Proof.} The sequence
\[
  0\longrightarrow Z^n(E)\longrightarrow E^n\longrightarrow E^{n+1}\longrightarrow\cdots\longrightarrow 0
\]
is exact, and \(E^i\in\Ob\mathcal C_0\) for \(i\geq n\). Since \(\mathcal C_0\) is locally quasi-stable under kernels of epimorphisms, it follows that \(Z^n(E)\) is of finite type, hence so is \(\Ho^n(E)\), which is a quotient of it.

\textbf{Proposition 2.2.} Let \(F\in\Ob\Dpx(\mathcal C_S)\) and \(n\in\mathbb Z\). Then:

\textbf{a)} If \(F\) is strictly \(n\)-pseudo-coherent (relative to \(\mathcal C_0\)) and if
\[
  u:E\longrightarrow F
\]
is a quasi-isomorphism, then the set of \(X\) over \(S\) such that there exists a quasi-isomorphism
\[
  v:F'\longrightarrow E_X
\]
with \(F'\) strictly \(n\)-pseudo-coherent is a refinement of \(S\).

\textbf{b)} The following conditions are equivalent:
\begin{enumerate}[label=(\roman*)]
\item the set of \(X\) over \(S\) such that there exists a quasi-isomorphism
\[
  u:E\longrightarrow F_X
\]
of \(\Cpx(\mathcal C_X)\), with \(E\) strictly \(n\)-pseudo-coherent, is a refinement of \(S\);
\item[\((\mathrm{i}')\)] the same condition as (i), but with \(u\) an isomorphism of \(\Dpx(\mathcal C_X)\);
\item the set of \(X\) over \(S\) such that there exists an \(n\)-quasi-isomorphism
\[
  u:E\longrightarrow F_X
\]
of \(\Cpx(\mathcal C_X)\), with \(E\) strictly perfect, is a refinement of \(S\);
\item[\((\mathrm{ii}')\)] the same condition as (ii), but with \(u\) an \(n\)-isomorphism of \(\Dpx(\mathcal C_X)\).
\end{enumerate}

\emph{Proof.} For a), there exists, for \(i\) sufficiently large, an \(i\)-quasi-isomorphism
\[
  v_i:F'_i\longrightarrow E
\]
with \(F'_i\) strictly \(n\)-pseudo-coherent (take \(F'_i=0\)). We show that, if \(i>n\), there exists, after refining \(S\), an \((i-1)\)-quasi-isomorphism
\[
  v_{i-1}:F'_{i-1}\longrightarrow E
\]
with \(F'_{i-1}\) strictly \(n\)-pseudo-coherent. By (1.4.4), it is enough to show that \(\Ho^{i-1}(M')\), where \(M'\) is the cone of \(v_i\), is of \(\mathcal C_0\)-finite type. Now if \(M''\) denotes the cone of \(uv_i\), the hypothesis that \(u\) is a quasi-isomorphism and (1.4.2) imply that
\[
  \Ho^j(M')\qis \Ho^j(M'')\qquad\text{for all }j,
\]
and
\[
  \Ho^j(M')=0\qquad\text{for }j\geq i.
\]
The cone formula
\[
  M''{}^j=F'_i{}^{\,j+1}\oplus F^j
\]
shows, on the other hand, that \(M''\) is strictly \((i-1)\)-pseudo-coherent. Hence, by (2.1.1),
\[
  \Ho^{i-1}(M'')=\Ho^{i-1}(M')
\]
is of finite type, whence the assertion. It follows, by iteration and after refining \(S\), that there is an \(n\)-quasi-isomorphism
\[
  v_n:F'_n\longrightarrow E
\]
where \(F'_n\) is strictly \(n\)-pseudo-coherent. But, by (1.4.5), one can extend \(v_n\) into a quasi-isomorphism
\[
  v:F'\longrightarrow E
\]
without changing anything in degrees \(\geq n\); in particular \(F'\) is strictly \(n\)-pseudo-coherent. This proves a).

For b), the equivalence of (i) and \((\mathrm{i}')\) follows at once from a) and from the definition of isomorphisms of \(\Dpx^-(\mathcal C_X)\). If \(E\in\Ob\Dpx^-(\mathcal C_S)\) is strictly \(n\)-pseudo-coherent, there is an exact sequence
\[
  0\longrightarrow E'\longrightarrow E\longrightarrow E''\longrightarrow 0
\]
where \(E'\) is strictly perfect and \(E''{}^i=0\) for \(i\geq n\). Consequently \(E'\to E\) is an \(n\)-quasi-isomorphism, which gives the implication \((\mathrm{i})\Rightarrow(\mathrm{ii})\). The implication \((\mathrm{ii})\Rightarrow(\mathrm{ii}')\) is trivial, so it remains to prove \((\mathrm{ii}')\Rightarrow(\mathrm{i})\). Let
\[
  u:E\longrightarrow F
\]
be an \(n\)-isomorphism of \(\Dpx^-(\mathcal C_S)\), with \(E\) strictly perfect. Then \(u\) is represented by a diagram in \(\Cpx^-(\mathcal C_S)\)
\[
\begin{tikzcd}
& E' \arrow[dl,"s"'] \arrow[dr,"u'"]\\
E && F ,
\end{tikzcd}
\]
where \(s\) is a quasi-isomorphism and \(u'\) an \(n\)-quasi-isomorphism. By a), after refining \(S\), there exists a quasi-isomorphism
\[
  t:E''\longrightarrow E',
\]
with \(E''\) strictly \(n\)-pseudo-coherent. Applying (1.4.5) to \(u't:E''\to F\), we extend it to a quasi-isomorphism
\[
  E'''\longrightarrow F
\]
with \(E'''\) strictly \(n\)-pseudo-coherent. This proves \((\mathrm{ii}')\Rightarrow(\mathrm{i})\) and completes the proof of (2.2).

\textbf{Definition 2.3.} Under the hypotheses of (2.2), we shall say that \(F\in\Ob\Dpx^-(\mathcal C_S)\) is \(n\)-\(\mathcal C_0\)-pseudo-coherent if \(F\) satisfies the equivalent conditions b) of (2.2). We shall say that \(F\) is \(\mathcal C_0\)-pseudo-coherent if \(F\) is \(n\)-\(\mathcal C_0\)-pseudo-coherent for every \(n\).

When this cannot cause confusion, we shall sometimes allow ourselves to say \(n\)-pseudo-coherent, pseudo-coherent, instead of \(n\)-\(\mathcal C_0\)-pseudo-coherent, \(\mathcal C_0\)-pseudo-coherent.

\textbf{Remark 2.3.1.} Let \(F\in\Ob\Dpx(\mathcal C_S)\) be a pseudo-coherent complex. The set of \(X\) over \(S\) such that there exists a quasi-isomorphism
\[
  E\longrightarrow F_X
\]
with \(E\) strictly pseudo-coherent (2.1) is not in general a refinement of \(S\). However, we shall see later (2.7 a) and Exposé II) important cases in which this is so.

The proof of (2.2 a) gives the following result.

\textbf{Proposition 2.4.} Let \(n,i\in\mathbb Z\), and let
\[
  u:E\longrightarrow F
\]
be an \(i\)-quasi-isomorphism of \(\Cpx(\mathcal C_S)\), with \(F\) \(n\)-pseudo-coherent and \(E\) strictly \(n\)-pseudo-coherent. Then, after refining \(S\), there exists a commutative triangle in \(\Cpx(\mathcal C_S)\)
\[
\begin{tikzcd}
E \arrow[r,"v"] \arrow[d,"u"'] & E' \arrow[dl,"u'"]\\
F &
\end{tikzcd}
\]
where \(E'\) is strictly \(n\)-pseudo-coherent, \(u'\) is a quasi-isomorphism, and \(v\) is a monomorphism inducing the identity in degree \(\geq i\).

\textbf{Proposition 2.5.} \textbf{a)} If \(F\in\Ob\Dpx(\mathcal C_S)\) is \(n\)-pseudo-coherent (respectively pseudo-coherent), then \(F[m]\) is \((n-m)\)-pseudo-coherent (respectively pseudo-coherent).

\textbf{b)} Let
\[
\begin{tikzcd}
E \arrow[rr] && F \arrow[dl]\\
& G \arrow[ul,"+1"] &
\end{tikzcd}
\tag{*}
\]
be a distinguished triangle of \(\Dpx(\mathcal C_S)\). If \(E\) and \(G\) are \(n\)-pseudo-coherent (respectively pseudo-coherent), then \(F\) is \(n\)-pseudo-coherent (respectively pseudo-coherent).

\emph{Proof.} a) is trivial from the definitions.

b) It is enough, of course, to prove the non-parenthesized assertion. The triangle \((*)\) can also be written
\[
\begin{tikzcd}
G[-1] \arrow[rr] && E \arrow[dl]\\
& F \arrow[ul,"+1"] &
\end{tikzcd}
\]
By a), \(G[-1]\) is \((n+1)\)-pseudo-coherent. Thus we are reduced to proving that, if \(E\) is \((n+1)\)-pseudo-coherent and \(F\) is \(n\)-pseudo-coherent, then \(G\) is \(n\)-pseudo-coherent. After localizing, we may assume \(E\) (respectively \(F\)) strictly \((n+1)\)-pseudo-coherent (respectively \(n\)-pseudo-coherent). The arrow
\[
  E\xrightarrow{u}F
\]
of \(\Dpx^-(\mathcal C_S)\) is represented by a pair of arrows of \(\Cpx^-(\mathcal C_S)\)
\[
\begin{tikzcd}
& E' \arrow[dl,"s"'] \arrow[dr,"u'"]\\
E && F ,
\end{tikzcd}
\]
where \(s\) is a quasi-isomorphism. By (2.2 a)), after localizing, there exists a quasi-isomorphism
\[
  t:E''\longrightarrow E',
\]
where \(E''\) is strictly \((n+1)\)-pseudo-coherent. We are therefore reduced to the case where \(u\) is defined by a true arrow of \(\Cpx^-(\mathcal C_S)\). The assertion then follows from the definitions and from the formula
\[
  C(u)^i=E^{i+1}\oplus F^i,
\]
where \(C(u)\) denotes the cone of \(u\).

\textbf{Scholium 2.6.} By rotation of triangles, one sees that pseudo-coherence properties for two of the objects of the triangle \((*)\) imply another such property for the third. The various cases may be summarized in the table below; the lower-case letters denote degrees of pseudo-coherence, and in each row the conclusion is underlined:
\[
\begin{array}{ccc}
E&F&G\\[0.2em]
n+1&n&\underline n\\
n&\underline n&n\\
\underline n&n&n-1 .
\end{array}
\]

We shall write
\[
  \Dpx(\mathcal C)_{\mathcal C_0\text{-}n\mathrm{coh}}
  \quad\text{respectively}\quad
  \Dpx(\mathcal C)_{\mathcal C_0\text{-}\mathrm{coh}}
\]
for the strictly full fibered subcategory of \(\Dpx(\mathcal C)\) whose fiber over each \(X\in\Ob S\) is formed by the objects of \(\Dpx(\mathcal C_X)\) that are \(n\)-pseudo-coherent (respectively pseudo-coherent) relative to \(\mathcal C_0\). By (2.5), the category
\[
  \Dpx(\mathcal C)_{\mathcal C_0\text{-}\mathrm{coh}}
\]
is a triangulated sub-\(S\)-category of \(\Dpx(\mathcal C)\) (1.1).

\textbf{Proposition 2.7.} Suppose that \(S\) is the punctual site, so that
\[
  \mathcal C=\mathcal C_S,\qquad \mathcal C_0=\mathcal C_{0,S}.
\]
Then:
\begin{enumerate}[label=\alph*)]
\item If \(F\in\Ob\Dpx(\mathcal C)\) is pseudo-coherent, there exists a quasi-isomorphism
\[
  E\longrightarrow F
\]
with \(E\) strictly pseudo-coherent (2.1).
\item Let \(\Kpx^{-,\mathrm{ac}}(\mathcal C_0)\) denote the full triangulated subcategory of \(\Kpx^-(\mathcal C_0)\) formed by acyclic complexes. The canonical functor
\[
  \Kpx^-(\mathcal C_0)/\Kpx^{-,\mathrm{ac}}(\mathcal C_0)
  \longrightarrow \Dpx(\mathcal C)
\]
is fully faithful and has essential image \(\Dpx(\mathcal C)_{\mathcal C_0\text{-}\mathrm{coh}}\).
\end{enumerate}

\emph{Proof.} a) Let \(F\in\Ob\Dpx(\mathcal C)_{\mathcal C_0\text{-}\mathrm{coh}}\), and suppose constructed an \(i\)-quasi-isomorphism
\[
  u_i:E_i\longrightarrow F
\]
of \(\Cpx(\mathcal C)\), with \(E_i\) strictly pseudo-coherent; this is possible for \(i\) sufficiently large, taking \(E_i=0\). Let \(G_i\) be the cone of \(u_i\). It follows from (1.4.2) and (2.1.1) that \(\Ho^{i-1}(G_i)\) is of \(\mathcal C_0\)-finite type. Hence, by (1.4.4), there exists a commutative triangle in \(\Cpx(\mathcal C)\)
\[
\begin{tikzcd}
E_i \arrow[r,"v_{i-1}"] \arrow[d,"u_i"'] & E_{i-1} \arrow[dl,"u_{i-1}"]\\
F &
\end{tikzcd}
\]
where \(E_{i-1}\) is strictly pseudo-coherent, \(u_{i-1}\) an \((i-1)\)-quasi-isomorphism, and \(v_{i-1}\) a monomorphism inducing the identity in degrees \(\geq i\). Passing to the limit (cf. the proof of (1.4.5)) gives a quasi-isomorphism
\[
  u:E\longrightarrow F
\]
with \(E\) strictly pseudo-coherent.

b) Apply the criterion of \([V]\), Chap. I, §2, 4.2 b)(i), together with a).

We note that the proof of a) gives the following more precise result, which generalizes (1.4.5).

\textbf{Proposition 2.7.1.} Under the hypotheses of (2.7), let \(u:E\to F\) be an \(n\)-quasi-isomorphism of \(\Cpx(\mathcal C)\), with \(F\) pseudo-coherent and \(E\) strictly pseudo-coherent. Then there exists a commutative triangle in \(\Cpx(\mathcal C)\)
\[
\begin{tikzcd}
E \arrow[r,"v"] \arrow[d,"u"'] & E' \arrow[dl,"u'"]\\
F &
\end{tikzcd}
\]
where \(E'\) is strictly pseudo-coherent, \(u'\) is a quasi-isomorphism, and \(v\) is a monomorphism inducing the identity in degree \(\geq n\).

\textbf{Definition 2.8.} Let \(n\in\mathbb N\) and \(F\in\Ob\mathcal C_S\). We shall say that \(F\) is of \(n\)-\(\mathcal C_0\)-finite presentation if the set of \(X\) over \(S\) such that there exists an exact sequence
\[
  E^{-n}\longrightarrow\cdots\longrightarrow E^0\longrightarrow F_X\longrightarrow 0,
\qquad E^i\in\Ob\mathcal C_{0,X}\quad(-n\leq i\leq 0),
\]
is a refinement of \(S\). We shall say that \(F\) is of \(\infty\)-\(\mathcal C_0\)-finite presentation if \(F\) is of \(n\)-\(\mathcal C_0\)-finite presentation for every \(n\in\mathbb N\).

Of course, we shall sometimes simply say \(n\)-finite presentation when there is no possible ambiguity about \(\mathcal C_0\). We shall say ``finite presentation'' for ``1-finite presentation,'' and ``finite type'' for ``0-finite presentation,'' the latter notion coinciding with the one already introduced in (1.2).

\textbf{Proposition 2.9.} Let \(n\in\mathbb N\) and \(F\in\Ob\mathcal C_S\). The following conditions are equivalent:
\begin{enumerate}[label=(\roman*)]
\item \(F\) is of \(n\)-finite presentation (respectively of \(\infty\)-finite presentation);
\item \(F\) is \((-n)\)-pseudo-coherent (respectively pseudo-coherent).
\end{enumerate}
In the case where \(S\) is the punctual site, \(F\) is pseudo-coherent if and only if there exists an exact sequence
\[
  \cdots\longrightarrow E^i\longrightarrow\cdots\longrightarrow E^0\longrightarrow F\longrightarrow 0,
\qquad E^i\in\Ob\mathcal C_0\quad\text{for all }i.
\]

\emph{Proof.} For the equivalence \((\mathrm{i})\Leftrightarrow(\mathrm{ii})\), it is plainly enough to treat the non-parenthesized case. The implication \((\mathrm{i})\Rightarrow(\mathrm{ii})\) is immediate from the definitions, and \((\mathrm{ii})\Rightarrow(\mathrm{i})\) follows at once from (2.4): the arrow \(0\to F\) is a \(1\)-quasi-isomorphism, and from this one deduces, after localizing, a quasi-isomorphism
\[
  E\longrightarrow F
\]
with \(E\) strictly \((-n)\)-pseudo-coherent and \(E^i=0\) for \(i\geq 1\), that is, an \(n\)-finite presentation of \(F\). The second assertion is proved in the same way, using (2.7.1).

\textbf{Proposition 2.10.} \textbf{a)} Let \(m,n\in\mathbb Z\), with \(n\leq m\). If \(F\in\Ob\Dpx^-(\mathcal C_S)\) is \(n\)-pseudo-coherent and acyclic in degrees \(\geq m\), then there exists locally a quasi-isomorphism
\[
  E\longrightarrow F
\]
with \(E\) strictly \(n\)-pseudo-coherent and \(E^i=0\) for \(i\geq m\). If \(S\) is punctual, and if \(F\) is pseudo-coherent and acyclic in degrees \(\geq m\), there exists a quasi-isomorphism
\[
  E\longrightarrow F
\]
with \(E\) strictly pseudo-coherent and \(E^i=0\) for \(i\geq m\).

\textbf{b)} Let \(F\in\Ob\Dpx^-(\mathcal C_S)\) be acyclic in degrees \(>n\). For \(F\) to be \(n\)-pseudo-coherent, it is necessary and sufficient that \(\Ho^n(F)\) be of finite type.

\emph{Proof.} a) Since \(F\) is acyclic in degrees \(\geq m\), the arrow \(0\to F\) is an \(m\)-quasi-isomorphism. The assertion follows from (2.4) and (2.7.1).

b) There exists a distinguished triangle in \(\Dpx^-(\mathcal C_S)\):
\[
\begin{tikzcd}
F \arrow[rr] && F' \arrow[dl]\\
& \Ho^n(F)[-n] \arrow[ul,"+1"] &
\end{tikzcd}
\]
(to see this, begin by truncating \(F\), replacing \(F^n\) by \(Z^n(F)\) and \(F^i\) by \(0\) for \(i>n\)). The exact cohomology sequence implies
\[
  \Ho^i(F')=0\quad\text{for }i\geq n-1,
\]
so \(F'\) is \((n-1)\)-pseudo-coherent. We conclude by (2.5) and (2.6).



% Additional notation for Expose I continuation
\providecommand{\RHom}{\mathrm R\!\Hom}
\providecommand{\Ext}{\operatorname{Ext}}
\providecommand{\coh}{\operatorname{coh}}
\providecommand{\PreCoh}{\operatorname{precoh}}
\providecommand{\Tr}{\operatorname{Tr}}
\providecommand{\Qis}{\operatorname{Qis}}
\providecommand{\Rad}{\operatorname{rad}}
\providecommand{\Spec}{\operatorname{Spec}}
\providecommand{\End}{\operatorname{End}}
\providecommand{\Id}{\operatorname{Id}}
\providecommand{\limdir}{\mathop{\mathrm{lim}}\limits_{\longrightarrow}}
\providecommand{\liminv}{\mathop{\mathrm{lim}}\limits_{\longleftarrow}}
\providecommand{\qcoh}{\operatorname{qcoh}}
\providecommand{\amp}{\operatorname{amp}}


\section*{Exposé I, continued: pseudo-coherent complexes and perfect complexes}

\textbf{Proposition 2.11.} Let \(n\in\mathbb Z\).

\textbf{a)} Let \(F\in\Ob\Cpx^-(\mathcal C_S)\). If, for every \(i\geq n\) (respectively for every \(i\)), \(F^i\) is of finite \((i-n)\)-presentation (respectively of finite \(\infty\)-presentation), then \(F\) is \(n\)-pseudo-coherent (respectively pseudo-coherent).

\textbf{b)} Let \(F\in\Ob\Dpx^-(\mathcal C_S)\). If, for every \(i\geq n\) (respectively for every \(i\)), \(\Ho^i(F)\) is of finite \((i-n)\)-presentation (respectively of finite \(\infty\)-presentation), then \(F\) is pseudo-coherent.

\emph{Proof.} It is enough to treat the non-parenthesized assertions.

\textbf{a)} We argue by induction on \(N\) such that \(F^i\neq0\) only for \(i\leq n+N\). If \(N=0\), the assertion follows at once from (2.10 b). Suppose \(N\geq1\). We have an exact sequence
\[
  0\longrightarrow F^{n+N}[-n-N]\longrightarrow F\longrightarrow F'\longrightarrow 0 .
\]
By the induction hypothesis, \(F'\) is \(n\)-pseudo-coherent; on the other hand \(F^{n+N}[-n-N]\) is \(n\)-pseudo-coherent by hypothesis, using (2.9 and 2.5 a). Hence, by (2.5 b), \(F\) is \(n\)-pseudo-coherent.

\textbf{b)} We argue by induction on \(N\) such that \(\Ho^i(F)=0\) for \(i>n+N\). If \(N=0\), the assertion follows from (2.10 b). Suppose \(N\geq1\). As in the proof of (2.10 b), we have a distinguished triangle in \(\Dpx^-(\mathcal C_S)\)
\[
  F'\longrightarrow F\longrightarrow \Ho^{n+N}(F)[-n-N]\longrightarrow F'[1].
\]
more plainly, with cone \(\Ho^{n+N}(F)[-n-N]\). The exact cohomology sequence implies
\[
  \Ho^i(F')=0\quad\text{for } i>n+N-1,
\]
that is, for \(i>(n-1)+(N-1)\), and
\[
  \Ho^i(F')\qis \Ho^{i+1}(F)\quad\text{for }i<n+N-1 .
\]
By the induction hypothesis, applied with \(n-1\) in place of \(n\), \(F'\) is therefore \((n-1)\)-pseudo-coherent. On the other hand, \(\Ho^{n+N}(F)[-n-N]\) is \(n\)-pseudo-coherent by (2.9 and 2.5 a). Hence, by (2.6), \(F\) is \(n\)-pseudo-coherent.

\textbf{Remark 2.11.1.} If \(F\) is pseudo-coherent, the \(F^i\) (respectively the \(\Ho^i(F)\)) are not in general pseudo-coherent, nor even of finite type. Nevertheless, in the following paragraph, in the excursions on the notions of coherence, we shall see important examples where pseudo-coherence of \(F\) implies that of the \(\Ho^i(F)\).

\textbf{Proposition 2.12.} Let \(F',F''\in\Ob\Dpx^-(\mathcal C_S)\), and let \(F=F'\oplus F''\). For \(F\) to be \(n\)-pseudo-coherent (respectively pseudo-coherent), it is necessary and sufficient that \(F'
\) and \(F''\) be \(n\)-pseudo-coherent (respectively pseudo-coherent).

\emph{Proof.} It is enough to prove the non-parenthesized assertion. Sufficiency is obvious; let us prove necessity. By induction we may suppose that it has already been proved that \(F'\) and \(F''\) are \((n+1)\)-pseudo-coherent. After localizing, we may even suppose that \(F'\) and \(F''\) are strictly \((n+1)\)-pseudo-coherent and that \(F=F'\oplus F''\) in \(\Cpx^-(\mathcal C_S)\). Let \(F'_{>n}\) (respectively \(F'_{\leq n}\)) denote the complex obtained from \(F'\) by replacing \(F'^i\) by \(0\) for \(i\leq n\) (respectively for \(i>n\)), and define \(F''_{>n}\) and \(F''_{\leq n}\) similarly. We have an exact sequence
\[
  0\longrightarrow F'_{>n}\oplus F''_{>n}\longrightarrow F
  \longrightarrow F'_{\leq n}\oplus F''_{\leq n}\longrightarrow 0 .
\]
The complex \(F'_{>n}\oplus F''_{>n}\) is strictly perfect; since \(F\) is \(n\)-pseudo-coherent, (2.6) implies that \(F'_{\leq n}\oplus F''_{\leq n}\) is \(n\)-pseudo-coherent. This means, by (2.10 b), that
\[
  \Ho^n(F'_{\leq n}\oplus F''_{\leq n})
\]
is of finite type. But
\[
  \Ho^n(F'_{\leq n}\oplus F''_{\leq n})
  =\Ho^n(F'_{\leq n})\oplus \Ho^n(F''_{\leq n}),
\]
so \(\Ho^n(F'_{\leq n})\) (respectively \(\Ho^n(F''_{\leq n})\)) is of finite type, and again by (2.10 b) \(F'_{\leq n}\) (respectively \(F''_{\leq n}\)) is \(n\)-pseudo-coherent. Applying (2.5) to the exact sequence
\[
  0\longrightarrow F'_{>n}\longrightarrow F'\longrightarrow F'_{\leq n}\longrightarrow 0
\]
(and similarly for \(F''\)), we deduce that \(F'\) and \(F''\) are \(n\)-pseudo-coherent.

\textbf{Proposition 2.13.} Let \((S',\mathcal C',\mathcal C'_0)\) be a triple satisfying the same hypotheses (2.0) as \((S,\mathcal C,\mathcal C_0)\). Let \(f:S'\to S\) be a morphism of sites, corresponding to an underlying functor \(f^{-1}:S\to S'\), and let
\[
  u:\Dpx^*(\mathcal C)\longrightarrow \Dpx^*(\mathcal C')\qquad (*=+,-,b)
\]
be a Cartesian functor above \(f\), exact on each fiber. Suppose that there exist \(a,b\in\mathbb Z\) such that, for every \(S\in\Ob S\), the following conditions hold:

\begin{enumerate}[label=(\roman*)]
\item if \(E\in\Ob\mathcal C_{0,S}\), then \(u(E)\) is \(a\)-\(\mathcal C'_0\)-pseudo-coherent (respectively \(\mathcal C'_0\)-pseudo-coherent);
\item if \(F\in\Ob\Dpx^*(\mathcal C_S)\) is acyclic in degrees \(>0\), then \(u(F)\) is acyclic in degrees \(>b\).
\end{enumerate}

Then, if \(F\in\Ob\Dpx^*(\mathcal C_S)\) is \(m\)-pseudo-coherent relative to \(\mathcal C_0\) and acyclic in degrees \(>n\), the object
\[
  u(F)\in\Ob\Dpx^*(\mathcal C'_{S'}),\qquad S'=f^{-1}(S),
\]
is \(\sup(m+b,n+a)\)-pseudo-coherent relative to \(\mathcal C'_0\) (respectively \((m+b)\)-pseudo-coherent), and is acyclic in degrees \(>n+b\). In particular, in the parenthesized case, if \(F\in\Ob\Dpx^*(\mathcal C_S)\) is pseudo-coherent relative to \(\mathcal C_0\), then \(u(F)\) is pseudo-coherent relative to \(\mathcal C'_0\).

\emph{Proof.} We may restrict to the non-parenthesized case. First suppose that \(F\in\Ob\Dpx^*(\mathcal C_S)\) is strictly perfect (2.1), with \(F^i=0\) for \(i>n\). Then \(u(F)\) is \((n+a)\)-\(\mathcal C'_0\)-pseudo-coherent. We argue by induction on the number \(N\) of indices \(i\) such that \(F^i\neq0\). If \(N=1\), then \(F=F^p[-p]\) for some \(p\leq n\), so \(u(F)\) is \((p+a)\)-\(\mathcal C'_0\)-pseudo-coherent, hence a fortiori \((n+a)\)-pseudo-coherent. If \(N>1\), there is an exact sequence
\[
  0\longrightarrow F'\longrightarrow F\longrightarrow F''\longrightarrow 0,
\]
where \(F'\) and \(F''\) are strictly perfect of length \(<N-1\), and where \(F'^i=F''^i=0\) for \(i>n\). Applying (2.5 b) to the distinguished triangle
\[
  u(F')\longrightarrow u(F)\longrightarrow u(F'')\longrightarrow u(F')[1].
\]
and using the induction hypothesis, we get that \(u(F)\) is \((n+a)\)-pseudo-coherent.

Now pass to the general case: \(F\) is \(m\)-\(\mathcal C_0\)-pseudo-coherent and acyclic in degrees \(>n\). After localizing, we may, by (2.10 a), suppose \(F\) strictly \(m\)-pseudo-coherent and such that \(F^i=0\) for \(i>n\). There is then an exact sequence
\[
  0\longrightarrow F'\longrightarrow F\longrightarrow F''\longrightarrow 0,
\]
where \(F''\) is zero in degrees \(\geq m\), and \(F'\) is strictly perfect and zero in degrees \(>n\). In the distinguished triangle
\[
  u(F')\longrightarrow u(F)\longrightarrow u(F'')\longrightarrow u(F')[1].
\]
the object \(u(F'')\) is acyclic in degrees \(>m+b\), hence \((m+b)\)-pseudo-coherent, while \(u(F')\) is \((n+a)\)-pseudo-coherent by what was just proved. We conclude by (2.5 b).

\textbf{Corollary 2.14.} Let \(\mathcal C_1\) be a strictly full additive sub-\(S\)-category of \(\mathcal C\) such that
\[
  \mathcal C_0\subset \mathcal C_1\subset \mathcal C.
\]
Assume that every object of \(\mathcal C_1\) is \(\mathcal C_0\)-pseudo-coherent. Then:

\begin{enumerate}[label=(\roman*)]
\item \(\mathcal C_1\) is locally quasi-liftable in \(\mathcal C\) (1.2), and locally quasi-stable under kernels of epimorphisms (1.2);
\item for \(S\in\Ob S\) and \(F\in\Ob\Dpx(\mathcal C_S)\), in order that \(F\) be \(n\)-pseudo-coherent (respectively pseudo-coherent) relative to \(\mathcal C_1\), it is necessary and sufficient that it be so relative to \(\mathcal C_0\).
\end{enumerate}

\emph{Proof.} For (i), let \(E\to F\) be an epimorphism in \(\mathcal C_S\), with \(F\in\Ob\mathcal C_{1,S}\). Since \(F\) is \(\mathcal C_0\)-pseudo-coherent, hence a fortiori of \(\mathcal C_0\)-finite type, there exists, after localizing, an epimorphism \(E'\to F\) with \(E'\in\Ob\mathcal C_{0,S}\). Using the quasi-liftability of \(\mathcal C_0\) in \(\mathcal C\), we may suppose, after localizing again, that this epimorphism factors through \(E\); this gives the quasi-liftability of \(\mathcal C_1\). On the other hand, if \(E\in\Ob\mathcal C_S\) admits a finite right resolution by objects of \(\mathcal C_1\), then \(E\) is \(\mathcal C_0\)-pseudo-coherent by (2.11 a), hence a fortiori of \(\mathcal C_1\)-finite type.

For (ii), it is enough to prove necessity. For this one may apply, at choice, (2.11 a) or (2.13), with \(S'=S\), \(\mathcal C'=\mathcal C\), \(f=\Id\), \(u=\Id\), and with \(\mathcal C'_0\) and \(\mathcal C_0\) replaced respectively by \(\mathcal C_0\) and \(\mathcal C_1\).

\textbf{Scholium 2.14.1.} There is therefore a largest strictly full additive sub-\(S\)-category of \(\mathcal C\) giving the same notion of \(n\)-pseudo-coherence (respectively pseudo-coherence) as \(\mathcal C_0\), namely the category formed by the \(\mathcal C_0\)-pseudo-coherent objects. This latter category is locally stable under kernels of epimorphisms.

\textbf{Example 2.15.} Let \((S,A)\) be a ringed topos, and let \(\mathcal C,\mathcal C_0\) be the \(S\)-categories considered in (1.3 d). The fibered category \(\Dpx(\mathcal C)\) will often be denoted by \(\Dpx(S,A)\), or \(\Dpx_A(S)\), when one wants to suggest that the modules are left \(A\)-modules, or even simply \(\Dpx(S)\) if there is no ambiguity about \(A\). We denote by \(S\) the final object of \(S\), and put \(\Dpx(\mathcal C_S)=\Dpx_A(S)\) (or \(\Dpx(S)\)). An object of \(\Dpx(S)\) will be said to be \(n\)-pseudo-coherent (respectively pseudo-coherent) if it is \(n\)-\(\mathcal C_0\)-pseudo-coherent (respectively \(\mathcal C_0\)-pseudo-coherent), and similarly for the strict notions. The category fibered in \(\Dpx(\mathcal C)\), denoted \(\Dpx(\mathcal C)_{n-\mathcal C_0-\coh}\) (respectively \(\Dpx(\mathcal C)_{\mathcal C_0-\coh}\)) in (2.6), will be denoted \(\Dpx_A(S)_{n-\coh}\) (respectively \(\Dpx_A(S)_{\coh}\)). One observes that, by (2.14), the same notion of pseudo-coherence is obtained if \(\mathcal C_0\) is replaced by the \(S\)-category \(\mathcal C_1\) defined by
\[
  U\longmapsto \text{the category of left }A_U\text{-modules locally direct summands of free finite-type modules}.
\]
If an \(A\)-module \(M\) is a direct summand of a free finite-type \(A\)-module, it does not follow in general that \(M\) is locally free of finite type. One does, however, have the following result.

\textbf{Lemma 2.15.1.} Suppose that \(S\) satisfies one of the following two hypotheses:
\begin{enumerate}[label=(\roman*)]
\item \(S\) has enough points (SGA 4 IV 6.4), and for every point \(s\) of \(S\), the left \(A_s\)-modules which are projective of finite type are free; this last condition is satisfied, for example, if \(A_s/\Rad(A_s)\) is a field or if \(\Rad(A_s)\) is nilpotent;
\item \((S,A)\) is a ``locally ringed'' topos in the sense of [1], which means that \(A\) is commutative and that, for every \(X\in\Ob S\) and every \(f\in\Gamma(X,A)\), one has
\[
  X=X_f\cup X_{1-f},
\]
where \(X_f\) (respectively \(X_{1-f}\)) denotes the largest subobject of \(X\) on which \(f\) (respectively \(1-f\)) is invertible.
\end{enumerate}
Then every direct summand of a free finite-type left \(A\)-module is locally free of finite type.

\emph{Proof.} Suppose first that (i) holds. We note that, if \(E\) and \(F\) are two \(A\)-modules and \(E\) is of finite presentation, the canonical arrow
\[
  \Hom(E,F)_s\longrightarrow \Hom(E_s,F_s)
\]
is an isomorphism for every point \(s\) of \(S\). It follows that, if \(E\) and \(F\) are both of finite presentation and if \(E_s\qis F_s\) at a point \(s\) of \(S\), then there is an \(s\)-pointed object \(U\) of \(S\) such that \(E|U\qis F|U\). The conclusion follows at once.

Suppose next that (ii) holds. Let \(n\in\mathbb N\), and let \(\mathbb D^n\) denote the functor on the category of locally ringed topoi belonging to a fixed universe, defined by
\[
  T\longmapsto \Gamma(T,\mathcal O_T)^n.
\]
It is known [1] that the functor \(\End_{\mathbb D}(\mathbb D^n)\) is represented by
\[
  M_n=\Spec\mathbb Z[T_{ij}]\quad(1\leq i\leq n,
  1\leq j\leq n),
\]
endowed with the endomorphism \(U_n\) of \((\mathcal O_{M_n})^n\) defined by the matrix \((T_{ij})\). Now let
\[
  p:A^n\longrightarrow A^n
\]
be a projector. We show that \(p(A^n)\) is locally free of finite type. From what has just been said, there exists a morphism
\[
  s:S\longrightarrow M_n
\]
of locally ringed topoi such that \(s^*(u_n)=p\). We may interpret \(s\) as a point of \(M_n\) with values in \(\Gamma(S,A)\), that is, as a matrix of order \(n\) with coefficients in \(\Gamma(S,A)\). Since this matrix is a projector, \(s\) factors through the closed subscheme
\[
  i:P_n\hookrightarrow M_n
\]
corresponding to projectors; say \(S\xrightarrow{s'}P_n\xrightarrow{i}M_n\). Hence
\[
  p=s'^* i^*(u_n),
\]
and we are done, because, by the assertion under hypothesis (i), the image of \(i^*(u_n)\) is a locally free finite-type sheaf.

\textbf{Corollary 2.16.} Let
\[
  f:(S',A')\longrightarrow (S,A)
\]
be a morphism of ringed topoi, defined by a functor \(f^{-1}:S\to S'\) and a morphism of sheaves of rings \(f^{-1}(A)\to A'\). Let \(\Dpx(_{A'}S'_A)\) denote the derived category of \(A'\)-\(A\)-bimodules, left over \(A'\) and right over \(A\), and let
\[
  \otimes_A^{\mathrm L}:\Dpx^-(_{A'}S'_A)\times_{S'}\Dpx^-(_A S)
  \longrightarrow \Dpx^-(_{A'}S')
\]
be the derived functor of the tensor product
\[
  (E',E)\longmapsto E'\otimes_A f^{-1}(E).
\]
Let \(X\in\Ob S\), put \(X'=u(X)\), and let \(E'\in\Ob\Dpx^-(_{A'}X'_A)\), \(E\in\Ob\Dpx^-(_A X)\). If \(E'\) is \(m'\)-pseudo-coherent (as an object of \(\Dpx^-(_A S')\)) and acyclic in degrees \(>n'\), and if \(E\) is \(m\)-pseudo-coherent and acyclic in degrees \(>n\), then
\[
  E'\otimes_A^{\mathrm L}E
\]
is \(\sup(m+n',m'+n)\)-pseudo-coherent and acyclic in degrees \(>n+n'\).

\emph{Proof.} After replacing \(S\) by \(S/X\), we may suppose \(X=S\), and apply (2.13) with
\[
  u=E'\otimes_A^{\mathrm L}:\Dpx^-(_A S)\longrightarrow \Dpx^-(_{A'}S').
\]
Let \(U\in\Ob S\), let \(U'=f^{-1}(U)\), and let \(E\) be a complex of left \(A_U\)-modules. If \(E\) is acyclic in degrees \(>n\), then \(E'\otimes_A^{\mathrm L}E\) is acyclic in degrees \(>n+n'\), by the definition of \(\otimes_A^{\mathrm L}\). On the other hand, if \(E=(A_U)^p\), then
\[
  E'\otimes_A^{\mathrm L}E=(E'|U')^p,
\]
which is \(m'\)-pseudo-coherent by hypothesis and (2.12). The conditions of (2.13) are therefore satisfied, with \(a=m'\) and \(b=n'\), whence the conclusion.

In particular:

\textbf{Corollary 2.16.1.} Under the hypotheses of (2.16), one has:

\textbf{a)} The functor
\[
  \mathrm L f^*: \Dpx^-(_A S)\longrightarrow \Dpx^-(_{A'}S')
\]
(which is none other than the functor \(A'\otimes_A^{\mathrm L}\)) induces a functor
\[
  \mathrm L f^*: \Dpx^-(_A S)_{\coh}\longrightarrow \Dpx^-(_{A'}S')_{\coh}
\]
(with the notation of (2.6)).

\textbf{b)} If \(A\) is commutative, the functor
\[
  \otimes_A^{\mathrm L}:\Dpx^-(S)\times_S\Dpx^-(S)
  \longrightarrow \Dpx^-(S)
\]
induces a functor
\[
  \otimes_A^{\mathrm L}:\Dpx^-(S)_{\coh}\times_S\Dpx^-(S)_{\coh}
  \longrightarrow \Dpx^-(S)_{\coh} .
\]

\textbf{Exercise 2.16.2.} Let \(S\) be a topos, and let \(A\to A'\) be a morphism of rings of \(S\). Assume one of the following hypotheses:
\begin{enumerate}[label=(\roman*)]
\item \(A\to A'\) is surjective, with nilpotent ideal as kernel;
\item the functor \(A'\otimes_A-\), on the category of left \(A\)-modules, is exact and faithful.
\end{enumerate}
Show that \(E\in\Ob\Dpx^-(_A S)\) is \(n\)-pseudo-coherent (respectively pseudo-coherent) if and only if
\[
  A'\otimes_A E\in\Ob\Dpx^-(_{A'}S)
\]
is \(n\)-pseudo-coherent (respectively pseudo-coherent).

Hints: first show that a left \(A\)-module \(L\) is of finite type if and only if \(A'\otimes_A L\) is of finite type (cf. Bourbaki, Commutative Algebra, Chap. I, §3, Prop. 11 and Chap. II, §3, Prop. 4), and then reduce to this case by truncation, using (2.5) and (2.10 b).

\section*{3. Link with the classical notion of coherence}

\textbf{3.0.} The hypotheses and notation are those of (2.0), except that \(\mathcal C_0\) is no longer assumed locally quasi-stable under kernels of epimorphisms.

\textbf{Definition 3.1.} Let \(F\in\Ob\mathcal C_S\). We shall say that \(F\) is \emph{precoherent} relative to \(\mathcal C_0\) if, for every \(U\) over \(S\) and every morphism
\[
  u:E\longrightarrow F|U,
\]
where \(E\in\Ob\mathcal C_{0,U}\), \(\Ker(u)\) is of \(\mathcal C_0\)-finite type (1.2). We shall say that \(F\) is \emph{coherent} relative to \(\mathcal C_0\) if \(F\) is precoherent and of finite type.

\emph{Technical note.} We introduce this notion here only as a technical intermediary; the only notion to retain is that of coherence.

\textbf{Remark 3.2.} Let \(F\in\Ob\mathcal C_S\). It is clear from the definitions that, if \(F\) is precoherent (respectively of finite type), then every subobject (respectively quotient) of \(F\) is precoherent (respectively of finite type).

\textbf{Lemma 3.3.}

\textbf{a)} Let
\[
  0\longrightarrow F\longrightarrow G\longrightarrow H
\]
be an exact sequence of \(\mathcal C_S\). If \(G\) is of finite type and \(H\) is precoherent, then \(F\) is of finite type.

\textbf{b)} Let
\[
  F\longrightarrow G\longrightarrow H\longrightarrow 0
\]
be an exact sequence of \(\mathcal C_S\). If \(F\) is of finite type and \(G\) is precoherent, then \(H\) is precoherent.

\textbf{c)} Let
\[
  0\longrightarrow F\longrightarrow G\longrightarrow H\longrightarrow 0
\]
be an exact sequence of \(\mathcal C_S\). If \(F\) and \(H\) are precoherent (respectively of finite type), then so is \(G\).

\emph{Proof.} a) After localizing, there is an epimorphism \(G'\to G\) with \(G'\in\Ob\mathcal C_0\). Hence there is a commutative diagram with exact rows
\[
\begin{tikzcd}
0\arrow[r]&F'\arrow[r]\arrow[d]&G'\arrow[r]\arrow[d]&H\arrow[d,"\Id"]\\
0\arrow[r]&F\arrow[r]&G\arrow[r]&H .
\end{tikzcd}
\]
Since \(H\) is precoherent, \(F'\) is of finite type; but \(F'\to F\) is an epimorphism, so \(F\) is of finite type.

b) Let \(U\) be over \(S\), and let \(H'\xrightarrow{f}H|U\) be a morphism of \(\mathcal C_U\), with \(H'\in\Ob\mathcal C_{0,U}\). After a change of notation, we may suppose \(U=S\). By the quasi-liftability of \(\mathcal C_0\) in \(\mathcal C\), there exists, after localizing, a commutative square
\[
\begin{tikzcd}
G' \arrow[r,"p"] \arrow[d] & H' \arrow[d,"f"]\\
G \arrow[r] & H,
\end{tikzcd}
\]
where the horizontal arrows are epimorphisms and \(G'\in\Ob\mathcal C_0\). Since \(\Ker(f)\) is a quotient of \(\Ker(fp)\), it is enough to show that \(\Ker(fp)\) is of finite type; hence we may suppose \(H'=G'\). The arrows \(F\to G\) and \(H'\to G\) define
\[
  g:F\oplus H'\longrightarrow G.
\]
Putting \(K=\Ker(G\to H)\), we have a commutative diagram, with exact rows and with \(q\) an epimorphism,
\[
\begin{tikzcd}
0\arrow[r]&F\arrow[r]\arrow[d,"q"]&F\oplus H'\arrow[r]\arrow[d,"g"]&H'\arrow[r]\arrow[d,"f"]&0\\
0\arrow[r]&K\arrow[r]&G\arrow[r]&H\arrow[r]&0 .
\end{tikzcd}
\]
The snake diagram shows that \(\Ker(g)\to\Ker(f)\) is an epimorphism. But, since \(F\oplus H'\) is of finite type and \(G\) is precoherent, \(\Ker(g)\) is of finite type by a); hence we are done.

c) First prove the non-parenthesized assertion. Up to a change of notation, it is a matter of showing that, if \(g:G'\to G\) is a morphism of \(\mathcal C_S\), with \(G'\in\Ob\mathcal C_0\), then \(\Ker(g)\) is of finite type. Let \(H'\) be the image of the composed arrow \(G'\to H\). We have a commutative diagram with exact rows
\[
\begin{tikzcd}
0\arrow[r]&F'\arrow[r]\arrow[d,"f"]&G'\arrow[r]\arrow[d,"g"]&H'\arrow[r]\arrow[d,"h"]&0\\
0\arrow[r]&F\arrow[r]&G\arrow[r]&H\arrow[r]&0 .
\end{tikzcd}
\]
Since \(h\) is a monomorphism, the snake diagram gives an isomorphism \(\Ker(f)\qis\Ker(g)\). But \(F'\) is of finite type by a), hence \(\Ker(f)\) is of finite type, again by a). This proves the non-parenthesized assertion. For the parenthesized assertion, strictly speaking one cannot invoke (2.5 b), since \(\mathcal C_0\) is not assumed locally quasi-stable under kernels of epimorphisms. After localizing, there are epimorphisms
\[
  F'\xrightarrow{u}F,
  \qquad H'\xrightarrow{v}H,
\]
with \(F',H'\in\Ob\mathcal C_0\). After localizing further, we may suppose that \(v\) factors through \(G\), whence a commutative diagram with exact rows
\[
\begin{tikzcd}
0\arrow[r]&F'\arrow[r]\arrow[d,"u"]&F'\oplus H'\arrow[r]\arrow[d,"w"]&H'\arrow[r]\arrow[d,"v"]&0\\
0\arrow[r]&F\arrow[r]&G\arrow[r]&H\arrow[r]&0 .
\end{tikzcd}
\]
Since \(u\) and \(v\) are epimorphisms, so is \(w\), and this proves the assertion. The proof of (3.3) is complete.

It follows immediately that:

\textbf{Corollary 3.4.} The full (respectively strictly full) subcategory of \(\mathcal C_S\) formed by coherent objects is stable under finite projective (respectively inductive) limits. In particular, the kernel, cokernel, and image of an arrow between two coherent objects are coherent. Let
\[
  E^0\longrightarrow E^1\longrightarrow E^2\longrightarrow E^3\longrightarrow E^4
\]
be an exact sequence of \(\mathcal C_S\). If \(E^0,E^1,E^3,E^4\) are coherent, then \(E^2\) is coherent as well.

\textbf{Corollary 3.5.} Suppose that the objects of \(\mathcal C_0\) are \(\mathcal C_0\)-coherent (which in particular implies that \(\mathcal C_0\) is locally quasi-stable under kernels of epimorphisms, that is, that we are under the conditions of (2.0)).

\textbf{a)} For \(F\in\Ob\mathcal C_S\), the following conditions are equivalent:
\begin{enumerate}[label=(\roman*)]
\item \(F\) is of finite presentation;
\item \(F\) is coherent;
\item \(F\) is of finite \(\infty\)-presentation (2.8).
\end{enumerate}

\textbf{b)} Let \(F\in\Ob\Dpx^-(\mathcal C_S)\), and let \(n\in\mathbb Z\). Then
\[
\begin{split}
F\text{ is }n\text{-pseudo-coherent}
\quad\Longleftrightarrow\quad&
\Ho^i(F)\text{ is coherent for }i\geq n+1,\\
&\text{and }\Ho^n(F)\text{ is of finite type.}
\end{split}
\]
In particular, for \(F\) to be pseudo-coherent it is necessary and sufficient that \(\Ho^n(F)\) be coherent for every \(n\).

\emph{Proof.} Assertion a) and the implication \(\Rightarrow\) in b) follow trivially from (3.4). The implication \(\Leftarrow\) in b) follows at once from a) and (2.11 b).

\textbf{Examples 3.6.} Let \((S,A)\) be a ringed topos, and let \(\mathcal C,\mathcal C_0\) be the \(S\)-categories of (1.3 d), also as in (2.15). In order that the objects of \(\mathcal C_0\) be coherent, it is necessary and sufficient that \(A\) be coherent, that is, precoherent. Here are examples where this is the case:
\begin{itemize}
\item \(A\) is constant with value a left noetherian ring (SGA 4 IX 2.1);
\item \(S\) is the topos of sheaves of sets, for the Zariski topology, on a locally noetherian scheme \(S\), and \(A=\mathcal O_S\);
\item \(S\) is the topos of sheaves of sets on a complex analytic space \(S\), and \(A=\mathcal O_S\).
\end{itemize}
On the other hand, the sheaf of continuous complex-valued functions on a topological space is not coherent in general.

Of course, if \((S,A)\) is defined by a ringed space, the notion of coherence considered here coincides with the usual notion of Serre (EGA \(0_I\), 5.3.1).

\textbf{Proposition 3.7.} \textbf{a)} Let
\[
  f:(S',A')\longrightarrow(S,A)
\]
be a morphism of ringed topoi. Assume that \(f\) is flat on the right, that is, that the functor \(f^*=A'\otimes_A-\) is exact. There is, for \(F\in\Ob\Dpx_A^-(S)\) and \(G\in\Ob\Dpx_A^+(S)\), a functorial canonical arrow
\[
  f^*\RHom(F,G)\longrightarrow \RHom(f^*(F),f^*(G)).
\]
This arrow is an isomorphism when \(F\) is pseudo-coherent.

\textbf{b)} Suppose \(A\) is commutative and coherent, and let \(F\in\Ob\Dpx^-(S)\), \(G\in\Ob\Dpx^+(S)\). If \(\Ho^i(F)\) and \(\Ho^i(G)\) are coherent for all \(i\), then the same is true of \(\Ext^i(F,G)\).

\emph{Proof.} a) To define the arrow (cf. SGA 4 XVII), one may suppose that \(G\) has injective components, and one takes the composite arrow
\[
  f^*\Hom(F,G)\longrightarrow \Hom(f^*F,f^*G)
  \longrightarrow \RHom(f^*F,f^*G).
\]
Let us show that it is an isomorphism for \(F\) pseudo-coherent. For brevity write \(T\) (respectively \(T'\)) for the functor \(f^*\RHom(\cdot,G)\) (respectively \(\RHom(f^*\cdot,f^*G)\)), and let \(a\in\mathbb Z\) be such that \(\Ho^i(G)=0\) for \(i<a\). It is clear that one has \(TA\qis T'A\), hence by dévissage \(TF\qis T'F\) for \(F\) strictly perfect (2.1). On the other hand, if \(F\) is acyclic in degrees \(>n\), then \(TF\) and \(T'F\) are acyclic in degrees \(<a-n\). To verify that \(TF\qis T'F\) for \(F\) pseudo-coherent, the question being local, one may suppose \(F\) strictly \(n\)-pseudo-coherent; that is, there is a distinguished triangle
\[
  F'\longrightarrow F\longrightarrow F''\longrightarrow F'[1].
\]
with \(F'\) strictly perfect and \(F''\) acyclic in degrees \(>n-1\). From this one obtains a morphism of distinguished triangles
\[
\begin{tikzcd}[column sep=large]
TF'' \arrow[r] \arrow[d] & TF \arrow[r] \arrow[d] & TF' \arrow[d] \\
T'F'' \arrow[r] & T'F \arrow[r] & T'F' .
\end{tikzcd}
\]
Now \(TF'\qis T'F'\), since \(F'\) is strictly perfect, and \(TF''\) and \(T'F''\) are acyclic in degrees \(\leq a-n\). Writing the induced morphism on the exact cohomology sequences of the preceding triangles, one deduces that \(TF\to T'F\) induces an isomorphism
\[
  \Ho^i(TF)\qis \Ho^i(T'F)
\]
for \(i\leq a-n\). Since \(n\) may be chosen arbitrarily negative, this proves the assertion.

b) The proof is left as an exercise to the reader; see, for example, [H] II 3.3.

\textbf{Remark 3.7.1.} We shall see below (7.1.2) another case of compatibility of \(\RHom\) with \(\mathrm L f^*\).

\textbf{Corollary 3.7.2.} (cf. EGA \(0_I\), 5.2.7.) Let \((S,A)\) be a ringed topos, let \(s\) be a point of \(S\) (SGA 4 VIII 7.8), and let
\[
  E,F\in\Ob\D^b_{A,\coh}(S).
\]
If \(E_s\qis F_s\), then there exists an \(s\)-pointed object \(U\) of \(S\) such that
\[
  E|U\qis F|U.
\]

\emph{Proof.} Let \(\alpha:E_s\to F_s\) and \(\beta:F_s\to E_s\) be reciprocal isomorphisms. By (2.22 a), one has
\[
  \Hom(E,F)_s\qis\Hom(E_s,F_s),
  \qquad
  \Hom(F,E)_s\qis\Hom(F_s,E_s).
\]
There are therefore a neighbourhood \(U\) of \(s\) and a section \(u\) (respectively \(v\)) of \(\Hom(E,F)\) (respectively of \(\Hom(F,E)\)) over \(U\) such that \(u_s=\alpha\) (respectively \(v_s=\beta\)). The morphisms \(vu:E|U\to E|U\) and \(uv:F|U\to F|U\) induce the identity at \(s\); hence on a suitable \(V\) over \(U\), as required.

\section*{4. Perfect complexes}

\textbf{4.0.} In this section, \(S\) denotes a site, \(S\) an object of \(S\), \(\mathcal C\) a flat abelian \(S\)-category, and \(\mathcal C_0\) a strictly full additive sub-\(S\)-category of \(\mathcal C\) (1.1). Unless otherwise mentioned, we suppose that \(\mathcal C_0\) is locally liftable in \(\mathcal C\) and locally stable under kernels of epimorphisms (1.2).

\textbf{Lemma 4.1.} Let
\[
  u:E\longrightarrow F
\]
be an arrow of \(\Cpx(\mathcal C_S)\). If \(E\) is strictly \(\mathcal C_0\)-perfect (2.1) and \(F\) is acyclic, then \(u\) is locally homotopic to zero.

\emph{Proof.} One has to construct locally a homotopy \(k\) such that
\[
  u=dk+kd.
\]
Since \(E\) is strictly perfect, there are \(a,b\in\mathbb Z\), with \(a\leq b\), such that \(E^i=0\) for \(i\notin[a,b]\). Thus one may take \(k^i=0\) for \(i\notin[a,b]\). We construct the \(k^i\), for \(i>a\), by descending induction on \(i\). Suppose \(k^i:E^i\to F^{i-1}\) has been constructed so that
\[
  u^i=k^{i+1}d_E^i+d_F^{i-1}k^i
\]
for \(i>n>a\). We show that, after localizing, there exists
\[
  k^{n-1}:E^{n-1}\longrightarrow F^{n-2}
\]
such that
\[
  u^{n-1}=k^n d_E^{n-1}+d_F^{n-2}k^{n-1}.
\]
By the induction hypothesis, the arrow
\[
  v^{n-1}=u^{n-1}-k^n d_E^{n-1}:E^{n-1}\longrightarrow F^{n-1}
\]
factors through \(Z^{n-1}(F)\). But \(Z^{n-1}(F)=B^{n-1}(F)\), since \(F\) is acyclic. Since \(\mathcal C_0\) is locally liftable in \(\mathcal C\), it follows, after localizing, that there exists \(k^{n-1}:E^{n-1}\to F^{n-2}\) such that \(v^{n-1}=d_F^{n-2}k^{n-1}\), as required.

\textbf{Corollary 4.2.} Let
\[
\begin{tikzcd}
P \arrow[d,"v"'] \\
F & E \arrow[l,"f"]
\end{tikzcd}
\]
be a diagram of \(\Cpx(\mathcal C_S)\), where \(P\) is strictly \(\mathcal C_0\)-perfect and \(f\) is a quasi-isomorphism. The set of \(X\) over \(S\) for which there exists an arrow
\[
  u:P_X\longrightarrow E_X
\]
of \(\Cpx(\mathcal C_X)\) such that the diagram
\[
\begin{tikzcd}
& P_X \arrow[d,"v_X"]\\
E_X \arrow[r,"f_X"'] \arrow[ur,"u"] & F_X
\end{tikzcd}
\]
is commutative in \(\Kpx(\mathcal C_X)\), is a refinement of \(S\).

\emph{Proof.} Let \(G\) be the cone of \(f\). There is an exact sequence
\[
  \longrightarrow \Hom_{\Kpx(\mathcal C_S)}(P,E)
  \longrightarrow \Hom_{\Kpx(\mathcal C_S)}(P,F)
  \longrightarrow \Hom_{\Kpx(\mathcal C_S)}(P,G)
  \longrightarrow .
\]
Since \(G\) is acyclic, the image of \(v\) in \(\Hom_{\Kpx(\mathcal C_S)}(P,G)\) is locally zero by (4.1). Therefore, after localizing, there exists \(u:P\to E\) such that \(v=fu\) in \(\Kpx(\mathcal C_S)\), as required.

\textbf{Corollary 4.3.} Let \(f:E\to F\) be a quasi-isomorphism of \(\Cpx(\mathcal C_S)\), with \(F\) strictly perfect relative to \(\mathcal C_0\). The set of \(X\) over \(S\) for which there exists a quasi-isomorphism
\[
  g:F_X\longrightarrow E_X
\]
such that \(f_Xg\) is homotopic to \(\Id_{F_X}\), is a refinement of \(S\).

\emph{Proof.} This is a special case of (4.2).

\textbf{Scholium 4.4.} Let \(F\in\Ob\Kpx(\mathcal C_S)\). One may define, for example with the aid of a cleavage of \(\mathcal C\) (SGA 1 VI 7.1), a category \((\Qis/F)\) fibered over \(S/S\), whose fiber over \(X\) over \(S\) is the category \((\Qis/F_X)^\circ\), opposite to the category of quasi-isomorphisms with target \(F_X\). One may then express (4.3) by saying that, if \(F\) is strictly perfect, the arrow \(\Id_F\) is locally cofinal in \((\Qis/F)\).

\textbf{Notation 4.5.} Let \(\mathcal M\) be a fibered \(S\)-category. Let \(X\in\Ob S\), and let \(E,F\in\Ob\mathcal M_X\). We shall denote by
\[
  \underline{\Hom}_{\mathcal M}(E,F)
\]
the sheaf on \(S/X\) associated to the presheaf of \(S\)-morphisms from \(E\) to \(F\) in \(\mathcal M\). For the definition of this presheaf, see [G] I 2.6. One must be careful that, even if \(\mathcal C\) is a stack ([G]), the corresponding presheaf for \(\mathcal M=\Kpx(\mathcal C)\) or \(\Dpx(\mathcal C)\) is not in general a sheaf.

\textbf{Corollary 4.6.} Let \(E,F\in\Ob\Cpx(\mathcal C_S)\), with \(E\) strictly perfect. The canonical arrow
\[
  \underline{\Hom}_{\Kpx(\mathcal C)}(E,F)
  \longrightarrow
  \underline{\Hom}_{\Dpx(\mathcal C)}(E,F)
\]
is an isomorphism.

\emph{Proof.} This is an immediate consequence of (4.4) and of the formula
\[
  \underline{\Hom}_{\Dpx(\mathcal C_X)}(E_X,F_X)
  = \lim_{(\Qis/E_X)^\circ}\underline{\Hom}_{\Kpx(\mathcal C_X)}(\,\cdot\,,F_X)
\]
for \(X\) over \(S\).

\textbf{Definition 4.7.} Let \(F\in\Ob\Dpx(\mathcal C_S)\), and let \(a,b\in\mathbb Z\), with \(a\leq b\). We shall say that \(F\) has \(\mathcal C_0\)-perfect amplitude contained in the interval \([a,b]\), and we shall write
\[
  \mathcal C_0\text{-}\Parf\text{-}\amp(F)\subset[a,b],
\]
if the set of \(X\) over \(S\) for which there exists a quasi-isomorphism
\[
  F'\longrightarrow F
\]
in \(\Cpx(\mathcal C_X)\), with \(F'\) strictly \(\mathcal C_0\)-perfect (2.1) and such that \(F'^i=0\) for \(i\notin[a,b]\), is a refinement of \(S\). By (4.3), this condition is equivalent to the one obtained by replacing ``quasi-isomorphism in \(\Cpx(\mathcal C_X)\)'' by ``isomorphism in \(\Dpx(\mathcal C_X)\).'' We shall say that \(F\) has finite \(\mathcal C_0\)-perfect amplitude if there are integers \(a\) and \(b\) such that \(\mathcal C_0\)-\(\Parf\)-\(\amp(F)\subset[a,b]\). We shall say that \(F\) is \(\mathcal C_0\)-perfect if \(F\) is locally of finite \(\mathcal C_0\)-perfect amplitude.

\textbf{Example 4.8.} Let \((S,A)\) be a ringed topos, let \(\mathcal C\) be the \(S\)-category defined by
\[
  U\longmapsto \text{the category of left }A_U\text{-modules},
\]
and let \(\mathcal C_0\) be the sub-\(S\)-category defined by
\[
  U\longmapsto \text{the category of left }A_U\text{-modules which are direct summands of free finite-type modules}.
\]
The pair \((\mathcal C,\mathcal C_0)\) satisfies the hypotheses (4.0) by (1.3.1). An object \(F\) of \(\Dpx_A(S)\) (cf. (2.15)) will be said to have perfect amplitude contained in \([a,b]\) (respectively finite perfect amplitude, respectively to be perfect) if \(F\) has \(\mathcal C_0\)-perfect amplitude contained in \([a,b]\) (respectively the corresponding property). When the direct summands of free finite-type modules are locally free (cf. (2.15.1)), to say that \(F\) is perfect means that \(F\) is locally isomorphic, in \(\Dpx(S)\), to a bounded complex whose components are free of finite type.

\textbf{Notation 4.9.} We shall denote by \(\Dpx(\mathcal C)_{\mathcal C_0\text{-}\Parf}\), or \(\Parf(\mathcal C,\mathcal C_0)\), the strictly full sub-\(S\)-category of \(\Dpx(\mathcal C)\) formed by the \(\mathcal C_0\)-perfect objects. In the case of Example (4.8), we shall write \(\Dpx_A(S)_{\Parf}\), or \(\Parf_A(S)\), or simply \(\Parf(S)\) if no ambiguity about \(A\) is possible.

\textbf{Proposition 4.10.} The category \(\Parf(\mathcal C,\mathcal C_0)\) is a triangulated sub-\(S\)-category of \(\Dpx(\mathcal C)\) (1.1.4). The canonical functor
\[
  \Kpx^b(\mathcal C_0)\longrightarrow \Parf(\mathcal C,\mathcal C_0)
\]
(respectively
\[
  \Tr(\Kpx^b(\mathcal C_0))\longrightarrow \Tr(\Parf(\mathcal C,\mathcal C_0))
\]
(1.1.3)) induces an \(S\)-equivalence on the associated stacks ([G] II 2.2.4), which means:

\begin{enumerate}[label=(\roman*)]
\item for every \(X\in\Ob S\) and all \(E,F\in\Ob\Kpx^b(\mathcal C_{0,X})\) (respectively \(\Tr(\Kpx^b(\mathcal C_{0,X}))\)), the canonical arrow
\[
  \Hom_{\Kpx^b(\mathcal C_0)}(E,F)
  \longrightarrow
  \Hom_{\Dpx(\mathcal C)}(E,F)
\]
(respectively
\[
  \Hom_{\Tr(\Kpx^b(\mathcal C_0))}(E,F)
  \longrightarrow
  \Hom_{\Tr(\Dpx(\mathcal C))}(E,F)
\]) is an isomorphism; and
\item every object of \(\Parf(\mathcal C,\mathcal C_0)\) (respectively of \(\Tr(\Parf(\mathcal C,\mathcal C_0))\)) is locally isomorphic in \(\Dpx(\mathcal C)\) (respectively in \(\Tr(\Dpx(\mathcal C))\)) to an object of \(\Kpx^b(\mathcal C_0)\) (respectively of \(\Tr(\Kpx^b(\mathcal C_0))\)).
\end{enumerate}

\emph{Proof.} The non-parenthesized assertion (i) is a special case of (4.6), and the non-parenthesized assertion (ii) is true by the definition (4.7) and (4.9). It is clear, moreover, that \(\Parf(\mathcal C,\mathcal C_0)\) is stable under translation of degrees. To prove that \(\Parf(\mathcal C,\mathcal C_0)\) is a triangulated sub-\(S\)-category of \(\Dpx(\mathcal C)\), it remains to show that, if
\[
  E\xrightarrow{u}F\longrightarrow G\longrightarrow E[1].
  \tag{*}
\]
is a distinguished triangle of \(\Dpx(\mathcal C)\) with \(E\) and \(F\) perfect, then \(G\) is perfect. After localizing, we may suppose that \(E\) and \(F\) are objects of \(\Kpx^b(\mathcal C_0)\). After localizing further, we may suppose, by the non-parenthesized assertion (i), that \(u\) is an arrow of \(\Kpx^b(\mathcal C_0)\); then \(G\) is perfect, since it is isomorphic to the cone of \(u\), which is strictly perfect. Hence \(\Parf(\mathcal C,\mathcal C_0)\) is indeed a triangulated sub-\(S\)-category of \(\Dpx(\mathcal C)\). The preceding argument also proves the parenthesized assertion (ii). It remains to prove the parenthesized assertion (i). Let \(X\in\Ob S\), and let
\[
E=(E_1\xrightarrow{a_1}E_2\xrightarrow{a_2}E_3\xrightarrow{a_3}E_1[+1]),
\quad
F=(F_1\xrightarrow{b_1}F_2\xrightarrow{b_2}F_3\xrightarrow{b_3}F_1[+1])
\]
be distinguished triangles of \(\Kpx^b(\mathcal C_{0,X})\). We must show that
\[
  \Hom_{\Tr(\Kpx^b(\mathcal C_0))}(E,F)
  \longrightarrow
  \Hom_{\Tr(\Dpx(\mathcal C))}(E,F)
  \tag{**}
\]
is an isomorphism. By the non-parenthesized assertion (i), we already know that \((**)\) is injective. Let \((u_i):E\to F\) be a morphism of triangles in \(\Dpx(\mathcal C)\). After localizing, the non-parenthesized assertion (i) permits us to suppose that \(u_i\), for \(1\leq i\leq3\), is an arrow of \(\Kpx^b(\mathcal C_0)\). The arrow
\[
  u_{i+1}a_i-b_iu_i,
\]
for \(1\leq i\leq3\) (with the convention \(u_4=u_1\)), is zero in \(\Dpx(\mathcal C)\), hence locally zero in \(\Kpx^b(\mathcal C_0)\) by the non-parenthesized assertion (i). In other words, \((u_i)\) locally defines a morphism of triangles in \(\Kpx^b(\mathcal C_0)\). This proves that \((**)\) is surjective and completes the proof of (4.10).

\textbf{Remark 4.10.1.} In (4.10), one may replace \(\Parf(\mathcal C,\mathcal C_0)\) by
\[
  \Dpx^*(\mathcal C)_{\mathcal C_0\text{-}\Parf}
  =\Dpx^*(\mathcal C)\cap\Parf(\mathcal C,\mathcal C_0)
  \qquad(*=-\text{ or }b),
\]
and nothing changes in the proof.

\textbf{Complement 4.11.} The triangulated structure of \(\Parf(\mathcal C,\mathcal C_0)\) can be made more precise in terms of perfect amplitude (4.7). If \(E\in\Ob\Dpx(\mathcal C)\) has perfect amplitude contained in \([a,b]\), then \(E[n]\) has perfect amplitude contained in \([a-n,b-n]\).

Let
\[
  E\longrightarrow F\longrightarrow G\longrightarrow E[1].
\]
be a distinguished triangle of \(\Dpx(\mathcal C)\). One may draw up a table analogous to (2.6), where the entries indicate perfect amplitudes and each row corresponds to an implication whose target is underlined:
\[
\begin{array}{ccc}
E&F&G\\[0.3em]
{}[a+1,b+1]&{}[a,b]&\underline{[a,b]}\\
{}[a,b]&\underline{[a,b]}&{}[a,b]\\
\underline{[a,b]}&{}[a,b]&{}[a-1,b-1].
\end{array}
\]
From (4.10) one immediately deduces:

\textbf{Corollary 4.12.} Suppose that \(S\) is a punctual site. For \(F\in\Ob\Dpx(\mathcal C)\) to have perfect amplitude relative to \(\mathcal C_0\) contained in \([a,b]\), it is necessary and sufficient that there exist a quasi-isomorphism \(F'\to F\), with \(F'\) strictly perfect and such that \(F'^i=0\) for \(i\notin[a,b]\). The canonical functors
\[
  \Kpx^b(\mathcal C_0)\longrightarrow\Parf(\mathcal C,\mathcal C_0)
\]
and
\[
  \Tr(\Kpx^b(\mathcal C_0))\longrightarrow\Tr(\Parf(\mathcal C,\mathcal C_0))
\]
are equivalences of categories.

\textbf{Remark 4.12.1.} When \(S\) is a punctual site, the objects of \(\mathcal C_0\) are projective. When \(\mathcal C_0\) contains all the projective objects of \(\mathcal C\), one says ``projective amplitude'' instead of ``\(\mathcal C_0\)-perfect amplitude.''



% Additional notation for Expose I continuation
\providecommand{\Cpx}{\mathrm C}
\providecommand{\Kpx}{\mathrm K}
\providecommand{\Dpx}{\mathrm D}
\providecommand{\Ho}{\mathrm H}
\providecommand{\Ob}{\operatorname{ob}}
\providecommand{\qis}{\xrightarrow{\sim}}
\providecommand{\RHom}{\mathrm R\!\Hom}
\providecommand{\Ext}{\operatorname{Ext}}
\providecommand{\Hom}{\operatorname{Hom}}
\providecommand{\Tr}{\operatorname{Tr}}

\providecommand{\Add}{\operatorname{Add}}
\providecommand{\Cart}{\operatorname{Cart}}
\providecommand{\perfamp}{\operatorname{perf.amp}}
\providecommand{\perfdim}{\operatorname{perf.dim}}
\providecommand{\toramp}{\operatorname{tor.amp}}
\providecommand{\torf}{\operatorname{torf}}
\providecommand{\injf}{\operatorname{injf}}
\providecommand{\rg}{\operatorname{rg}}
\providecommand{\is}{\mathrm{is}}
\providecommand{\Mod}{\operatorname{Mod}}


\section*{Exposé I, continued: perfect amplitude, finite Tor-dimension, and rank}

\textbf{Remark 4.12.2.} Let \(\mathcal C\) be an abelian category (which may be regarded as fibered over a punctual site), and let \(\mathcal C_0\) be a strictly full additive subcategory. Suppose that \(\mathcal C_0\) is stable under kernels of epimorphisms and is quasi-liftable in \(\mathcal C\). Then an acyclic object of \(\Kpx^b(\mathcal C_0)\) is not in general zero; hence there is no hope that the functor
\[
  \Kpx^b(\mathcal C_0)\longrightarrow \Dpx(\mathcal C)
\]
should be faithful. However, if \(\Kpx^{b,\mathrm{ac}}(\mathcal C_0)\) denotes the thick subcategory of \(\Kpx^b(\mathcal C_0)\) formed by the objects which are acyclic in \(\mathcal C\), the canonical functor
\[
  \Kpx^b(\mathcal C_0)/\Kpx^{b,\mathrm{ac}}(\mathcal C_0)
  \longrightarrow \Dpx(\mathcal C)
\]
is fully faithful.

Indications: one already knows from (2.7) that the functor
\[
  \Kpx^-(\mathcal C_0)/\Kpx^{-,\mathrm{ac}}(\mathcal C_0)
  \longrightarrow \Dpx^-(\mathcal C)
\]
is fully faithful. Show, for example with the aid of the criterion of [V], 4.2 b)(ii), that the functor
\[
  \Kpx^b(\mathcal C_0)/\Kpx^{b,\mathrm{ac}}(\mathcal C_0)
  \longrightarrow
  \Kpx^-(\mathcal C_0)/\Kpx^{-,\mathrm{ac}}(\mathcal C_0)
\]
is fully faithful.

\textbf{Lemma 4.13.} Let \(F\in\Ob \Dpx(\mathcal C_S)\), and let \(a,b,c\in\mathbb Z\) be such that \(a\leq b\) and \(a\leq c\). If \(F\) is of perfect amplitude contained in \([a,c]\) (4.7), and if \(\Ho^i(F)=0\) for \(i>b\), then \(F\) is of perfect amplitude contained in \([a,d]\), where \(d=\inf(b,c)\).

\emph{Proof.} We may suppose \(b\leq c\), and \(F\) strictly perfect with \(F^i=0\) for \(i\notin [a,c]\). Then \(F\) is isomorphic to its truncation \(F'\) defined by
\[
  F'^i=F^i\quad (i<b),\qquad F'^b=Z^b(F),\qquad F'^i=0\quad (i>b).
\]
But \(Z^b(F)\) admits a finite right resolution by objects of \(\mathcal C_0\), hence is locally an object of \(\mathcal C_0\), whence the assertion.

\textbf{Corollary 4.14.} Let \(F\in\Ob\mathcal C_S\) and \(n\in\mathbb N\). The following conditions are equivalent:
\begin{enumerate}[label=(\roman*)]
\item \(\perfamp(F)\subset[-n,0]\);
\item \(F\) locally admits a left resolution of length \(n\) by objects of \(\mathcal C_0\), that is, the set of objects \(X\) over \(S\) such that there exists an exact sequence
\[
  0\longrightarrow L^{-n}\longrightarrow L^{-n+1}\longrightarrow \cdots
  \longrightarrow L^0\longrightarrow F_X\longrightarrow 0,
\]
with \(L^i\in\Ob\mathcal C_{0,X}\) for all \(i\), is a refinement of \(S\).
\end{enumerate}

\emph{Proof.} This is an immediate consequence of (4.7) and (4.13).

\textbf{Definition 4.14.1.} We shall say that \(F\) is of perfect dimension \(\leq n\), and shall write \(\perfdim(F)\leq n\), if \(F\) satisfies the equivalent conditions of (4.14). We shall call the perfect dimension of \(F\) the smallest integer \(n\) such that \(\perfdim(F)\leq n\); if there is no such integer, we shall say that \(F\) is of infinite perfect dimension.

When \(S\) is the punctual site and \(\mathcal C_0\) is the full subcategory of \(\mathcal C\) formed by the projective objects, the perfect dimension is just the usual projective dimension.

\textbf{Proposition 4.15.} Let \(a,b,n\in\mathbb Z\), with \(a\leq b\) and \(n\geq0\), and let \(F\in\Ob\Dpx(\mathcal C)\) be such that \(F^i=0\) (respectively \(\Ho^i(F)=0\)) for \(i\notin[a,b]\). If, for every \(i\in[a,b]\), \(F^i\) (respectively \(\Ho^i(F)\)) is of perfect dimension \(\leq n+(i-a)\), then \(F\) is of perfect amplitude contained in \([a-n,b]\). Consequently, if for every \(i\in[a,b]\), \(F^i\) (respectively \(\Ho^i(F)\)) is perfect, then \(F\) is perfect.

\emph{Proof.} Proceed by induction on \(b\) (with \(a\) and \(n\) fixed), by truncating \(F\) and applying (4.11), cf. the proof of (2.11).

\textbf{Lemma 4.16.} Let \(a,b\in\mathbb Z\), \(a\leq b\), and let \(F\in\Ob\Dpx(\mathcal C)\) be such that \(F^i=0\) for \(i\notin[a,b]\). Suppose that \(\perfamp(F)\subset[a,b]\), and that \(F^i\in\Ob\mathcal C_0\) for \(i>a\). Then \(F^a\) is locally an object of \(\mathcal C_0\).

\emph{Proof.} There is an exact sequence, obtained by truncating \(F\),
\[
  0\longrightarrow F'\longrightarrow F\longrightarrow F^a[-a]\longrightarrow 0,
\]
with \(\perfamp(F')\subset[a+1,b+1]\). From (4.11) and (4.13) one deduces that \(\perfdim(F^a)=0\), i.e. that \(F^a\) is locally an object of \(\mathcal C_0\).

\textbf{Proposition 4.17.} Let \(F',F''\in\Ob\Dpx(\mathcal C_S)\), and let \(a,b\in\mathbb Z\), with \(a\leq b\). Consider the conditions:
\begin{enumerate}[label=(\roman*)]
\item \(\perfamp(F')\subset[a,b]\) and \(\perfamp(F'')\subset[a,b]\);
\item \(\perfamp(F'\oplus F'')\subset[a,b]\).
\end{enumerate}
Then (i) implies (ii). If \(\mathcal C_0\) is locally stable under direct factors, i.e. if the relation \(E'\oplus E''\in\Ob\mathcal C_0\), for \(E',E''\in\Ob\mathcal C_X\) and \(X\in\Ob S\), implies that locally \(E'\) and \(E''\) are objects of \(\mathcal C_0\), then (i) and (ii) are equivalent.

\emph{Proof.} The implication (i) \(\Rightarrow\) (ii) is a particular case of (4.11). Let us prove the converse. Put \(F=F'\oplus F''\). Since \(F\) is pseudo-coherent, \(F'\) and \(F''\) are also pseudo-coherent (2.12). On the other hand, the relation \(\Ho^i(F)=0\) for \(i\notin[a,b]\) implies \(\Ho^i(F')=\Ho^i(F'')=0\) for \(i\notin[a,b]\). By (2.10 a), we may therefore suppose, after localizing, that \(F'^i\) and \(F''^i\) are in \(\Ob\mathcal C_0\) for \(i\geq a+1\), and that \(F'^i=F''^i=0\) for \(i>b\). Then, by truncating, we get \(F'^i=F''^i=0\) for \(i<a\). Since \(\perfamp(F)\subset[a,b]\), it follows from (4.16) that \(F^a\) is locally an object of \(\mathcal C_0\). But \(F^a=F'^a\oplus F''^a\), hence \(F'^a\) and \(F''^a\) are locally objects of \(\mathcal C_0\), since \(\mathcal C_0\) is locally stable under direct factors.

\textbf{Proposition 4.18.} Let \((S',\mathcal C',\mathcal C'_0)\) be a triple satisfying the same hypotheses (4.0) as \((S,\mathcal C,\mathcal C_0)\). Let \(f:S'\to S\) be a morphism of sites, defined by \(f^{-1}:S\to S'\), and let
\[
  u:\Dpx^*(\mathcal C)\longrightarrow \Dpx^*(\mathcal C')\qquad (*=-,+,b)
\]
be a functor over \(f\) (cf. (2.13)). If \(u\) sends objects of \(\mathcal C_0\) to \(\mathcal C'_0\)-perfect complexes, then \(u\) defines a functor
\[
  \Dpx^*(\mathcal C)_{\mathcal C_0\text{-}\mathrm{parf}}
  \longrightarrow
  \Dpx^*(\mathcal C')_{\mathcal C'_0\text{-}\mathrm{parf}}.
\]

\emph{Proof.} Let \(F\in\Ob\Dpx^*(\mathcal C)_{\mathcal C_0\text{-}\mathrm{parf}}\). We show that \(u(F)\in\Ob\Dpx^*(\mathcal C')_{\mathcal C'_0\text{-}\mathrm{parf}}\). Since the question is local, one may suppose \(F\) strictly perfect. It remains only to use dévissage, taking into account the fact that \(\Dpx^*(\mathcal C')_{\mathcal C'_0\text{-}\mathrm{parf}}\) is a triangulated subcategory of \(\Dpx^*(\mathcal C')\) (4.10).

\textbf{Remark 4.18.1.} Suppose that there exist \(m,n\in\mathbb Z\), \(m\leq n\), such that \(F\in\Ob\mathcal C_0\) implies
\[
  \mathcal C'_0\text{-}\perfamp(u(F))\subset[m,n].
\]
Then, if \(F\in\Ob\Dpx^*(\mathcal C)\), one has
\[
  \mathcal C_0\text{-}\perfamp(F)
  \subset [a,b]
  \quad\Rightarrow\quad
  \mathcal C'_0\text{-}\perfamp(u(F))
  \subset [a+m,b+n].
\]
The proof is left as an exercise to the reader (use (4.11)).

\textbf{Corollary 4.19.} Under the hypotheses of (2.16), if \(E'\in\Ob\Dpx^-(A,S')\) is perfect as an object of \(\Dpx^-(A,S')\), and \(E\in\Ob\Dpx^-(A,S)\) is perfect, then
\[
  E'\overset{\mathrm L}{\otimes}_A E
\]
is perfect.

\emph{Proof.} We may suppose that \(E'\in\Ob\Dpx^-(A,S'_A)\), where \(S'\) is the final object of \(S'\), and apply (4.18) to
\[
  u=E'\overset{\mathrm L}{\otimes}_A:
  \Dpx^-(A,S)\longrightarrow \Dpx^-(A,S').
\]
In particular:

\textbf{Corollary 4.19.1.} \textbf{a)} The functor
\[
  \mathrm L f^*:\Dpx^-(A,S)\longrightarrow \Dpx^-(A,S')
\]
induces a functor
\[
  \mathrm L f^*:\Dpx^-(A,S)_{\mathrm{parf}}
  \longrightarrow
  \Dpx^-(A,S')_{\mathrm{parf}}.
\]

\textbf{b)} If \(A\) is commutative, the functor
\[
  \overset{\mathrm L}{\otimes}_A:
  \Dpx^-(S)\times_S\Dpx^-(S)
  \longrightarrow \Dpx^-(S)
\]
induces a functor
\[
  \overset{\mathrm L}{\otimes}_A:
  \Dpx^-(S)_{\mathrm{parf}}\times_S\Dpx^-(S)_{\mathrm{parf}}
  \longrightarrow \Dpx^-(S)_{\mathrm{parf}}.
\]

\textbf{Remark 4.19.2.} One can refine (4.19) in terms of perfect amplitude:
\[
  \perfamp(E)\subset[a,b],\qquad
  \perfamp(E')\subset[a',b']
  \quad\Rightarrow\quad
  \perfamp(E'\overset{\mathrm L}{\otimes}_A E)\subset[a+a',b+b'].
\]
In particular,
\[
  \perfamp(E)\subset[a,b]
  \quad\Rightarrow\quad
  \perfamp(\mathrm L f^*(E))\subset[a,b].
\]

\section*{5. Finite Tor-dimension and perfection}

\textbf{5.0.} The reader will surely have asked the question: what is missing from a pseudo-coherent complex \(E\) in order that it be perfect? Is it enough, for example, that \(E\) have bounded cohomology? The answer is no. Indeed, take for \(\mathcal C\) the category of modules over a commutative noetherian ring \(A\), and for \(\mathcal C_0\) the full subcategory of projective \(A\)-modules of finite type. Any \(A\)-module of finite type \(E\) is pseudo-coherent and has bounded cohomology, but \(E\) is perfect only if, in addition, it is of finite projective dimension (or, equivalently, of finite Tor-dimension). This example suggests that the difference between a pseudo-coherent complex and a perfect complex lies in a question of cohomological dimension. This is what we shall now examine in the case of the category of modules of a ringed topos (cf. (3.3)). We shall need a few preliminary sorties on the notion of finite Tor-dimension. We shall be brief, since this question already appears in the literature: cf. [H] II 4 and SGA 4 XVII 4.1.9.

\textbf{Proposition 5.1.} Let \((S,A)\) be a ringed topos, \(X\in\Ob S\), \(F\in\Ob\Dpx^-(A,X)\), and let \(m,n\in\mathbb Z\), with \(m\leq n\). The following conditions are equivalent:
\begin{enumerate}[label=(\roman*)]
\item \(F\) is isomorphic, in \(\Dpx^-(A,X)\), to a complex \(F'\) such that \(F'^i\) is flat for every \(i\), and zero for \(i\notin[m,n]\);
\item for every right \(A_X\)-module \(E\), one has
\[
  \Ho^i(E\overset{\mathrm L}{\otimes}_A F)=0\quad\text{for }i\notin[m,n].
\]
\end{enumerate}

\emph{Proof.} The implication (i) \(\Rightarrow\) (ii) is trivial; the implication (ii) \(\Rightarrow\) (i) follows, by truncation, from the following lemma.

\textbf{Lemma 5.1.1.} (cf. (4.16)). Let \(F\in\Ob\Dpx^-(A,X)\). Suppose that \(F\) satisfies hypothesis (ii) of (5.1), that \(F^i=0\) for \(i\notin[m,n]\), and that \(F^i\) is flat for \(i\ne m\). Then \(F^m\) is flat.

\emph{Proof.} There is a distinguished triangle
\[
\begin{tikzcd}
F' \arrow[rr] && F \arrow[dl]\\
& F^m[-m] \arrow[ul,"+1"] &
\end{tikzcd}
\]
where \(F'\) has flat components and is zero outside the interval \([m+1,n]\) (suppose \(m<n\)). Apply the functor \(E\overset{\mathrm L}{\otimes}_A\cdot\), for an arbitrary right \(A_X\)-module \(E\), and write the exact cohomology sequence.

\textbf{Definition 5.2.} An object \(F\) of \(\Dpx^-(A,S)\) satisfying the equivalent conditions of (5.1) will be said to be of flat amplitude contained in the interval \([m,n]\), and we shall write
\[
  \toramp(F)\subset[m,n].
\]
We shall say that \(F\) is of finite flat amplitude, or of finite Tor-dimension, if there exist integers \(m\) and \(n\) such that the preceding condition is satisfied.

\textbf{5.3.} Let \(X\in\Ob S\). We shall write \(\Dpx^-(A,X)_{\torf}\) for the full subcategory of \(\Dpx^-(A,X)\) formed by the objects of finite Tor-dimension. By criterion (ii) of (5.1), this is a triangulated subcategory of \(\Dpx^-(A,X)\). As \(X\) varies, the categories \(\Dpx^-(A,X)_{\torf}\) form a triangulated sub-\(S\)-category (1.1.4) of \(\Dpx^-(A,S)\), which we shall denote by \(\Dpx^-(S)_{A\text{-}\torf}\). The functor \(\overset{\mathrm L}{\otimes}_A\) induces an exact \(S\)-functor
\[
  \overset{\mathrm L}{\otimes}_A:
  \Dpx^b(S_A)\times_S \Dpx(A/S)_{\torf}
  \longrightarrow \Dpx^b(S_{\mathbb Z}).
\]

\textbf{5.4.} In order that a left \(A\)-module \(F\) be of flat amplitude contained in \([-n,0]\), \(n\in\mathbb N\), it is necessary and sufficient that \(F\) be of Tor-dimension \(\leq n\) in the ordinary sense, that is, that \(F\) admit a left resolution of length \(n\) by flat \(A\)-modules.

\textbf{5.5.} Let \(X\in\Ob S\), and let \(m,n\in\mathbb Z\), with \(m\leq n\). In order that \(F\in\Ob\Dpx^-(A,X)\) be of flat amplitude contained in \([m,n]\), it is necessary and sufficient that there exist a covering family \((X_i\to X)\) such that \(F|X_i\) is of flat amplitude contained in \([m,n]\). In particular, if \(S/X\) is of finite type (SGA 4 VI 1.7), and if \(F\) is locally of finite Tor-dimension, then \(F\) is of finite Tor-dimension.

\textbf{Proposition 5.6.} Under the hypotheses of (2.16), if \(E'\) is of flat amplitude contained in \([m',n']\) as an object of \(\Dpx^-(A,S')\), and if \(E\) is of flat amplitude contained in \([m,n]\), then
\[
  E'\overset{\mathrm L}{\otimes}_A E
\]
is of flat amplitude contained in \([m+m',n+n']\).

\emph{Proof.} This is immediate: use form (5.1 (i)) of the notion of finite Tor-dimension and the fact that, if \(E\) is a flat left \(A\)-module, then \(f^{-1}(E)\) is a flat left \(f^{-1}(A)\)-module (SGA 4 IV, App.).

In particular:

\textbf{Corollary 5.6.1.} \textbf{a)} The functor
\[
  \mathrm L f^*:\Dpx^-(A,S)\longrightarrow \Dpx^-(A,S')
\]
induces a functor
\[
  \mathrm L f^*:\Dpx(A,S)_{\torf}\longrightarrow \Dpx(A,S')_{\torf}.
\]

\textbf{b)} If \(A\) is commutative, the functor
\[
  \overset{\mathrm L}{\otimes}_A:
  \Dpx^-(S)\times_S\Dpx^-(S)\longrightarrow \Dpx^-(S)
\]
induces a functor
\[
  \overset{\mathrm L}{\otimes}_A:
  \Dpx(S)_{\torf}\times_S\Dpx(S)_{\torf}
  \longrightarrow \Dpx(S)_{\torf}.
\]

\textbf{Remark 5.6.2.} Let \(E\in\Ob\Dpx^-(A,S)\), where \(S\) is the final object of \(S\), and let \(m,n\in\mathbb Z\), with \(m\leq n\). If \(E\) is of flat amplitude contained in \([m,n]\), then so is \(E_s\), for every point \(s\) of \(S\). The converse is true if \(S\) has enough points; use form (5.1(ii)) of the notion of flat amplitude and the fact that \(\overset{\mathrm L}{\otimes}_A\) commutes with passage to the fibers.

\textbf{Exercises 5.7.} \textbf{a)} Suppose that \(A\) is commutative and that \((S',A')\) is faithfully flat over \((S,A)\) (cf. (2.16.2)). In order that \(E\in\Ob\Dpx^-(A,S)\) be of flat amplitude contained in \([m,n]\), it is necessary and sufficient that \(f^*(E)\) have that property.

\textbf{b)} Let \(F\in\Ob\Dpx^+(A,S)\). Show the equivalence of the following conditions:
\begin{enumerate}[label=(\roman*)]
\item \(F\) is isomorphic, in \(\Dpx^+(A,S)\), to a complex \(F'\) such that \(F'^i\) is injective for every \(i\) and zero for \(i\notin[m,n]\);
\item for every left \(A_S\)-module \(E\), one has \(\Ext^i(E,F)=0\) for \(i\notin[m,n]\).
\end{enumerate}
If \(F\) satisfies the preceding conditions, one says that \(F\) is of injective amplitude contained in \([m,n]\). The full subcategory of \(\Dpx^+(A,S)\) formed by the objects of finite injective amplitude is a triangulated subcategory, denoted \(\Dpx^+(A,S)_{\injf}\). Show that, if \(A\) is commutative, the functor
\[
  \mathrm R\Hom:\Dpx(S)^\circ\times\Dpx^+(S)\longrightarrow \Dpx(S)
\]
induces a functor
\[
  \mathrm R\Hom:\Dpx(S)^\circ_{\torf}\times\Dpx(S)_{\injf}
  \longrightarrow \Dpx(S)_{\injf}.
\]

We can now answer the question of (5.0) and state the main result of this section.

\textbf{Theorem 5.8.} Let \((S,A)\) be a ringed topos, \(E\in\Ob\Dpx(A,S)\), and let \(m,n\in\mathbb Z\), with \(m\leq n\). The following conditions are equivalent:
\begin{enumerate}[label=(\roman*)]
\item \(\perfamp(E)\subset[m,n]\) (4.8);
\item \(E\) is \((m-1)\)-pseudo-coherent (2.15) and \(\toramp(E)\subset[m,n]\) (5.2).
\end{enumerate}
Since ``perfect'' plainly implies ``pseudo-coherent'', one deduces:

\textbf{Corollary 5.8.1.} In order that \(E\) be perfect (4.8), it is necessary and sufficient that \(E\) be pseudo-coherent (2.15) and locally of finite Tor-dimension (5.2).

For the proof, we shall need a few lemmas.

\textbf{Lemma 5.8.2.} Let \(B\) be a ring of \(S\). If \(E\) is a left \(A\)-module, \(G\) a right \(B\)-module, and \(F\) an \(A\)-\(B\)-bimodule (left over \(A\), right over \(B\)), there is a canonical functorial homomorphism
\begin{equation}\tag{5.8.2.1}
  \Hom_B(F,G)\otimes_A E
  \longrightarrow
  \Hom_B(\Hom_A(E,F),G).
\end{equation}
It is an isomorphism for \(G\) injective and \(E\) of finite presentation.

\emph{Proof.} First define the arrow. By the adjunction formula dear to Cartan, it amounts to defining a homomorphism of \(B\)-modules:
\begin{equation}\tag{*}
  \Hom_B(F,G)\otimes_A E\otimes_{\mathbb Z_S}\Hom_A(E,F)
  \longrightarrow G.
\end{equation}
There is a canonical homomorphism of \(A\)-\(B\)-bimodules
\begin{equation}\tag{**}
  E\otimes_{\mathbb Z_S}\Hom_A(E,F)
  \longrightarrow F
\end{equation}
coming, by adjunction, from the identity arrow of \(\Hom_A(E,F)\). Similarly, there is a canonical homomorphism of \(B\)-modules
\begin{equation}\tag{***}
  \Hom_B(F,G)\otimes_A F\longrightarrow G.
\end{equation}
We then define (*) as the composite of (**) and (***). To see that the arrow (5.8.2.1) is an isomorphism when \(G\) is injective and \(E\) is of finite presentation, it is enough to note that, if \(G\) is injective, the source and the target of the arrow are right-exact functors in \(E\), and that the arrow induces the identity for \(E=A\).

Here is now a particular case of (5.8).

\textbf{Lemma 5.8.3.} Let \(E\) be a left \(A\)-module. The following conditions are equivalent:
\begin{enumerate}[label=(\roman*)]
\item \(E\) is locally a direct factor of a free module of finite type;
\item \(E\) is flat and of finite presentation.
\end{enumerate}

\emph{Proof} (after Bourbaki, Alg. comm., Chap. I, §2, Exer. 15). It is clear that (i) implies (ii). Let us prove (ii) \(\Rightarrow\) (i). By (1.3.1), it is enough to show that the functor \(\Hom_A(E,\cdot)\) is exact. Let
\[
  F\longrightarrow F''\longrightarrow 0
\]
be an exact sequence of left \(A\)-modules. We must show that the sequence
\[
  \Hom_A(E,F)\longrightarrow \Hom_A(E,F'')\longrightarrow 0
\]
is exact. For this it is enough to prove that, for every injective \(\mathbb Z_S\)-module \(I\), the sequence obtained by applying the functor \(\Hom_{\mathbb Z_S}(\cdot,I)\) is exact. But there is a commutative diagram, where the vertical arrows are the canonical arrows of (5.8.2):
\[
\begin{array}{ccccccccc}
0&\to&\Hom_{\mathbb Z_S}(F'',I)\otimes_A E&\to&\Hom_{\mathbb Z_S}(F,I)\otimes_A E\\
&&\downarrow&&\downarrow\\
0&\to&\Hom_{\mathbb Z_S}(\Hom_A(E,F''),I)&\to&\Hom_{\mathbb Z_S}(\Hom_A(E,F),I).
\end{array}
\]
The vertical arrows are isomorphisms by (5.8.2). The top row is exact because \(E\) is flat. Hence the bottom row is exact, q.e.d.

\emph{Proof of (5.8).} The implication (i) \(\Rightarrow\) (ii) follows trivially from the definitions. Let us prove (ii) \(\Rightarrow\) (i). Since the question is local, one may suppose \(E\) strictly \((m-1)\)-pseudo-coherent. Since, on the other hand, \(E\) is acyclic in degrees \(<m\), the truncation arrow
\[
\begin{array}{ccccccc}
\cdots&\to&E^{m-1}&\to&E^m&\to&E^{m+1}\to\cdots\\
&&\downarrow&&\downarrow&&\downarrow\\
\cdots&\to&0&\to&E^m/B^m&\to&E^{m+1}\to\cdots
\end{array}
\]
is a quasi-isomorphism. Let \(E'\) be the truncated complex. By construction, \(E'\) is of flat amplitude contained in \([m,n]\), and \(E'^i\) is free of finite type for \(i>m\). Therefore, by (5.1.1), \(E'^m\) is flat. On the other hand there is a distinguished triangle
\[
\begin{tikzcd}
E'' \arrow[rr] && E' \arrow[dl]\\
& E'^m[-m] \arrow[ul,"+1"] &
\end{tikzcd}
\]
where \(E''\) is strictly perfect. By (2.6), \(E'^m[-m]\) is \((m-1)\)-pseudo-coherent; i.e. by (2.9), \(E'^m\) is of finite presentation. We conclude from (5.8.3) that \(E'^m\) is locally a direct factor of a free module of finite type, hence that \(\perfamp(E)\subset[m,n]\).

\textbf{Remarks 5.8.4.} \textbf{a)} Using (2.5), one recovers the fact that perfect complexes form a triangulated subcategory of \(\Dpx(A,S)\).

\textbf{b)} A perfect complex, even with bounded cohomology, is not necessarily of finite Tor-dimension. It is so, however, if \(S\) is of finite type (5.5). Write \(\Dpx(A,S)_{\mathrm{parf}}\) for the full sub-\(S\)-category of \(\Dpx^b(A,S)\) formed by perfect complexes of finite Tor-dimension. It follows at once from (2.16.1 b) and (5.3) that, if \(A\) is commutative, the functor \(\overset{\mathrm L}{\otimes}_A\) induces a functor
\[
  \overset{\mathrm L}{\otimes}_A:
  \Dpx(S)_{\mathrm{parf}}\times_S \Dpx^b(S)_{\mathrm{coh}}
  \longrightarrow \Dpx^b(S)_{\mathrm{coh}}.
\]
This formula is the starting point for the theory of intersections on schemes. Note, on the other hand, that the derived tensor product of two pseudo-coherent complexes with bounded cohomology is not, in general, of bounded cohomology.

\textbf{Corollary 5.9.} Suppose that \(S=(\mathrm{Ens})\), so that \(A\) is an ordinary ring, and that every left \(A\)-module is of finite Tor-dimension; this is the case, for example, if \(A\) is of finite global cohomological dimension, cf. (MULT VI 2.6) and (EGA \(0_{\mathrm{IV}}\) 17.2.8). Then
\[
\text{(i)}\quad \Dpx^b(A,S)=\Dpx(A,S)_{\torf},
\qquad
\text{and}
\qquad
\text{(ii)}\quad \Dpx^b(A,S)_{\mathrm{coh}}=\Dpx(A,S)_{\mathrm{parf}}.
\]
In other words, in order that \(E\in\Ob\Dpx(A,S)\) be of finite Tor-dimension (respectively finite Tor-dimension and perfect), it is necessary and sufficient that \(E\) have bounded cohomology (respectively bounded and pseudo-coherent cohomology).

\emph{Proof.} By (5.8), it is enough to prove (i). Let \(E\in\Ob\Dpx^b(A,S)\). We may suppose, after extending \(E\) by \(0\), that \(E\in\Ob\Dpx^b(A,S)\), where \(S\) is the final object of \(S\). By induction on the number of non-zero cohomology objects of \(E\), one is reduced, by dévissage, to the case where \(E\) is an ordinary \(A\)-module; then one wins.

More generally, one may state:

\textbf{Corollary 5.10.} Let \(X\in\Ob S\). Suppose that \(S/X\) has enough points (SGA 4 IV 6.4.1) and that, for every point \(s\) of \(S/X\), every left \(A_s\)-module is of finite Tor-dimension. Then
\[
  \Dpx^b(A,X)_{\mathrm{coh}}=\Dpx^b(A,X)_{\mathrm{parf}}.
\]
Equivalently, in order that \(E\in\Ob\Dpx^b(A,X)\) be pseudo-coherent, it is necessary and sufficient that \(E\) be perfect.

\emph{Proof.} Let \(E\in\Ob\Dpx^b(A,X)_{\mathrm{coh}}\) and let \(s\) be a point of \(S/X\). One must show that \(E\) is perfect in a neighbourhood of \(s\), that is, that there exists an \(s\)-pointed object \(U\) of \(S\) over \(X\) such that \(E|U\) is perfect. Now, since the functor \(s^*\) is exact, (2.16.1 a) gives \(E_s\in\Ob\Dpx^b(A_s)_{\mathrm{coh}}\). Hence, by (5.9), \(E_s\) is perfect. The assertion is therefore a consequence of (3.7.2) and of the following lemma.

\textbf{Lemma 5.10.1.} Let \(s\) be a point of \(S\), and let \(F\in\Ob\Dpx(A_s)_{\mathrm{parf}}\). There exist an \(s\)-pointed object \(U\) of \(S\) and \(F'\in\Ob\Dpx(A_U)_{\mathrm{parf}}\) such that \(F'_s=F\).

\emph{Proof.} This is evident if \(F\) is a free \(A_s\)-module of finite type. Suppose now that \(F\) is projective of finite type, hence the image of a projector
\[
  p:L\longrightarrow L,
\]
where \(L\) is free of finite type. There is then a neighbourhood \(U\) of \(s\), a free \(A_U\)-module \(L'\), and a homomorphism \(p':L'\to L'\) such that \(p'_s=p\). Since \(p_s^2=p_s\), we may suppose, after shrinking \(U\), that \(p'^2=p'\). Thus \(F'=\Im(p')\) is a direct factor of \(L'\) and satisfies \(F'_s=F\). The case where \(F\) is strictly perfect, that is, a bounded complex made up of projective modules of finite type, follows immediately.

\textbf{Examples 5.11.} Corollary (5.10) applies in particular when \((S,A)\) is defined by a locally ringed space in regular local rings, for example a regular scheme or a smooth complex analytic space over \(\mathbb C\).

\section*{6. Rank of a perfect complex}

\textbf{6.0.} Let \((X,\mathcal O_X)\) be a ringed space in local rings. It is well known that one can attach, to every locally free sheaf of finite type \(L\) on an open set \(U\) of \(X\), a function \(\rg_U(L)\), locally constant on \(U\) and with integral values, called the rank of \(L\). The family of functions \(\rg_U(L)\) is characterized by the following properties:
\begin{enumerate}[label=(\roman*)]
\item if \(V\subset U\) is an inclusion of open sets, then \(\rg_V(L|V)=\rg_U(L)|V\);
\item for every exact sequence of locally free \(\mathcal O_U\)-modules of finite type
\[
  0\longrightarrow L'\longrightarrow L\longrightarrow L''\longrightarrow 0,
\]
one has \(\rg_U(L)=\rg_U(L')+\rg_U(L'')\);
\item \(\rg_X(\mathcal O_X)=1\).
\end{enumerate}
We propose in this section to extend the notion of rank to perfect complexes.

\textbf{Definition 6.1} (cf. SGA 5 VIII 2). Let \(\mathcal C\) be an additive (respectively triangulated) category, and let \(F\) be an abelian group. A function
\[
  f:\Ob(\mathcal C)\longrightarrow F
\]
is said to be an additive function from \(\mathcal C\) to \(F\) if, for all \(L,L',L''\in\Ob(\mathcal C)\) such that \(L\simeq L'\oplus L''\) (respectively such that there is a distinguished triangle
\[
\begin{tikzcd}
L' \arrow[rr] && L \arrow[dl]\\
& L'' \arrow[ul] &
\end{tikzcd}
\]
), one has
\[
  f(L)=f(L')+f(L'').
\]
The additive functions from \(\mathcal C\) to \(F\) form an abelian group, denoted \(\Add(\mathcal C,F)\).

\textbf{6.2.} Let \(\mathcal C\) be a triangulated category, and let \(f\) be an additive function from \(\mathcal C\) to an abelian group \(F\). Then \(f\) induces an additive function on the underlying additive category \(\mathcal C^{\mathrm{ad}}\) of \(\mathcal C\); hence an inclusion
\[
  \Add(\mathcal C,F)\subset \Add(\mathcal C^{\mathrm{ad}},F).
\]
In particular, \(f(0)=0\) and \(f(L)=f(M)\) if \(L\simeq M\). On the other hand,
\[
  f(L[n])=(-1)^n f(L)
\]
for all \(L\in\Ob(\mathcal C)\) and \(n\in\mathbb Z\).

\textbf{6.3.} Let \(\mathcal C\) be an additive (respectively triangulated) category. The group \(\Add(\mathcal C,F)\) depends functorially on \(F\). The functor \(\Add(\mathcal C,\cdot)\) is representable (cf. SGA 5 VIII 2) by an abelian group, denoted \(k(\mathcal C)\), endowed with an additive application
\[
  \cl_{\mathcal C}:\Ob(\mathcal C)\longrightarrow k(\mathcal C).
\]
If \(\mathcal C\) is additive, one takes for \(k(\mathcal C)\) the symmetrized monoid of isomorphism classes of objects of \(\mathcal C\). If \(\mathcal C\) is triangulated, one takes for \(k(\mathcal C)\) the quotient of the free abelian group generated by \(\Ob(\mathcal C)\) by the relations identifying \(L\) with \(L'+L''\) whenever there is a distinguished triangle
\[
  L'\longrightarrow L\longrightarrow L''\longrightarrow L'[1].
\]

The group \(k(\mathcal C)\) depends functorially on \(\mathcal C\), with respect to additive (respectively exact) functors. If \(u,u':\mathcal C\to\mathcal D\) are isomorphic additive (respectively exact) functors, then \(k(u)=k(u')\).

If \(\mathcal C\) is a triangulated category, one has a canonical epimorphism
\[
  k(\mathcal C^{\mathrm{ad}})\longrightarrow k(\mathcal C)
\]
(cf. (6.2)).

\textbf{Lemma 6.4.} Let \(\mathcal C\) be an additive category. The canonical functor
\[
  \mathcal C\longrightarrow \Kpx^b(\mathcal C)
\]
induces an isomorphism, functorial in \(\mathcal C\),
\[
  k(\mathcal C)\xrightarrow{\sim} k(\Kpx^b(\mathcal C)).
\]

\emph{Proof.} This is an immediate consequence of the following formula, which follows trivially from (6.2):
\begin{equation}\tag{6.4.1}
  f(L)=\sum_i (-1)^i f(L^i)
\end{equation}
for \(L\in\Ob\Kpx^b(\mathcal C)\) and \(f\) an additive function on \(\Kpx^b(\mathcal C)\).

\emph{Remark.} We shall not use here the usual notation \(K(\mathcal C)\), since \(K(\mathcal C)\) already denotes, when \(\mathcal C\) is additive, the category of complexes of \(\mathcal C\) ``up to homotopy.''

\textbf{Definition 6.5.} Let \(S\) be a category, \(\mathcal C\) an additive (respectively triangulated) \(S\)-category (1.1.1), and \(F\) an abelian presheaf on \(S\). A family of functions
\[
  f_X:\Ob(\mathcal C_X)\longrightarrow F(X),
  \qquad X\in\Ob S,
\]
is said to be an additive \(S\)-function from \(\mathcal C\) to \(F\) if the following conditions are satisfied:
\begin{enumerate}[label=(\roman*)]
\item for every \(X\in\Ob S\), \(f_X\) is an additive function (6.1);
\item for every arrow \(u:X\to Y\) of \(S\), the following diagram is commutative:
\[
\begin{tikzcd}
\Ob\mathcal C_Y \arrow[r,"f_Y"] \arrow[d,"u^*"'] & F(Y) \arrow[d,"F(u)"]\\
\Ob\mathcal C_X \arrow[r,"f_X"'] & F(X).
\end{tikzcd}
\]
\end{enumerate}
The additive \(S\)-functions from \(\mathcal C\) to \(F\) form an abelian group, denoted \(\Add_S(\mathcal C,F)\).

\textbf{Remarks 6.6.} \textbf{a)} Let \(\mathcal C^{\is}\) denote the fibered \(S\)-category whose fiber over \(X\in\Ob S\) is the category with the same objects as \(\mathcal C_X\) but whose arrows are the isomorphisms of \(\mathcal C_X\). Let \(\underline F\) denote the fibered \(S\)-category whose fiber over \(X\in\Ob S\) is the discrete category (i.e. the only arrows are identities) \(F(X)\), the inverse-image functor by \(u\in\mathrm{Fl}(S)\) being \(F(u)\). Then an additive \(S\)-function from \(\mathcal C\) to \(F\) is simply a Cartesian \(S\)-functor from \(\mathcal C^{\is}\) to \(\underline F\), inducing on each fiber an additive function in the sense of (6.1). It will be convenient for us to interpret additivity as follows. Suppose, for definiteness, that \(\mathcal C\) is triangulated. The category
\[
  \Cart_S(\mathcal C^{\is},\underline F)
\]
(respectively \(\Cart_S(\Tr(\mathcal C)^{\is},\underline F)\) (1.1.2)) of Cartesian \(S\)-functors from \(\mathcal C^{\is}\) (respectively \(\Tr(\mathcal C)^{\is}\)) to \(\underline F\) is the discrete category associated to an abelian group, and there is a homomorphism
\[
  \rho:\Cart_S(\mathcal C^{\is},\underline F)
  \longrightarrow
  \Cart_S(\Tr(\mathcal C)^{\is},\underline F),
\]
defined by
\[
  \rho(f)(T)=f(L)-f(L')-f(L'')
\]
for \(f\in\Cart_S(\mathcal C^{\is},\underline F)\) and \(T=(L'\to L\to L''\to L'[1])\in\Ob\Tr(\mathcal C)\). To say that \(f\in\Cart_S(\mathcal C^{\is},\underline F)\) is additive means that \(\rho(f)=0\), or equivalently that one has a functorial canonical exact sequence
\begin{equation}\tag{6.6.1}
  0\longrightarrow \Add_S(\mathcal C,F)
  \longrightarrow \Cart_S(\mathcal C^{\is},\underline F)
  \xrightarrow{\rho}
  \Cart_S(\Tr(\mathcal C)^{\is},\underline F).
\end{equation}

\textbf{b)} Let \(u:X\to Y\) be an arrow of \(S\). The inverse-image functors by \(u\), being all isomorphic to one another, define by (6.3) a unique homomorphism
\[
  k(Y)\longrightarrow k(X).
\]
Consequently, \(X\mapsto k(X)\) defines an abelian presheaf on \(S\), denoted \(k(\mathcal C)\). An additive \(S\)-function from \(\mathcal C\) to \(F\) is then simply interpreted as a homomorphism of presheaves
\[
  k(\mathcal C)\longrightarrow F,
\]
or in other words
\[
  \Add_S(\mathcal C,F)=\Hom(k(\mathcal C),F).
\]
The presheaf \(k(\mathcal C)\) of course depends functorially on \(\mathcal C\) with respect to additive (respectively exact) \(S\)-functors.

\textbf{Proposition 6.7.} Let \(S\) be a site, \(\mathcal C\) a flat abelian \(S\)-category, and \(\mathcal C_0\) a strictly full additive sub-\(S\)-category of \(\mathcal C\) satisfying the hypotheses (4.0). Then the morphism of abelian presheaves
\[
  k(\mathcal C_0)\longrightarrow k(\Parf(\mathcal C,\mathcal C_0))
\]
induced by the inclusion of \(\mathcal C_0\) in \(\Parf(\mathcal C,\mathcal C_0)\), defines an isomorphism on the associated sheaves. In other words, for every abelian sheaf \(F\) on \(S\), the restriction homomorphism
\[
  \Add_S(\Parf(\mathcal C,\mathcal C_0),F)
  \longrightarrow
  \Add_S(\mathcal C_0,F)
\]
is an isomorphism.

\emph{Proof.} Let \(F\) be an abelian sheaf on \(S\). We show that the arrow
\[
  \Add_S(\Parf(\mathcal C,\mathcal C_0),F)
  \longrightarrow
  \Add_S(\mathcal C_0,F)
\]
is an isomorphism. It factors as
\[
\Add_S(\Parf(\mathcal C,\mathcal C_0),F)
 \xrightarrow{a}
\Add_S(\Kpx^b(\mathcal C_0),F)
 \xrightarrow{b}
\Add_S(\mathcal C_0,F).
\]
By (6.4), the arrow \(b\) is an isomorphism. We are reduced to proving that \(a\) is an isomorphism.

Consider the following commutative diagram, whose rows are exact by (6.6.1), and whose vertical arrows are the restriction arrows:
\[
\begin{array}{ccccccccc}
0&\to&\Add_S(\Parf(\mathcal C),F)&\to&\Cart_S(\Parf(\mathcal C)^{\is},\underline F)&\to&\Cart_S(\Tr(\Parf(\mathcal C))^{\is},\underline F)\\
&&\downarrow a&&\downarrow&&\downarrow\\
0&\to&\Add_S(\Kpx^b(\mathcal C_0),F)&\to&\Cart_S(\Kpx^b(\mathcal C_0)^{\is},\underline F)&\to&\Cart_S(\Tr(\Kpx^b(\mathcal C_0))^{\is},\underline F).
\end{array}
\]
Since \(F\) is a sheaf, \(\underline F\) is a stack (G II 2.2.4); hence, by (4.10), the two right vertical arrows are isomorphisms. The same is therefore true of \(a\), q.e.d.

\textbf{Remark 6.7.1} (cf. 4.10.1). In (6.7), one may replace \(\Parf(\mathcal C,\mathcal C_0)\) by \(\Dpx^*(\mathcal C)_{\mathcal C_0\text{-}\mathrm{parf}}\) (\(*=-\) or \(b\)).

\textbf{6.8.} Let \((S,A)\) be a ringed topos, and let \(\mathcal C\) and \(\mathcal C_0\) be the fibered \(S\)-categories considered in (4.8). There are homomorphisms of abelian presheaves
\[
  k(\mathcal C_0)
  \xrightarrow{\alpha}
  k(\Dpx^-(A,S)_{\mathrm{parf}})
  \xrightarrow{\beta}
  k(\Parf(A,S)),
\]
which, by (6.7), induce isomorphisms on the associated sheaves. The homomorphism \(\alpha\) is compatible with inverse images: that is, if
\[
  u:(S',A')\longrightarrow(S,A)
\]
is a morphism of ringed topoi, there is a commutative diagram, whose vertical arrows are induced by the functor \(\mathrm L u^*\) (cf. (4.19.1)):
\[
\begin{tikzcd}
 k(\mathcal C_0) \arrow[r] \arrow[d] & k(\Dpx^-(A,S)_{\mathrm{parf}}) \arrow[d]\\
 k(\mathcal C'_0) \arrow[r] & k(\Dpx^-(A',S')_{\mathrm{parf}}),
\end{tikzcd}
\]
where \(\mathcal C'_0\) is defined analogously to \(\mathcal C_0\). If \(A\) is commutative, the derived tensor product endows \(k(\mathcal C_0)\) and \(k(\Dpx^-(A,S)_{\mathrm{parf}})\) with structures of presheaves of rings (cf. (4.19.1)), and \(\alpha\) is a homomorphism of presheaves of rings.

\textbf{6.9.} In the situation of (6.8), suppose that \((S,A)\) is locally ringed (cf. (2.15.1)). Then, for every \(X\in\Ob S\), \(\mathcal C_{0,X}\) is the category of locally free \(A_X\)-modules of finite type. Let \(\mathcal C_1\) be the fibered \(S\)-category defined by
\[
  X\longmapsto\text{the category of free }A_X\text{-modules of finite type}.
\]
It is clear that \(\mathcal C_1\) satisfies the hypotheses (4.0). An object of \(\Dpx(S)\) is perfect if and only if it is perfect relative to \(\mathcal C_1\). There are homomorphisms of abelian presheaves
\[
  k(\mathcal C_1)\longrightarrow k(\mathcal C_0)\longrightarrow
  k(\Dpx^-(S)_{\mathrm{parf}})\longrightarrow k(\Parf(S))
\]
(the two on the left being homomorphisms of presheaves of rings), which, by (6.7), induce isomorphisms on the associated sheaves. We shall show that the associated sheaves are in fact canonically isomorphic to the constant sheaf \(\mathbb Z_S\).

\textbf{Lemma 6.10.} Let \((S,A)\) be a locally ringed topos, let \(U\) be an object of \(S\) distinct from the empty object, and let \(p,q\in\mathbb N\). If
\[
  A_U^p\simeq A_U^q,
\]
then \(p=q\).

\emph{Proof.} By an argument similar to that given in the proof of (2.15.1), one is reduced to the case where \((S,A)\) is the locally ringed topos defined by a prescheme; in this case the assertion is trivial: take a point in \(U\).

\textbf{Definition 6.11.} The rank, denoted \(\rg\), is the additive \(S\)-function from \(\mathcal C_1\) to \(\mathbb Z_S\) defined by
\[
  \rg(L)=p\quad\text{if } L\simeq A_X^p,
  \qquad X\ne\varnothing,
\]
\[
  \rg(L)=0\quad\text{if } L\in\Ob\mathcal C_{1,X},
  \qquad X=\varnothing.
\]
By (6.7), the rank function extends uniquely to an additive \(S\)-function from \(\Parf(S)\) to \(\mathbb Z_S\), which we shall again call rank and denote by \(\rg\).

\textbf{Proposition 6.12.} \textbf{a)} For \(L,M\in\Ob\Dpx^-(X)_{\mathrm{parf}}\), with \(X\in\Ob S\), one has
\[
  \rg(L\overset{\mathrm L}{\otimes} M)=\rg(L)\,\rg(M).
\]

\textbf{b)} Let
\[
  u:(S',A')\longrightarrow(S,A)
\]
be a morphism of locally ringed topoi. For \(M\in\Ob\Dpx^-(S)_{\mathrm{parf}}\), one has
\[
  \rg(\mathrm L u^*(M))=u^*(\rg(M)).
\]
In other words, the rank homomorphism
\[
  \rg:k(\Dpx^-(S)_{\mathrm{parf}})\longrightarrow \mathbb Z_S
\]
is a homomorphism of presheaves of rings, compatible with inverse images.



% Additional notation for Exposé I
\providecommand{\Cpx}{\mathrm C}
\providecommand{\Kpx}{\mathrm K}
\providecommand{\Dpx}{\mathrm D}
\providecommand{\Ho}{\mathrm H}
\providecommand{\Ob}{\operatorname{ob}}
\providecommand{\Hom}{\operatorname{Hom}}
\providecommand{\RHom}{\mathrm R\!\operatorname{Hom}}
\providecommand{\Ext}{\operatorname{Ext}}
\providecommand{\Tor}{\operatorname{Tor}}
\providecommand{\Tr}{\operatorname{Tr}}
\providecommand{\Add}{\operatorname{Add}}
\providecommand{\Cart}{\operatorname{Cart}}
\providecommand{\Parf}{\operatorname{Parf}}
\providecommand{\coh}{\operatorname{coh}}
\providecommand{\amp}{\operatorname{amp}}
\providecommand{\cl}{\operatorname{cl}}
\providecommand{\rg}{\operatorname{rg}}
\providecommand{\Id}{\operatorname{Id}}
\providecommand{\Ker}{\operatorname{Ker}}
\providecommand{\Im}{\operatorname{Im}}
\providecommand{\cH}{\mathcal H}
\providecommand{\Zs}{\mathbb Z_S}
\providecommand{\A}{\mathcal A}
\providecommand{\cC}{\mathcal C}
\providecommand{\cCzero}{\mathcal C_0}
\providecommand{\otL}{\mathop{\otimes}^{\mathrm L}}
\providecommand{\tamp}{\operatorname{tor.amp}}
\providecommand{\pamp}{\operatorname{parf.amp}}
\providecommand{\pdim}{\operatorname{parf.dim}}
\providecommand{\bon}{\mathrm b}
\providecommand{\parf}{\operatorname{parf}}
\providecommand{\parff}{\operatorname{parf}}
\providecommand{\ac}{\mathrm{ac}}


% Additional notation for Exposé I, completion pages 154--166
\providecommand{\End}{\operatorname{End}}
\providecommand{\tr}{\operatorname{tr}}

\emph{Proof.} It is enough to show that the homomorphism induced by rank on the sheaf associated to \(k(\Dpx^-(S)_{\Parf})\) is a homomorphism of sheaves of rings compatible with inverse images. By (6.9), one may replace \(k(\Dpx^-(S)_{\Parf})\) by \(k(\mathcal C_1)\). We are thus reduced to verifying a) and b) for free modules of finite type, which is trivial.

\textbf{6.13.} With \((S,A)\) still denoting a locally ringed topos, let \(X\in\Ob S\) and \(n\in H^0(X,\mathbb Z_S)\). We define an object of \(\Parf(X)\), denoted \(A_X^n\), as follows. If \(n\) is constant of value an integer \(n\geq0\) (respectively \(n<0\)), put \(A_X^n=A_X^{\oplus n}\) (respectively \(A_X^n=A_X^{-n}[1]\)). In the general case, decompose \(X\) into disjoint pieces \(X_i\) on which \(n|X_i\) is constant, and define \(A_X^n\) as the unique object of \(\Parf(X)\) inducing \(A_{X_i}^{n|X_i}\) on each \(X_i\). More generally, for \(L\in\Ob\Parf(X)\), one may define \(L^n\) by putting \(L^n=A_X^n\otimes_A L\).

The map
\[
  i_X:H^0(X,\mathbb Z_S)\longrightarrow k(\Parf(X)),
  \qquad i_X(n)=\cl(A_X^n),
\]
is a homomorphism of rings, with \(i_X(1)=1\), and defines, as \(X\) varies, a homomorphism of presheaves of rings
\[
  i:\mathbb Z_S\longrightarrow k(\Parf(S)).
\]

\textbf{Proposition 6.14.} Under the hypotheses of (6.13), one has
\[
  \rg\circ i=\Id(\mathbb Z_S).
\]
In other words, \(k(\Parf(S))\), considered as a presheaf of \(\mathbb Z_S\)-algebras by means of \(i\), is augmented toward \(\mathbb Z_S\) by the rank. Moreover, the rank
\[
  \rg:k(\Parf(S))\longrightarrow \mathbb Z_S
\]
defines an isomorphism on the associated sheaves.

\emph{Proof.} The first assertion follows trivially from the relations
\[
  \cl(L[1])=-\cl(L),\qquad \cl(L^n)=n\cl(L),
  \qquad L\in\Ob\Parf(S).
\]
Let us prove the second. First observe that \(i\) factors in the evident way through \(k(\mathcal C_1)\). Since \(\rg\circ i\) is the identity, we are reduced, by (6.9), to proving that the homomorphism
\[
  i\circ\rg:k(\mathcal C_1)\longrightarrow k(\mathcal C_1)
\]
induces an isomorphism on the associated sheaves; equivalently, for \(L\in\Ob\mathcal C_1\), that \(i(\rg(L))\) is locally isomorphic to \(L\). This is immediate.

One deduces, by (6.6 b):

\textbf{Corollary 6.15.} Under the hypotheses of (6.13), for every abelian sheaf \(F\) on \(S\), the group of \(S\)-additive functions from \(\Parf(S)\) to \(F\) is canonically identified with \(H^0(S,F)\), where \(S\) is the final object of \(S\). The rank function corresponds, for \(F=\mathbb Z_S\), to the unit section of \(\mathbb Z_S\).

\section*{7. Duality of perfect complexes}

\textbf{7.0.} In this section, \((S,A)\) denotes a topos ringed in commutative rings, and \(S\) also denotes the final object of \(S\). We propose to generalize to perfect complexes the duality notions that are well known for locally free sheaves of finite type.

\textbf{Proposition 7.1.} Let \(E\in\Ob\Dpx^-(S)_{\Parf}\) and \(F\in\Ob\Dpx^+(S)\).

\textbf{a)} If \(F\) has locally bounded cohomology, then so does \(\RHom(E,F)\). If \(F\) is pseudo-coherent (respectively locally of finite Tor-dimension, respectively perfect), then the same is true of \(\RHom(E,F)\).

\textbf{b)} Suppose, moreover, that \(E\) is of finite Tor-dimension, a condition satisfied if \(S\) is of finite type. Then, if \(F\) has bounded cohomology (respectively finite Tor-dimension), the same is true of \(\RHom(E,F)\). More precisely, let \(a,b,m,n\) be integers such that
\[
  \tamp(E)\subset[a,b]
\]
(or, equivalently by (5.8), \(\pamp(E)\subset[a,b]\)), and such that \(\Ho^i(F)=0\) for \(i\notin[m,n]\) (respectively \(\tamp(F)\subset[m,n]\)). Then
\[
  \Ho^i(\RHom(E,F))=0\quad(i\notin[m-b,n-a]),
\]
respectively
\[
  \tamp(\RHom(E,F))\subset[m-b,n-a].
\]

\emph{Proof.} Assertion a) is local, so we may suppose \(E\) strictly perfect, and then conclude by devissage on \(E\), reducing to \(E=A\). For assertion b), we may suppose \(E\) strictly perfect with \(E^i=0\) for \(i\notin[a,b]\), and \(F^i=0\) for \(i\notin[m,n]\) (respectively \(F^i\) flat for every \(i\) and \(F^i=0\) for \(i\notin[m,n]\)). It is then enough to apply the following lemma.

\textbf{Lemma 7.1.1.} Let \(E\in\Ob\Dpx^-(S)\) and \(F\in\Ob\Dpx^+(S)\). If \(E\) is strictly pseudo-coherent, that is, in the sense of (2.1), with components direct factors of free finite-type modules, then the canonical arrow of \(\Dpx^+(S)\)
\[
  \Hom^\bullet(E,F)\longrightarrow \RHom(E,F)
\]
is an isomorphism.

\emph{Proof.} The spectral sequence
\[
  E_2^{p0}=\Ho^p(\Hom^\bullet(E,F))
  \Longrightarrow \Ext^p(E,F)
\]
where \(\Ext^q(E^i,F^j)=0\) for every \(q>0\) by (1.3.1), degenerates. One obtains isomorphisms
\[
  E_2^{p0}=\Ho^p(\Hom^\bullet(E,F))
  \xrightarrow{\sim}\Ext^p(E,F)=\Ho^p(\RHom(E,F)),
\]
which prove the assertion.

\textbf{Proposition 7.1.2.} Let
\[
  f:(S',A')\longrightarrow(S,A)
\]
be a morphism of ringed topoi. Let \(E\in\Ob\Dpx^-(S)\) and \(F\in\Ob\Dpx^b(S)\) be such that \(\RHom(E,F)\in\Ob\Dpx^b(S)\) and \(\mathrm Lf^*(F)\in\Ob\Dpx^b(S')\). There is a functorial canonical arrow
\[
  \mathrm Lf^*\RHom(E,F)
  \longrightarrow
  \RHom(\mathrm Lf^*E,\mathrm Lf^*F).
\]
This is an isomorphism if \(E\) is perfect.

\emph{Proof.} By the adjunction of \(\mathrm Lf^*\) and \(\mathrm Rf_*\), it is enough to define an arrow
\[
  \RHom(E,F)
  \longrightarrow
  \mathrm Rf_*\RHom(\mathrm Lf^*E,\mathrm Lf^*F).
\]
For this we use the adjunction arrow \(F\to\mathrm Rf_*\mathrm Lf^*F\) and the duality isomorphism
\[
  \mathrm Rf_*\RHom(\mathrm Lf^*E,\mathrm Lf^*F)
  \xrightarrow{\sim}
  \RHom(E,\mathrm Rf_*\mathrm Lf^*F).
\]
To see that the arrow of (7.1.2) is an isomorphism for \(E\) perfect, one reduces trivially, by localization and devissage, to the case \(E=A\), where it is a tautology.

\textbf{Proposition 7.2.} There exists (cf. SGA 5 I 1.6), for \(E\in\Ob\Dpx(S)\) and \(F\in\Ob\Dpx^+(S)\), a functorial canonical arrow
\[
  E\longrightarrow\RHom(\RHom(E,F),F).
\]
This is an isomorphism if \(E\in\Ob\Dpx(S)_{\Parf}\) and \(F=A\).

\emph{Proof.} Define the arrow by resolving \(F\) injectively and using the canonical arrow in \(\Kpx(S)\)
\[
  E\longrightarrow \Hom^\bullet(\Hom^\bullet(E,F),F).
\]
To see that the arrow in question is an isomorphism if \(E\) is perfect and \(F=A\), one reduces, by localization and devissage, to \(E=A\), where it is evident.

\textbf{Definition 7.3.} Let \(E\in\Ob\Dpx(S)\). We shall put
\[
  \RHom(E,A)=E^\vee,
\]
and we shall say that \(E^\vee\) is the dual of \(E\).

By (7.1) and (7.2), if \(E\in\Ob\Dpx(S)\) is perfect (respectively perfect and of finite Tor-dimension), then so is \(E^\vee\), and the canonical arrow \(E\to E^{\vee\vee}\) is an isomorphism. In other words, the functor \(\RHom(\,\cdot\,,A)\) defines an anti-auto-equivalence of the category \(\Dpx(S)_{\Parf}\) (respectively \(\Dpx(S)_{\parf}\), cf. (5.8.4 b)) with itself.

\textbf{Lemma 7.4} (Cartan isomorphism) (cf. ([H] II 5.15)). For \(E\in\Ob\Dpx(S)\), \(F\in\Ob\Dpx^-(S)\), and \(G\in\Ob\Dpx^+(S)\), there are functorial canonical isomorphisms
\[
\begin{aligned}
\RHom(E\otL F,G)&\xrightarrow{\sim}\RHom(E,\RHom(F,G)),\\
\RHom(E\otimes F,G)&\xrightarrow{\sim}\RHom(E,\RHom(F,G)),\\
\Hom(E\otimes F,G)&\xrightarrow{\sim}\Hom(E,\RHom(F,G)).
\end{aligned}
\]

\emph{Proof.} This is standard: resolve \(G\) injectively and \(F\) flatly.

\textbf{Lemma 7.5.} For \(E\in\Ob\Dpx^-(S)\) and \(F\in\Ob\Dpx^+(S)\), there is a functorial canonical homomorphism
\[
  E\otL \RHom(E,F)\longrightarrow F.
\]

\textbf{Proposition 7.6} (cf. ([H] II 5.13)). Let \(E\in\Ob\Dpx(S)\), \(F\in\Ob\Dpx^+(S)\), and \(G\in\Ob\Dpx^+(S)\). Under either of the following hypotheses:
\begin{enumerate}[label=(\roman*)]
\item \(G\) is of finite Tor-dimension;
\item \(E\in\Ob\Dpx^-(S)\), \(F\) is of finite Tor-dimension, and \(G\in\Ob\Dpx^b(S)\),
\end{enumerate}
there is a functorial canonical homomorphism
\[
  \RHom(E,F)\otL G
  \longrightarrow
  \RHom(E,F\otL G).
\]
This is an isomorphism when \(E\) is perfect.

\emph{Proof.} Let us define the arrow. Under hypothesis (i), it is enough to resolve \(F\) injectively and \(G\) flatly. Under hypothesis (ii), by (7.4) it is equivalent to define an arrow
\[
  E\otL\RHom(E,F)\otL G\longrightarrow F\otL G.
\]
Apply the functor \(-\otL G\) to the arrow of (7.5).

To see that the arrow is an isomorphism when \(E\) is perfect, use the standard method (cf. the proof of (7.2)): reduce, by localization and devissage, to the case \(E=A\).

\textbf{Corollary 7.7.} There exists, for \(E\in\Ob\Dpx(S)_{\Parf}\) and \(F\in\Ob\Dpx^+(S)\), a functorial isomorphism
\[
  E^\vee\otL F\xrightarrow{\sim}\RHom(E,F),
\]
with the notation of (7.3).

\emph{Proof.} Apply (7.6), with \(F\) and \(G\) replaced respectively by \(A\) and \(F\), and use the fact (7.1) that \(E^\vee\in\Ob\Dpx(S)_{\Parf}\).

\textbf{Corollary 7.8.} Let \(E,F\in\Ob\Dpx(S)_{\Parf}\). There is a natural arrow
\[
  \RHom(E,F)\otL \RHom(F,E)\longrightarrow A,
\]
which is a ``perfect duality,'' that is, modulo (7.4), it identifies each of the complexes \(\RHom(E,F)\) and \(\RHom(F,E)\) with the dual of the other, in the sense of (7.3).

\emph{Proof.} By (7.5), one has natural pairings \(E^\vee\otL E\to A\) and \(F^\vee\otL F\to A\). Hence, by tensor product, one obtains a natural pairing
\[
  E^\vee\otL F\otL F^\vee\otL E\longrightarrow A,
\]
which, in view of (7.7), gives the announced pairing on the \(\RHom\)'s. Let us show that it is a perfect duality. There are canonical isomorphisms
\[
\begin{aligned}
\RHom(F^\vee\otL E,A)&\xrightarrow{\sim}\RHom(F^\vee,E^\vee) &&(7.4)\\
&\xrightarrow{\sim} F^{\vee\vee}\otL E^\vee &&(7.6)\\
&\xrightarrow{\sim} F\otL E^\vee &&(7.2).
\end{aligned}
\]
Thus, up to a compatibility verification, the result follows. One could also localize and devissage to reduce to \(E=F=A\).


\textbf{Exercise 7.9.} Let \(E,F\in\Ob\Dpx(S)_{\Parf}\). Show that there is a canonical isomorphism
\[
  E^\vee\otL F^\vee \xrightarrow{\sim} (E\otL F)^\vee .
\]

\section*{8. Traces and cup-products}

\textbf{8.0.} In this section we generalize to perfect complexes the notion of the trace of an endomorphism of a projective module of finite type (see SGA 5 III for applications). The hypotheses and notation are those of (7.0).

\textbf{8.1.} Let \(E\in\Ob\Dpx(S)_{\parf}\) (5.8.4 b). There is a natural arrow
\[
  \tau:\RHom(E,E)\longrightarrow A,
\]
obtained by composing the inverse of the isomorphism (7.7)
\[
  \RHom(E,E)\xrightarrow{\sim}E^\vee\otL E
\]
with the canonical arrow (7.5) \(E^\vee\otL E\to A\).

Let \(U\in\Ob S\). Applying to \(\tau\) the functor \(H^0(U,\cdot)\), and using the fact that \(H^0(U,\RHom(E,E))\) identifies with \(\Hom_{\Dpx(U)}(E|U,E|U)\), one obtains a homomorphism of \(H^0(U,A)\)-modules
\[
  \Tr_{E|U}:\Hom_{\Dpx(U)}(E|U,E|U)\longrightarrow H^0(U,A),
\]
which we shall call the trace homomorphism. The finite Tor-dimension hypothesis on \(E\) is in fact superfluous for the definition of \(\Tr_{E|U}\); see (8.5).

Before going further, we give some properties of the trace that follow immediately from the definition.

\textbf{8.1.1.} If \(E=A\), then
\[
  \Tr_{A|U}:\Hom(A|U,A|U)\longrightarrow H^0(U,A)
\]
is just the well-known isomorphism. If \(E\) is a direct factor of a free module of finite type, an endomorphism of \(E\) over \(U\) is locally defined by
\[
  e\longmapsto \sum_i\langle e,x_i^\vee\rangle x_i,
\]
where \(x_i\) (respectively \(x_i^\vee\)) is a local section of \(E\) (respectively \(E^\vee\)); the corresponding trace is \(\sum_i\langle x_i,x_i^\vee\rangle\). It follows easily that, if \(E\) and \(F\) are direct factors of free modules of finite type, and if
\[
  u:(E+F)|U\longrightarrow(E+F)|U
\]
is defined by a matrix \(\begin{pmatrix}a&b\\ c&d\end{pmatrix}\), then
\[
  \Tr_{(E+F)|U}(u)=\Tr_{E|U}(a)+\Tr_{F|U}(d).
\]

\textbf{8.1.2.} Let \(E\) be a strictly perfect complex, let \(U\in\Ob S\), and let
\[
  u:E|U\longrightarrow E|U
\]
be a homomorphism of complexes, with components \(u^i:E^i|U\to E^i|U\). Then
\[
  \Tr_{E|U}(u)=\sum_i(-1)^i\Tr_{E^i|U}(u^i).
\]
Indeed, the question being local, one may suppose that \(u\) is defined by a section of \(H^0(E^\vee\otimes E)\), and even that this section is of the form
\[
  \sum_{i,j}x_{ij}^\vee\otimes x_{ij},
\]
where \(x_{ij}\) (respectively \(x_{ij}^\vee\)) is a section of \(E^i\) (respectively \(E^{i\vee}\)). Making explicit the arrow (7.5) \(E^\vee\otL E\to A\), one finds
\[
  \Tr_{E|U}(u)=\sum_{i,j}(-1)^i\langle x_{ij},x_{ij}^\vee\rangle
  =\sum_i(-1)^i\Tr_{E^i|U}(u^i).
\]

\textbf{8.2.} Let \(E,F\in\Ob\Dpx(S)_{\Parf}\). We have seen in (7.8) that there is a canonical homomorphism
\[
  \chi:\RHom(E,F)\otL\RHom(F,E)\longrightarrow A.
\]
Let \(U\in\Ob S\), and let
\[
  u\in H^0(U,\RHom(E,F)),\qquad v\in H^0(U,\RHom(F,E)).
\]
One may interpret \(u\) (respectively \(v\)) as a homomorphism, in \(\Dpx(U)\), from \(A|U\) to \(\RHom(E,F)|U\) (respectively \(\RHom(F,E)|U\)), and form
\[
  u\otL v:A|U\longrightarrow(\RHom(E,F)\otL\RHom(F,E))|U.
\]
The map \((u,v)\mapsto\langle u,v\rangle=\chi\circ(u\otL v)\) defines an \(H^0(U,A)\)-bilinear pairing
\[
  \Hom_{\Dpx(U)}(E|U,F|U)\otimes_{H^0(U,A)}\Hom_{\Dpx(U)}(F|U,E|U)
  \longrightarrow H^0(U,A),
\]
which we shall call the cup-product. The same comment as for the definition of \(\Tr_{E|U}\) applies here.

The relation between trace and cup-product is expressed in the following proposition.

\textbf{Proposition 8.3.} In the situation of (8.2), for \(u\in\Hom(E|U,F|U)\) and \(v\in\Hom(F|U,E|U)\), one has
\[
  \Tr_{E|U}(vu)=\Tr_{F|U}(uv)=\langle u,v\rangle.
\]
In particular, if \(E=F\), then
\[
  \Tr_{E|U}(u)=\langle u,\Id\rangle=\langle\Id,u\rangle.
\]

\emph{Proof.} For simplicity we suppose \(U=S\). We therefore have two arrows
\[
  u:A\longrightarrow E^\vee\otL F,
  \qquad
  v:A\longrightarrow F^\vee\otL E,
\]
and hence
\[
  u\otL v:A\longrightarrow E^\vee\otL F\otL F^\vee\otL E.
\]
On the other hand, we have seen that \(vu:A\to E^\vee\otL E\) and \(uv:A\to F^\vee\otL F\). The proposition follows from the commutativity of the following diagrams, whose verification is purely mechanical:
\[
\begin{tikzcd}[column sep=large,row sep=large]
 A \arrow[r,"u\otL v"] \arrow[dr,"vu"'] & E^\vee\otL F\otL F^\vee\otL E \arrow[d,"\Id\otimes\tau_F"] \\
 & E^\vee\otL E
\end{tikzcd}
\qquad
\begin{tikzcd}[column sep=large,row sep=large]
 A \arrow[r,"u\otL v"] \arrow[dr,"uv"'] & E^\vee\otL F\otL F^\vee\otL E \arrow[d,"\tau_E\otimes\Id"] \\
 & F^\vee\otL F
\end{tikzcd}
\]
\[
\begin{tikzcd}[column sep=large,row sep=large]
E^\vee\otL F\otL F^\vee\otL E \arrow[r,"\Id\otimes\tau_E"] \arrow[d,"\Id\otimes\tau_F"'] & F^\vee\otL F \arrow[d,"\tau_F"]\\
E^\vee\otL E \arrow[r,"\tau_E"'] & A .
\end{tikzcd}
\]

\textbf{Corollary 8.4.} Let \(E,F,U\) be as in (8.2), let \(u\in\Hom(E|U,E|U)\), and let \(s:E|U\to F|U\) be an isomorphism. Then \(sus^{-1}\in\Hom(F|U,F|U)\), and
\[
  \Tr_{F|U}(sus^{-1})=\Tr_{E|U}(u).
\]

\emph{Proof.} Immediate.

\textbf{Remark 8.5.} If \(\mathcal C\) is a fibered \(S\)-category, denote by \(\mathcal E(\mathcal C)\) the fibered \(S\)-category whose fiber over \(U\in\Ob S\) is the category whose objects are pairs \((L,u)\), with \(L\in\Ob\mathcal C_U\) and \(u\in\operatorname{End}_S(L)\), and whose arrows from \((L,u)\) to \((M,v)\) are the arrows \(f:L\to M\) such that \(fu=vf\). By (8.4) one may regard the trace homomorphism as an \(S\)-functor from \(\mathcal E(\Dpx(S)_{\Parf})^{\mathrm{is}}\) to \(\mathcal A\), where the exponent \(\mathrm{is}\) means that only isomorphisms are considered, and where \(\mathcal A\) is the \(S\)-category with discrete fibers defined by \(A\) (cf. (6.6 a)). Since every perfect complex is locally of finite Tor-dimension, the forgetful \(S\)-functor
\[
  \mathcal E(\Dpx(S)_{\Parf})^{\mathrm{is}}
  \longrightarrow
  \mathcal E(\Dpx(S)_{\parf})^{\mathrm{is}}
\]
induces an isomorphism on the associated stacks. Consequently, the trace homomorphism extends uniquely to an \(S\)-functor from \(\mathcal E(\Dpx(S)_{\Parf})^{\mathrm{is}}\) to \(\mathcal A\).

For \(E\in\Ob\Parf(U)\), \(U\in\Ob S\), we again write
\[
  \Tr_E:\Hom(E,E)\longrightarrow H^0(U,A)
\]
for the corresponding homomorphism, which is an \(H^0(U,A)\)-module homomorphism. One sees in the same way that the definition of the cup-product (8.2), as well as assertions (8.3) and (8.4), generalize to perfect complexes, not necessarily of finite Tor-dimension.

\textbf{Proposition 8.6.} Let \(U\in\Ob S\), let \(E,F\in\Parf(U)\), and let \(u\in\Hom(E,E)\), \(v\in\Hom(F,F)\); hence \(u\otL v\in\Hom(E\otL F,E\otL F)\). Then
\[
  \Tr_{E\otL F}(u\otL v)=\Tr_E(u)\,\Tr_F(v).
\]

\emph{Proof.} The question being local, by (8.4) and (4.10) one may suppose that \(u\) and \(v\) are arrows of \(K^b(\mathcal C_0)\). Using (8.1.2), one is reduced to the case where \(u\) and \(v\) are arrows of \(\mathcal C_0\). One finishes by (8.1.1).

\section*{Bibliography}
\begin{description}[leftmargin=2.5cm,style=nextline]
\item[{[H]}] Hartshorne, R. \emph{Residues and Duality}, Lecture Notes no. 20, Springer, 1966.
\item[{[V]}] Verdier, J.-L. \emph{Catégories dérivées}, I.H.E.S., 1963.
\item[{[G]}] Giraud, J. \emph{Cohomologie non abélienne}, to appear in North-Holland Publishing Company.
\item[{SGA 4}] \emph{Cohomologie étale des schémas}, by M. Artin, A. Grothendieck, and J.-L. Verdier, to appear in North-Holland Publishing Company.
\item[{SGA 5}] \emph{Cohomologie \(\ell\)-adique et fonctions L}, by A. Grothendieck, to appear in North-Holland Publishing Company.
\item[{[1]}] Hakim, M. \emph{Topos annelés et schémas relatifs}, thesis (in preparation).
\item[{[2]}] Verdier, J.-L. \emph{Les catégories dérivées des catégories abéliennes}, to appear in North-Holland Publishing Company.
\item[{[MUL]}] Cartan, H. and Eilenberg, S. \emph{Homological Algebra}, Princeton, 1956.
\end{description}



\providecommand{\cS}{\mathcal S}
\providecommand{\cCz}{\mathcal C_0}
\providecommand{\cCone}{\mathcal C_1}
\providecommand{\cA}{\mathcal A}
\providecommand{\cB}{\mathcal B}
\providecommand{\cL}{\mathcal L}
\providecommand{\cI}{\mathcal I}
\providecommand{\cJ}{\mathcal J}
\providecommand{\Db}{\mathrm D}
\providecommand{\Kb}{\mathrm K}
\providecommand{\Dminus}{\mathrm D^{-}}
\providecommand{\Kminus}{\mathrm K^{-}}
\providecommand{\Kbd}{\mathrm K^{b}}
\providecommand{\Mod}{\operatorname{Mod}}
\providecommand{\Lib}{\operatorname{Lib}}
\providecommand{\Loclib}{\operatorname{Loclib}}
\providecommand{\Qcoh}{\operatorname{Qcoh}}
\providecommand{\Coh}{\operatorname{Coh}}
\providecommand{\Ext}{\operatorname{Ext}}
\providecommand{\Hom}{\operatorname{Hom}}
\providecommand{\Ob}{\operatorname{ob}}
\providecommand{\im}{\operatorname{Im}}
\providecommand{\coker}{\operatorname{Coker}}
\providecommand{\coh}{\operatorname{coh}}
\providecommand{\parf}{\operatorname{parf}}
\providecommand{\amp}{\operatorname{amp}}
\providecommand{\RHom}{\operatorname{RHom}}


\section*{Exposé II. Existence of global resolutions}
\begin{center}
by L. Illusie
\end{center}

\section*{1. General globalization criteria}

\subsection*{1.0}
In this section we are given a site $S$, a flat abelian $S$-category $\mathcal C$ (I 1.1), and a strictly full additive sub-$S$-category $\mathcal C_0$ (I 1.1). We assume that $\mathcal C_0$ is locally quasi-liftable in $\mathcal C$ (I 1.2) and locally quasi-stable under kernels of epimorphisms (I 1.2).

\textbf{Proposition 1.1.} Let $S\in \Ob(S)$, and suppose that:
\begin{enumerate}[label=(\roman*)]
\item $\mathcal C_{0S}$ is quasi-liftable in $\mathcal C_S$ (I 1.2), where $\mathcal C_S$ is regarded as a category fibered over a punctual site;
\item every object of $\mathcal C_S$ which is of finite $\mathcal C_0$-type (I 1.2), respectively of finite $\mathcal C_0$-presentation (I 2.8), is of finite $\mathcal C_{0S}$-type.
\end{enumerate}
Then:

\begin{enumerate}[label=\alph*)]
\item $\mathcal C_{0S}$ is quasi-stable under kernels of epimorphisms (I 1.2), where $\mathcal C_{0S}$ is regarded as a category fibered over a punctual site.
\item Let $E\in\Ob \mathrm D^{-}(\mathcal C_S)$ and $n\in\mathbb Z$. If $E$ is $n$-$\mathcal C_0$-pseudo-coherent (I 2.2), then $E$ is $n$-$\mathcal C_{0S}$-pseudo-coherent, respectively $(n+1)$-$\mathcal C_{0S}$-pseudo-coherent. In other words,
\[
\mathrm D^{-}(\mathcal C_S)_{n-\mathcal C_0\text{-coh}}
=
\mathrm D^{-}(\mathcal C_S)_{n-\mathcal C_{0S}\text{-coh}},
\]
and, respectively,
\[
\mathrm D^{-}(\mathcal C_S)_{n-\mathcal C_0\text{-coh}}
\subset
\mathrm D^{-}(\mathcal C_S)_{(n+1)-\mathcal C_{0S}\text{-coh}}
\subset
\mathrm D^{-}(\mathcal C_S)_{(n+1)-\mathcal C_0\text{-coh}},
\]
inclusions of strictly full subcategories of $\mathrm D(\mathcal C_S)$, the inclusion on the right being trivial.
\item For an object $E\in\Ob\mathrm D(\mathcal C_S)$ to be pseudo-coherent relative to $\mathcal C_0$ (I 2.3), it is necessary and sufficient that it be pseudo-coherent relative to $\mathcal C_{0S}$; equivalently,
\[
\mathrm D^{-}(\mathcal C_S)_{\mathcal C_0\text{-coh}}
=
\mathrm D^{-}(\mathcal C_S)_{\mathcal C_{0S}\text{-coh}}.
\]
Moreover, the canonical functor
\[
\mathrm K^{-}(\mathcal C_{0S})/\mathrm K^{-,\delta}(\mathcal C_{0S})
\longrightarrow \mathrm D(\mathcal C_S),
\]
where $\mathrm K^{-,\delta}(\mathcal C_{0S})$ denotes the thick subcategory of $\mathrm K^{-}(\mathcal C_{0S})$ formed by acyclic complexes, is fully faithful and has as essential image $\mathrm D^{-}(\mathcal C_S)_{\mathcal C_0\text{-coh}}$. If $\mathcal C_{0S}$ is liftable in $\mathcal C_S$ (I 1.2), that is, if the objects of $\mathcal C_{0S}$ are projective, then $\mathrm K^{-,\delta}(\mathcal C_{0S})$ is reduced to zero objects, and consequently the functor
\[
\mathrm K^{-}(\mathcal C_{0S})\longrightarrow \mathrm D(\mathcal C_S)
\]
is fully faithful and has essential image $\mathrm D^{-}(\mathcal C_S)_{\mathcal C_0\text{-coh}}$.
\end{enumerate}

\emph{Proof.} a) Let
\[
0\longrightarrow L\longrightarrow L^0\longrightarrow\cdots\longrightarrow L^m\longrightarrow0
\]
be an exact sequence in $\mathcal C_S$, with $L^i\in\Ob\mathcal C_{0S}$ for all $i$. It is enough to prove that $L$ is of finite $\mathcal C_{0S}$-type. By (I 2.2), $L$ is $\mathcal C_0$-pseudo-coherent, hence, by (I 2.9), of finite $\mathcal C_0$-presentation; the assertion follows from (ii).

b) We prove by descending induction on $i$ that $E$ is $i$-$\mathcal C_{0S}$-pseudo-coherent for $i\ge n$, respectively for $i\ge n+1$. One may already suppose that $E$ is strictly $(n+1)$-$\mathcal C_{0S}$-pseudo-coherent, respectively strictly $(n+2)$-$\mathcal C_{0S}$-pseudo-coherent (I 2.1); the point is then to prove that $E$ is $n$-$\mathcal C_{0S}$-pseudo-coherent, respectively $(n+1)$-$\mathcal C_{0S}$-pseudo-coherent. Thus there is an exact sequence
\[
0\longrightarrow E'\longrightarrow E\longrightarrow E''\longrightarrow0,
\]
where $E'$ is strictly $\mathcal C_{0S}$-perfect (I 2.1) and $E''{}^i=0$ for $i>n$, respectively for $i>n+1$. Replacing $E$ by $E''$, one may therefore suppose (I 2.6) that $E^i=0$ for $i>n$, respectively for $i>n+1$. Then, by (I 2.10 a)), $H^n(E)$, respectively $H^{n+1}(E)$, is of finite $\mathcal C_0$-type, respectively of finite $\mathcal C_0$-presentation. Hence, by (ii), $H^n(E)$, respectively $H^{n+1}(E)$, is of finite $\mathcal C_{0S}$-type. Applying (I 2.10 b)) again, one obtains that $E$ is $n$-$\mathcal C_{0S}$-pseudo-coherent, respectively $(n+1)$-$\mathcal C_{0S}$-pseudo-coherent. The hypothesis (i) and part a) have been used implicitly in the references to (I 2.6 and 2.10).

c) The first assertion follows from b), the second from (I 2.7), and the last is evident.

\textbf{Proposition 1.2.} Suppose that $\mathcal C_0$ is locally liftable in $\mathcal C$ (I 1.2) and locally stable under kernels of epimorphisms (I 1.2), and make the hypotheses (1.1 (i) and (ii), respectively). Let $\mathcal C_{1S}$ denote the strictly full subcategory of $\mathcal C_S$ formed by the objects which are locally in $\mathcal C_0$. Then:
\begin{enumerate}[label=\alph*)]
\item $\mathcal C_{1S}$ is stable under kernels of epimorphisms (I 1.2) and is quasi-liftable in $\mathcal C_S$ (I 1.2), as a category fibered over a punctual site.
\item Let $E\in\Ob\mathrm D(\mathcal C_S)$ and let $[m,n]$ be an interval of $\mathbb Z$. For $E$ to have $\mathcal C_0$-perfect amplitude contained in $[m,n]$, it is necessary and sufficient that $E$ be isomorphic, in $\mathrm D(\mathcal C_S)$, to a complex $F$ such that $F^i=0$ for $i\notin[m,n]$, $F^i\in\Ob\mathcal C_{0S}$ for $i>m$, and $F^m\in\Ob\mathcal C_{1S}$.
\item The canonical functor
\[
\mathrm K^b(\mathcal C_{1S})/\mathrm K^{b,\delta}(\mathcal C_{1S})
\longrightarrow \mathrm D(\mathcal C_S)
\]
is fully faithful and has as essential image the full subcategory of $\mathrm D(\mathcal C_S)$ formed by complexes of finite $\mathcal C_0$-perfect amplitude.
\item If $\mathcal C_{1S}$ is not merely quasi-liftable but liftable in $\mathcal C_S$, one may replace, in b), the expression ``$E$ is isomorphic to $F$ in $\mathrm D(\mathcal C_S)$'' by ``there exists a quasi-isomorphism $F\to E$''; and, in c), $\mathrm K^{b,\delta}(\mathcal C_{1S})$ is reduced to zero objects.
\end{enumerate}

\emph{Proof.} a) It is clear that $\mathcal C_{1S}$ is stable under kernels of epimorphisms. On the other hand, the objects of $\mathcal C_{1S}$ are $\mathcal C_0$-pseudo-coherent, hence $\mathcal C_{0S}$-pseudo-coherent by (1.1 c)); the second assertion follows from (I 2.14), taking $S$ to be the punctual site.

b) It suffices to prove necessity. By (1.1 c)), $E$ is already $\mathcal C_{0S}$-pseudo-coherent. Since $E$ is acyclic in degree $>n$, one may, after replacing $E$ by an isomorphic complex in $\mathrm D(\mathcal C_S)$, suppose (I 2.10) that $E$ is strictly $(m+1)$-$\mathcal C_{0S}$-pseudo-coherent and that $E^i=0$ for $i>n$. Since $E$ is acyclic in degree $<m$, one may also replace $E$ by its truncation
\[
0\longrightarrow E^m/B^m\longrightarrow E^{m+1}\longrightarrow\cdots\longrightarrow E^n\longrightarrow0,
\]
so that $E^i=0$ for $i<m$. Then, by (I 4.16), $E^m$ is locally an object of $\mathcal C_0$, and therefore an object of $\mathcal C_{1S}$.

c) In view of (1.2 a)), the assertion follows from (1.2 b)) and (I 4.12.2). d) The fact that $\mathrm K^{b,\delta}(\mathcal C_{1S})$ is reduced to zero objects is evident. The other assertion follows from (I 4.6), applied over the punctual site.

\section*{2. Application to certain categories of modules}

\subsection*{2.0. Notation}
Let $(S,A)$ be a ringed topos. For $U\in\Ob S$, write $\Mod(U)$ for the category of left $A_U$-modules, and $\Lib(U)$, respectively $\Loclib(U)$, for the full subcategory of $\Mod(U)$ formed by free $A_U$-modules, respectively modules locally direct factors of free modules, of finite type.

\subsection*{2.1. Lifting lemmas}

\textbf{Lemma 2.1.1} (cf. EGA $0_{\mathrm I}$, 5.2.3). Let $(S,A)$ be a ringed topos, and let $S\in\Ob S$. Suppose that $S$ is quasi-compact, that is, every covering family $(S_i\to S)$ is dominated by a finite covering family.
\begin{enumerate}[label=\alph*)]
\item Let $F$ be a left $A_S$-module which is the filtered inductive limit, over a filtered category $\Lambda$, of $A_S$-modules $F_\lambda$, and let
\[
u=\varinjlim u_\lambda:F\longrightarrow G
\]
be an epimorphism, where $G$ is an $A_S$-module of finite type. Then there exists $\lambda\in\Lambda$ such that $u_\lambda$ is an epimorphism.
\item Let $F$ be a left $A_S$-module generated by its sections, that is, such that there exists an epimorphism $A_S^{(I)}\to F$ for a suitable set of indices $I$. If $F$ is of finite type, then there exists an epimorphism $A_S^p\to F$ for some $p\in\mathbb N$.
\end{enumerate}
The proof of a) is analogous to the proof of the cited result and is left to the reader; b) follows immediately from a).

\textbf{Lemma 2.1.2.} Let $(S,A)$ be a ringed topos, let $\mathcal C$ be the $S$-category fibered in left $A$-modules, let $\mathcal C_0$ be the fibered subcategory of free $A$-modules of finite type (I 1.3 d)), let $S\in\Ob S$, and let $\mathcal C'_S$ be a strictly full additive subcategory of $\mathcal C_S$ containing $\mathcal C_{0S}$ and stable under kernels, cokernels, and extensions. Suppose that
\[
H^1(S,F)=0\quad\text{for every }F\in\Ob\mathcal C'_S.
\]
Then:
\begin{enumerate}[label=\alph*)]
\item Every exact sequence in $\mathcal C_S$
\[
0\longrightarrow E\longrightarrow F\longrightarrow G\longrightarrow0,
\]
such that $\operatorname{Hom}(G,E)\in\Ob\mathcal C'_S$ and such that $G$ is locally a direct factor of a free module of finite type, is split. The category $\mathcal C_{0S}$ is liftable in $\mathcal C'_S$ (I 2.2).
\item Let $\Mod(A(S))$ denote the category of $A(S)$-modules. The functor
\[
\Gamma_S:\mathcal C'_S\longrightarrow\Mod(A(S))
\]
is exact. Moreover:
\begin{enumerate}[label=(\roman*)]
\item if every object $F$ of $\mathcal C'_S$ admits a finite global presentation
\[
A_S^q\longrightarrow A_S^p\longrightarrow F\longrightarrow0,
\]
then $\Gamma_S$ is fully faithful and has as essential image the full subcategory of $\Mod(A(S))$ formed by $A(S)$-modules of finite presentation;
\item if $\mathcal C'_S$ contains the free $A_S$-modules, if every object $F$ of $\mathcal C'_S$ admits a global presentation
\[
A_S^{(J)}\longrightarrow A_S^{(I)}\longrightarrow F\longrightarrow0,
\]
and if $\Gamma_S$ commutes with filtered inductive limits, then $\Gamma_S$ is an equivalence of categories.
\end{enumerate}
\end{enumerate}

\emph{Proof.} a) Since $G$ is locally a direct factor of a free module of finite type, the spectral sequence
\[
E_2^{pq}=H^p(S,\mathcal{E}xt^q(G,E))\Longrightarrow \Ext^{p+q}(G,E)
\]
degenerates (I 1.3.1) and gives an isomorphism
\[
H^1(S,\operatorname{Hom}(G,E))\simeq \Ext^1(G,E),
\]
whence the first assertion. The second is a particular case, since for $G\in\Ob\mathcal C_{0S}$ and $E\in\Ob\mathcal C'_S$ one has $\operatorname{Hom}(G,E)\in\Ob\mathcal C'_S$.

b) The first assertion is trivial. Let us prove, in case (i) or (ii), the full faithfulness of $\Gamma_S:\mathcal C'_S\to\Mod(A(S))$. Let $F,G\in\Ob\mathcal C'_S$ and let
\[
A_S^{(J)}\longrightarrow A_S^{(I)}\longrightarrow F\longrightarrow0
\tag{$*$}
\]
be a global presentation of $F$, with $I$ and $J$ finite in case (i). Applying $\Gamma_S$, which is exact on $\mathcal C'_S$ and, in case (ii), commutes with filtered inductive limits, one obtains an exact sequence
\[
A(S)^{(J)}\longrightarrow A(S)^{(I)}\longrightarrow\Gamma_S(F)\longrightarrow0.
\tag{$**$}
\]
Applying the left exact functors $\operatorname{Hom}_{\mathcal C_S}(-,G)$ and $\operatorname{Hom}_{A(S)}(-,\Gamma_S(G))$ to $(*)$ and $(**)$ gives a commutative diagram with exact rows:
\[
\begin{tikzcd}[column sep=small]
0 \arrow[r] & \operatorname{Hom}_{\mathcal C_S}(F,G) \arrow[r] \arrow[d] &
\operatorname{Hom}_{\mathcal C_S}(A_S^{(I)},G) \arrow[r] \arrow[d,"\sim"] &
\operatorname{Hom}_{\mathcal C_S}(A_S^{(J)},G) \arrow[d,"\sim"] \\
0 \arrow[r] & \operatorname{Hom}_{A(S)}(\Gamma_SF,\Gamma_SG) \arrow[r] &
\operatorname{Hom}_{A(S)}(A(S)^{(I)},\Gamma_SG) \arrow[r] &
\operatorname{Hom}_{A(S)}(A(S)^{(J)},\Gamma_SG).
\end{tikzcd}
\]
The two vertical arrows on the right are visibly isomorphisms, hence so is the arrow on the left. The description of the essential image in cases (i) and (ii) is immediate and is left to the reader.

\textbf{Lemma 2.1.3.} Let $\mathcal A$ and $\mathcal B$ be abelian categories, and let $\varphi:\mathcal A\to\mathcal B$ be an exact functor.
\begin{enumerate}[label=\alph*)]
\item For $M\in\Ob\mathcal A$, the following conditions are equivalent:
\begin{enumerate}[label=(\roman*)]
\item for every $L\in\Ob\mathcal A$, the canonical map induced by $\varphi$
\[
\Ext^n_{\mathcal A}(L,M)\longrightarrow\Ext^n_{\mathcal B}(\varphi L,\varphi M)
\]
is an isomorphism for all $n$;
\item for every $L\in\Ob\mathcal A$, the canonical map
\[
\operatorname{Hom}_{\mathcal A}(L,M)\longrightarrow\operatorname{Hom}_{\mathcal B}(\varphi L,\varphi M)
\]
is an isomorphism, and for every $n\ge1$ and every $x\in\Ext^n_{\mathcal B}(\varphi L,\varphi M)$ there exists an epimorphism $L'\to L$ in $\mathcal A$ which kills $x$.
\end{enumerate}
\item The following conditions are equivalent:
\begin{enumerate}[label=(\roman*)]
\item for all $L,M\in\Ob\mathcal A$ and all $n\in\mathbb Z$, the canonical map
\[
\Ext^n_{\mathcal A}(L,M)\longrightarrow\Ext^n_{\mathcal B}(\varphi L,\varphi M)
\]
is an isomorphism;
\item the canonical functor induced by $\varphi$
\[
\mathrm D^b(\mathcal A)\longrightarrow\mathrm D^b(\mathcal B)
\]
is fully faithful and has as essential image the full subcategory of $\mathrm D^b(\mathcal B)$ formed by complexes $F$ such that $H^i(F)$ is in the essential image of $\varphi$ for every $i$.
\end{enumerate}
\end{enumerate}

\emph{Proof.} For (i)$\Rightarrow$(ii) in a), it is enough to show that, if $n\ge1$, $L\in\Ob\mathcal A$, and $y\in\Ext^n_{\mathcal A}(L,M)$, then there exists an epimorphism $L'\to L$ killing $y$. Using Yoneda's interpretation ([18], III 3.2), $y$ is the class of an exact sequence
\[
0\longrightarrow M\longrightarrow E_{n-1}\longrightarrow\cdots\longrightarrow E_0\longrightarrow L\longrightarrow0,
\]
and the epimorphism $E_0\to L$ kills $y$ ([18], III 3.2.5 and 3.2.8). Conversely, put
\[
F^n(L)=\Ext^n_{\mathcal A}(L,M),\qquad
G^n(L)=\Ext^n_{\mathcal B}(\varphi L,\varphi M).
\]
The effaceability hypothesis says, for $n\ge1$ and $L\in\Ob\mathcal A$, that
\[
G^n(L)=\varinjlim\operatorname{Coker}igl(G^{n-1}(L_0)\to G^{n-1}(L_1)\bigr),
\]
the limit being taken over the exact sequences $L_1\to L_0\to L\to0$. As above, one has the same formula for $F^n(L)$. Since $F^0\to G^0$ is an isomorphism, induction on $n$ shows that the morphism $F^n\to G^n$ induced by $\varphi$ is an isomorphism for every $n$. The equivalence in b) follows by an immediate dévissage, left to the reader.

\textbf{Remarks 2.1.3.1.} There is of course a dual version of a), which we leave to the reader to state. On the other hand, if $\mathcal A$ and $\mathcal B$ have enough injectives, the conditions (i) and (ii) of b) are equivalent to:
\begin{enumerate}[label=(\roman*),start=3]
\item for $L\in\Ob\mathrm D^{-}(\mathcal A)$ and $M\in\Ob\mathrm D^{+}(\mathcal A)$, the canonical map
\[
\mathrm R\!\operatorname{Hom}(L,M)\longrightarrow
\mathrm R\!\operatorname{Hom}(\varphi L,\varphi M)
\]
is an isomorphism,
\end{enumerate}
as is seen at once by a spectral sequence argument.

\subsection*{2.2. Noetherian schemes and divisorial schemes}

\textbf{Notation 2.2.1.} Let $S$ be a scheme. We write $\Qcoh(S)$, respectively $\Coh(S)$, for the abelian category of quasi-coherent, respectively coherent, $\mathcal O_S$-modules. Let $[m,n]$ be an interval of $\mathbb Z$. We shall say that an object $E$ of $\mathrm D(\Qcoh(S))$ is $n$-pseudo-coherent, respectively pseudo-coherent, respectively of perfect amplitude contained in $[m,n]$, respectively perfect, if it is so relative to the category fibered over the Zariski site of $S$ defined by
\[
U\longmapsto \Loclib(U)
\]
((I 2.3), (I 4.7), and (2.0)). We denote by $\mathrm D^{-}(\Qcoh(S))_{\coh}$, respectively $\mathrm D(\Qcoh(S))_{\parf}$, the full triangulated subcategory of $\mathrm D(\Qcoh(S))$ formed by pseudo-coherent, respectively perfect, complexes.

\textbf{Proposition 2.2.2.} Let $S$ be a noetherian scheme. The canonical functor
\[
\mathrm D^{-}(\Coh(S))\longrightarrow \mathrm D(\Qcoh(S))
\]
is fully faithful and has as essential image $\mathrm D^{-}(\Qcoh(S))_{\coh}$.

\emph{Proof.} Apply (1.1), taking for the site the Zariski site of $S$, and for $\mathcal C$, respectively $\mathcal C_0$, the fibered categories $U\mapsto\Qcoh(U)$, respectively $U\mapsto\Coh(U)$. Since $\mathcal O_S$ is coherent, the objects of $\mathcal C_0$ are pseudo-coherent (I 3.5); hence an object of $\mathrm D(\Qcoh(U))$ is pseudo-coherent, in the sense of (2.2.1), if and only if it is pseudo-coherent relative to $\mathcal C_0$ (I 2.14.1). By (1.1 c)), it remains to verify conditions (i) and (ii) of (1.1). For (ii) this is trivial, since a quasi-coherent module which is a quotient of a coherent module is coherent, and coherence is local. For (i), one uses (2.1.1 a)) and the fact that every quasi-coherent $\mathcal O_S$-module is a filtered inductive limit of coherent $\mathcal O_S$-modules (EGA I 9.4.9).

Granting (3.5) and (3.7), which are independent of the present section, one easily deduces the following.

\textbf{Corollary 2.2.2.1.} Let $S$ be a noetherian scheme, respectively a noetherian scheme of finite dimension. The canonical functor
\[
\mathrm D^b(\Coh(S))\longrightarrow \mathrm D(S)
\]
respectively
\[
\mathrm D^{-}(\Coh(S))\longrightarrow \mathrm D(S)
\]
is fully faithful and has as essential image $\mathrm D^b(S)_{\coh}$, respectively $\mathrm D^{-}(S)_{\coh}$ (I 2.15).

\textbf{Proposition 2.2.3} (cf. EGA II 4.5.2 and 4.5.5). Let $S$ be a quasi-compact and quasi-separated scheme, and let $(L_i)_{i\in I}$ be a family of invertible sheaves on $S$. For every $\mathcal O_S$-module $F$ and every pair $(i,n)\in I\times\mathbb Z$, let $F_{(i,n)}$ denote the sub-$\mathcal O_S$-module of $F\otimes L_i^{\otimes n}$ generated by the sections of $F\otimes L_i^{\otimes n}$ over $S$. The following conditions are equivalent:
\begin{enumerate}[label=(\roman*)]
\item the open sets $S_f$, for $f\in\Gamma(S,L_i^{\otimes n})$, $i\in I$ and $n\in\mathbb N$, form a base for the topology of $S$;
\item those of the $S_f$ which are affine cover $S$;
\item for every quasi-coherent $\mathcal O_S$-module $F$ of finite type, there exists a finite family $(i_\alpha,n_\alpha)$ with $n_\alpha\ge0$ and an epimorphism
\[
\bigoplus_\alpha L_{i_\alpha}^{\otimes(-n_\alpha)}\longrightarrow F;
\]
\item for every $F\in\Ob\Qcoh(S)$, one has
\[
F=\sum_{(i,n)\in I\times\mathbb N}F_{(i,n)}\otimes L_i^{\otimes(-n)},
\]
a sum, not necessarily direct, of sub-sheaves of $F$;
\item same condition as (iv), but with $F$ a quasi-coherent ideal of $\mathcal O_S$.
\end{enumerate}

\emph{Proof.} This is a simple paraphrase of the cited passages; we recall the argument briefly for the reader's convenience. For (i)$\Rightarrow$(ii), choose for $x\in S$ an affine open $U$ containing $x$. An open $S_f$ with $x\in S_f\subset U$, where $f\in\Gamma(S,L_i^{\otimes n})$ and $n\ge0$, is necessarily affine, by the following lemma.

\textbf{Lemma 2.2.3.1} (EGA II 5.5.8). Let $X$ be a scheme, $L$ an invertible $\mathcal O_X$-module, and $f\in\Gamma(X,L)$. The canonical immersion $X_f\to X$ is affine.

\emph{Proof.} This is trivial: the question is local on $X$, so one may suppose $X$ affine and $L\simeq\mathcal O_X$.

For (ii)$\Rightarrow$(iii), choose a finite covering of $S$ by affine opens $S_{f_{ij}}$, with $f_{ij}\in\Gamma(S,L_i^{\otimes m_{ij}})$. Replacing the $f_{ij}$ by suitable powers, one may suppose that the $m_{ij}$ with fixed $i$ are all equal to an integer $m_i$. If $F$ is quasi-coherent of finite type, then, for each $i,j$, there is an integer $q_i\ge0$ such that every section of $F$ over $S_{f_{ij}}$ extends to a section of $F\otimes L_i^{\otimes q_im_i}$ over $S$. Thus $F\otimes L_i^{\otimes q_im_i}$ is generated over $S_{f_{ij}}$ by finitely many global sections. Proceeding as in the proof of EGA II 4.5.5, one obtains, for each $i$, an integer $n_i\ge0$ such that $F\otimes L_i^{\otimes n_i}$ is generated over the union of the corresponding $S_{f_{ij}}$ by finitely many sections over $S$. Hence there are integers $r_i\ge0$ and a morphism
\[
L_i^{-n_i\, r_i}\longrightarrow F
\]
which induces an epimorphism over that union; the assertion follows.

For (iii)$\Rightarrow$(iv), since every quasi-coherent sheaf on $S$ is the inductive limit of its quasi-coherent subsheaves of finite type (EGA I 9.4.9 and IV 1.7.7), it is enough to treat the case where $F$ is of finite type. There exist morphisms
\[
u_i:L_i^{-n_i\, r_i}\longrightarrow F
\]
whose sum is an epimorphism. But
\[
\operatorname{Im}(u_i)=\operatorname{Im}\bigl(u_i\otimes\operatorname{Id}(L_i^{n_i})\bigr)\otimes L_i^{-n_i},
\]
and $\operatorname{Im}\bigl(u_i\otimes\operatorname{Id}(L_i^{n_i})\bigr)$ is a subsheaf of $F_{(i,n_i)}$, which gives the assertion. The implication (iv)$\Rightarrow$(v) is trivial. Finally, for (v)$\Rightarrow$(i), let $x\in S$, let $U$ be an open containing $x$, and let $F$ be a quasi-coherent ideal defining $S-U$. By (v), there are $(i,n)\in I\times\mathbb N$ and a section $f$ of $F\otimes L_i^{\otimes n}$ over $S$ with $f(x)\ne0$. Then $x\in S_f\subset U$, proving (i).

\textbf{Definition 2.2.4.} Let $S$ be a quasi-compact and quasi-separated scheme, and let $(L_i)_{i\in I}$ be a family of invertible sheaves on $S$. We say that the $L_i$ form an \emph{ample family} if the equivalent conditions of (2.2.3) hold. In particular, a family consisting of a single sheaf $L$ is ample if and only if $L$ is ample in the usual sense (EGA II 4.5.3).

\textbf{Definition 2.2.5.} A scheme $S$ is called \emph{divisorial} if $S$ is quasi-compact and quasi-separated and if there exists an ample family (2.2.4) of invertible $\mathcal O_S$-modules.

\textbf{Lemma 2.2.6.} Let $X$ be a locally noetherian scheme. Let $U$ be an open of $X$ such that the inclusion $U\to X$ is affine, for example $U$ affine and $X$ separated. Then $X-U$ is purely of codimension $0$ or $1$; that is, at every maximal point $x$ of $X-U$ one has $\dim\mathcal O_{X,x}=0$ or $1$. It is purely of codimension $1$ if $U$ contains the maximal points of $X$.

\emph{Proof.} Put $Y=X-U$. Let $x$ be a maximal point of $Y$ and make the base change $X'=\operatorname{Spec}\mathcal O_{X,x}\to X$; one obtains a cartesian diagram
\[
\begin{tikzcd}
U' \arrow[r] \arrow[d] & U \arrow[d] \\
X' \arrow[r] & X .
\end{tikzcd}
\]
Let $Y'=X'-U'$. Since $x$ is a maximal point of $Y$, $Y'$ is just the closed point of $X'$. Since $X'$ is affine, $U'$ is an affine open of $X'$. We show that the dimension $d$ of $X'$ is $0$ or $1$. Indeed, one knows (SGA 2 V 5.7) that $H^d_{Y'}(X',\mathcal O_{X'})\ne0$. But if $d\ge2$, then
\[
H^d_{Y'}(X',\mathcal O_{X'})=H^{d-1}(U',\mathcal O_{X'}),
\]
and the latter group is zero because $U'$ is affine and $d-1\ge1$, a contradiction.

\textbf{Remark 2.2.6.1.} Without the separation hypothesis on $X$, (2.2.6) would be false for an affine open $U$ of $X$, as shown by the example of the affine plane with the origin doubled (EGA I 5.5.11).

\textbf{Proposition 2.2.7.} Let $S$ be a noetherian separated locally factorial scheme, that is, a scheme whose local rings are factorial. Then $S$ is divisorial (2.2.5). In particular (EGA IV 21.11.1):

\textbf{Corollary 2.2.7.1.} Every noetherian regular separated scheme is divisorial.

\emph{Proof of 2.2.7.} One may suppose that $S$ is integral. Let $(S_i)_{i\in I}$ be a finite covering of $S$ by nonempty affine opens. By (2.2.6), $S-S_i$ is purely of codimension $1$; hence, by EGA IV 21.7.4, it is the support of a positive divisor. In other words, for each $i\in I$, there is an invertible sheaf $L_i$ and a section $f_i\in\Gamma(S,L_i)$ such that $S_i=S_{f_i}$. The family $(L_i)$ is then ample by (2.2.4), using (2.2.3 (ii)).

\textbf{Remark 2.2.7.2.} As in (2.2.6), the separation hypothesis on $S$ is essential. Let $S$ again be the affine plane with the origin doubled. Recall that $S$ is obtained by gluing two copies of the affine plane $\mathbb A^2_k$ along the identity on the complement $U$ of the origin $O$; hence there is a canonical projection
\[
p:S\longrightarrow\mathbb A^2_k
\]
which is an isomorphism over $U$, while $p^{-1}(O)$ consists of two points $O_1$ and $O_2$. An invertible sheaf on $S$ is obtained by gluing two invertible sheaves on $\mathbb A^2_k$ by an isomorphism over $U$. Since every invertible sheaf on $\mathbb A^2_k$ is trivial, an invertible sheaf on $S$ is defined by an element $g\in\Gamma(U,\mathcal O_U^*)$. But $g$ is necessarily constant, since the restriction $\Gamma(\mathbb A^2_k,\mathcal O_{\mathbb A^2})\to\Gamma(U,\mathcal O_U)$ is an isomorphism, the origin being of depth $2$ in the plane. Thus every invertible sheaf on $S$ is trivial. It is immediate, on the other hand, that $p$ induces an isomorphism
\[
\Gamma(\mathbb A^2_k,\mathcal O_{\mathbb A^2})\xrightarrow{\sim}\Gamma(S,\mathcal O_S).
\]
Consequently every global section of $\mathcal O_S$ which vanishes at $O_1$ also vanishes at $O_2$, and the opens $S_f$, for $f\in\Gamma(S,\mathcal O_S)$, do not form a base for the topology of $S$.

\textbf{Theorem 2.2.8.} Let $S$ be a divisorial scheme (2.2.5), and let $(L_i)_{i\in I}$ be an ample family of invertible $\mathcal O_S$-modules. Let $L(S)$ be the strictly full additive subcategory of $\Qcoh(S)$ generated by the $L_i^{\otimes(-n)}$, for $(i,n)\in I\times\mathbb N$. Let $E\in\Ob\mathrm D(\Qcoh(S))$.
\begin{enumerate}[label=\alph*)]
\item The categories $L(S)$ and $\Loclib(S)$ are quasi-liftable in $\Qcoh(S)$ (I 1.2) and quasi-stable under kernels of epimorphisms (I 1.2).
\item Let $n\in\mathbb Z$. For $E$ to be $n$-pseudo-coherent, respectively pseudo-coherent, in the sense of (2.2.1), it is necessary and sufficient that $E$ be so relative to $L(S)$ or to $\Loclib(S)$ (I 2.3). The canonical functors
\[
\mathrm K^{-}(L(S))/\mathrm K^{-,\delta}(L(S))\longrightarrow \mathrm D(\Qcoh(S)),
\]
\[
\mathrm K^{-}(\Loclib(S))/\mathrm K^{-,\delta}(\Loclib(S))\longrightarrow \mathrm D(\Qcoh(S))
\]
are fully faithful and have as essential image $\mathrm D^{-}(\Qcoh(S))_{\coh}$.
\item Let $[m,n]$ be an interval of $\mathbb Z$. For $E$ to have perfect amplitude contained in $[m,n]$ (2.2.1), it is necessary and sufficient that $E$ be isomorphic, in $\mathrm D(\Qcoh(S))$, to a complex $F$ such that $F^i=0$ for $i\notin[m,n]$, $F^i\in\Ob L(S)$ for $i>m$, and $F^m$ is locally free of finite type. The canonical functor
\[
\mathrm K^b(\Loclib(S))/\mathrm K^{b,\delta}(\Loclib(S))\longrightarrow \mathrm D(\Qcoh(S))
\]
is fully faithful and has as essential image $\mathrm D(\Qcoh(S))_{\parf}$.
\end{enumerate}

\emph{Proof.} Let $F\to G$ be an epimorphism in $\Qcoh(S)$, with $G$ of finite type. Since $F$ is the inductive limit of its quasi-coherent subsheaves of finite type (EGA I 9.4.9), there exists, by (2.1.1 a)), a quasi-coherent subsheaf $F'$ of finite type in $F$ such that $F'\to F\to G$ is an epimorphism. By (2.2.3 (iii)), there is an epimorphism $E\to F'$, with $E\in\Ob L(S)$; hence $E\to F'\to F\to G$ is an epimorphism. This proves that the category of quasi-coherent $\mathcal O_S$-modules of finite type is quasi-liftable in $\Qcoh(S)$, and a fortiori the same is true for $\Loclib(S)$ and for $L(S)$.

For the second assertion of a), apply (1.1) on the Zariski site of $S$, with $\mathcal C$ the fibered category $U\mapsto\Qcoh(U)$ and $\mathcal C_0$ the fibered category $U\mapsto L(U)$. Since every object of $L$ is pseudo-coherent, the notions of $n$-pseudo-coherence and pseudo-coherence in the fibered category $\mathrm D(\Qcoh)$ coincide with the same notions relative to $L$ (I 2.14). Condition (ii) of (1.1) is satisfied by (2.2.3 (iii)), and thus (1.1 a)) gives the part of a) concerning $L(S)$. The assertion for $\Loclib(S)$ is proved similarly. Part b) follows immediately from (1.1 b), c)), taking into account that, if $E\in\Ob\mathrm D(\Qcoh(S))$ is $n$-pseudo-coherent, then $E$ lies in $\mathrm D^{-}(\Qcoh(S))$, since $S$ is quasi-compact. An object of the fibered category $\Qcoh$ is locally in $L$ if and only if it is locally free of finite type; hence $L$ is locally liftable in $\Qcoh$ and locally stable under kernels of epimorphisms. Thus c) follows from (1.2), with $\mathcal C=\Qcoh$, $\mathcal C_0=L$, and $\mathcal C_1=\Loclib(S)$.

\textbf{Proposition 2.2.9.} The hypotheses and notation are those of (2.2.8), but assume moreover that $S$ is separated or noetherian.
\begin{enumerate}[label=\alph*)]
\item The canonical functors
\[
\mathrm K^{-}(L(S))/\mathrm K^{-,\delta}(L(S))\longrightarrow \mathrm D(S),
\qquad
\mathrm K^{-}(\Loclib(S))/\mathrm K^{-,\delta}(\Loclib(S))\longrightarrow \mathrm D(S)
\]
are fully faithful and have as essential image $\mathrm D^{-}(S)_{\coh}$. Moreover, if $S$ is of finite cohomological dimension in the sense of (3.7), for example if the underlying space of $S$ is noetherian of finite Krull dimension, then the functors
\[
\mathrm K(L(S))/\mathrm K^{\delta}(L(S))\longrightarrow
\mathrm K(\Loclib(S))/\mathrm K^{\delta}(\Loclib(S))\longrightarrow
\mathrm D^{-}(S)_{\coh}
\]
are equivalences of categories.
\item The canonical functor
\[
\mathrm K^b(\Loclib(S))/\mathrm K^{b,\delta}(\Loclib(S))\longrightarrow \mathrm D(S)
\]
is fully faithful and has as essential image $\mathrm D(S)_{\parf}$ (I 4.9). For $E\in\Ob\mathrm D(S)$ to have perfect amplitude contained in $[m,n]$ (I 4.8), it is necessary and sufficient that $E$ be isomorphic, in $\mathrm D(S)$, to a complex $F$ such that $F^i=0$ for $i\notin[m,n]$, $F^i\in\Ob L(S)$ for $i>m$, and $F^m$ is locally free of finite type.
\end{enumerate}

\emph{Proof.} In view of (3.5) and (3.7), this is an immediate corollary of (2.2.8). The only point to check, and it follows trivially from (3.5), respectively (3.7), is that for an object $E$ of $\mathrm D^{-}(\Qcoh(S))$, respectively of $\mathrm D(\Qcoh(S))$ when $S$ has finite cohomological dimension in the sense of (3.7), the notions of pseudo-coherence and perfection of (2.2.1) coincide with the usual notions (I 2.15 and 4.9) after regarding $E$ as an object of $\mathrm D(S)$.

\textbf{Remark 2.2.10.} a) If $S=\operatorname{Spec}A$ is affine, the family consisting of the single sheaf $\mathcal O_S$ is ample, and for it $L(S)=\Lib(S)$. It is also clear that the categories $\mathrm K^{-,\delta}(\Lib(S))$ and $\mathrm K^{-,\delta}(\Loclib(S))$ are reduced to zero objects. It follows from (2.2.9) that, writing $\Mod(A)$ for the category of $A$-modules and $\mathrm D(A)=\mathrm D(\Mod(A))$, the functor $R\Gamma_S:\mathrm D(S)\to\mathrm D(A)$ induces fully faithful functors
\[
\mathrm D^{-}(S)_{\coh}\longrightarrow \mathrm D(A),
\qquad
\mathrm D(S)_{\parf}\longrightarrow \mathrm D(A),
\]
whose essential images are respectively the triangulated subcategories $\mathrm D^{-}(A)_{\coh}$ and $\mathrm D(A)_{\parf}$, formed by complexes with bounded-above cohomology and pseudo-coherent, respectively perfect, relative to the subcategory $\Lib(A)$, respectively $\Loclib(A)$, of free, respectively projective, modules of finite type. There are commutative diagrams of equivalences of categories
\[
\begin{tikzcd}
\mathrm K^{-,b}(\Lib(S)) \arrow[r] \arrow[d,"\Gamma_S"'] &
\mathrm K^{-,b}(\Loclib(S)) \arrow[r] \arrow[d,"\Gamma_S"'] &
\mathrm D^b(S)_{\coh} \arrow[d,"R\Gamma_S"] \\
\mathrm K^{-,b}(\Lib(A)) \arrow[r] &
\mathrm K^{-,b}(\Loclib(A)) \arrow[r] &
\mathrm D^b(A)_{\coh},
\end{tikzcd}
\]
\[
\begin{tikzcd}
\mathrm K^b(\Loclib(S)) \arrow[r] \arrow[d,"\Gamma_S"'] &
\mathrm D(S)_{\parf} \arrow[d,"R\Gamma_S"] \\
\mathrm K^b(\Loclib(A)) \arrow[r] &
\mathrm D(A)_{\parf}.
\end{tikzcd}
\]
In the upper diagram, one may replace $b$ by $-$ when $S$ has finite cohomological dimension in the sense of (3.7).

b) For an arbitrary scheme $S$, it follows from a) that a complex of $\mathcal O_S$-modules which is pseudo-coherent and has bounded cohomology is locally isomorphic, in $\mathrm D(S)$, to a strictly pseudo-coherent complex, which was by no means evident a priori (cf. I 2.3.1).

\subsection*{2.3. Soft sheaves}

\textbf{2.3.1.} Recall a few definitions and results from [T.F.]. Let $X$ be a paracompact space. A sheaf of sets $F$ on $X$ is called \emph{soft} if, for every closed subset $Y$ of $X$, every section of $F$ over $Y$ extends to a section of $F$ over $X$. Soft abelian sheaves are acyclic: if $F$ is soft, then $H^i(X,F)=0$ for $i>0$ (T.F. II 4.4.3).

Let $A$ be a sheaf of rings on $X$. If $F$ is a soft $A$-module, then $F$ is generated by its sections (cf. 2.1.1): indeed, the homomorphism $A_X^{(I)}\to F$ defined by $I=H^0(X,F)$ is an epimorphism. Moreover, if $A$ is soft, every $A$-module is soft (T.F. II 3.7.1). Examples of soft sheaves of rings include the sheaf of continuous real or complex-valued functions on a paracompact space and the sheaf of $C^\infty$ functions on a $C^\infty$ manifold.

\textbf{Proposition 2.3.2.} Let $S$ be a compact space equipped with a soft sheaf of rings $A$. Write $\Mod(S)$ for the category of left $A$-modules and $\mathrm D(S)=\mathrm D(\Mod(S))$.
\begin{enumerate}[label=\alph*)]
\item The category $\Lib(S)$, respectively $\Loclib(S)$ (2.0), is liftable in $\Mod(S)$ and quasi-stable, respectively stable, under kernels of epimorphisms (I 1.2). This assertion holds more generally for $S$ paracompact.
\item Let $E\in\Ob\mathrm D(S)$ and $n\in\mathbb Z$. For $E$ to be $n$-pseudo-coherent, respectively pseudo-coherent, it is necessary and sufficient that it be so relative to $\Lib(S)$ or to $\Loclib(S)$. The canonical functors
\[
\mathrm K^{-}(\Lib(S))\longrightarrow \mathrm D(S),
\qquad
\mathrm K^{-}(\Loclib(S))\longrightarrow \mathrm D(S)
\]
are fully faithful and have as essential image $\mathrm D^{-}(S)_{\coh}$ (I 2.15).
\item Let $E\in\Ob\mathrm D(S)$ and let $[m,n]$ be an interval of $\mathbb Z$. For $E$ to have perfect amplitude contained in $[m,n]$ (I 4.8), it is necessary and sufficient that there exist a quasi-isomorphism $F\to E$, where $F$ is a complex such that $F^i=0$ for $i\notin[m,n]$, $F^i\in\Ob\Lib(S)$ for $i>m$, and $F^m\in\Ob\Loclib(S)$. The canonical functor
\[
\mathrm K^b(\Loclib(S))\longrightarrow \mathrm D(S)
\]
is fully faithful and has as essential image $\mathrm D(S)_{\parf}$ (I 4.9).
\item Let $\Mod(A(S))$ be the category of left $A(S)$-modules and $\mathrm D(A(S))$ the corresponding derived category. The functor
\[
H^0(S,-):\Mod(S)\longrightarrow\Mod(A(S))
\]
is exact and is an equivalence of categories. Moreover it defines equivalences of categories
\[
\mathrm D(S)\xrightarrow{\sim}\mathrm D(A(S)),\qquad
\mathrm D^{-}(S)_{\coh}\xrightarrow{\sim}\mathrm D^{-}(A(S))_{\coh},\qquad
\mathrm D(S)_{\parf}\xrightarrow{\sim}\mathrm D(A(S))_{\parf},
\]
where pseudo-coherence, respectively perfection, in $\mathrm D(A(S))$ is understood relative to the category of free, respectively projective, modules of finite type. It also gives an equivalence from the full subcategory of $\mathrm D(S)$ formed by $n$-pseudo-coherent objects, respectively by objects of perfect amplitude contained in $[m,n]$, onto the analogous full subcategory of $\mathrm D(A(S))$.
\end{enumerate}

\emph{Proof.} The first assertion of a) follows from (2.1.2 a)), applied with $S$ the topos of sheaves of sets on $S$ and $\mathcal C'_S=\Mod(S)$, using (2.3.1). The second follows from (1.1 a)), respectively (1.2 a)), with $\mathcal C=\Mod$ and $\mathcal C_0=\Lib$, respectively $\Loclib$; the condition (ii) of (1.1) is verified by (2.3.1) and (2.1.1 b)). Part b) follows from (1.1 b), c)) with the same $\mathcal C$ and $\mathcal C_0$. Part c) follows from (1.2 b), c), d)), with $\mathcal C=\Mod$ and $\mathcal C_0=\Lib$. The first assertion of d) follows, using (2.3.1), from (2.1.2 b) (ii)) with $\mathcal C'_S=\mathcal C_S=\Mod(S)$. The other assertions in d) follow from it, together with b) and c), and from the fact that $\Gamma_S$ carries free modules, respectively direct factors of free modules, of finite type to free, respectively projective, modules of finite type.

\subsection*{2.4. Stein spaces}

\textbf{Definition 2.4.1.} Let $S$ be a locally ringed topos (I 2.15.1 (ii)), with final object also denoted $S$. We say that $S$ is \emph{Stein} if the structural sheaf $\mathcal O_S$ is coherent (I 3.1) and if the following properties hold:

\emph{Theorem A:} every coherent $\mathcal O_S$-module (I 3.1) is generated by its sections;

\emph{Theorem B:} for every coherent $\mathcal O_S$-module $F$, one has $H^i(S,F)=0$ for $i>0$.

\textbf{Remark 2.4.1.1.} If $S$ has enough closed points $s$ such that the ideal $\mathfrak m_s$ defined by $s$ is coherent, then the relation $H^1(S,F)=0$ for every coherent sheaf $F$ implies Theorem A. Indeed, let $F$ be a coherent $\mathcal O_S$-module and let $s$ be a closed point of $S$. There is an exact sequence
\[
0\longrightarrow \mathfrak m_sF\longrightarrow F\longrightarrow F/\mathfrak m_sF\longrightarrow0,
\]
where $F/\mathfrak m_sF$ is a coherent sheaf concentrated at $s$. Applying $H^0(S,-)$ and writing the cohomology exact sequence, one finds that
\[
H^0(S,F)\longrightarrow H^0(S,F/\mathfrak m_sF)
\]
is an epimorphism. Since
\[
H^0(S,F/\mathfrak m_sF)=F_s/\mathfrak m_sF_s,
\]
where $\mathfrak m_s$ is the maximal ideal of $\mathcal O_s$, Nakayama's lemma gives the conclusion. The editor does not know whether the preceding implication is true in general when one assumes only that $\mathcal O_S$ is a coherent sheaf of local rings. The editor also does not know whether the relation $H^1(S,F)=0$ for every coherent sheaf $F$ implies Theorem B.

\textbf{Examples 2.4.2.} A complex analytic Stein space, a compact convex subset of $\mathbb C^n$ equipped with the sheaf induced by the sheaf of holomorphic functions, the spectrum of a noetherian ring, and an affine rigid analytic space define Stein topoi.

\textbf{Lemma 2.4.3.} Let $S$ be a Stein topos (2.4.1). Suppose that the final object $S$ is quasi-compact (2.1.1), and let $\Coh(S)$ be the abelian category of coherent $\mathcal O_S$-modules, that is, modules of finite presentation. Then:
\begin{enumerate}[label=\alph*)]
\item $\Lib(S)$, respectively $\Loclib(S)$, is liftable in $\Coh(S)$ and quasi-stable, respectively stable, under kernels of epimorphisms (I 1.2).
\item The canonical functors
\[
\mathrm K^{-}(\Lib(S))\longrightarrow \mathrm K^{-}(\Loclib(S))\longrightarrow \mathrm D^{-}(\Coh(S))
\]
are equivalences of categories.
\item Put $A=\Gamma(S,\mathcal O_S)$, and write $\Mod(A)$ for the category of $A$-modules and $\mathrm D(A)$ for the corresponding derived category. The functor
\[
\Gamma_S:\Coh(S)\longrightarrow\Mod(A)
\]
is exact, fully faithful, and has as essential image the strictly full additive subcategory of $\Mod(A)$ formed by $A$-modules of finite presentation. It induces a fully faithful functor
\[
\mathrm D^{-}(\Coh(S))\longrightarrow\mathrm D(A),
\]
whose essential image is the full subcategory $\mathrm D^{-}(A)_{\coh}$ formed by complexes pseudo-coherent relative to the subcategory of free modules of finite type.
\end{enumerate}
Here ``enough closed points'' means enough fiber functors (SGA 4 IV 6.4.1) of the form $i^*:S\to s$, with $i_*$ the inclusion of a closed subtopos of $S$; and $\mathfrak m_s$ denotes the sections of $\mathcal O_S$ whose fiber at $s$ belongs to the maximal ideal of $\mathcal O_s$.

\emph{Proof.} For the first assertion of a), apply (2.1.2 a)) with $\mathcal C'_S=\Coh(S)$; the hypothesis of (2.1.2) is verified by Theorem B, since if $E$ is coherent and $G$ is locally free of finite type, then $\operatorname{Hom}(G,E)$ is coherent. Thus $\Loclib(S)$, and a fortiori $\Lib(S)$, is liftable in $\Coh(S)$. The second assertion of a) follows from (1.1 a)), respectively (1.2 a)), applied with $\mathcal C_0=\Lib$, respectively $\Loclib$; condition (ii) of (1.1) is verified by Theorem A and (2.1.1 b)). Part b) follows from (1.1 c)), applied successively with $\mathcal C_0=\Lib$ and $\mathcal C_0=\Loclib$. The first assertion of c) follows from (2.1.2 b)(i)), applied with $\mathcal C'_S=\Coh(S)$. To prove that $\Gamma_S:\mathrm D^{-}(\Coh(S))\to\mathrm D(A)$ is fully faithful, it is enough, by b), to prove that the composite
\[
\mathrm K^{-}(\Lib(S))\longrightarrow\mathrm D^{-}(\Coh(S))\longrightarrow\mathrm D(A)
\]
is so. But this composite factors as
\[
\mathrm K^{-}(\Lib(S))\xrightarrow{\Gamma_S}\mathrm K^{-}(\Lib(A))\longrightarrow\mathrm D(A),
\]
where $\Lib(A)$ denotes the category of free $A$-modules of finite type; the first arrow is an equivalence of categories, and the second is fully faithful by (I 2.7 b)), applied with $\mathcal C=\Mod(A)$ and $\mathcal C_0=\Lib(A)$, since $\mathrm K^{-,\delta}(\Lib(A))=0$. This proves the last assertion of c).

\textbf{Lemma 2.4.4.} Under the hypotheses of (2.4.3), the canonical functor
\[
\mathrm D^b(\Coh(S))\longrightarrow\mathrm D(S)
\]
is fully faithful and has as essential image $\mathrm D^b(S)_{\coh}$ (I 2.15).

\emph{Proof.} Apply (2.1.3) with $\mathcal A=\Coh(S)$, $\mathcal B=\Mod(S)$, and $\varphi:\mathcal A\to\mathcal B$ the inclusion functor. Since this functor is fully faithful, it is enough, by (2.1.3 a)(ii)), to check that for fixed $M\in\Ob\Coh(S)$ and $n>0$ the functor
\[
L\longmapsto \Ext^n_{\Mod(S)}(L,M)
\]
is coeffaceable on $L\in\Ob\Coh(S)$. By Theorem A and (2.1.1 b)), every coherent $\mathcal O_S$-module is a quotient of a free module of finite type; and if $L$ is free of finite type, then $\Ext^n_{\Mod(S)}(L,M)=0$ by Theorem B. The assertion follows. Hence the functor $\mathrm D^b(\Coh(S))\to\mathrm D(S)$ is fully faithful and has as essential image the full subcategory of $\mathrm D(S)$ formed by complexes with bounded coherent cohomology, which is precisely $\mathrm D^b(S)_{\coh}$ since $\mathcal O_S$ is coherent (I 3.5).

\textbf{Proposition 2.4.5.} With the hypotheses and notation of (2.4.3), one has:
\begin{enumerate}[label=\alph*)]
\item The canonical functors
\[
\mathrm K^{-,b}(\Lib(S))\longrightarrow\mathrm K^{-,b}(\Loclib(S))\longrightarrow \mathrm D^b(S)_{\coh}
\]
are equivalences of categories. Moreover, the functor $R\Gamma_S:\mathrm D^b(S)\to\mathrm D^b(A)$ induces a fully faithful functor
\[
R\Gamma_S:\mathrm D^b(S)_{\coh}\longrightarrow \mathrm D^b(A),
\]
whose essential image is $\mathrm D^b(A)_{\coh}$, the triangulated subcategory of $\mathrm D(A)$ formed by complexes with bounded cohomology and pseudo-coherent relative to the category of free modules of finite type.
\item Let $E\in\Ob\mathrm D(S)$, and let $[m,n]$ be an interval of $\mathbb Z$. If $E$ has perfect amplitude contained in $[m,n]$ (I 4.5), then $E$ is isomorphic to a complex $F$ such that $F^i=0$ for $i\notin[m,n]$, $F^i\in\Ob\Lib(S)$ for $i>m$, and $F^m\in\Ob\Loclib(S)$.
\item The canonical functor
\[
\mathrm K^b(\Loclib(S))\longrightarrow \mathrm D(S)_{\parf}
\]
is an equivalence of categories (I 4.9). Moreover, $R\Gamma_S:\mathrm D^b(S)\to\mathrm D^b(A)$ induces a fully faithful functor
\[
R\Gamma_S:\mathrm D(S)_{\parf}\longrightarrow\mathrm D^b(A),
\]
whose essential image is $\mathrm D^b(A)_{\parf}$, the triangulated subcategory of $\mathrm D(A)$ formed by complexes perfect relative to the category of projective modules of finite type.
\end{enumerate}

\emph{Proof.} a) The functors
\[
\mathrm K^{-,b}(\Lib(S))\longrightarrow\mathrm K^{-,b}(\Loclib(S))\longrightarrow
\mathrm D^b(\Coh(S))\longrightarrow\mathrm D^b(A)_{\coh}
\]
are equivalences of categories by (2.4.3 b), c). Since the fully faithful functor $\Gamma_S:\mathrm D^b(\Coh(S))\to\mathrm D(A)$ factors as
\[
\mathrm D^b(\Coh(S))\longrightarrow\mathrm D^b(S)_{\coh}\xrightarrow{R\Gamma_S}\mathrm D(A),
\]
the conclusion follows from (2.4.4).

b) Suppose $\operatorname{parf.amp}(E)\subset[m,n]$. In particular $E$ is pseudo-coherent; by a), after replacing $E$ by an isomorphic object, we may assume that $E^i\in\Ob\Lib(S)$ for $i>m$ and $E^i=0$ for $i>n$. Truncating, we may also suppose that $E^i=0$ for $i<m$. By (I 4.16), applied with $\mathcal C=\Mod$ and $\mathcal C_0=\Lib$, one then has $E^m\in\Ob\Loclib(S)$.

c) The hypotheses (1.1 (i)) and, respectively, (1.1 (ii)) are verified for $\mathcal C=\Coh(S)$ and $\mathcal C_0=\Lib$, by (2.4.3 a)) and by (2.1.1 b)) together with Theorem A. Hence (1.2 c), d)) shows that the functor
\[
\mathrm K^b(\Loclib(S))\longrightarrow\mathrm D^b(\Coh(S))
\]
is fully faithful. Therefore, by (2.4.4), the same is true for
\[
\mathrm K^b(\Loclib(S))\longrightarrow\mathrm D(S),
\]
and its essential image is $\mathrm D(S)_{\parf}$ by b). By an argument analogous to that used in the proof of (2.4.3 c)), one sees that the functor $\mathrm K^b(\Loclib(S))\to\mathrm D(S)$ induced by $\Gamma_S$ is fully faithful and has essential image $\mathrm D^b(A)_{\parf}$; one finishes the proof as in a).

\textbf{Remark 2.4.6.} The editor does not know whether in (2.4.4) one may replace $\mathrm D^b$ by $\mathrm D^{-}$, and consequently whether (2.4.5 a)) remains true when one replaces $\mathrm K^{-,b}$ by $\mathrm K^{-}$ and $\mathrm D^b$ by $\mathrm D^{-}$.

We finish with an interesting complement to (2.4.5). First we give a definition.

\textbf{Definition 2.4.7.} A \emph{complex pseudo-analytic space} is a ringed space locally isomorphic to a ringed space of the form $(Y,\mathcal O_X|Y)$, where $X$ is a complex analytic space and $Y$ is a closed subset of $X$.

The structural sheaf of a complex pseudo-analytic space is therefore coherent. We say that a complex pseudo-analytic space is Stein if the corresponding topos is Stein (2.4.1).



% Exposé II continuation macros
\providecommand{\Qcoh}{\operatorname{Qcoh}}
\providecommand{\Mod}{\operatorname{Mod}}
\providecommand{\Coh}{\operatorname{Coh}}
\providecommand{\Ob}{\operatorname{Ob}}
\providecommand{\Hom}{\operatorname{Hom}}
\providecommand{\Ext}{\operatorname{Ext}}
\providecommand{\Spec}{\operatorname{Spec}}
\providecommand{\qcoh}{\mathrm{qcoh}}
\providecommand{\parf}{\mathrm{parf}}
\providecommand{\coh}{\mathrm{coh}}
\providecommand{\Loclib}{\operatorname{Loclib}}
\providecommand{\Lib}{\operatorname{Lib}}
\providecommand{\im}{\operatorname{Im}}

\section*{Exposé II. Global resolutions (continued)}

\textbf{Proposition 2.4.8.} Let $S$ be a compact Stein complex pseudo-analytic space (2.4.7). Suppose that, for every open $U$ of $S$ and every reduced closed subspace $X$ of $U$, the normalization of $X$ is locally connected. This condition is verified, for example, when $S$ is locally defined by real analytic equations or inequalities in a space $\mathbb R^n$. Then
\[
 A=\Gamma(S,\mathcal O_S)
\]
is a noetherian ring.

\emph{Proof.} It is enough to show that, if $M$ is an $A$-module of finite presentation, then the set of submodules of $M$, or equivalently the set of submodules of finite type of $M$, satisfies the maximal condition. Indeed, if $N$ is a submodule of finite type of $M$, then $M/N$ is of finite presentation by I 2.6 and 2.9, so that $N$ is itself of finite presentation by the equivalence (2.4.3 c). Using again the equivalence (2.4.3 c) between the category of coherent $\mathcal O_S$-modules and the category of $A$-modules of finite presentation, together with the compactness of $S$, one is reduced to proving the following lemma.

\textbf{Lemma 2.4.8.1.} Let $S$ be a complex pseudo-analytic space satisfying the local connectedness condition of (2.4.8). Let $F$ be a coherent $\mathcal O_S$-module, let $(E_n)$ be an increasing sequence of coherent sub-$\mathcal O_S$-modules of $F$, and let $s$ be a point of $S$. Then there exists an open neighbourhood $U$ of $s$ such that the sequence $(E_n|U)$ is stationary.

\emph{Proof.} Since the ring $\mathcal O_{S,s}$ is noetherian, there exists an integer $n_0$ such that $(E_n)_s=(E_{n_0})_s$ for $n\geq n_0$. Replacing the sequence $(E_n)$ by the sequence $(E_n/E_{n_0})$, one may therefore suppose that $(E_n)_s=0$ for every $n$. On the other hand $F_s=M$ is a finite $\mathcal O_{S,s}$-module; hence it has a finite filtration
\[
 M=M_0\supset M_1\supset\cdots\supset M_q=0
\]
such that
\[
 M_i/M_{i+1}\simeq \mathcal O_s/\mathfrak p_i,
\]
where $\mathfrak p_i$ is a prime ideal of $\mathcal O_s$. Since the category of finite-type $\mathcal O_s$-modules is equivalent to the category of germs at $s$ of coherent $\mathcal O_S$-modules, after shrinking around $s$ one may suppose that the $M_i$, respectively the $\mathfrak p_i$, are the fibres at $s$ of coherent sub-$\mathcal O_S$-modules $F_i$, respectively coherent ideals $\mathfrak p_i$, of $F$, respectively of $\mathcal O_S$, and that one has
\[
 F=F_0\supset F_1\supset\cdots\supset F_q=0,
 \qquad F_i/F_{i+1}\simeq\mathcal O_S/\mathfrak p_i.
\]
To prove that the sequence $(E_n)$ is stationary in a neighbourhood of $s$ is therefore equivalent to proving the same assertion for the induced sequences on the graded pieces of this filtration. Thus one may suppose that $F=\mathcal O_S$ and that $\mathcal O_{S,s}$ is an integral domain; the $E_n$ then form an increasing sequence of ideals of $\mathcal O_S$ with $(E_n)_s=0$.

Since $\mathcal O_{S,s}$ is an integral domain, hence in particular reduced, one may moreover suppose, by Oka's theorem, after localizing at $s$, that $S$ is reduced. We have to show that there is an open neighbourhood $U$ of $s$ such that $E_n|U=0$ for $n$ large enough. Consider, to do this, the normalization $\widetilde S$ of $S$. It is the analytic spectrum of a finite $\mathcal O_S$-algebra $\widetilde{\mathcal O}_S$ whose fibre at each $t\in S$ is the integral closure of $\mathcal O_{S,t}$. Since $E_n\subset E_n\widetilde{\mathcal O}_S$, it is enough to prove that $E_n\widetilde{\mathcal O}_S=0$ for $n$ large. This reduces us to the case where $S$ is normal.

In that case, since $(E_n)_s=0$, the closed immersion
\[
 i_n:S_n\longrightarrow S
\]
defined by $E_n$ is étale at $s$. But, $S$ being normal, the set of points at which $i_n$ is étale is both open and closed by the lemma below. Hence its image in $S$ is simultaneously open and closed. The restriction of $E_n$ to the connected component of $s$ in $S$ is therefore zero; this component is open because $S$ is locally connected. The lemma is proved.

\textbf{Lemma 2.4.8.2} (cf. EGA IV 18.10.3). Let $Y$ be a normal pseudo-analytic space, and let $f:X\to Y$ be a non-ramified morphism. The set of points of $X$ at which $f$ is étale is both open and closed.

\emph{Proof.} It is immediate that the set of points at which $f$ is étale is open. It remains to show that the set of points at which $f$ is not étale is open. Since a non-ramified morphism is locally an immersion, one may suppose that $f$ is a closed immersion defined by a coherent ideal $I$. Let $x$ be a point of $X$ such that $I_x\ne0$. Let $\mathfrak p$ be a prime ideal of $\mathcal O_{Y,x}$ containing $I_x$. One has
\[
 \dim(\mathcal O_{Y,x}/\mathfrak p)>0,
\]
for otherwise $\mathfrak p=0$, since $\mathcal O_{Y,x}$ is integral. Thus $Y$ has positive codimension at $x$, hence positive codimension in a neighbourhood of $x$; the assertion follows.

\section*{3. Complements on quasi-coherent sheaves on schemes}

\textbf{3.0.} In this section we fix a quasi-compact and quasi-separated scheme $S$. The canonical, fully faithful inclusion functor
\[
 \iota:\Qcoh(S)\longrightarrow\Mod(S)
\]
extends to a functor, again denoted $\iota$, from $\mathrm D(\Qcoh(S))$ to $\mathrm D(S)=\mathrm D(\Mod(S))$, which we shall study.

\textbf{Lemma 3.1.} The category $\Qcoh(S)$ is an abelian category satisfying AB5, and the quasi-coherent sheaves of finite type form a small family of generators. In particular $\Qcoh(S)$ has enough injectives.

\emph{Proof.} The assertion about generators follows at once from the fact that every quasi-coherent sheaf on $S$ is the filtered inductive limit of its quasi-coherent subsheaves of finite type (EGA I 9.4.9 and IV 1.7.7). The other assertions are evident.

\textbf{Remark 3.1.1.} The editor does not know whether $\Qcoh(S)$ still has enough injectives when one assumes only that $S$ is quasi-separated.

\textbf{Lemma 3.2.} The functor
\[
 \iota:\Qcoh(S)\longrightarrow\Mod(S)
\]
has a right adjoint
\[
 Q:\Mod(S)\longrightarrow\Qcoh(S),
\]
called the coherer.

\emph{Proof.} If $S$ is affine, the existence of $Q$ is immediate: one takes $Q=\widetilde{\Gamma_S}$ (EGA I 1.3.7). Let now $f:U\to S$ be an affine open of $S$. It is clear that, for $E\in\Ob\Qcoh(S)$ and $F\in\Ob\Mod(U)$, one has a canonical isomorphism
\begin{equation}\tag{1}
 \Hom(E,f_*F)\simeq \Hom(E,f_*Q(F)).
\end{equation}
Let $(f_i:U_i\to S)$ be a finite covering of $S$ by affine opens, and, for each pair $(i,j)$, let $(U_{ijk}\to U_i\cap U_j)_k$ be a finite covering of $U_i\cap U_j$ by affine opens, which is possible since $S$ is quasi-separated. For $F\in\Ob\Mod(S)$ one has an exact sequence
\begin{equation}\tag{2}
 0\longrightarrow F\longrightarrow \prod_i f_{i*}(F|U_i)
 \longrightarrow \prod_{i,j,k}f_{ijk*}(F|U_{ijk}),
\end{equation}
where $f_{ijk}$ denotes the inclusion of $U_{ijk}$ into $S$. If $g:V\to U$ is an inclusion of affine opens of $S$, there is an evident natural arrow
\[
 Q(F|U)\longrightarrow g_*Q(F|V).
\]
Consequently (2) gives an arrow
\begin{equation}\tag{3}
 \prod_i f_{i*}Q(F|U_i)\longrightarrow
 \prod_{i,j,k}f_{ijk*}Q(F|U_{ijk}),
\end{equation}
whose kernel we denote by $QF$. If $E\in\Ob\Qcoh(S)$, applying the functor $\Hom(E,\cdot)$ to (2), and using (1), one obtains a functorial isomorphism in $E$,
\[
 \Hom(E,F)\simeq\Hom(E,QF).
\]
This proves the existence of $Q$; its uniqueness is evident, since it is an adjoint.

\textbf{Lemma 3.3.} \begin{enumerate}[label=\alph*)]
\item The coherer is an additive functor and transforms injectives into injectives.
\item Let
\[
 RQ:\mathrm D^+(S)\longrightarrow \mathrm D^+(\Qcoh(S))
\]
be the right derived functor of the coherer. For $E\in\Ob\mathrm D(\Qcoh(S))$ and $F\in\Ob\mathrm D^+(S)$ there are canonical functorial isomorphisms
\begin{align}
 R\Hom_{\Mod}(E,F)&\simeq R\Hom_{\Qcoh}(E,RQ(F)), \tag{i}\\
 \Hom_{\Mod}(E,F)&\simeq \Hom_{\Qcoh}(E,RQ(F)). \tag{ii}
\end{align}
\end{enumerate}

\emph{Proof.} Assertion a) is an immediate consequence of the definition. For b), it is enough to establish the isomorphism (i), since (ii) follows by passing to cohomology. For this purpose one may suppose that $F$ has injective components. Then $QF$ has injective components by a), and the isomorphism (i) follows directly from the definition of $Q$.

\textbf{Lemma 3.4.} Let $f:X\to Y$ be a morphism of quasi-compact and quasi-separated schemes. For $E\in\Ob\Mod(X)$ there is a canonical functorial isomorphism
\[
 Q_Y f_*(E)\simeq f_*Q_X(E).
\]

\emph{Proof.} Since $f$ is quasi-compact and quasi-separated (EGA IV 1.2.3 and 1.2.4), $f_*$ transforms quasi-coherent sheaves into quasi-coherent sheaves (EGA III 1.4.10 and IV 1.7.21). The functors $Q_Yf_*$ and $f_*Q_X$ from $\Mod(X)$ to $\Qcoh(Y)$ are therefore both right adjoints to the functor
\[
 f^*: \Qcoh(Y)\longrightarrow \Mod(X),
\]
and hence are canonically isomorphic.

\textbf{Proposition 3.5.} If $S$ is noetherian, or quasi-compact and separated, the functor
\[
 \iota:\mathrm D^+(\Qcoh(S))\longrightarrow \mathrm D^+(S)
\]
is fully faithful and has as essential image the full subcategory $\mathrm D^+(S)_{\qcoh}$ of $\mathrm D^+(S)$ formed by complexes with quasi-coherent cohomology.

\emph{Proof.} One may restrict to proving the assertion in the case where $S$ is separated, since the noetherian case is already known. It is plainly equivalent to prove the following two assertions.

\begin{enumerate}[label=(3.5.\arabic*)]
\item For $E\in\Ob\mathrm D^+(\Qcoh(S))$, the adjunction arrow
\[
 E\longrightarrow RQ(E)
\]
is an isomorphism.
\item For $E\in\Ob\mathrm D^+(S)_{\qcoh}$, the adjunction arrow
\[
 RQ(E)\longrightarrow E
\]
is an isomorphism.
\end{enumerate}

\emph{Proof of (3.5.1).} By an immediate dévissage one is reduced to proving that quasi-coherent sheaves are acyclic for the coherer, that is, for $E\in\Ob\Qcoh(S)$, one has $R^iQ(E)=0$ for $i>0$, and that, moreover, the adjunction arrow $E\to QE$ is an isomorphism. If $S$ is affine, this follows from the formula $Q=\widetilde{\Gamma_S}$ used in the proof of (3.2). In the general case, cover $S$ by a finite number of affine opens $U_i$. Since $S$ is separated, the inclusions
\[
 f_{i_0\ldots i_p}:U_{i_0}\cap\cdots\cap U_{i_p}\longrightarrow S
\]
are affine morphisms. Let $E\in\Ob\Qcoh(S)$, and consider the resolution of $E$ defined by this covering,
\begin{equation}\tag{*}
 0\longrightarrow E\longrightarrow \prod_i f_{i*}f_i^*E\longrightarrow\cdots\longrightarrow
 \prod_{i_0\ldots i_p}f_{i_0\ldots i_p*}f_{i_0\ldots i_p}^*E\longrightarrow\cdots .
\end{equation}
It gives a spectral sequence
\begin{equation}\tag{**}
 E_1^{pq}=\prod_{i_0\ldots i_p}R^qQ\bigl(f_{i_0\ldots i_p*}f_{i_0\ldots i_p}^*E\bigr)
 \Longrightarrow R^{p+q}Q(E).
\end{equation}
Since $f_{i_0\ldots i_p}$ is affine, (3.4) gives, for every $q$,
\begin{equation}\tag{***}
 R^qQ\bigl(f_{i_0\ldots i_p*}f_{i_0\ldots i_p}^*E\bigr)
 \simeq f_{i_0\ldots i_p*}R^qQ\bigl(f_{i_0\ldots i_p}^*E\bigr).
\end{equation}
In particular $E_1^{pq}=0$ for $q>0$, by the affine case already treated. From (**) one deduces isomorphisms
\[
 R^nQ(E)\simeq H^n(QE')\quad(n\geq0),
\]
where $E'$ is the resolution of $E$ defined by the exact sequence (*). But by (***) one has $QE^n\simeq E^n$ for $n\geq0$; hence finally
\[
 R^nQ(E)=0\quad(n>0),
\]
and $Q(E)\simeq E$. This proves (3.5.1).

\emph{Proof of (3.5.2).} The spectral sequence
\[
 R^pQ(H^q(E))\Longrightarrow R^{p+q}Q(E)
\]
reduces us immediately to the case where $E\in\Ob\Qcoh(S)$. The composite
\[
 E\longrightarrow RQE\longrightarrow E
\]
of the adjunction arrows is the identity, and the assertion follows from (3.5.1). This completes the proof of (3.5).

\textbf{Remark 3.6.} The conclusion of (3.5) would be false if one omitted the hypothesis ``separated or noetherian'', as is shown by a neat counterexample of J.-L. Verdier; see Appendix I to the present exposé.

\textbf{Proposition 3.7.} Under the hypothesis of (3.5), suppose moreover that there is an integer $N$ such that, for every $F\in\Ob\Mod(S)$ and every quasi-compact open $U$ of $S$, one has
\[
 H^i(U,F)=0\quad\text{for }i>N.
\]
Then:
\begin{enumerate}[label=\alph*)]
\item The coherer has finite cohomological dimension.
\item The functor
\[
 \iota:\mathrm D(\Qcoh(S))\longrightarrow \mathrm D(S)
\]
is fully faithful and has as essential image the full subcategory $\mathrm D(S)_{\qcoh}$ of $\mathrm D(S)$ formed by complexes with quasi-coherent cohomology.
\end{enumerate}

\emph{Proof of a).} If $S$ is affine, the formula $Q=\widetilde{\Gamma_S}$ implies that $Q$ is of cohomological dimension $\leq N$. We pass to the general case. Let us show that there is an integer $a$ such that $R^qQ(E)=0$ for every $E\in\Ob\Mod(S)$ and every $q>a$. The hypothesis implies that every $E\in\Ob\Mod(S)$ admits a finite right resolution $E'$ of length $N$ such that the $E^i$ are acyclic, that is, $H^q(U,E^i)=0$ for every open $U$ of $S$ and every $q>0$. Thus it is enough to prove that there is an integer $b$ such that
\[
 R^qQ(E)=0
\]
for every acyclic $\mathcal O_S$-module $E$ and every $q>b$; one may then take $a=b+N$.

Let $(U_i)$ be a finite covering of $S$ by $N'$ affine opens. Let $E$ be an acyclic $\mathcal O_S$-module and let $C'(E)$ be the resolution of $E$ defined by the covering $(U_i)$. There are two cases.

If $S$ is separated, the components of $C'(E)$ are sums of terms of the form $f_*f^*E$, where $f:U\to S$ is the inclusion of an affine open, hence an affine morphism. One has
\[
 RQ(Q_Xf_*)f^*E\simeq RQ\,Rf_*f^*E\simeq RQ\,f_*f^*E,
\]
since $f^*E$ is acyclic. On the other hand, by (3.4),
\[
 R(Q_Xf_*)f^*E\simeq R(f_*Q_X)f^*E\simeq f_*RQf^*E,
\]
since $f$ is affine. Hence
\[
 R^nQ(f_*f^*E)\simeq f_*R^nQ(f^*E)\quad(n\geq0).
\]
By the affine case,
\[
 R^nQ(f_*f^*E)=0\quad(n>N).
\]
Since $C'(E)$ has length $N'$, it follows that
\[
 R^nQ(E)=0\quad(n>N+N').
\]

If $S$ is noetherian, the components of $C'(E)$ are sums of terms of the same form $f_*f^*E$, where now $f:U\to S$ is the inclusion of a quasi-compact and separated open. As in the first case, one has
\[
 R(Q_Xf_*)f^*E\simeq RQ\,f_*f^*E.
\]
Since $U$ is noetherian, $\Qcoh(U)$ has enough injective objects in $\Mod(U)$, and hence the right derived functor of $f_*:\Qcoh(U)\to\Qcoh(S)$ is induced by $Rf_*:\mathrm D^+(U)\to\mathrm D^+(S)$ when $\mathrm D^+(\Qcoh(U))$ and $\mathrm D^+(\Qcoh(S))$ are identified with full subcategories of $\mathrm D^+(U)$ and $\mathrm D^+(S)$ by (3.5). Thus
\[
 R(f_*Q_X)f^*E\simeq Rf_*RQf^*E.
\]
One deduces a spectral sequence
\[
 E_2^{pq}=R^pf_*R^qQ(f^*E)\Longrightarrow R^{p+q}(f_*Q)f^*E.
\]
By the first case there is an integer $n(U)$ depending only on $U$ such that $Q$ is of cohomological dimension $<n(U)$ on $\Mod(U)$; hence $E_2^{pq}=0$ for $q>n(U)$. Also $E_2^{pq}=0$ for $p>N$, by the cohomological-dimension hypothesis on $S$. Thus
\[
 R^n(f_*Q)f^*E=0\quad(n>N+N''),
\]
where $N''$ is an integer greater than all the $n(U)$, with $U$ running over the intersections $U_{i_0}\cap\cdots\cap U_{i_p}$. Taking account of the preceding displayed isomorphism, one obtains
\[
 R^nQ(f_*f^*E)=0\quad(n>N+N''),
\]
and hence finally
\[
 R^nQ(E)=0\quad(n>N+N'+N'').
\]
This proves a).

\emph{Proof of b).} The coherer, being of finite cohomological dimension by a), admits a right derived functor
\[
 RQ:\mathrm D(S)\longrightarrow\mathrm D(\Qcoh(S))
\]
whose restriction to $\mathrm D^+(S)$ coincides with the functor $RQ$ defined above. For $E\in\Ob\mathrm D(\Qcoh(S))$ there is a canonical functorial isomorphism
\begin{equation}\tag{1}
 E\xrightarrow{\sim}RQ(E),
\end{equation}
since, by (3.5.1), the $E^i$ are acyclic for $Q$ and satisfy $E^i\xrightarrow{\sim}Q(E^i)$. On the other hand, for $E\in\Ob\mathrm D(S)$ one has a canonical functorial homomorphism
\begin{equation}\tag{2}
 RQ(E)\longrightarrow E,
\end{equation}
which, in the case where the $E^i$ are $Q$-acyclic, comes from the ordinary adjunction homomorphism between $\iota$ and $Q$. It is checked directly that these homomorphisms make
\[
 RQ:\mathrm D(S)\longrightarrow\mathrm D(\Qcoh(S))
\]
a right adjoint of
\[
 \iota:\mathrm D(\Qcoh(S))\longrightarrow\mathrm D(S).
\]
The fact that (1) is an isomorphism implies that $\iota$ is fully faithful. Finally, using the spectral sequence
\[
 R^pQ(H^q(E))\Longrightarrow R^{p+q}Q(E),
\]
which converges since $Q$ has finite cohomological dimension, one sees that (2) is an isomorphism for $E\in\Ob\mathrm D(S)_{\qcoh}$; hence $\iota$ has image $\mathrm D(S)_{\qcoh}$. This proves b) and completes the proof of (3.7).

\section*{Appendix I. A counterexample of Verdier}

\subsection*{0. Introduction}
If $S$ is a scheme, write $\Mod(S)$, respectively $\Qcoh(S)$, for the category of $\mathcal O_S$-modules, respectively of quasi-coherent $\mathcal O_S$-modules. The purpose of this appendix is to draw attention to certain pathological properties of the category $\mathrm D(\Qcoh(S))$, and in particular to prove the following facts.

\textbf{0.1.} Let $X$ be an affine scheme with ring $A$, and let $I$ be an injective $A$-module. We know that, if $A$ is noetherian, the sheaf $\widetilde I$ is an injective $\mathcal O_X$-module. If $A$ is not noetherian, $\widetilde I$ is not necessarily injective, nor even flasque, even if the underlying space of $X$ is noetherian of finite Krull dimension. Moreover, the restriction of $\widetilde I$ to an open $U$ of $X$ is not, in general, an injective object of $\Qcoh(U)$.

\textbf{0.2.} Let $f:X\to Y$ be a morphism of quasi-compact and quasi-separated schemes. The functor
\[
 Rf_*:\mathrm D^+(X)\longrightarrow\mathrm D^+(Y)
\]
does not in general induce, on $\mathrm D^+(\Qcoh(X))$ and $\mathrm D^+(\Qcoh(Y))$, the right derived functor of
\[
 f_*:\Qcoh(X)\longrightarrow\Mod(Y).
\]
More precisely, an injective object of $\Qcoh(X)$ is not necessarily acyclic for $f_*:\Mod(X)\to\Mod(Y)$.

\textbf{0.3.} Let $S$ be a quasi-compact and quasi-separated scheme. An object of $\Qcoh(S)$, even an injective object, is not necessarily acyclic for the coherer (II 3.2). The functor
\[
 \iota:\mathrm D^b(\Qcoh(S))\longrightarrow\mathrm D^b(S)
\]
is not in general fully faithful.

\subsection*{1. The counterexamples}
We first recall the following lemma (SGA 2 II 9).

\textbf{Lemma 1.1.} Let $A$ be a ring, $X=\Spec(A)$, let $\mathbf f=(f_\alpha)$ be a finite system of elements of $A$, and let $Y$ be the closed subset defined by $\mathbf f(x)=0$. For $i>0$, the following conditions are equivalent:
\begin{enumerate}[label=(\roman*)]
\item for every injective $A$-module $I$, one has $H^i_Y(X,\widetilde I)=0$;
\item the pro-object $n\mapsto H_i(\mathbf f^n,A)$ is zero, that is, for every $n$ there exists $n'\geq n$ such that
\[
 H_i(\mathbf f^{n'},A)\longrightarrow H_i(\mathbf f^n,A)
\]
is zero.
\end{enumerate}

\textbf{1.2.} We shall use (1.1) to prove (0.1). Let $k$ be a field,
\[
 B=k[x,y],\qquad \mathbf f=(x,y),\qquad
 M=\bigoplus_{n>0}B/(\mathbf f^n),
\]
and let $A$ be the $B$-algebra $D_B(M)$ (EGA $0_{\mathrm{IV}}$ 18.2.3). One has trivially
\[
 H_2^A(\mathbf f^n,A)\simeq H_2^B(\mathbf f^n,A).
\]
On the other hand,
\[
 H_2^B(\mathbf f^n,A)\simeq H_2^B(\mathbf f^n,B)\oplus H_2^B(\mathbf f^n,M),
\]
and
\[
 H_2^B(\mathbf f^n,B)=0,
\]
since $\mathbf f$ is a regular sequence. Thus
\[
 H_2^B(\mathbf f^n,A)\simeq H_2^B(\mathbf f^n,M).
\]
Finally,
\[
 H_2^B(\mathbf f^n,M)\simeq H_B^0(\mathbf f^n,M),
\]
where $H_B^0(\mathbf f^n,M)$ is canonically identified with the submodule $(\mathbf f^n,M)$ of $M$ annihilated by $\mathbf f^n$. Hence
\begin{equation}\tag{*}
 H_2^A(\mathbf f^n,A)\simeq (\mathbf f^n,M).
\end{equation}
Since the transition morphism
\[
 H_2^A(\mathbf f^{n'},A)\longrightarrow H_2^A(\mathbf f^n,A)
\]
is induced by multiplication by $\mathbf f^{\,n'-n}$, one checks immediately from (*) that the pro-object $n\mapsto H_2^A(\mathbf f^n,A)$ is not zero. Hence, by (1.1), there is an injective $A$-module $I$ such that
\[
 H^2_Y(X,\widetilde I)\ne0,
\]
and therefore $\widetilde I$ is not flasque. Notice also that the projection $\Spec(A)\to\Spec(B)$ is a homeomorphism, so that the underlying space of $X$ is noetherian of dimension $2$. If $U$ denotes the complement of $Y$ in $X$, then by SGA 2 II 4,
\[
 H^2_Y(X,I)\simeq H^1(U,I)\ne0.
\]
But, since $U$ is separated, the functor $\mathrm D^+(\Qcoh(U))\to\mathrm D^+(U)$ is fully faithful by II 3.5, so $\widetilde I|U$ is not injective in $\Qcoh(U)$. This proves assertion (0.1).

Let $Z$ be the quasi-compact and quasi-separated scheme obtained by gluing two copies of $X$ along $U$; let $j_1,j_2$ be the two canonical open immersions of $X$ into $Z$. Thus one has a Cartesian square
\[
\begin{tikzcd}
 U \arrow[r,"i"] \arrow[d,"i"'] & X \arrow[d,"j_1"]\\
 X \arrow[r,"j_2"'] & Z .
\end{tikzcd}
\]
One has
\[
 j_2^*R^1j_{1*}(\widetilde I)\simeq R^1i_*(i^*\widetilde I),
\]
and hence
\[
 H^0(X,R^1j_{1*}(\widetilde I))\simeq H^1(U,\widetilde I)\ne0,
\]
so that $R^1j_{1*}(\widetilde I)\ne0$, which proves (0.2).

On the other hand, by II 3.4 and 3.5 one has
\[
 RQ\,Rj_{1*}(\widetilde I)\simeq j_{1*}(\widetilde I).
\]
It follows from this a spectral sequence whose exact sequence in low degrees gives the isomorphism
\[
 R^1j_{1*}(\widetilde I)\simeq R^2Q(j_{1*}\widetilde I).
\]
Since the sheaf $j_{1*}\widetilde I$ is injective in $\Qcoh(Z)$, the first part of (0.3) is proved.

Finally, let $E\in\Ob\Qcoh(Z)$. We have the duality isomorphism
\[
 R\Hom_{\Mod}(E,Rj_{1*}\widetilde I)
 \simeq R\Hom_{\Mod}(j_1^*E,\widetilde I),
\]
but, since $X$ is affine, II 3.5 gives
\[
 R\Hom_{\Mod}(j_1^*E,\widetilde I)
 \simeq R\Hom_{\Qcoh}(j_1^*E,I)
 \simeq \Hom(j_1^*E,I).
\]
Consequently
\[
 R\Hom_{\Mod}(E,Rj_{1*}\widetilde I)
 \simeq \Hom(j_1^*E,I).
\]
A spectral sequence then gives, in low degree, the isomorphism
\[
 \Ext^2_{\Mod}(E,j_{1*}\widetilde I)
 \simeq \Hom(E,R^1j_{1*}\widetilde I).
\]
Thus the identity morphism of $E=R^1j_{1*}\widetilde I$ supplies a non-zero element of
\[
 \Ext^2_{\Mod}(E,j_{1*}\widetilde I).
\]
It follows that the functor
\[
 \mathrm D^b(\Qcoh(Z))\longrightarrow \mathrm D^b(Z)
\]
is not fully faithful. This completes the proof of (0.3).




% Macros for Exposé II, Appendix II
\providecommand{\Diff}{\operatorname{Diff}}
\providecommand{\Hom}{\operatorname{Hom}}
\providecommand{\im}{\operatorname{Im}}
\providecommand{\ind}{\operatorname{ind}}
\providecommand{\Cone}{\operatorname{Cone}}
\providecommand{\Loclib}{\operatorname{Loclib}}
\providecommand{\ob}{\operatorname{ob}}
\providecommand{\Parf}{\operatorname{Parf}}
\providecommand{\Mod}{\operatorname{Mod}}
\providecommand{\Ob}{\operatorname{Ob}}
\providecommand{\Coker}{\operatorname{Coker}}
\providecommand{\Ker}{\operatorname{Ker}}
\providecommand{\cH}{\mathcal H}
\providecommand{\cO}{\mathcal O}
\providecommand{\cP}{\mathcal P}
\providecommand{\cE}{\mathcal E}
\providecommand{\cF}{\mathcal F}
\providecommand{\cG}{\mathcal G}
\providecommand{\cD}{\mathcal D}
\providecommand{\cT}{\mathcal T}
\providecommand{\cL}{\mathcal L}
\providecommand{\RHom}{\mathrm R\!\operatorname{Hom}}
\providecommand{\RGamma}{\mathrm R\Gamma}
\providecommand{\dlim}{\operatorname*{lim}}
\providecommand{\plim}{\operatorname*{lim}}


\newpage
\section*{Appendix II. Definition of the analytic index of a relative elliptic complex}

In this appendix we propose to show that the notion of a perfect complex leads to a simple and natural definition of the ``analytic index'' of a continuous family of elliptic differential operators, or more generally of a relative elliptic complex. For the initiated reader, let us say that the definition given here has the advantage, compared with the classical definitions (see for example [2] and [17]), of not requiring any deformation of the elliptic complex under consideration.

Sections 1 and 2 consist of reminders on the classical notion of a ``mixed structure''.
\footnote{The informed reader is advised to skip them.}
We give, of course, only the minimum of definitions and results required in order to state and prove, in Section 3, the central theorem we have in view. It says, more precisely, that the direct image, by the projection (assumed proper) onto the base (assumed locally compact), of a relative elliptic complex is a perfect complex. When the base $S$ is compact, the definition of the analytic index as an element of $K(S)$, the Grothendieck group of complex vector bundles of finite rank on $S$, then follows from the globalization sorite of II 2.3.2.

Throughout what follows, if $S$ is a topological space, $\mathcal O_S$ will denote the sheaf of continuous real-valued functions on $S$.

\subsection*{1. The sorite of mixed structures}

In this section a topological space $S$ is fixed.

\subsubsection*{1.1. Definition of mixed varieties}

Let $V$ be an open subset of $\mathbb R^n$. For $r\in\mathbb N$ one writes $C^r(V)$ for the Fréchet space of real-valued functions of class $C^r$ on $V$. Put
\[
 C^\infty(V)=\plim_r C^r(V).
\]
Let $U$ be an open subset of $S$. For $r\in\mathbb N\cup\{\infty\}$, denote by ${}^r\mathcal O_{U\times V}$ the sheaf of rings on $U\times V$ defined by
\[
 \Gamma(U'\times V',{}^r\mathcal O_{U\times V})
 =C(U',C^r(V')),
\]
where $U'$ (respectively $V'$) is an open subset of $U$ (respectively of $V$), and $C(U',C^r(V'))$ denotes the space of continuous maps from $U'$ into $C^r(V')$. There is a canonical injection
\[
 \operatorname{pr}_1^{-1}(\mathcal O_U)\longrightarrow {}^r\mathcal O_{U\times V}
\]
which makes $(U\times V,{}^r\mathcal O_{U\times V})$ into a ringed space over $U$.

A ringed space $X$ over $S$ is said to be a \emph{mixed variety of class $C^r$} ($r\in\mathbb N\cup\{\infty\}$) if $X$ is locally isomorphic, as a ringed space over $S$, to a space of the type just described. If $S$ is reduced to one point, a mixed variety of class $C^r$ over $S$ is just a variety of class $C^r$ in the usual sense.

Let $f:X\to S$ be a mixed variety of class $C^r$. For $s\in S$, the fiber $X_s=f^{-1}(s)$, endowed with the sheaf $i_s^{-1}(\mathcal O_X)/I_s$, where $i_s:X_s\to X$ is the injection and $I_s$ is the ideal of sections which vanish on $X_s$, is a variety of class $C^r$.

\subsubsection*{1.2. Relative differential operators}

It is easy, on the model of algebraic geometry (EGA IV 16, SGA 3 VII), to develop an infinitesimal study of mixed varieties. We shall only review rapidly the notions we shall need.

\textbf{1.2.1.} Let $f:X\to S$ be a mixed variety of class $C^\infty$. For each $n\in\mathbb N$ one defines the sheaf $P^n_{X/S}$ of principal parts of order $n$ on $X$. In two different ways, corresponding to the two projections $\operatorname{pr}_1$ and $\operatorname{pr}_2$ of $X\times_SX$ onto $X$, it is an $\mathcal O_X$-algebra augmented toward $\mathcal O_X$, and is locally free of finite type as an $\mathcal O_X$-module. When the $\mathcal O_X$-module structure on $P^n_{X/S}$ is not specified, it will be understood to be the one coming from $\operatorname{pr}_1$. One has the canonical identifications
\[
 P^0_{X/S}=\mathcal O_X,
 \qquad
 P^1_{X/S}/\mathcal O_X=\Omega^1_{X/S}.
\]
The sheaf $\Omega^1_{X/S}$ is the sheaf of relative differentials of $X$ over $S$, or relative cotangent sheaf. Its dual is denoted $T_{X/S}$ and called the relative tangent sheaf. The sheaf $T_{X/S}$, respectively $\Omega^1_{X/S}$, is the sheaf of sections of a vector bundle of finite rank on $X$, called the relative tangent, respectively cotangent, bundle; these bundles are mixed varieties over $X$.

The $f^{-1}(\mathcal O_S)$-linear homomorphism, corresponding to the $\mathcal O_X$-algebra structure defined by $\operatorname{pr}_2$,
\[
 d^n_{X/S}:\mathcal O_X\longrightarrow P^n_{X/S},
\]
which associates to a section of $\mathcal O_X$ its principal part of order $n$, is called the universal differential operator of order $n$ on $X$ relative to $S$.

\textbf{1.2.2.} More generally, if $E$ is a locally free $\mathcal O_X$-module of finite type, one may consider the sheaf of principal parts of order $n$ of $E$,
\[
 P^n_{X/S}(E)=P^n_{X/S}\otimes_{\mathcal O_X}E,
\]
where the tensor product is defined by means of the second $\mathcal O_X$-module structure on $P^n_{X/S}$. One has a universal differential operator
\[
 d^n_{X/S}(E):E\longrightarrow P^n_{X/S}(E).
\]
Let $E,F$ be locally free $\mathcal O_X$-modules of finite type. By definition, a differential operator from $E$ to $F$ of order $\le n$ is an $f^{-1}(\mathcal O_S)$-linear homomorphism from $E$ to $F$ which factors through $d^n_{X/S}(E)$ as an $\mathcal O_X$-linear homomorphism $P^n_{X/S}(E)\to F$; the factorization is then unique. In other words, $d^n_{X/S}(E)$ gives a bijection
\begin{equation}\tag{1.2.2.1}
 \Hom_{\mathcal O_X}(P^n_{X/S}(E),F)\xrightarrow{\sim}\Diff^n(E,F),
\end{equation}
where $\Diff^n(E,F)$ denotes the set of differential operators from $E$ to $F$ of order $\le n$.

\textbf{1.2.3.} We shall also have to consider differential operators with complex coefficients. Let $\mathcal O_X\otimes\mathbb C$ (respectively $\mathcal O_S\otimes\mathbb C$) be the complexification of $\mathcal O_X$ (respectively $\mathcal O_S$). Let $E,F$ be locally free $(\mathcal O_X\otimes\mathbb C)$-modules of finite type. By a differential operator from $E$ to $F$ of order $\le n$ one means an $f^{-1}(\mathcal O_S\otimes\mathbb C)$-linear homomorphism from $E$ to $F$ which factors through $d^n_{X/S}(E)$ as an $(\mathcal O_X\otimes\mathbb C)$-linear homomorphism $P^n_{X/S}(E)\to F$. If $\Diff^n_{\mathbb C}(E,F)$ denotes the set of such operators, $d^n_{X/S}(E)$ therefore defines a bijection
\begin{equation}\tag{1.2.3.1}
 \Hom_{\mathcal O_X\otimes\mathbb C}(P^n_{X/S}(E),F)
 \xrightarrow{\sim}\Diff^n_{\mathbb C}(E,F).
\end{equation}

\textbf{1.2.4.} For $n\ge1$ there is a canonical exact sequence
\[
 0\longrightarrow S^n(\Omega^1_{X/S})\longrightarrow P^n_{X/S}
 \longrightarrow P^{n-1}_{X/S}\longrightarrow 0.
\]
Hence, tensoring over $\mathcal O_X$ with a locally free $\mathcal O_X$-module, respectively $(\mathcal O_X\otimes\mathbb C)$-module, $E$ of finite type, one obtains an exact sequence
\begin{equation}\tag{1.2.4.1}
0\longrightarrow S^n(\Omega^1_{X/S})\otimes E\longrightarrow P^n_{X/S}(E)
\longrightarrow P^{n-1}_{X/S}(E)\longrightarrow0.
\end{equation}
Here $S^n(\Omega^1_{X/S})$ denotes the $n$-th symmetric power.

Let $F$ be a locally free $\mathcal O_X$-module, respectively $(\mathcal O_X\otimes\mathbb C)$-module, of finite type. Applying $\Hom(\ ,F)$, respectively $\Hom_{\mathcal O_X\otimes\mathbb C}(\ ,F)$, to (1.2.4.1), and using (1.2.2.1) and (1.2.3.1), one obtains an exact sequence
\begin{equation}\tag{1.2.4.2}
0\longrightarrow \Diff^{n-1}(E,F)\longrightarrow \Diff^n(E,F)
\xrightarrow{\sigma_n}
\Hom_{\mathcal O_X}\bigl(S^n(\Omega^1_{X/S})\otimes E,F\bigr),
\end{equation}
and analogously in the complex case. If $X$ is paracompact and $\mathcal O_X$ is soft, the sequence may be completed on the right by adding a zero. The homomorphism $\sigma_n$ is called the \emph{symbol}. It associates to a differential operator from $E$ to $F$ of order $\le n$ a homogeneous polynomial map of degree $n$ from $T^*_{X/S}$ into $\Hom(E,F)$, or equivalently a homomorphism $p^*E\to p^*F$, where $p$ is the projection of the relative cotangent bundle onto $X$, satisfying the usual homogeneity condition.

Let $E,F,G$ be locally free $\mathcal O_X$-modules of finite type, $u\in\Diff^n(E,F)$ and $v\in\Diff^m(F,G)$. Then $vu\in\Diff^{m+n}(E,G)$ and one checks at once the relation
\begin{equation}\tag{1.2.4.3}
 \sigma_{m+n}(vu)=\sigma_m(v)\sigma_n(u).
\end{equation}
The same relation holds for differential operators with complex coefficients. In particular, if the composite of two differential operators is zero, then the composite of their symbols is zero.

\textbf{1.2.5.} Let $E^\bullet$ be a bounded complex of differential operators with complex coefficients on $X$ relative to $S$, and let $d^i:E^i\to E^{i+1}$ be of order $n_i$. By (1.2.4.3),
\[
 \sigma_{n_{i+1}}(d^{i+1})\sigma_{n_i}(d^i)=0.
\]
The complex $E^\bullet$ is called an \emph{elliptic complex of order $(n_i)$} if the symbol complex on the relative cotangent bundle, pulled back to that bundle, is acyclic outside the zero section.

When $S$ is reduced to one point, this notion of elliptic complex on $X$ coincides with the usual one, studied in detail in [4] and [16].

\textbf{1.2.6.} Let $s\in S$, and let $i_s$ be the inclusion of the fiber $X_s$ in $X$, considered as a monomorphism of ringed spaces (1.1). For $n\in\mathbb N$ one has a canonical isomorphism
\[
 i_s^*P^n_{X/S}\simeq P^n_{X_s},
\]
and hence a restriction homomorphism
\[
 i_s^*:\Diff^n(E,F)\longrightarrow \Diff^n(i_s^*E,i_s^*F),
\]
for $E,F$ locally free $\mathcal O_X$-modules of finite type, and similarly for operators with complex coefficients. If $E^\bullet$ is a bounded complex of differential operators with complex coefficients on $X$, then $E^\bullet$ is elliptic if and only if, for every $s\in S$, $i_s^*(E^\bullet)$ is an elliptic complex. The verification is left to the reader.

\subsection*{2. Rigidity lemmas}

\textbf{2.0.} It is ``well known'' that compact differentiable varieties are rigid, that is, they admit no continuous deformation other than the trivial one. Since this result will be used essentially in Section 3, and references seem rather scarce, it is perhaps useful to give precise statements and a few proofs.

A topological space $S$ is fixed. The mixed varieties under consideration are assumed to be of class $C^r$, with $r\gg1$; possibly $r=\infty$.

\textbf{2.1. Notation.} Let $X,Y$ be mixed $S$-varieties. We denote by $\Hom(X,Y)$ the set of $S$-morphisms from $X$ to $Y$, and by $\mathcal Hom(X,Y)$ the sheaf on $X$ of germs of $S$-morphisms from $X$ to $Y$, defined by
\[
 \mathcal Hom(X,Y)(U)=\Hom(U,Y)
\]
for $U$ open in $X$.

Let $s\in S$, and write $i_s$ for the inclusion of $X_s$ in $X$. There is a canonical induction homomorphism
\begin{equation}\tag{2.1.1}
 i_s^{-1}\mathcal Hom(X,Y)\longrightarrow \mathcal Hom(X_s,Y_s)
\end{equation}
which is an epimorphism, as is immediate. In the case $Y=S\times\mathbb R$, the sheaf $\mathcal Hom(X,Y)$ is canonically identified with $\mathcal O_X$, and the epimorphism (2.1.1) with the epimorphism
\[
 i_s^{-1}(\mathcal O_X)\longrightarrow i_s^*(\mathcal O_X)
\]
considered in (1.1).

\textbf{Proposition 2.2.} Let $X,Y$ be mixed $S$-varieties, and let $s\in S$. Suppose that the fiber $X_s$ is paracompact and that $s$ has a fundamental system of paracompact neighborhoods. Let $A$ be a closed subset of $X_s$, and let
\[
 f\in H^0(A,\mathcal Hom(X_s,Y_s)).
\]
The subsheaf $\mathcal F$ of $i_s^{-1}\mathcal Hom(X,Y)$ consisting of the germs which extend $f$, that is, whose image by (2.1.1) coincides with $f$ over $A$, is soft.

\emph{Proof.} Let $X'$ be an open neighborhood in $X$ which is isomorphic, as a mixed variety, to $S'\times U$, where $S'$ is a paracompact neighborhood of $s$ and $U$ is a relatively compact open subset of $\mathbb R^n$, and such that the image by $f$ of $X'_s\cap A$ is contained in an open subset of $Y$ isomorphic to $S'\times V$, where $V$ is a convex open subset of $\mathbb R^q$ containing the origin. By T.F. II 3.4.1 it is enough to show that every section of $\mathcal F$ over a closed subset of $X_s$ contained in $X'_s$ extends to $X'_s$. Since $X'_s=\{s\}\times U$ is relatively compact, a closed subset of $X_s$ contained in $X'_s$ is compact. Identifying $X'$ with $S'\times U$ and $Y'$ with $S'\times V$, putting $A'=A\cap X'_s$, and using T.F. II 3.3.1, the assertion is reduced to the following elementary lemma.\footnote{Communicated to the writer by A. Douady. A closed subset of a paracompact space has a fundamental system of paracompact neighborhoods.}

\textbf{Lemma 2.2.1.} Let $U$ be an open subset of $\mathbb R^n$, and let $V$ be a convex open subset of $\mathbb R^q$ containing $0$. Let $A_1,A_2$ be closed subsets of $\{s\}\times U$. Let $U_1$, respectively $U_2$, be an open subset of $S\times U$, respectively of $\{s\}\times U$, containing $A_1$, respectively $A_2$. Let $f_1:U_1\to V$ be a morphism and let $f_2:U_2\to S\times V$ be an $S$-morphism which, on $U_1\cap(\{s\}\times U)$, defines the same germ as $f_1$ near $A_1$. Then there exists an open subset $U'$ of $S\times U$ containing $\{s\}\times U$ and an $S$-morphism
\[
 F:U'\longrightarrow S\times V
\]
such that the germ of $F$ along $\{s\}\times U$ extends the germ defined by $f_1$, and the morphism induced by $F$ on $\{s\}\times U$ defines the same germ as $f_2$ near $A_2$.

\emph{Proof.} Replacing $f_1$ by $af_1$, where $a$ is a $C^\infty$ function on $U$ equal to $1$ near $A_1$ and with support in $U_1$, one may suppose $U_1=U$. Let $b$ be a $C^\infty$ function on $U$, equal to $1$ on $A_2$ and with support in $U_2\cap(\{s\}\times U)$. Then
\[
 U' = U_1\cup\bigl(S\times (U-\operatorname{Supp}b)\bigr)
\]
is an open subset of $S\times U$ containing $\{s\}\times U$. The morphism $F:U'\to S\times V$ defined by
\[
 F(t,x)=\bigl(t,\,b(x)\operatorname{pr}_2 f_2(x)+(1-b(x))f_1(t,x)\bigr)
\]
visibly satisfies the required conditions.

\textbf{Corollary 2.3.} Under the hypotheses of (2.2), the sheaves $i_s^{-1}\mathcal Hom(X,Y)$ and $\mathcal Hom(X_s,Y_s)$ are soft. In particular, the sheaves $i_s^{-1}(\mathcal O_X)$ and $\mathcal O_{X_s}$ are soft. For every closed subset $A$ of $X_s$, the map
\[
 H^0(X_s,i_s^{-1}\mathcal Hom(X,Y))\longrightarrow H^0(A,\mathcal Hom(X_s,Y_s))
\]
deduced from (2.1.1) is surjective.

\emph{Proof.} The softness of $i_s^{-1}\mathcal Hom(X,Y)$ follows from (2.2) applied with $A=\varnothing$. Taking $S=\{s\}$, one obtains the softness of $\mathcal Hom(X_s,Y_s)$. The last assertion follows from (2.2).\footnote{For $\mathcal O_{X_s}$ this is very well known. The writer does not know whether $\mathcal O_X$ is soft when the hypotheses of (2.2) are satisfied for every $s\in S$ and $X$ is paracompact. This would trivially imply, under those hypotheses, the softness of $i_s^{-1}(\mathcal O_X)$ as well as (2.5) below.}

\textbf{Corollary 2.4.} Under the hypotheses of (2.2), let $E$ be an $\mathcal O_X$-module. For every closed subset $A$ of $X_s$, the canonical homomorphism
\[
 H^0(X_s,i_s^{-1}E)\longrightarrow H^0(A,i_s^*E)
\]
deduced from the epimorphism
\[
 i_s^{-1}E\longrightarrow i_s^*E=i_s^{-1}E\otimes_{i_s^{-1}\mathcal O_X}\mathcal O_{X_s}
\]
is an epimorphism.

\emph{Proof.} Since $i_s^{-1}\mathcal O_X$ is soft by (2.3), the functor $H^0(A,\ )$ is exact on the category of $i_s^{-1}\mathcal O_X$-modules (T.F. II 3.7.1, 3.4.2, 4.4.3). Hence
\[
 H^0(A,i_s^{-1}E)\longrightarrow H^0(A,i_s^*E)
\]
is an epimorphism. Since $i_s^{-1}E$ is soft, the homomorphism
\[
 H^0(X_s,i_s^{-1}E)\longrightarrow H^0(A,i_s^{-1}E)
\]
is an epimorphism, and this proves the assertion.

\textbf{Corollary 2.5.} Let $p:X\to S$ be a mixed variety and let $s\in S$. Suppose that $s$ has a fundamental system of paracompact neighborhoods $S'$ such that the $p^{-1}(S')$ form a fundamental system of paracompact neighborhoods of $X_s$. This condition is satisfied, for example, if $S$ is locally compact and $p$ is a proper map of topological spaces. Then, for every $\mathcal O_X$-module $E$,
\[
 (R^ip_*E)_s=0\qquad(i>0).
\]

\emph{Proof.} By T.F. II 4.17.1,
\[
 (R^ip_*E)_s=H^i(X_s,i_s^{-1}E).
\]
But $i_s^{-1}\mathcal O_X$ is soft by (2.3), hence $H^i(X_s,i_s^{-1}E)=0$ for $i>0$.

\textbf{Corollary 2.6.} Let $X,Y$ be mixed $S$-varieties, and let $s\in S$. Suppose that $X$ satisfies the hypotheses of (2.5). For every morphism $f:X_s\to Y_s$, there exists an open neighborhood $S'$ of $s$ and a morphism
\[
 F:X|S'\longrightarrow Y|S'
\]
such that $F_s=f$.

\emph{Proof.} This is an immediate consequence of (2.3), with $A=X_s$, and of the fact that the hypotheses on $X$ give
\[
 H^0(X_s,i_s^{-1}\mathcal Hom(X,Y))
 =\dlim_{S'\ni s}\Hom(X|S',Y|S')
\]
for open $S'$ containing $s$ (T.F. II 8.5.5).

\textbf{Proposition 2.7.} Let $X,Y$ be mixed $S$-varieties, let $f:X\to Y$ be a morphism, let $s\in S$, and let $x\in X_s$. If the tangent linear map of
\[
 f_s:X_s\longrightarrow Y_s
\]
is invertible at $x$, then $f$ is a local isomorphism at $x$; that is, there is an open neighborhood $U$ of $x$ in $X$ and an open subset $V$ of $Y$ such that $f$ induces an isomorphism $U\simeq V$.

\emph{Proof.} This is an easy consequence of the implicit function theorem. The details are left to the reader.

\textbf{Corollary 2.8.} Let $X,Y$ be mixed $S$-varieties and let $s\in S$. Suppose that $s$ has a fundamental system of neighborhoods $S'$ such that the $X|S'$, respectively $Y|S'$, form a fundamental system of neighborhoods of $X_s$, respectively $Y_s$; this is the case, for example, if $S$ is locally compact and $X$, respectively $Y$, is proper over $S$. Let $f:X\to Y$ be an $S$-morphism such that
\[
 f_s:X_s\longrightarrow Y_s
\]
is an isomorphism. Then there exists a neighborhood $S'$ of $s$ such that
\[
 f|X|S':X|S'\longrightarrow Y|S'
\]
is an isomorphism.

\emph{Proof.} By (2.7), there exists an open neighborhood $U$ of $X_s$ in $X$ such that $f|U$ is étale, that is, a local isomorphism. We may suppose that $U=X|S'$ for a neighborhood $S'$ of $s$; after shrinking $S$, we may therefore assume that $f$ is étale. Then the fiber product $Z=X\times_YX$ is naturally a mixed $S$-variety, étale over $X$ for each projection, and the diagonal morphism $\Delta:X\to Z$ is an open immersion. Moreover the sets $Z|S'$ form a fundamental system of neighborhoods of $Z_s$. Since
\[
 \Delta_s:X_s\longrightarrow Z_s
\]
is an isomorphism, it follows, after shrinking $S$, that $\Delta$ is an isomorphism. Thus $f$ is an open immersion, and after shrinking $S$ again it is an isomorphism.

\textbf{Corollary 2.9.} Let $X$ be a mixed $S$-variety and let $s\in S$. Suppose that the hypothesis of (2.5) is satisfied. Then there exists a neighborhood $S'$ of $s$ in $S$ such that $X|S'$ is $S'$-isomorphic to $S'\times X_s$.

\emph{Proof.} Apply (2.6) and (2.8).

\textbf{Corollary 2.10.} Let $s\in S$ and let $V$ be a paracompact variety. Put $X=S\times V$, and let $i_s$ be the injection of $V$ into $X$ defined by $s$. Suppose that the hypothesis of (2.5) is satisfied. Let $E$ be a locally free $\mathcal O_X$-module, respectively $\mathcal O_X\otimes\mathbb C$-module, of finite type. Then there exists a neighborhood $S'$ of $s$ in $S$ such that $E|X|S'$ is isomorphic to $\operatorname{pr}_2^*(i_s^*E)$, where $\operatorname{pr}_2:S'\times V\to V$ is the second projection.

\emph{Proof.} Put $F=\operatorname{pr}_2^*(i_s^*E)$. By (2.4), applied to $\mathcal Hom(E,F)$, the induction homomorphism
\[
 \Hom(i_s^{-1}E,i_s^{-1}F)\longrightarrow \Hom(i_s^*E,i_s^*F)
\]
is an epimorphism. On the other hand, by T.F. II 3.3.1,
\[
 \Hom(i_s^{-1}E,i_s^{-1}F)
 =\dlim_{S'\ni s}\Hom(E|X|S',F|X|S') .
\]
Thus the canonical isomorphism $i_s^*E\simeq i_s^*F$ is induced, after shrinking $S$, by a homomorphism $f:E\to F$. Since $f_s$ is an isomorphism, $f$ is an isomorphism near $X_s$ by Nakayama, hence over $X|S'$ for a suitable neighborhood $S'$ of $s$.

\subsection*{3. The finiteness theorem}

\subsubsection*{3.0. Conventions and notation}

\textbf{3.0.1.} Let $X$ and $Y$ be topological spaces. We write $C(X,Y)$ for the set of continuous maps from $X$ to $Y$, and $\mathcal C(X,Y)$ for the sheaf on $X$ defined by $U\mapsto C(U,Y)$.

\textbf{3.0.2.} Let $E,F$ be Banach spaces, real or complex. We write $L(E,F)$ for the Banach space of continuous linear maps from $E$ to $F$.

\textbf{3.0.3.} The mixed varieties occurring in this section are assumed to be of class $C^\infty$. If $X$ is a mixed $S$-variety, we shall often identify a locally free $\mathcal O_X$-module, respectively $\mathcal O_X\otimes\mathbb C$-module, of finite type with the corresponding real, respectively complex, vector bundle.

\textbf{3.0.4.} Let $X$ be a compact ordinary variety, and let $E$ be a real or complex vector bundle of finite rank on $X$. Write $C^\infty(X,E)$ for the Fréchet space of $C^\infty$ sections of $E$ on $X$. For $r\in\mathbb Z$, write $C_r(X,E)$ for the Hilbert space of sections of $E$ which are $r$ times differentiable in $L^2$ ([16] IX). For $s\le r$, $C_r(X,E)$ maps to $C_s(X,E)$ by a continuous linear injective map with dense image, and
\[
 C^\infty(X,E)=\plim_r C_r(X,E).
\]

\textbf{3.0.5.} Let $S$ be a topological space. We shall need the notion of a Hilbert bundle over $S$, for which we refer the reader to [12] or [14]. There is an additive functor from the category of real, respectively complex, Hilbert bundles over $S$ to the category of $\mathcal O_S$-, respectively $\mathcal O_S\otimes\mathbb C$-, modules: it sends a Hilbert bundle $E$ to the sheaf of continuous sections of $E$.

\subsubsection*{3.1. Auxiliary technical results}

\textbf{Lemma 3.1.1.} Let $X,Y$ be topological spaces, with $X$ locally compact. Let $\mathcal C'(X,Y)$ be the presheaf, on the basis of relatively compact open subsets of $X$, defined by
\[
 \mathcal C'(X,Y)(U)=C(\overline U,Y).
\]
The canonical restriction morphism
\[
 \mathcal C'(X,Y)\longrightarrow \mathcal C(X,Y)
\]
induces an isomorphism on the associated sheaves.

\emph{Proof.} Immediate.

\textbf{Lemma 3.1.2.} Let $E,F$ be locally convex topological vector spaces over $\mathbb C$, and let $u:E\to F$ be a continuous linear map, respectively a continuous linear map with dense image. Let $X$ be compact. Endow $C(X,E)$ and $C(X,F)$ with the topology of uniform convergence. Then the induced linear map
\[
 u_X:C(X,E)\longrightarrow C(X,F)
\]
is continuous, respectively continuous with dense image.

\emph{Proof.} It is clear that $u_X$ is continuous. On the other hand, by Bourbaki, \emph{Intégration}, chapter II, §2, lemmas 1 and 2, the subspace $C(X,\mathbb C)\otimes E$ of $C(X,E)$, formed by the continuous maps from $X$ to $E$ whose values lie in a finite-dimensional vector subspace, is dense; and similarly $C(X,\mathbb C)\otimes F$ is dense in $C(X,F)$. If $u$ has dense image, then $u_X(C(X,E))$ is dense in $C(X,\mathbb C)\otimes F$, and taking closures gives the assertion.

\textbf{3.1.3.} Let $V$ be a compact variety, let $E$ and $F$ be complex vector bundles of finite rank on $V$, and let $n\in\mathbb N$. The space $\Diff^n_{\mathbb C}(E,F)$, identified by (1.2.3.1) with
\[
 \Hom(P^n(E),F)\simeq C^\infty(V,\Hom(P^n(E),F)),
\]
is naturally a Fréchet space. Let $u\in\Diff^n_{\mathbb C}(E,F)$. One checks easily that the linear map $C^\infty(V,E)\to C^\infty(V,F)$ defined by $u$ is continuous. If $L(C^\infty(V,E),C^\infty(V,F))$ denotes the space of continuous linear maps, one therefore has a linear map
\[
 \delta:\Diff^n_{\mathbb C}(E,F)
 \longrightarrow L(C^\infty(V,E),C^\infty(V,F)).
\]
Moreover, for $r\in\mathbb Z$, $u$ acts by extension of scalars on the $C_r$ sections of $E$ and defines a linear map
\[
 \delta_r:\Diff^n_{\mathbb C}(E,F)
 \longrightarrow L(C_r(V,E),C_{r-n}(V,F)).
\]
The maps $\delta_r$ are continuous, and for $u\in\Diff^n_{\mathbb C}(E,F)$, $\delta(u)$ is the inverse-limit element of the projective system of the $\delta_r(u)$. These assertions follow easily from [16] IX: one reduces immediately to $n=0$; the statement on the $C_r$ follows directly from the definitions, while the statement for $C^\infty$ follows, after localizing by a partition of unity, from the formula for successive derivatives of a product. The analogous real statements are also valid.

\textbf{Lemma 3.1.4.} Let $S$ be a topological space, let $V$ be a compact variety, let $E,F$ be complex vector bundles of finite rank on $V$, and let $n\in\mathbb N$. Put $X=S\times V$. For
\[
 u\in\Diff^n_{\mathbb C}(\operatorname{pr}_2^*E,\operatorname{pr}_2^*F),
\]
the map $s\mapsto i_s^*u$ from $S$ into $\Diff^n_{\mathbb C}(E,F)$ is continuous. The linear map
\[
 \Diff^n_{\mathbb C}(\operatorname{pr}_2^*E,\operatorname{pr}_2^*F)
 \longrightarrow C(S,\Diff^n_{\mathbb C}(E,F))
\]
so defined is an isomorphism.

\emph{Proof.} By (1.2.3.1) one may restrict to the case $n=0$. Putting $G=\Hom(E,F)$, one is reduced to showing that for $u\in\Gamma(X,\operatorname{pr}_2^*G)$, the map $s\mapsto u_s=i_s^*u$ from $S$ into $C^\infty(V,G)$ is continuous and that the resulting map
\[
 \Gamma(X,\operatorname{pr}_2^*G)\longrightarrow C(S,C^\infty(V,G))
\]
is an isomorphism. This follows immediately from the definition of mixed varieties.

\textbf{3.1.5.} Under the hypotheses of (3.1.4), let
\[
 u\in\Diff^n_{\mathbb C}(\operatorname{pr}_2^*E,\operatorname{pr}_2^*F).
\]
Identifying $u$, by (3.1.4), with a continuous map from $S$ into $\Diff^n(E,F)$, one sees that, for $U$ open in $S$, the morphism
\[
 \operatorname{pr}_{1*}(u)(U):C^\infty(U,E)\longrightarrow C^\infty(U,F)
\]
is canonically identified with the linear map
\[
 C(U,C^\infty(V,E))\longrightarrow C(U,C^\infty(V,F))
\]
deduced from
\[
 \delta u:U\longrightarrow L(C^\infty(V,E),C^\infty(V,F)).
\]
Suppose $S$ locally compact. By (3.1.1), the morphism of sheaves
\[
 \operatorname{pr}_{1*}(u):
 \operatorname{pr}_{1*}\operatorname{pr}_2^*E
 \longrightarrow
 \operatorname{pr}_{1*}\operatorname{pr}_2^*F
\]
is the morphism associated to the morphism of presheaves
\[
 \mathcal C'(S,C^\infty(V,E))\longrightarrow \mathcal C'(S,C^\infty(V,F))
\]
defined by $\delta u$. On the other hand, for $r\in\mathbb Z$,
\[
 \delta_r u:S\longrightarrow L(C_r(V,E),C_{r-n}(V,F))
\]
defines a morphism of presheaves
\[
 u_r:\mathcal C'(S,C_r(V,E))\longrightarrow \mathcal C'(S,C_{r-n}(V,F)).
\]
By (3.1.3), the $u_r$ form a projective system whose limit is the previous morphism. It will be convenient to interpret $\delta_r u$ as a morphism of trivial Hilbert bundles
\[
 S\times C_r(V,E)\longrightarrow S\times C_{r-n}(V,F),
\]
inducing on sheaves of sections the sheaf morphism associated to $u_r$.

\subsubsection*{3.2. Definition of the index}

\textbf{Theorem 3.2.1.} Let $S$ be a locally compact space, let $p:X\to S$ be a proper mixed variety, and let $E^\bullet$ be an elliptic complex of order $(n_i)$ of complex vector bundles on $X$ (1.2.5). Then the complex $Rp_*(E^\bullet)$, which by (2.5) is isomorphic to $p_*(E^\bullet)$ in the derived category of $\mathcal O_S\otimes\mathbb C$-modules, is perfect.

\emph{Proof.} The question being local on $S$, one may suppose by (2.9) that $X=S\times V$, where $V$ is a compact variety, and that $p=\operatorname{pr}_1$. Moreover, by (2.10), one may suppose that for all $i$,
\[
 E^i=\operatorname{pr}_2^*(F^i),
\]
where $F^i$ is a complex vector bundle of finite rank on $V$. By (3.1.5), $p_*(E^\bullet)$ is associated to the complex of presheaves
\[
 \mathcal C'(S,C^\infty(V,F^\bullet)),
\]
where the differentials are induced by the given differential operators. Put, for each $i\in\mathbb Z$,
\[
 m_i=\sum_{j<i}n_j,
\]
with the convention $n_i=0$ if $E^i=0$. For each $r\in\mathbb Z$ one has, again by (3.1.5), a complex of presheaves
\[
 \mathcal C'(S,C_{r-m_i}(V,F^i)),
\]
which we denote by $\mathcal C'(S,C_r(V,F^\bullet))$. These complexes form a projective system, and
\begin{equation}\tag{3.2.1.1}
 \mathcal C'(S,C^\infty(V,F^\bullet))
 =\plim_r \mathcal C'(S,C_r(V,F^\bullet)).
\end{equation}
Finally, denote by $F_r^\bullet$ the complex of trivial Hilbert bundles over $S$ defined by the corresponding $\delta_r(d^i)$, and by $\mathcal F_r^\bullet$ the complex of sheaves of sections; by (3.1.5), it is just the complex of sheaves associated to $\mathcal C'(S,C_r(V,F^\bullet))$. To prove the theorem it is enough to prove the following assertions:

(a) For $r\in\mathbb Z$, the canonical morphism
\[
 \mathcal C'(S,C^\infty(V,F^\bullet))\longrightarrow \mathcal C'(S,C_r(V,F^\bullet))
\]
is a quasi-isomorphism of complexes of presheaves, that is, it induces isomorphisms on the presheaves of cohomology.

(b) For $r\in\mathbb Z$, the complex $\mathcal F_r^\bullet$ is perfect.

Proof of (b). The ellipticity of $E^\bullet$ implies ([16] XI, theorem 6) that, for every $s\in S$, the differential of $(F_r^\bullet)_s$ has closed image and the spaces $H^i((F_r^\bullet)_s)$ are finite-dimensional vector spaces. In the terminology of [11], $F_r^\bullet$ is a direct quasi-acyclic complex of Hilbert bundles. Hence, by [12] 2.3, $F_r^\bullet$ is locally homotopically equivalent to a bounded complex of vector bundles of finite rank. Since the sheaf-of-sections functor is additive (3.0.5), $\mathcal F_r^\bullet$ is locally homotopically equivalent to a bounded complex of locally free sheaves of finite type; hence it is perfect.

For the proof of (a) we shall need the following intermediate assertion:

(a') For $p\gg q$, the canonical morphism
\[
 F_p^\bullet\longrightarrow F_q^\bullet
\]
is a quasi-isomorphism.

Proof of (a'). Let $M_{pq}^\bullet$ be the complex of Hilbert bundles which is the cone of the morphism $F_p^\bullet\to F_q^\bullet$. The ellipticity of $E^\bullet$ implies that, for every $s\in S$, the morphism $(F_p^\bullet)_s\to(F_q^\bullet)_s$ is a quasi-isomorphism ([16] XI, corollary 2 and theorem 6). In other words, $(M_{pq}^\bullet)_s$ is a direct acyclic complex. By [11] 2.3, $M_{pq}^\bullet$ is locally homotopically trivial; hence its sheaf complex, the cone of $\mathcal F_p^\bullet\to\mathcal F_q^\bullet$, is acyclic.

Proof of (a). It is enough to show that, for every relatively compact open subset $U$ of $S$, the morphism
\[
 C^\infty(U,C^\infty(V,F^\bullet))\longrightarrow C^\infty(U,C_r(V,F^\bullet))
\]
is a quasi-isomorphism. Replacing $S$ by $\overline U$ and $E^\bullet$ by the induced complex over $\overline U$, one is reduced to proving that
\[
 C(S,C^\infty(V,F^\bullet))\longrightarrow C(S,C_r(V,F^\bullet))
\]
is a quasi-isomorphism, with $S$ compact. By (3.2.1.1),
\[
 C(S,C^\infty(V,F^\bullet))=\plim_i C(S,C_i(V,F^\bullet)).
\]
For each $i$, $C(S,C_i(V,F^j))$ is a Banach space for the topology of uniform convergence, and the transition morphisms
\[
 C(S,C_i(V,F^j))\longrightarrow C(S,C_{i-1}(V,F^j))
\]
are continuous linear maps with dense image by (3.0.4) and (3.1.2). On the other hand, applying to the quasi-isomorphisms (a') the functor ``sections'', which is exact since $\mathcal O_S$ is soft, one sees that the transition morphisms are quasi-isomorphisms. Assertion (a) follows from EGA $0_{\mathrm I}$ 13.2.3, topological variant. This completes the proof of (3.2.1).

\textbf{3.2.2.} If $S$ is a topological space, we denote by $\Parf(S)$ the triangulated category of perfect complexes of $(\mathcal O_S\otimes\mathbb C)$-modules, and by $K^0(S)$ the corresponding Grothendieck group.\footnote{See I 6.3. For a general review of $K$-functors, see Exposé IV.}

Let $S,X,E^\bullet$ be as in (3.2.1). The \emph{analytic index} of $E^\bullet$ is the element
\[
 \ind_{\mathrm{an}}(E^\bullet)=[Rp_*(E^\bullet)]\in K^0(S).
\]

If $S$ is compact, $\ind_{\mathrm{an}}(E^\bullet)$ canonically defines an element in the usual group $K(S)$ of topologists. Indeed, in general, for every topological space $S$, let $\Loclib(S)$ be the additive category of locally free $(\mathcal O_S\otimes\mathbb C)$-modules of finite type, and let $K(S)$ be the Grothendieck group of the triangulated category
\[
 K^b(\Loclib(S))/K^b_{\mathrm{ac}}(\Loclib(S)).
\]
This category is in fact identified with $K^b(\Loclib(S))$ when $S$ is paracompact (II 2.2.2(a)), so in this case, by I 6.4, $K(S)$ is canonically isomorphic to the Grothendieck group of the additive category $\Loclib(S)$. There is a natural homomorphism
\[
 K(S)\longrightarrow K^0(S).
\]
By II 2.3.2 this is an isomorphism for compact $S$, proving the assertion. More generally, $\ind_{\mathrm{an}}(E^\bullet)$ defines an element in the completed group
\[
 \widehat K(S)=\plim_A K(A),
\]
the projective limit being taken over the compact subsets $A$ of $S$.

\textbf{3.2.3.} The writer does not know whether the finiteness theorem (3.2.1) remains valid if, instead of assuming $S$ locally compact, one merely assumes that each point of $S$ has a fundamental system of paracompact neighborhoods. Moreover, the local compactness hypothesis on $S$ does not seem very natural. The space of ``universal local parameters'' for elliptic complexes of fixed order is in fact a subspace of infinite dimension in a Fréchet space. More precisely, let $V$ be a compact variety, let $F=(F^i)$ be a sequence of complex vector bundles of finite rank on $V$, zero except for finitely many indices, and let $n=(n_i)$ be a sequence of integers $\ge0$. Let $\operatorname{Comp.diff}_n(F)$ be the space of complexes of differential operators of base $F$ and order $n$, that is, the closed subspace of
\[
 \prod_i \Diff^{n_i}(F^i,F^{i+1})
\]
defined by the equations $d^{i+1}d^i=0$. This product is a Fréchet space. The subset $\operatorname{Ell}_n(F)$ formed by the elliptic complexes of order $n$ is open. It is this space, in general not locally compact, which plays the role of a space of universal local parameters.

\textbf{3.2.4.} Nevertheless, even if the generalization of (3.2.1) just mentioned is impossible, one can always associate to $E^\bullet$, under the sole hypothesis that the fibration $p$ is proper and locally trivial, an element of $K^0(S)$ deserving the name analytic index of $E^\bullet$. Indeed, suppose, for simplicity, that we are in the local situation where $X=S\times V$ with $V$ compact, and $E^\bullet$ is given by a continuous map from $S$ into $\operatorname{Ell}_n(F)$ (3.2.3). Assertions (a') and (b) in the proof of (3.2.1) are in fact valid without any hypothesis on $S$; it is therefore reasonable to put, with the notation of that proof,
\[
 \ind_{\mathrm{an}}(E^\bullet)=[\mathcal F_r^\bullet]\in K^0(S),
\]
this element being independent of $r$.

\textbf{3.2.5.} Suppose $S$ paracompact. Let $H$ be a complex Hilbert space of infinite Hilbert dimension, and let $I(H)$ be the subspace of $L(H,H)$ formed by the Fredholm operators. The complexes $F_r^\bullet$ define [12] a canonical element in the group $[S,I(H)]$ of homotopy classes of continuous maps from $S$ to $I(H)$. But the writer does not know, for non-compact $S$ (even when $S$ is locally compact), what the relations are between $[S,I(H)]$ and $K^0(S)$: is there a natural arrow from one to the other? are they isomorphic? A fortiori, the relation between the two indices thus defined remains unclear.

For a generalization, under the hypotheses of (3.2.1), of the Atiyah--Singer index theorem, we refer to the recent work of Shih [17], or to work in preparation of M. F. Atiyah and I. M. Singer.\footnote{For the general case, use the notion of ``weak Hilbert bundle'', or continuous field of Hilbert spaces, of [12].}

\subsection*{Bibliography}

\begin{enumerate}[label={[\,\arabic*\,]}]
\item Artin, M., Grothendieck, A., and Verdier, J.-L., \emph{Cohomologie étale des schémas}, to appear with North-Holland Publishing Company (cited as SGA 4).
\item Atiyah, M., and Segal, G., secret papers.
\item Cartan, H., \emph{Funzioni e varietà complesse. Faisceaux analytiques cohérents}, course at the Centro Internazionale Matematico Estivo, 1964, Rome.
\item Cartan, H., and Schwartz, L., \emph{Théorème d'Atiyah--Singer}, Seminar 1963--64, Institut Henri Poincaré, Paris.
\item Frisch, J., Fonctions analytiques sur un ensemble semi-analytique, note in C.R.A.S. Paris, t. 260, pp. 2974--2976 (1965).
\item Godement, R., \emph{Théorie des faisceaux}, Hermann (1958), cited as T.F.
\item Grothendieck, A., \emph{Cohomologie locale des faisceaux cohérents et théorèmes de Lefschetz locaux et globaux}, Séminaire de Géométrie Algébrique de l'I.H.E.S. 1962, to appear with North-Holland Publishing Company, cited as SGA 2.
\item Grothendieck, A., \emph{Fondements de la géométrie analytique}, Séminaire Cartan 1959--60, Institut Henri Poincaré, Paris.
\item Hartshorne, R., \emph{Residues and Duality}, Lecture Notes in Mathematics 20, Springer, cited as [H].
\item Houzel, C., \emph{Géométrie analytique locale}, Séminaire Cartan 1959--60, Institut Henri Poincaré, Paris.
\item Illusie, L., \emph{Complexes de fibrés et définitions du foncteur $K$}, Séminaire Shih Weishu, I.H.E.S. 1964--65.
\item Illusie, L., Complexes quasi-acycliques directs de fibrés banachiques, note in C.R.A.S. Paris, t. 260, pp. 6499--6502.
\item Kiehl, R., Theorem A und Theorem B in der nichtarchimedischen Funktionentheorie, \emph{Inventiones Mathematicae} 2, pp. 256--273 (1967).
\item Lang, S., \emph{Introduction to differentiable manifolds}.
\item Łojasiewicz, S., \emph{Annali della Scuola Normale Superiore di Pisa}, Serie III, 18, fasc. IV, 1964.
\item Palais, R., lectures at the Princeton Seminar 1963--64 on the index theorem.
\item Shih, W., \emph{Fiber cobordism and the index of a family of elliptic differential operators}.
\item Verdier, J.-L., \emph{Les catégories dérivées des catégories abéliennes}, to appear with North-Holland Publishing Company.
\end{enumerate}


\end{document}
